\documentclass[11pt]{book}

%---------------------------------------
% Packages
%---------------------------------------
% === Document Setup and Language ===
\usepackage[english]{babel}        % English language
\usepackage[utf8]{inputenc}        % Unicode support
\usepackage[T1]{fontenc}           % Font encoding
\usepackage{subfiles}              % Allow for subfiles which can be compiled standalone
\setcounter{secnumdepth}{2}

% === Math Packages ===
\usepackage{amsfonts, amsmath, amssymb, amsthm}  % AMS math packages
\usepackage{dsfont}                % For \mathds{1}
\usepackage{mathtools}             % Enhanced math tools
\usepackage{accents}               % Mathematical accents

% === Typography and Fonts ===
%\usepackage{mathptmx}              % Times New Roman font
\usepackage{newtxtext,newtxmath}   % Enhanced Times New Roman with math support
\usepackage{csquotes}              % Smart quotes and quotation management
\linespread{1.05}                  % Line spacing

% === Graphics and Diagrams ===
\usepackage{tikz-cd, tikz}         % Diagrams and drawings
\usetikzlibrary{
	positioning, 	% For relative positioning (below=of, right=of)
	calc,        	% For coordinate calculations (finding midpoints)
	fit,			% Needed for the bounding box (fit)
	patterns,
	decorations.markings,
	decorations.pathmorphing,	% Needed for squiggly arrows
	backgrounds
}
\usepackage{graphicx}              % Image handling
\usepackage{booktabs}              % \toprule, \midrule, \bottomrule for tables

% === Page Layout ===
\usepackage{geometry}
\geometry{left=1.3in,right=1.2in}
\setlength{\parskip}{\smallskipamount}
\setlength{\parindent}{0pt}

% === Lists and Enumeration ===
\usepackage{enumerate}             % Enhanced enumeration
\renewcommand{\labelenumi}{\normalfont{(\arabic{enumi})}} % Default to [(1)]
\usepackage{etoolbox}              % Programming tools for LaTeX

% === Hyperlinks and References ===
\usepackage[hyphens]{url}
\urlstyle{same}
\usepackage{bookmark}
\usepackage{hyperref}              % Must be loaded after most other packages
\hypersetup{
	colorlinks=true,
	linkcolor=black,
	filecolor=black,
	urlcolor=black,
	citecolor=black,
	bookmarksdepth=subsubsection
}

% === Bibliography ===
\usepackage[
backend=biber,
natbib,
style=alphabetic,
maxcitenames=4,
maxalphanames=5,
maxbibnames=99
]{biblatex}
\renewbibmacro{in:}{}              % Suppresses "In: " before journal names
\addbibresource{C:/Users/LocalAdmin/OneDrive/Documenten OneDrive/A - Regensburg/Research and notes/Main_Bibliography.bib}
%\addbibresource{Bibliography.bib}

% === Optional Settings (uncomment if needed) ===
% \setcounter{secnumdepth}{2}
% \setcounter{tocdepth}{2}

%---------------------------------------
% Gestalten
%---------------------------------------
\DeclareMathOperator{\fpqc}{fpqc}

%---------------------------------------
% Algebraic geometry / motivic cohomology
%---------------------------------------
\DeclareMathOperator{\KGL}{KGL}
\DeclareMathOperator{\Nis}{Nis}
\DeclareMathOperator{\Et}{Et}
\DeclareMathOperator{\cdh}{cdh}
\DeclareMathOperator{\vdim}{vdim}
\DeclareMathOperator{\Krdim}{Kr.dim}
\DeclareMathOperator{\qcqs}{qcqs}
\DeclareMathOperator{\Sm}{Sm}
\DeclareMathOperator{\SSeq}{SSeq}
\DeclareMathOperator{\Sym}{Sym}
\DeclareMathOperator{\KH}{KH}

%---------------------------------------
% Topos theory
%---------------------------------------
\newcommand{\ssquare}{\,\square\,}
\DeclareMathOperator{\Aff}{Aff}
\DeclareMathOperator{\Mono}{Mono}
\DeclareMathOperator{\Gerb}{Gerb}
\DeclareMathOperator{\hyp}{hyp}
\DeclareMathOperator{\mono}{mono}
\DeclareMathOperator{\cons}{cons}
\DeclareMathOperator{\wcons}{wcons}
\DeclareMathOperator{\quottext}{quot}
\DeclareMathOperator{\epi}{epi}
\DeclareMathOperator{\MonoCong}{MonoCong}
\DeclareMathOperator{\EpiCong}{EpiCong}
\DeclareMathOperator{\Cong}{Cong}
\DeclareMathOperator{\HypCong}{HypCong}
\DeclareMathOperator{\Mdl}{Mdl}
\DeclareMathOperator{\Acyc}{Acyc}
\DeclareMathOperator{\Sat}{Sat}
\DeclareMathOperator{\SSat}{SSat}
\DeclareMathOperator{\Conn}{Conn}
\DeclareMathOperator{\EffEpi}{EffEpi}
\DeclareMathOperator{\Epi}{Epi}
\DeclareMathOperator{\Cotr}{Cotr}
\DeclareMathOperator{\All}{All}
\DeclareMathOperator{\Pos}{Pos}
\DeclareMathOperator{\post}{post}
\DeclareMathOperator{\El}{El}
\DeclareMathOperator{\cocompl}{cocompl}
\newcommand{\Topos}{\catname{Topos}}
\newcommand{\Logos}{\catname{Logos}}
\newcommand{\bbCat}{\bbC \mathrm{at}}
\newcommand{\bbLog}{\bbL \mathrm{ogos}}
\newcommand{\bbTop}{\bbT \mathrm{opos}}
%\renewcommand{\top}{\mathrm{top}}
\DeclareMathOperator{\Geom}{Geom}
\DeclareMathOperator{\Pt}{Pt}
\DeclareMathOperator{\im}{im}				% Image
\DeclareMathOperator{\coim}{coim}			% Coimage
\DeclareMathOperator{\Act}{Act}
\DeclareMathOperator{\Bun}{Bun}
\DeclareMathOperator{\Site}{Site}
\DeclareMathOperator{\cont}{cont}
\DeclareMathOperator{\cocont}{cocont}
\DeclareMathOperator{\Cont}{Cont}
\DeclareMathOperator{\wccolim}{w.c.colim}

%---------------------------------------
% Miscellaneous
%---------------------------------------
\DeclareMathOperator{\Atl}{Atl}
\DeclareMathOperator{\DiffStk}{DiffStk}
\DeclareMathOperator{\SepStk}{SepStk}
\renewcommand{\H}{\mathrm{H}}
\newcommand{\SH}{\mathrm{SH}}
\DeclareMathOperator{\htp}{htp}
\DeclareMathOperator{\rep}{rep}
\DeclareMathOperator{\Loc}{Loc}
\DeclareMathOperator{\Thick}{Thick}
\DeclareMathOperator{\modtext}{mod}
\newcommand{\absmod}[1]{\abs{#1}_{\modtext}}
\DeclareMathOperator{\sep}{sep}
\DeclareMathOperator{\PB}{PB}
\DeclareMathOperator{\PF}{PF}
\DeclareMathOperator{\NerveDiff}{N^{\mathrm{Diff}}}
\newcommand{\Deltaalg}[1]{\Delta^{#1}_{\mathrm{alg}}}
%\newcommand{\Act}[1]{#1 \textup{Act}}
\DeclareMathOperator{\shear}{shear}
\DeclareMathOperator{\Sph}{Sph}
\DeclareMathOperator{\Sub}{Sub}
\DeclareMathOperator{\Open}{Open}
\DeclareMathOperator{\open}{open}
\DeclareMathOperator{\Corr}{Corr}
\DeclareMathOperator{\bbMap}{\mathbb{M} ap}
\DeclareMathOperator{\bbHom}{\mathbb{H} om}
\DeclareMathOperator{\gen}{gen}
\DeclareMathOperator{\Bgl}{B_{gl}}
\newcommand{\PrnBdl}[2]{\textup{PrnBdl}_{#1}(#2)}
\newcommand{\catnameb}[1]{{\normalfont\textbf{#1}}}
\DeclareMathOperator{\bLieGrpd}{\catnameb{LieGrpd}}
\newcommand{\cFun}{\Ff \textup{un}}
\DeclareMathOperator{\Tr}{Tr}
\DeclareMathOperator{\tr}{tr}
\DeclareMathOperator{\Shp}{Shp}
\DeclareMathOperator{\THH}{THH}

%---------------------------------------
% Higher Algebra
%---------------------------------------

\newcommand{\Assoc}{\mathcal{A}\mathit{ssoc}}
\newcommand{\Comm}{\mathcal{C}\!\,\mathit{omm}}
\newcommand{\Lie}{\mathcal{L}\mathit{ie}}
\newcommand{\Triv}{\mathcal{T}\!\!\mathit{riv}}
\newcommand{\oTens}{\mathcal{T}\!\!\mathit{ens}}
\newcommand{\oLMod}{\mathcal{L}\!\mathcal{M}\mathit{od}}
\newcommand{\oRMod}{\mathcal{R}\!\mathcal{M}\mathit{od}}
\newcommand{\oBMod}{\mathcal{B}\!\mathcal{M}\mathit{od}}
\newcommand{\oMod}{\mathcal{M}\mathit{od}}
\newcommand{\oAlg}{\mathcal{A}\mathit{lg}}
\newcommand{\oEnd}{\mathcal{E}\mathit{nd}}
\newcommand{\oCMon}{\mathcal{C}\!\mathcal{M}\mathit{on}}
\newcommand{\oCGrp}{\mathcal{C}\mathcal{G}\mathit{rp}}
\newcommand{\oSp}{\mathcal{S}\mathit{p}}
\newcommand{\Mult}{\mathcal{M}}
\DeclareMathOperator{\oCat}{\mathcal{C}\!\mathit{at}}
\DeclareMathOperator{\oFun}{\mathcal{F}\!\mathit{un}}
%\newcommand{\Mult}{\mathcal{M}\mathit{ult}}
\DeclareMathOperator{\Env}{Env}
\DeclareMathOperator{\Op}{Op}
\DeclareMathOperator{\AdTrip}{AdTrip}
\DeclareMathOperator{\rev}{rev}
\DeclareMathOperator{\tiso}{iso}
\DeclareMathOperator{\oTw}{\overline{\Tw}}
\DeclareMathOperator{\lax}{lax}
\DeclareMathOperator{\NTw}{NTw}
\newcommand{\OpCocart}{\Oo^{\amalg}}
\newcommand{\OpCart}{\Oo^{\times}}
\DeclareMathOperator{\Day}{Day}
\newcommand{\oDay}{\mathcal{D}\mathit{ay}}
\newcommand{\dlax}{\textup{-lax}}
\newcommand{\textred}{\textup{red}}
\newcommand{\exc}{\textup{exc}}
\DeclareMathOperator{\BiFun}{BiFun}
\DeclareMathOperator{\BV}{BV}
\renewcommand{\Bar}{\mathrm{Bar}}
\newcommand{\dnil}{\textup{-nil}}
\DeclareMathOperator{\Ext}{Ext}
\DeclareMathOperator{\Ex}{Ex}
\DeclareMathOperator{\Tor}{Tor}
\newcommand{\enr}{\textup{enr}}
\newcommand{\heart}{\heartsuit}
\newcommand{\dsm}{\textup{-sm}}
\DeclareMathOperator{\Free}{Free}

%---------------------------------------
% Stable Homotopy Theory
%---------------------------------------
\DeclareMathOperator{\ku}{ku}
\DeclareMathOperator{\KU}{KU}
\DeclareMathOperator{\MU}{MU}
\DeclareMathOperator{\MO}{MO}
\DeclareMathOperator{\MSO}{MSO}
\DeclareMathOperator{\BO}{BO}
\DeclareMathOperator{\bGr}{\mathbf{Gr}}
\DeclareMathOperator{\BUP}{BUP}
\DeclareMathOperator{\bBUP}{\mathbf{BUP}}
\DeclareMathOperator{\Th}{Th}
\DeclareMathOperator{\Pic}{Pic}
\DeclareMathOperator{\pic}{pic}
\DeclareMathOperator{\pre}{pre}
\DeclareMathOperator{\Line}{Line}
%\DeclareMathOperator{\Triv}{Triv}
\DeclareMathOperator{\PT}{PT}
\newcommand{\interiorsymb}[1]{ {\kern0pt#1}^{\mathrm{o}}}
\newcommand{\PSp}{\catname{PSp}}
\DeclareMathOperator{\Emb}{Emb}
\newcommand{\Ntop}{N^{\textup{top}}}
\newcommand{\Homtop}{\Hom^{\textup{top}}}
\DeclareMathOperator{\fw}{fw}
\DeclareMathOperator{\Shift}{Shift}
\DeclareMathOperator{\Der}{Der}

%---------------------------------------
% Category Theory Operators
%---------------------------------------

% Maps and functors
\DeclareMathOperator{\Hom}{Hom}
\DeclareMathOperator{\Aut}{Aut}
\DeclareMathOperator{\End}{End}
\DeclareMathOperator{\iHom}{\underline{\operatorname{Hom}}}			% Internal Hom
\DeclareMathOperator{\ihom}{\underline{\operatorname{hom}}}			% Internal Hom
\DeclareMathOperator{\uHom}{\underline{\operatorname{Hom}}}			% Internal Hom
\DeclareMathOperator{\Map}{Map}
\DeclareMathOperator{\map}{map}
\DeclareMathOperator{\Path}{Path}
\DeclareMathOperator{\Fun}{Fun}
\newcommand{\FunL}{\Fun^{\textup{L}}}
\newcommand{\FunR}{\Fun^{\textup R}}
\newcommand{\ulFun}{\ul{\Fun}}
\DeclareMathOperator{\Nat}{Nat}

% Fibrations and Related Concepts
\DeclareMathOperator{\Lift}{Lift}
\DeclareMathOperator{\lift}{lift}
\DeclareMathOperator{\cocart}{cocart}
\DeclareMathOperator{\CoCart}{CoCart}
\DeclareMathOperator{\Cocart}{Cocart}
\DeclareMathOperator{\isCocart}{\texttt{isCocart}}
\DeclareMathOperator{\Cart}{Cart}
\DeclareMathOperator{\RFib}{RFib}
\DeclareMathOperator{\LFib}{LFib}
\DeclareMathOperator{\Un}{Un}
\DeclareMathOperator{\cart}{cart}
\newcommand{\dcart}{\textup{-cart}}
\DeclareMathOperator{\cc}{cc}
\DeclareMathOperator{\ct}{ct}
\DeclareMathOperator{\Str}{Str}
\DeclareMathOperator{\el}{el}

%---------------------------------------
% Algebraic Geometry
%---------------------------------------
\DeclareMathOperator{\Spec}{Spec}
\DeclareMathOperator{\QCoh}{QCoh}
\DeclareMathOperator{\IndCoh}{IndCoh}
\DeclareMathOperator{\Perf}{Perf}
\DeclareMathOperator{\Proj}{Proj}
\DeclareMathOperator{\Zar}{Zar}
\DeclareMathOperator{\Sch}{Sch}
\newcommand{\ad}{\mathrm{ad}}
\newcommand{\corp}{\mathrm{corp}}
\newcommand{\loc}{\mathrm{loc}}
\newcommand{\mot}{\mathrm{mot}}
\newcommand{\ultra}{\mathrm{ultra}}

%---------------------------------------
% Other math-operators
%---------------------------------------
\DeclareMathOperator{\Ar}{Ar}
\DeclareMathOperator{\BC}{BC}
\DeclareMathOperator{\GL}{GL}
\DeclareMathOperator{\Gr}{Gr}
\DeclareMathOperator{\Ho}{Ho}
\let\Im\relax
\DeclareMathOperator{\Im}{Im}
\DeclareMathOperator{\Ind}{Ind}
\DeclareMathOperator{\Iso}{Iso}
\DeclareMathOperator{\LConst}{LConst}
\DeclareMathOperator{\Nm}{Nm}
\DeclareMathOperator{\Pro}{Pro}
\DeclareMathOperator{\Lan}{Lan}
\DeclareMathOperator{\Ran}{Ran}
\DeclareMathOperator{\Rex}{Rex}
\DeclareMathOperator{\PSh}{PSh}
\DeclareMathOperator{\Shape}{Shape}
\DeclareMathOperator{\Shv}{Shv}
\usepackage{relsize}
\newcommand{\simp}{\mathlarger{\boldsymbol{\Delta}}}
\newcommand{\bOmega}{{\boldsymbol{\Omega}}}
\DeclareMathOperator{\Span}{Span}
\newcommand{\SpanTwo}{\Span_2}
\newcommand{\SpanTwoP}[4]{\SpanTwo(#1,#2)_{#4,#3}}
\DeclareMathOperator{\NSpan}{NSpan}
\DeclareMathOperator{\Tw}{Tw}
\DeclareMathOperator{\LM}{LM}
\let\Re\relax
\DeclareMathOperator{\Re}{Re}
\DeclareMathOperator{\Fact}{Fact}

\DeclareMathOperator{\SD}{SD}				% Suave dual
\DeclareMathOperator{\PD}{PD}				% Prim dual
\DeclareMathOperator{\Suave}{Suave}			% Suave objects
\DeclareMathOperator{\Prim}{Prim}			% Prim objects

\DeclareMathOperator{\shift}{shift}
\DeclareMathOperator{\act}{act}
\DeclareMathOperator{\supp}{supp}
\DeclareMathOperator{\ex}{ex}
\DeclareMathOperator{\lex}{lex}
\DeclareMathOperator{\rex}{rex}
\DeclareMathOperator{\chain}{chain}
\DeclareMathOperator{\sd}{sd}
\DeclareMathOperator{\qc}{qc}
\DeclareMathOperator{\Poly}{Poly}
\DeclareMathOperator{\Mot}{Mot}

%---------------------------------------
% Abbreviations
%---------------------------------------

\newcommand{\can}{\mathrm{can}}             % Canonical
\newcommand{\ho}{\mathrm{ho}}               % Homotopy category
\newcommand{\nondeg}{\mathrm{nondeg}}       % Non-degenerate
\newcommand{\univ}{\mathrm{univ}}           % Universal
\newcommand{\prop}{\textup{pr}}
\newcommand{\inv}{\textup{inv}}
\DeclareMathOperator{\ac}{ac}				% Associated category
\DeclareMathOperator{\pb}{pb}				% Pullback

\DeclareMathOperator{\fgt}{fgt}				% Forgetful functor
\DeclareMathOperator{\id}{id}				% Identity
\DeclareMathOperator{\pr}{pr}				% Projection
\DeclareMathOperator{\incl}{incl}				% Inclusion
\DeclareMathOperator{\const}{const}			% Forgetful
\DeclareMathOperator{\sk}{sk}				% Skeleton
\DeclareMathOperator{\cosk}{cosk}			% Coskeleton
\DeclareMathOperator{\gl}{gl}               % Global
\DeclareMathOperator{\grp}{grp}             % Group completion
\DeclareMathOperator{\qis}{qis}             % Quasi-isomorphism
\DeclareMathOperator{\cyl}{cyl}             % Cylinder
\DeclareMathOperator{\surj}{surj}             % Surjective
\newcommand{\et}{\textup{ét}}				% Étale
\newcommand{\coh}{\mathrm{coh}}				% Coherent
\newcommand{\acc}{\mathrm{acc}}				% Accessible
\renewcommand{\small}{\mathrm{small}}

\DeclareMathOperator{\free}{free}
\DeclareMathOperator{\fin}{fin}
\DeclareMathOperator{\pt}{pt}
\DeclareMathOperator{\aug}{{aug}}
\DeclareMathOperator{\st}{st}
\DeclareMathOperator{\dual}{dual}
\DeclareMathOperator{\Calk}{Calk}
\DeclareMathOperator{\idem}{idem}
\DeclareMathOperator{\perf}{perf}
\newcommand{\rig}{\mathrm{rig}}
\DeclareMathOperator{\Nuc}{Nuc}

\DeclareMathOperator{\cofib}{cofib}
\DeclareMathOperator{\fib}{fib}
\DeclareMathOperator{\colim}{colim}
\let\lim\relax
\DeclareMathOperator{\lim}{lim}
\DeclareMathOperator{\coeq}{coeq}
\DeclareMathOperator{\coker}{coker}
\newcommand{\ev}{\textup{ev}}
\newcommand{\coev}{\textup{coev}}
\DeclareMathOperator{\res}{res}
\DeclareMathOperator{\ind}{ind}
\DeclareMathOperator{\coind}{coind}

%---------------------------------------
% Categories
%---------------------------------------
\newcommand{\catname}[1]{{\textup{#1}}}
\newcommand{\Set}{\catname{Set}}
\newcommand{\An}{\catname{An}}
\newcommand{\CondAn}{\catname{CondAn}}
\newcommand{\Ab}{\catname{Ab}}
\newcommand{\Alg}{\catname{Alg}}
\newcommand{\CAlg}{\catname{CAlg}}
\newcommand{\Cat}{\catname{Cat}}
\newcommand{\CAT}{\catname{CAT}}
\newcommand{\CRing}{\catname{CRing}}
\newcommand{\qCat}{\catname{qCat}}
\newcommand{\CGrp}{\catname{CGrp}}
\newcommand{\CMon}{\catname{CMon}}
\newcommand{\Fin}{\catname{Fin}}
\newcommand{\FinGrpd}{\catname{FinGrpd}}
\newcommand{\Grp}{\catname{Grp}}
\newcommand{\Grpd}{\catname{Grpd}}
\newcommand{\Mod}{\catname{Mod}}
\newcommand{\coMod}{\catname{coMod}}
\newcommand{\Mon}{\catname{Mon}}
\newcommand{\LMod}{\catname{LMod}}
\newcommand{\RMod}{\catname{RMod}}
\newcommand{\BMod}{\catname{BMod}}
\newcommand{\Rep}{\catname{Rep}}
\renewcommand{\Top}{\catname{Top}}
\newcommand{\TopGrpd}{\catname{TopGrpd}}
\newcommand{\bTopGrpd}{\mathbf{TopGrpd}}
\newcommand{\Spc}{\catname{Spc}}
\newcommand{\Sp}{\catname{Sp}}
\newcommand{\Vect}{\catname{Vect}}
\newcommand{\Orb}{\textup{Orb}}
\newcommand{\Orbit}{\textup{Orbit}}
\newcommand{\Glo}{\textup{Glo}}

\newcommand{\sAn}{\s\An}
\newcommand{\dAn}{\mathrm{d}\An}
\newcommand{\bVect}{\mathbf{V}\catname{ect}}
\newcommand{\LCHaus}{\catname{LCHaus}}
\newcommand{\CHaus}{\catname{CHaus}}
\newcommand{\ProFin}{\catname{ProFin}}
\newcommand{\AnRing}{\catname{AnRing}}
\newcommand{\AnStck}{\catname{AnStck}}
\newcommand{\Ring}{\catname{Ring}}
\newcommand{\twoCAlg}{\catname{2CAlg}}
\newcommand{\Idem}{\catname{Idem}}
\newcommand{\coIdem}{\catname{coIdem}}
\newcommand{\StRing}{\catname{StRing}}
\newcommand{\PreStRing}{\catname{PreStRing}}
\newcommand{\Gest}{\catname{Gest}}
\newcommand{\RingGest}{\catname{RingGest}}
\newcommand{\LieGrpd}{\catname{LieGrpd}}
\newcommand{\sTop}{\catname{sTop}}
\newcommand{\Diff}{\catname{Diff}}
\newcommand{\sSet}{\catname{sSet}}
\newcommand{\Poset}{\catname{Poset}}

\renewcommand{\Pr}{\textup{Pr}}
\newcommand{\PrL}{\textup{Pr}^{\textup{L}}}
\newcommand{\PrR}{\textup{Pr}^{\textup{R}}}
\newcommand{\PrLst}{\PrL_{\st}}

\def\ulSp{\smash{\ul{{\setbox0=\hbox{\Sp}\dp0=1pt \ht0=0pt \box0\relax}}}}
\def\ulSpc{\smash{\ul{{\setbox0=\hbox{\Spc}\dp0=1pt \ht0=0pt \box0\relax}}}}
\def\ulOrbSp{\smash{\ul{{\setbox0=\hbox{\Orb\Sp}\dp0=1pt \ht0=0pt \box0\relax}}}}
\def\ulOrbSpc{\smash{\ul{{\setbox0=\hbox{\Orb\Spc}\dp0=1pt \ht0=0pt \box0\relax}}}}

%---------------------------------------
% Miscellaneous
%---------------------------------------

\newcommand{\h}{\textup{h}}
\newcommand{\CW}{\mathrm{CW}}
\DeclareMathOperator{\hocolim}{hocolim}
\DeclareMathOperator{\holim}{holim}
\DeclareMathOperator{\Exc}{Exc}
\newcommand{\s}{\textup{s}}
\DeclareMathOperator{\Sing}{Sing}
\DeclareMathOperator{\Kan}{Kan}
\newcommand{\bKan}{\mathbf{Kan}}
\newcommand{\bqCat}{\mathbf{qCat}}
\DeclareMathOperator{\Ch}{\catname{Ch}}
\DeclareMathOperator{\inj}{inj}
\DeclareMathOperator{\all}{all}
\DeclareMathOperator{\isotext}{iso}
\DeclareMathOperator{\Cut}{Cut}
\DeclareMathOperator{\LCut}{LCut}
\DeclareMathOperator{\BCut}{BCut}
\DeclareMathOperator{\Red}{Red}
\DeclareMathOperator{\Segal}{Segal}
\DeclareMathOperator{\Seg}{Seg}
\DeclareMathOperator{\CSeg}{CSeg}
\DeclareMathOperator{\add}{add}
\DeclareMathOperator{\sadd}{sadd}
\newcommand{\Piinfty}[1]{\Pi_{\infty}(#1)}
\DeclareMathOperator{\oplax}{oplax}

\newcommand{\qednow}{\pushQED{\qed}\qedhere\popQED} % Use this 

\newcommand{\quadtext}[1]{\quad\textrm{#1}\quad}
\newcommand{\qquadtext}[1]{\qquad\textrm{#1}\qquad}
\newcommand{\qin}{\quad\in\quad}
\newcommand{\ul}[1]{\underline{#1}}
\newcommand{\abs}[1]{\lvert #1 \rvert}
\newcommand{\geom}[1]{\| #1 \|}
\newcommand{\floor}[1]{\lfloor #1 \rfloor}

\newcommand{\catop}{^{\textup{op}}}
\let\op\relax
\DeclareMathOperator{\op}{op}

\DeclareMathOperator{\unit}{\mathds{1}}    % unit for \otimes
\newcommand{\blank}{{\textup{ --}}}
\newcommand{\ceil}[1]{\lceil #1 \rceil}
\newcommand{\lra}[1]{\langle #1 \rangle}

\newcommand{\iso}{\xrightarrow{\;\smash{\raisebox{-0.5ex}{\ensuremath{\scriptstyle\sim}}}\;}}
% \newcommand{\mapsfrom}{\mathrel{\reflectbox{$\mapsto$}}}
\newcommand{\xrightrightarrows}[2]{\begin{tikzcd}[ampersand replacement=\&, cramped] \vphantom{x}\rar[shift left=0.45ex, "#1"] \rar[shift right=0.45ex, swap, "#2"]\& \vphantom{x} \end{tikzcd}}
\newcommand{\xadjto}[2]{\begin{tikzcd}[ampersand replacement=\&, cramped] \vphantom{x}\rar[shift left=0.45ex, "#1"]\& \lar[shift left=0.45ex, "#2"] \vphantom{x} \end{tikzcd}}

%---------------------------------------
% Alphabets
%---------------------------------------
\let\S\relax
\let\P\relax
\DeclareMathOperator{\C}{\mathbb C}
\DeclareMathOperator{\F}{\mathbb F}
\DeclareMathOperator{\N}{\mathbb N}
\DeclareMathOperator{\P}{\mathbb P}
\DeclareMathOperator{\Q}{\mathbb Q}
\DeclareMathOperator{\R}{\mathbb R}
\DeclareMathOperator{\S}{\mathbb S}
\DeclareMathOperator{\Z}{\mathbb Z}
\DeclareMathOperator{\D}{\mathrm{D}}
\renewcommand{\phi}{\varphi}
\renewcommand{\epsilon}{\varepsilon}
\renewcommand{\aa}{\mathfrak{a}}
\newcommand{\pp}{\mathfrak{p}}
\newcommand{\fa}{\mathfrak{a}}
\newcommand{\fb}{\mathfrak{b}}
\newcommand{\fc}{\mathfrak{c}}
\newcommand{\fm}{\mathfrak{m}}
\newcommand{\mm}{\mathfrak{m}}
\newcommand{\fn}{\mathfrak{n}}
\newcommand{\fp}{\mathfrak{p}}

% Caligraphic font
\newcommand{\Aa}{\mathcal{A}}
\newcommand{\Bb}{\mathcal{B}}
\newcommand{\Cc}{\mathcal{C}}
\newcommand{\Dd}{\mathcal{D}}
\newcommand{\Ee}{\mathcal{E}}
\newcommand{\Ff}{\mathcal{F}}
\newcommand{\Gg}{\mathcal{G}}
\newcommand{\Hh}{\mathcal{H}}
\newcommand{\Ii}{\mathcal{I}}
\newcommand{\Jj}{\mathcal{J}}
\newcommand{\Kk}{\mathcal{K}}
\newcommand{\Ll}{\mathcal{L}}
\newcommand{\Mm}{\mathcal{M}}
\newcommand{\Nn}{\mathcal{N}}
\newcommand{\Oo}{\mathcal{O}}
\newcommand{\Pp}{\mathcal{P}}
\newcommand{\Qq}{\mathcal{Q}}
\newcommand{\Rr}{\mathcal{R}}
\newcommand{\Ss}{\mathcal{S}}
\newcommand{\Tt}{\mathcal{T}}
\newcommand{\Uu}{\mathcal{U}}
\newcommand{\Vv}{\mathcal{V}}
\newcommand{\Ww}{\mathcal{W}}
\newcommand{\Xx}{\mathcal{X}}
\newcommand{\Yy}{\mathcal{Y}}
\newcommand{\Zz}{\mathcal{Z}}

% Bold font
\newcommand{\bA}{\mathbf{A}}
\newcommand{\bB}{\mathbf{B}}
\newcommand{\bC}{\mathbf{C}}
\newcommand{\bD}{\mathbf{D}}
\newcommand{\bE}{\mathbf{E}}
\newcommand{\bF}{\mathbf{F}}
\newcommand{\bG}{\mathbf{G}}
\newcommand{\bH}{\mathbf{H}}
\newcommand{\bI}{\mathbf{I}}
\newcommand{\bJ}{\mathbf{J}}
\newcommand{\bK}{\mathbf{K}}
\newcommand{\bL}{\mathbf{L}}
\newcommand{\bM}{\mathbf{M}}
\newcommand{\bN}{\mathbf{N}}
\newcommand{\bO}{\mathbf{O}}
\newcommand{\bP}{\mathbf{P}}
\newcommand{\bQ}{\mathbf{Q}}
\newcommand{\bR}{\mathbf{R}}
\newcommand{\bS}{\mathbf{S}}
\newcommand{\bT}{\mathbf{T}}
\newcommand{\bU}{\mathbf{U}}
\newcommand{\bV}{\mathbf{V}}
\newcommand{\bW}{\mathbf{W}}
\newcommand{\bX}{\mathbf{X}}
\newcommand{\bY}{\mathbf{Y}}
\newcommand{\bZ}{\mathbf{Z}}

% Blackboard bold
\newcommand{\bbA}{\mathbb{A}}
\newcommand{\bbB}{\mathbb{B}}
\newcommand{\bbC}{\mathbb{C}}
\newcommand{\bbD}{\mathbb{D}}
\newcommand{\bbE}{\mathbb{E}}
\newcommand{\bbF}{\mathbb{F}}
\newcommand{\bbG}{\mathbb{G}}
\newcommand{\bbH}{\mathbb{H}}
\newcommand{\bbI}{\mathbb{I}}
\newcommand{\bbJ}{\mathbb{J}}
\newcommand{\bbK}{\mathbb{K}}
\newcommand{\bbL}{\mathbb{L}}
\newcommand{\bbM}{\mathbb{M}}
\newcommand{\bbN}{\mathbb{N}}
\newcommand{\bbO}{\mathbb{O}}
\newcommand{\bbP}{\mathbb{P}}
\newcommand{\bbQ}{\mathbb{Q}}
\newcommand{\bbR}{\mathbb{R}}
\newcommand{\bbS}{\mathbb{S}}
\newcommand{\bbT}{\mathbb{T}}
\newcommand{\bbU}{\mathbb{U}}
\newcommand{\bbV}{\mathbb{V}}
\newcommand{\bbW}{\mathbb{W}}
\newcommand{\bbX}{\mathbb{X}}
\newcommand{\bbY}{\mathbb{Y}}
\newcommand{\bbZ}{\mathbb{Z}}

%---------------------------------------
% Environments
%---------------------------------------

% Theorem-style environments
\theoremstyle{plain}
\newtheorem{theorem}{Theorem}[chapter]

\newtheorem{conjecture}[theorem]{Conjecture}
\newtheorem{corollary}[theorem]{Corollary}
\newtheorem{lemma}[theorem]{Lemma}
\newtheorem{proposition}[theorem]{Proposition}
\newtheorem{thmx}{Theorem}
\renewcommand{\thethmx}{\Alph{thmx}}

% Definition-style environments
\theoremstyle{definition}
\newtheorem{axiom}{Axiom}
\newtheorem{claim}[theorem]{Claim}
\newtheorem{construction}[theorem]{Construction}
\newtheorem{convention}[theorem]{Convention}
\newtheorem{definition}[theorem]{Definition}
\newtheorem{discussion}[theorem]{Discussion}
\newtheorem{example}[theorem]{Example}
\newtheorem{exercise}[theorem]{Exercise}
\newtheorem{fact}[theorem]{Fact}
\newtheorem{notation}[theorem]{Notation}
\newtheorem{observation}[theorem]{Observation}
\newtheorem{question}[theorem]{Question}
\newtheorem{remark}[theorem]{Remark}
\newtheorem{terminology}[theorem]{Terminology}
\newtheorem{warning}[theorem]{Warning}

% Unnumbered environments
\newtheorem*{remark*}{Remark}
\newtheorem*{example*}{Example}
\newtheorem*{theorem*}{Theorem}
\newtheorem*{recognition*}{Recognition Principle for Infinite Loop Spaces}
\newtheorem*{claim*}{Claim}
\newtheorem*{convention*}{Convention}
\newtheorem*{definition*}{Definition}
\newtheorem*{exercise*}{Exercise}
\newtheorem*{fact*}{Fact}
\newtheorem*{question*}{Question}
\newtheorem*{warning*}{Warning}
\renewcommand{\theaxiom}{\Alph{axiom}}
\newtheorem{subaxiom}{Axiom}
\renewcommand{\thesubaxiom}{\Alph{axiom}.\arabic{subaxiom}}
\renewcommand{\theHsubaxiom}{\theaxiom.\thesubaxiom}  % Makes sure the hyperlinks remember in which axiom the subaxiom appears
\newtheorem{subaxiomprime}{Axiom}
\renewcommand{\thesubaxiomprime}{\Alph{axiom}.\arabic{subaxiomprime}'}
\renewcommand{\theHsubaxiomprime}{\theaxiom.\thesubaxiomprime}  % Makes sure the hyperlinks remember in which axiom the subaxiom appears

%---------------------------------------
% Margin notes
%---------------------------------------
\newcommand{\red}[1]{{\color{red}{\textbf #1}}}
\setlength{\parindent}{0cm}
\setlength{\parskip}{0.8ex}
\setlength\marginparwidth{57pt}
\newcommand{\notehelper}[3]{\textcolor{#3}{$\blacksquare$}\marginpar{\ifodd\thepage\raggedright\else\raggedleft\fi\color{#3}\tiny \textbf{#2:} #1}}
\newcommand\bastiaan[1]{\notehelper{#1}{Bastiaan}{blue}}
\newcommand\bastiaanlong[1]{{\itshape \color{blue} #1}}

%---------------------------------------
% Tikzcd
%---------------------------------------

\tikzcdset{pullback/.style = {phantom, "\lrcorner"{very near start}}}
\tikzcdset{pullbackdl/.style = {phantom, "{\reflectbox{$\lrcorner$}}"{very near start}}}
\tikzcdset{pushout/.style = {phantom, "{\rotatebox{180}{$\lrcorner$}}"{very near end}}}

\tikzset{curve/.style={settings={#1},to path={(\tikztostart)
			.. controls ($(\tikztostart)!\pv{pos}!(\tikztotarget)!\pv{height}!270:(\tikztotarget)$)
			and ($(\tikztostart)!1-\pv{pos}!(\tikztotarget)!\pv{height}!270:(\tikztotarget)$)
			.. (\tikztotarget)\tikztonodes}}}
\tikzcdset{scale cd/.style={every label/.append style={scale=#1},cells={nodes={scale=#1}}}}
\newcommand{\pullback}{\mathrel{\text{\tikz{
				\draw[line width=0.6pt] (0,0) -- (0.5em,0);
				\draw[line width=0.6pt] (0.5em,0) -- (0.5em,0.6em);
		}}\,}}
\newcommand{\pullbackto}{\mathrel{\text{\tikz{
				\draw[line width=0.6pt] (0,0) -- (0.5em,0);
				\draw[line width=0.6pt] (0.5em,0) -- (0.5em,0.6em);
		}}\kern0.5em\to}}
\newcommand{\pushout}{\mathrel{\text{\tikz{
				\draw[line width=0.4pt] (0,0) -- (0,0.5em);
				\draw[line width=0.4pt] (0,0.5em) -- (0.5em,0.5em);
		}}\,}}
\newcommand{\pushoutto}{\mathrel{\text{\tikz{
				\draw[line width=0.4pt] (0,0) -- (0,0.5em);
				\draw[line width=0.4pt] (0,0.5em) -- (0.5em,0.5em);
		}}\kern0.5em\to}}

%---------------------------------------
% Other
%---------------------------------------

\newcommand{\quot}{/\!\!/}
\newcommand{\quotcolim}{\text{``}\!\colim\!\text{''}}	% formal colimit in Ind

% Use this when you want to write a \colon but the arrow goes from right to left, as in $Y \leftarrow X\noloc f$.
\newcommand\noloc{%
	\nobreak
	\mspace{6mu plus 1mu}
	{:}
	\nonscript\mkern-\thinmuskip
	\mathpunct{}
	\mspace{2mu}
}

\newcommand{\pvto}[2]{{%
		\kern1pt
		\big(
		\begin{matrix}
			\text{\tiny $#1$} \\[-7pt]
			\scriptstyle\downarrow \\[-7pt]
			\text{\tiny $#2$}
		\end{matrix}
		\big)
}}
\newcommand{\puto}[3]{{%
		\kern1pt
		\big(
		\begin{matrix}
			\text{\tiny $#1$ \hspace{-6pt}\phantom{$#3$}} \\[-9pt]
			\scriptstyle\downarrow \raisebox{2pt}{\text{\tiny $#3$}} \\[-9pt]
			\text{\tiny $#2$ \hspace{-6pt}\phantom{$#3$}}
		\end{matrix}
		\big)
}}
\newcommand{\apvto}[2]{\raisebox{0pt}[1.3\ht\strutbox][1.3\dp\strutbox]{$\smash{\pvto{#1}{#2}}$}}
\newcommand{\pwto}[2]{{%
		\kern1pt
		\big(
		\begin{matrix}
			\text{\tiny $#1$} \\[-5pt]
			\text{\tiny $#2$}
		\end{matrix}
		\big)
}}
\newcommand{\rbox}[3]{\raisebox{0pt}[#1\ht\strutbox][#2\dp\strutbox]{$#3$}}

%---------------------------------------
% Cleverref
%---------------------------------------
\usepackage[capitalise]{cleveref}  % Smart cross-referencing

\crefname{axiom}{axiom}{axioms}
\crefname{claim}{claim}{claims}
\crefname{conjecture}{conjecture}{conjectures}
\crefname{construction}{construction}{constructions}
\crefname{convention}{convention}{conventions}
\crefname{corollary}{corollary}{corollaries}
\crefname{definition}{definition}{definitions}
\crefname{discussion}{discussion}{discussions}
\crefname{example}{example}{examples}
\crefname{exercise}{exercise}{exercises}
\crefname{fact}{fact}{facts}
\crefname{lemma}{lemma}{lemmas}
\crefname{notation}{notation}{notations}
\crefname{observation}{observation}{observations}
\crefname{proposition}{proposition}{propositions}
\crefname{question}{question}{questions}
\crefname{remark}{remark}{remarks}
\crefname{subaxiom}{axiom}{axioms}
\crefname{terminology}{terminology}{terminologies}
\crefname{theorem}{theorem}{theorems}
\crefname{warning}{warning}{warnings}

\Crefname{axiom}{Axiom}{Axioms}
\Crefname{claim}{Claim}{Claims}
\Crefname{conjecture}{Conjecture}{Conjectures}
\Crefname{construction}{Construction}{Constructions}
\Crefname{convention}{Convention}{Conventions}
\Crefname{corollary}{Corollary}{Corollaries}
\Crefname{definition}{Definition}{Definitions}
\Crefname{discussion}{Discussion}{Discussions}
\Crefname{example}{Example}{Examples}
\Crefname{exercise}{Exercise}{Exercises}
\Crefname{fact}{Fact}{Facts}
\Crefname{lemma}{Lemma}{Lemmas}
\Crefname{notation}{Notation}{Notations}
\Crefname{observation}{Observation}{Observations}
\Crefname{proposition}{Proposition}{Propositions}
\Crefname{question}{Question}{Questions}
\Crefname{remark}{Remark}{Remarks}
\Crefname{subaxiom}{Axiom}{Axioms}
\Crefname{terminology}{Terminology}{Terminologies}
\Crefname{theorem}{Theorem}{Theorems}
\Crefname{warning}{Warning}{Warnings}

% Appendix chapters are labelled \label[appendix]{...} so that \Cref calls them
% ``Appendix'' rather than ``Chapter'' (main-matter chapters keep ``Chapter'').
\crefname{appendix}{appendix}{appendices}
\Crefname{appendix}{Appendix}{Appendices}

% Invisible source metadata shared by PDF and web builds.
% Place immediately before a proof; the optional argument names a proof variant.
\NewDocumentCommand{\prooffor}{m o}{}


\usepackage{float}
\usepackage{etoolbox}
\usepackage[T1]{fontenc}
\usepackage[utf8]{inputenc}

\begin{document}
	\pagenumbering{roman}

	\title{Efimov K-theory\\ \vspace{10pt}
	\Large University of Regensburg, Winter term 2026/27}
	\author{\large Bastiaan Cnossen \\ Notes from the seminar organized by Marc Hoyois}

	\maketitle

	%	\newgeometry{top=1in,bottom=1in,left=1.5in,right=1.5in,headsep=1.5cm,marginparwidth=1.25in}
	\newgeometry{left=1.3in,right=1.2in, top=3.8cm, headsep=1.3cm}

	\begin{KeepFromToc}
		\tableofcontents
	\end{KeepFromToc}

	\titlespacing*{\chapter}{0pt}{20pt}{15pt}

	\chapter*{Preface}

	These notes grew out of the research seminar on Efimov K-theory organized by Marc Hoyois at the University of Regensburg in the winter term 2026/27. They are organized by topic rather than by talk and include material from the workshop \emph{Dualizable categories and continuous K-theory} (MPIM Bonn, September 2024) and from lectures by Dustin Clausen (July 2023).

	I thank Marc Hoyois for organizing the seminar and all the speakers whose lectures contributed to these notes.

	Generative AI has been used to convert lecture notes into LaTeX, reorganize the material, and draft additional exposition and proofs based on the literature. This includes material not covered in the lectures, as well as completions of arguments that were only sketched or omitted there. The source notes below distinguish lecture-based material from additions from the literature. Both kinds of material have received AI assistance. These remain working notes, and some of the added arguments have not yet been checked in detail.

	\section*{Sources by chapter}

	Chapters~1--7 and Appendix~A contain the material written so far. Chapters~8--11 are planned and are visibly marked as such. The following paragraphs record the sources of the written material. References for individual results and proofs also appear throughout the text.

	\subsection*{\Cref{chap:Introduction_And_Overview}}

	This chapter draws on Marc Hoyois's first seminar talk (13 October 2026) and Dustin Clausen's first three lectures (17--19 July 2023).

	\subsection*{\Cref{chap:Localization_Sequences}}

	Marc Hoyois's first seminar talk (13 October 2026) contributes to \Cref{sec:Localization_Sequences_PrL,sec:Localizing_Invariants}. Akhil Mathew's first Bonn lecture (9 September 2024) is the source for \Cref{sec:Calkin_Small} and parts of \Cref{sec:Karoubi_Sequences,sec:Localization_Sequences_PrL,sec:Semiorthogonal_Decompositions,sec:Localizing_Invariants}.

	\Cref{sec:Verdier_Sequences} and most of \Cref{sec:Karoubi_Sequences,sec:Localization_Sequences_PrL,sec:Semiorthogonal_Decompositions} were added from the literature, using~\cite{NikolausScholze2018Cyclic,Hermitian_II,lurie2009HTT,Efimov_Large_Categories}. The characterization of semiorthogonal decompositions in terms of adjoints (\Cref{prop:SOD_Characterization}) comes from Jacopo Mazzei's seminar talk (17 November 2026), supplemented by \cite[Remark~2.75]{Efimov_Large_Categories}. The sketch of the proof of Tamme's excision theorem (\Cref{thm:Tamme_Excision}) via continuous extensions follows \cite[Proposition~8.11]{Efimov_Large_Categories}; Tamme's original argument is summarized in \Cref{rem:Tamme_Lax_Pullback_Proof}. The explanation of \Cref{thm:K_Unit_Day_Convolution} was added from the stated non-finitary universal property of K-theory.

	\subsection*{\Cref{chap:Compactly_Assembled_Dualizable}}

	This chapter is based on Marc Hoyois's second seminar talk (20 October 2026). Akhil Mathew's second Bonn lecture (9 September 2024) contributes to \Cref{sec:Dualizable_Categories,sec:Omega_1_Presentability}. Dustin Clausen's first two lectures (17--18 July 2023) contribute to \Cref{sec:Localization_Sequences_Dualizable,sec:Omega_1_Presentability}.

	The introductory exposition was expanded with AI assistance. The example on continuous posets (\Cref{ex:Continuous_Posets}) was added from the literature, following Johnstone--Joyal and Gierz et al.~\cite{Johnstone_Joyal_Continuous_Categories,Gierz_Continuous_Lattices}. The proof of \Cref{prop:Strongly_Continuous_Compact_Morphisms} was added using \cite{Efimov_Large_Categories}. The proof of \Cref{thm:Compactly_Assembled_Omega_1} was revised using countable directed subposets, with the indexing lemma in Appendix~A.

	\subsection*{\Cref{chap:Compact_Morphisms_AB6}}

	This chapter is based on Jiapeng Liu's seminar talk (3 November 2026). The converse implication in \Cref{prop:Compact_Morphisms_Stable} and the proofs of \Cref{prop:Dualizable_Compactly_Exhaustible,cor:Compact_Iff_Identity_Weakly_Compact,cor:Yhat_Universal_Property,lem:Products_In_Ind} were added using~\cite{Efimov_Large_Categories,Krause_Nikolaus_Putzstuck_Sheaves}. A proof of (AB6) for sets and animae (\Cref{ex:AB6_Sets_And_Animae}) was added when preparing these notes. The interpolation results (\Cref{lem:Interpolation,prop:Compact_Interpolation}) were also added from the literature, following~\cite{Krause_Nikolaus_Putzstuck_Sheaves}. \Cref{sec:Automatic_Dualizability} is based on Anton Zakrewski's seminar talk (10 November 2026) and his notes for it, following \cite[Section~2.8]{Efimov_Large_Categories}; the second proof of \Cref{prop:Generation_Dualizable} was added using compact exhaustions. The alternative compact-morphism definitions are collected in \Cref{rem:Compact_Morphisms_Other_Definitions} after the criteria establishing their agreement.

	\subsection*{\Cref{chap:Examples}}

	Max Glaser's seminar talk (27 October 2026), Dustin Clausen's second lecture (18 July 2023), and Akhil Mathew's third Bonn lecture (10 September 2024) contribute to \Cref{sec:Homological_Epimorphisms,sec:Almost_Ring_Theory,sec:Continuous_Posets}. The proofs and references were supplemented from \cite{Neeman_Chromatic_Tower,Bhatt_Scholze_Witt_Grassmannian,Bhatt_Scholze_Prisms,Krause_Nikolaus_Putzstuck_Sheaves}. The proofs of \Cref{prop:Every_Dualizable_Kernel,prop:Operator_Ideal_Models} were added when preparing these notes. The latter supplies the Hilbert-space models for the compact-operator and $S_{0+}$ kernels, including their symmetric monoidal structures; the identification of the latter with liquid rigidification is still quoted from Clausen.

	The material on semiorthogonal decompositions and gluing in \Cref{sec:Extensions_And_Gluing} comes from Jacopo Mazzei's seminar talk (17 November 2026), following \cite[Propositions~2.76--2.78]{Efimov_Large_Categories}; parts of the proof of \Cref{prop:SOD_Functor_Categories} were completed when preparing these notes. The Tate-extension material comes from Alexander Efimov's fourth Bonn lecture (10 September 2024); the proofs of \Cref{prop:Extension_Dualizable_Tate} and \Cref{ex:Extension_Not_Dualizable} were completed using \cite{Efimov_Large_Categories,Hennion_Tate,Hermitian_II}.

	\subsection*{\Cref{chap:Rigidity_And_Nuclear_Objects}}

	This chapter draws on Max Glaser's seminar talk (27 October 2026) and Dustin Clausen's second and third lectures (18--19 July 2023). Much of the treatment of rigidity and nuclearity was added from \cite{GaitsgoryRozenbluym2017StudyDAG,Ramzi_Locally_Rigid,Efimov_Inverse_Limits,Scholze_Analytic_Geometry,Andreychev_Thesis}. The operator-algebra presentation and the Hochschild-homology obstruction in \Cref{ex:Nuclear_Solid} were supplied from Efimov's paper. The formula in \Cref{ex:Nuclear_Liquid_Gaseous} remains explicitly attributed to the unpublished lecture notes.

	\subsection*{\Cref{chap:Calkin_Limits_Colimits}}

	\Cref{sec:Colimits_Dualizable,sec:Continuous_Calkin,sec:Filtered_Colimits_Dualizable} are based on Anton Zakrewski's seminar talk (10 November 2026) and his notes for it. \Cref{sec:Finite_Limits_Dualizable} is based on Jacopo Mazzei's seminar talk (17 November 2026). Christoph Winges's seminar talk (24 November 2026) is the source for \Cref{sec:General_Limits_Dualizable,sec:Infinite_Products_Dualizable}. The proofs in these two sections follow \cite[Sections~2.13 and~2.15]{Efimov_Large_Categories} and were completed when preparing these notes; the products of short exact sequences of small categories are treated as in \cite[Lemma~5.2]{Kasprowski_Winges_Products}. The first two talks follow \cite[Sections~2.9--2.12]{Efimov_Large_Categories}, which was also used to complete the proofs. \Cref{prop:Epimorphisms_Localizations}, \Cref{ex:Dualizable_Kernel_DVR}, \Cref{cor:Monomorphisms_Prdual}, \Cref{rem:Dualizable_Coreflection} and part~(2) of \Cref{prop:Dualizable_Pullbacks} were added from \cite{Efimov_Large_Categories}, and the proof of \Cref{prop:Calkin_Short_Exact} uses \cite[Proposition~A.26]{Ramzi_Dualizable_Presentable}. The categorical facts preceding \Cref{prop:Compact_Objects_Filtered_Colimits} use \cite[Section~1.1.4]{lurie2016HA}; the explanations of coherence, adjunction units and conservativity of the core were added when checking these proofs. \Cref{rem:Calkin_Small_Vs_Continuous} identifies the continuous construction with the small Calkin construction in the compactly generated case, following \cite[Proposition~2.59]{Efimov_Large_Categories}.

	\subsection*{\Cref{app:Categorical_Background}}

	Marc Hoyois's second seminar talk (20 October 2026) contributes to \Cref{sec:Idempotent_Completeness,sec:PrL_Tensor_Products}. \Cref{sec:Filtered_Categories,sec:Accessible_Presentable,sec:Mates} and most of \Cref{sec:Idempotent_Completeness,sec:PrL_Tensor_Products} were added from the literature, using~\cite{lurie2009HTT,lurie2016HA,KellyStreet1974Mates,Riehl_Verity}.

	\titlespacing*{\chapter}{0pt}{50pt}{40pt}

	\restoregeometry
	\openany

	\mainmatter
	\pagestyle{plain}
	\setcounter{secnumdepth}{2}
	% Theorem numbers for Efimov_Large_Categories refer to arXiv:2405.12169v4.
	\chapter{Introduction and overview}
	\label[chapter]{chap:Introduction_And_Overview}

	% Source: Clausen, Lecture 1 (17.07.2023); Talk 1 (Marc Hoyois, 13.10.2026)
	These notes are about a recent extension, due to Efimov \cite{Efimov_Large_Categories}, of the domain of definition of algebraic K-theory: from small stable categories to a class of \emph{large} stable categories, the \emph{dualizable} ones. The history of algebraic K-theory can be read as a sequence of such extensions,
	\[
		\begin{tikzcd}[column sep=large, row sep=0pt]
			\Vect(X) \rar[squiggly] & \Perf(X) \rar[squiggly] & \Nuc(X) \\
			\text{exact categories} & \text{small stable categories} & \text{dualizable categories,}
		\end{tikzcd}
	\]
	and the guiding principle behind each of them is the same:
	\begin{center}
		\emph{the input of K-theory should already have the properties that K-theory is supposed to have.}
	\end{center}
	In this introductory chapter we explain this principle, motivate why one would want to go beyond small stable categories, and give an overview of the main definitions and results that will be developed in detail in the remaining chapters. Nothing in this chapter will be proved here; we give forward references to the places in the notes where the proofs can be found.

	\begin{convention}
		\label[convention]{con:Categories_And_K_Theory}
		Throughout, \emph{category} means $\infty$-category, and $K$ denotes \emph{nonconnective} algebraic K-theory, regarded as a functor $K\colon \Cat_{\perf} \to \Sp$ on small idempotent complete stable categories. If $C$ is a compactly generated stable category, we often abuse notation and write $K(C) := K(C^{\omega})$.
	\end{convention}

	\section{Why change the input of K-theory?}
	\label[section]{sec:Why_Change_Input}

	% Source: Clausen, Lecture 1 (17.07.2023)
	To illustrate the guiding principle, let us look at the first expansion: from vector bundles to perfect complexes.

	\subsection*{Pushforwards and the Riemann--Roch theorem}
	A crucial ingredient in Grothendieck's relative version of the Riemann--Roch theorem \cite{Borel_Serre_Riemann_Roch} is that, for every proper map $f\colon X \to Y$ of smooth quasi-projective varieties, there is a pushforward map
	\[
		f_*\colon K(\Vect(X)) \to K(\Vect(Y)).
	\]
	At the level of vector bundles, however, it is not at all obvious where such a map should come from: the pushforward $f_*\mathcal{E}$ of a vector bundle is in general neither a vector bundle nor concentrated in a single degree, and one is forced to use the \emph{derived} pushforward together with resolutions. If instead we allow perfect complexes as input, the map simply comes from the functor
	\[
		f_*\colon \Perf(X) \to \Perf(Y),
	\]
	and in this generality also the annoying assumption of quasi-projectivity disappears. In other words: K-theory admits proper pushforwards because the \emph{categories} $\Perf(-)$ admit proper pushforwards.

	\subsection*{Further examples}
	The same pattern appears everywhere once one works with stable categories:
	\begin{itemize}
		\item \emph{Descent.} The functor $X \mapsto K(\Perf(X))$ satisfies Zariski (and even Nisnevich) descent by the work of Thomason--Trobaugh \cite{Thomason_Trobaugh}, which fails for $X \mapsto K(\Vect(X))$ in general. The categorical reason is that for a quasi-compact open $U \subseteq X$ with closed complement $Z$, the restriction functor $\Perf(X) \to \Perf(U)$ is a localization (up to idempotent completion) with kernel $\Perf_Z(X)$, the perfect complexes supported on $Z$. K-theory sends such sequences to fiber sequences, and descent follows formally.
		\item \emph{Group completion is automatic.} The suspension functor $\Sigma$ on a stable category $C$ induces multiplication by $-1$ on $K(C)$: from the cofiber sequence $x \to 0 \to \Sigma x$ we get $[x] + [\Sigma x] = 0$. So the inverses that one has to formally adjoin in the group completion are already present in the category.
		\item \emph{Projective bundle formulas.} For a commutative ring $R$, there is a semiorthogonal decomposition (in the sense of \Cref{sec:Semiorthogonal_Decompositions})
		\[
			\Perf(\P^1_R) \simeq \lra{\Perf(R) \otimes \Oo(-1),\; \Perf(R) \otimes \Oo},
		\]
		and K-theory turns this into a splitting $K(\P^1_R) \simeq K(R) \oplus K(R)$.
	\end{itemize}
	In each of these cases, a property of K-theory that is hard to establish by hand becomes a formal consequence of a structural property of the input category, combined with the single fundamental property of K-theory: it is a \emph{localizing invariant}, i.e.\ it sends short exact sequences of small stable categories to fiber sequences of spectra (\Cref{def:Localizing_Invariant}).

	\section{Where small stable categories fall short}
	\label[section]{sec:Small_Categories_Fall_Short}

	% Source: Clausen, Lecture 1 (17.07.2023)
	Why should one want to go further, beyond small stable categories? Let us give two reasons, which we will see resolved by Efimov's theory later in this chapter.

	\subsection*{K-theory of analytic spaces}
	Suppose we want to develop a theory of algebraic K-theory for \emph{analytic} spaces, such as formal schemes, rigid analytic varieties, or complex analytic spaces, with good descent properties and a Grothendieck--Riemann--Roch theorem. Then the naive approach of taking $K(\Perf(-))$ does not work. The basic reason is that in analytic geometry, restriction to an open subset does not give rise to a localization in the categorical sense: for instance, if $U \subseteq X$ is a rational open subset of an affinoid adic space, the restriction functor on perfect complexes is not a Verdier quotient. Hence the Thomason--Trobaugh argument sketched above breaks down at its very first step.

	\subsection*{The K-theory of the compact operators}
	Let $\Hh$ be a separable infinite-dimensional Hilbert space, let $B = B(\Hh)$ be the C$^*$-algebra of bounded operators, and let $\Kk = \Kk(\Hh) \subseteq B$ be the closed two-sided ideal of compact operators. Since $\Kk$ is a non-unital ring, its algebraic K-theory is \emph{defined} using excision:
	\[
		K(\Kk) := \fib\big(K(B) \to K(B/\Kk)\big).
	\]
	A theorem of Suslin--Wodzicki \cite{Suslin_Wodzicki_Excision} shows that this is independent of the choice of $B$: we have $K(\Kk) \simeq \fib(K(R) \to K(R/\Kk))$ for every unital ring $R$ containing $\Kk$ as a two-sided ideal. They also proved \emph{Karoubi's conjecture}, which in particular implies that
	\[
		K(\Kk) \simeq \KU
	\]
	is the periodic complex topological K-theory spectrum. This is a beautiful statement, but it raises a question: K-theory is supposed to be an invariant of categories, so \emph{where is the category} whose K-theory is $K(\Kk)$? We will answer this question in \Cref{ex:Compact_Operators_Intro} below.

	\section{From small to large categories}
	\label[section]{sec:Small_To_Large}

	% Source: Clausen, Lecture 1 (17.07.2023); Talk 1 (Marc Hoyois, 13.10.2026)
	The first step is to reformulate K-theory of small categories in terms of large categories. Every small stable category $C_0$ gives rise to a compactly generated presentable stable category $\Ind(C_0)$, and if $C_0$ is idempotent complete, we can recover it as $C_0 \simeq \Ind(C_0)^{\omega}$. (Restricting to idempotent complete categories is harmless: these have the best behaved (co)fiber sequences, and passing to the idempotent completion only affects $K_0$, and does not affect nonconnective K-theory at all.) This gives an equivalence of categories
	\[
		\Cat_{\perf} \iso \PrL_{\omega,\st} = \left\{
		\begin{array}{c}
			\text{compactly generated categories in $\PrLst$,} \\
			\text{strongly continuous functors}
		\end{array}
		\right\},
		\qquad C_0 \mapsto \Ind(C_0),
	\]
	where a functor $F\colon C \to D$ in $\PrLst$ is called \emph{strongly continuous} if its right adjoint again preserves colimits. Between compactly generated stable categories, this is equivalent to $F$ preserving compact objects. So K-theory may equally well be regarded as a functor on compactly generated categories and strongly continuous functors.

	From this perspective, the obvious next question is: \emph{which presentable stable categories, beyond the compactly generated ones, admit a reasonable K-theory?} The answer is given by the following theorem, which provides many equivalent characterizations of a single class of categories.

	\begin{theorem}[{Lurie, \cite[Proposition~D.7.3.1]{SAG}}]
		\label[theorem]{thm:Dualizable_Characterizations}
		Consider $C \in \PrLst$. The following are equivalent:
		\begin{enumerate}
			\item There exists a localization $\phi\colon D \to E$ in $\PrL_{\omega,\st}$ such that $C = \ker(\phi)$;
			\item $C$ is dualizable in $(\PrLst,\otimes)$;
			\item The functor $\colim\colon \Ind(C) \to C$ has a left adjoint (``compactly assembled'');
			\item $C$ is a retract in $\PrLst$ of a compactly generated category;
			\item (Efimov, \cite[Theorem~5.2]{Efimov_Large_Categories}) $C \otimes (-)\colon \PrLst \to \PrLst$ preserves fully faithful functors (``flat'');
			\item (Clausen--Scholze, see \cite[Proposition~2.51]{Efimov_Large_Categories}) arbitrary products distribute over filtered colimits (``AB6'').
		\end{enumerate}
	\end{theorem}

	The equivalence of (1)--(4) will be proved in \Cref{thm:Lurie_Dualizable_Characterization}, the axiom (AB6) is discussed in \Cref{chap:Compact_Morphisms_AB6}, and flatness in \Cref{chap:Flatness}. Characterization (1) is the one most relevant for K-theory: a dualizable category is precisely something that can be described as the ``difference'' of two compactly generated categories. Characterization (3) is the one that is easiest to work with in practice. In a compactly generated category, the presentation of an object $x$ as the formal colimit of all compact objects mapping to it is the \emph{finest} presentation of $x$: it maps uniquely into every ind-object $X \in \Ind(C)$ with $\colim X \simeq x$, and as a consequence maps out of $x$ into filtered colimits can be computed at the level of ind-objects. Characterization (3) says precisely that every object $x$ of $C$ still admits such a finest presentation $\hat{y}(x) \in \Ind(C)$, whose terms are however no longer required to be compact. This heuristic is explained in more detail at the beginning of \Cref{chap:Compactly_Assembled_Dualizable}.

	In particular, compactly generated categories are dualizable, with dual
	\[
		\Ind(C_0)^{\vee} = \Ind(C_0\catop).
	\]
	But there are many more examples.

	\begin{example}
		\label[example]{ex:Dualizable_Categories}
		The following categories are dualizable but in general not compactly generated:
		\begin{itemize}
			\item Keller's example of a localization of a compactly generated category with a non-compactly generated kernel (\Cref{ex:Keller});
			\item the category $\Shv(X;\Sp)$ of sheaves of spectra on a locally compact Hausdorff space $X$ (see \Cref{chap:Sheaves}); for instance $\Shv(\R^n;\Sp)$ with $n \geq 1$ has no nonzero compact objects at all (for sheaves of abelian groups on a non-compact connected manifold this is due to Neeman \cite[Theorem~0.1]{Neeman_Sheaves_Manifold}; for sheaves of spectra it follows from Efimov's description of the compact objects \cite[Proposition~10.17]{Efimov_Large_Categories});
			\item the categories of nuclear modules arising in analytic geometry: $\Nuc(\Z_p^{\square})$, $\Nuc(\R^{p\text{-}\mathrm{liq}})$, $\Nuc(\C^{\mathrm{gas}})$ \cite{Scholze_Analytic_Geometry, Clausen_Scholze_Complex_Geometry}.
		\end{itemize}
		All of the latter examples are actually self-dual. The reason is: they are even \emph{rigid}.
	\end{example}

	\begin{definition}[{cf.\ \cite[Chapter~1, Definition~9.1.2]{GaitsgoryRozenbluym2017StudyDAG}}]
		\label[definition]{def:Rigid_Category}
		A \emph{rigid category} is $C \in \CAlg(\PrLst)$ such that for every finite set $I$, the multiplication functor $\mu\colon C^{\otimes I} \to C$ admits a right adjoint in $\Mod_{C^{\otimes I}}(\PrLst)$.
	\end{definition}

	\begin{example}
		\label[example]{ex:Rigid_Ind_Categories_Intro}
		Let $C_0$ be a small idempotent complete symmetric monoidal stable category in which every object is dualizable. Then $\Ind(C_0)$ is rigid. For instance, $\D(R) = \Ind(\Perf(R))$ is rigid for every commutative ring $R$.
	\end{example}

	\begin{proposition}[{\cite[Chapter~1, Section~9.2.1]{GaitsgoryRozenbluym2017StudyDAG}}]
		\label[proposition]{prop:Rigid_Self_Dual}
		If $C$ is rigid, then $C$ is self-dual via
		\[
			C \otimes C \xrightarrow{\mu} C \xrightarrow{\Gamma} \Sp,
		\]
		where $\Gamma = \hom_C(\unit,-)$.
	\end{proposition}

	Rigid categories, trace-class morphisms and nuclear modules are discussed in detail in \Cref{sec:Rigid_Categories,sec:Nuclear_Modules}.

	\begin{remark}
		\label[remark]{rem:Compact_Objects_And_Morphisms_Intro}
		There is a useful dictionary between the compactly generated and the dualizable world, which will be developed in \Cref{chap:Compact_Morphisms_AB6}: compact \emph{objects} are replaced by compact \emph{morphisms}. A dualizable category is generated under colimits by sequential colimits $\colim(x_0 \to x_1 \to x_2 \to \cdots)$ along compact morphisms, and in fact its $\omega_1$-compact objects are precisely of this form. In a symmetric monoidal setting, the analogy continues: just as a compactly generated category with compact unit is rigid if its compact objects are precisely its dualizable objects, a dualizable symmetric monoidal category is rigid if its compact morphisms are precisely its trace-class morphisms.
	\end{remark}

	The morphisms between dualizable categories which we will consider are the strongly continuous functors.

	\begin{definition}
		\label[definition]{def:Cat_Dual}
		We denote by
		\[
			\Pr^{\dual} \subseteq \PrLst
		\]
		the subcategory consisting of the dualizable categories and the strongly continuous functors, i.e.\ those functors whose right adjoint again preserves colimits (see \Cref{def:Dualizable_Category}).
	\end{definition}

	By the above discussion, Ind-completion defines a fully faithful functor
	\[
		\Ind\colon \Cat_{\perf} \hookrightarrow \Pr^{\dual}.
	\]
	The category $\Pr^{\dual}$ has excellent formal properties: it is closed under tensor products in $\PrLst$, the inclusion into $\PrLst$ creates colimits, and $\Pr^{\dual}$ admits all limits and internal homs, although these are more subtle (see \Cref{chap:Calkin_Limits_Colimits}).

	\section{Continuous K-theory}
	\label[section]{sec:Continuous_K_Theory_Intro}

	% Source: Talk 1 (Marc Hoyois, 13.10.2026); Clausen, Lecture 1 (17.07.2023)
	We would now like to extend the K-theory functor from $\Cat_{\perf}$ to $\Pr^{\dual}$. Characterization (1) of \Cref{thm:Dualizable_Characterizations} suggests the following: if $C$ is dualizable, write it as the kernel of a localization $D \to E$ of compactly generated categories along strongly continuous functors, and define
	\[
		K^{\cont}(C) := \fib\big(K(D^{\omega}) \to K(E^{\omega})\big).
	\]
	If $C$ is itself compactly generated, then $C^{\omega} \to D^{\omega} \to E^{\omega}$ is a short exact sequence of small stable categories by the Thomason--Neeman localization theorem (\Cref{thm:Thomason_Neeman}), so this gives back $K(C^{\omega})$. But it is not clear that this is well-defined in general: we would like to have a \emph{canonical} way of writing $C$ as the fiber of a map between compactly generated categories.

	Such a canonical presentation is provided by characterization (3). Let $C \in \Pr^{\dual}$ and let $\kappa$ be a regular cardinal which is large enough for the left adjoint $\hat{y}$ of $\colim\colon \Ind(C^{\kappa}) \to C$ to exist; in fact, $\kappa = \omega_1$ always suffices (\Cref{thm:Compactly_Assembled_Omega_1}). We get a localization sequence in $\PrLst$ with strongly continuous functors
	\[
		C \xhookrightarrow{\hat{y}} \Ind(C^{\kappa}) \to \Ind(\Calk^{\cont}_{\kappa}(C)),
	\]
	where the \emph{(continuous) Calkin category} $\Calk^{\cont}_{\kappa}(C)$ is defined as the category of compact objects of the quotient $\Ind(C^{\kappa})/\hat{y}(C)$.

	\begin{definition}
		\label[definition]{def:Continuous_Extension}
		Let $F$ be a localizing invariant. Then we define
		\[
			F^{\cont}(C) := \fib\big(F(C^{\kappa}) \to F(\Calk^{\cont}_{\kappa}(C))\big)
		\]
		for $\kappa = \omega_1$. When $F$ is K-theory, $K^{\cont}(C)$ is called the \emph{continuous K-theory} of $C$.
	\end{definition}

	For many localizing invariants, such as K-theory, the category $C^{\omega_1}$ is invisible: it admits countable coproducts, so by an Eilenberg swindle we have $K(C^{\omega_1}) = 0$, and hence
	\[
		K^{\cont}(C) \simeq \Omega K(\Calk^{\cont}_{\omega_1}(C)).
	\]

	At first sight this looks rather crazy: in order to compute $K^{\cont}_0(C)$, we have to compute $K_1$ of a huge and inexplicit category. Fortunately, this is not how one should think about continuous K-theory in practice. The following theorem gives a more conceptual characterization, which also makes much more flexible computational techniques available.

	\begin{theorem}[{Efimov, \cite[Theorem~8.10]{Efimov_Large_Categories}}]
		\label[theorem]{thm:Efimov_Continuous_Extension}
		Restriction along $\Ind$ induces an equivalence of categories
		\[
			\Fun_{\loc}(\Pr^{\dual},\Sp) \iso \Fun_{\loc}(\Cat_{\perf},\Sp),
		\]
		with inverse given by $F \mapsto F^{\cont}$. Moreover, $F$ commutes with filtered colimits if and only if $F^{\cont}$ does.
	\end{theorem}

	Here, on the left-hand side, a localizing invariant is a functor $\Pr^{\dual} \to \Sp$ that sends short exact sequences of dualizable categories, i.e.\ localization sequences in $\PrLst$ whose functors are strongly continuous, to fiber sequences. The key input in the proof is that the construction $C \mapsto \Calk^{\cont}_{\omega_1}(C)$ sends short exact sequences of dualizable categories to short exact sequences of small categories. In particular, $K^{\cont}$ is the \emph{unique} localizing invariant of dualizable categories extending K-theory, which makes the definition independent of all choices involved. This theorem will be proved in \Cref{chap:Localizing_Invariants_Dualizable}.

	\begin{example}
		\label[example]{ex:THH_Categorical_Trace}
		The continuous extension of $\THH$ admits a very clean description: it is the \emph{categorical trace}
		\[
			\THH\colon \Pr^{\dual} \to \Sp, \qquad C \mapsto \tr(\id_C) \in \End_{\PrLst}(\Sp) = \Sp,
		\]
		also known as the dimension of $C$ as a dualizable object of $\PrLst$. That this is a localizing invariant is a formal property of traces.
	\end{example}

	\section{First computations}
	\label[section]{sec:First_Computations}

	% Source: Clausen, Lectures 1--3 (17.--19.07.2023); Talk 1 (Marc Hoyois, 13.10.2026)
	Let us now see Efimov's theory in action, and in particular come back to the two problems raised in \Cref{sec:Small_Categories_Fall_Short}.

	\subsection*{Sheaves on locally compact Hausdorff spaces}
	The first computation shows that continuous K-theory is a genuinely new invariant: the category $\Shv(\R^n;\Sp)$ has no nonzero compact objects, so the only K-theory one could have assigned to it naively is zero.

	\begin{theorem}[{Efimov, \cite[Theorem~10.10]{Efimov_Large_Categories}}]
		\label[theorem]{thm:Kcont_Sheaves}
		Let $C$ be dualizable and $X$ a locally compact Hausdorff space. Then
		\[
			K^{\cont}(\Shv(X;C)) \simeq \Gamma_c(X;K^{\cont}(C)).
		\]
	\end{theorem}

	The idea of the proof is as follows. For $X$ compact, the functor $X \mapsto K^{\cont}(\Shv(X;C))$ commutes with cofiltered limits of compact Hausdorff spaces, sends $\emptyset$ to $0$, and sends a cover of $X$ by two closed subsets $Y,Z \subseteq X$ to a pullback square
	\[
		K^{\cont}(\Shv(X;C)) \iso K^{\cont}(\Shv(Y;C)) \times_{K^{\cont}(\Shv(Y \cap Z;C))} K^{\cont}(\Shv(Z;C)),
	\]
	both of which follow from the fact that $K^{\cont}$ is a localizing invariant commuting with filtered colimits. But every such functor is of the form $X \mapsto \Gamma(X;A)$ for the spectrum $A$ given by its value on the point. The case of general $X$ follows by passing to the one-point compactification. See \Cref{chap:Sheaves} for details.

	\begin{example}
		\label[example]{ex:Kcont_Euclidean_Sheaves}
		Taking $X = \R^n$, we get $K^{\cont}(\Shv(\R^n;C)) \simeq \Omega^n K^{\cont}(C)$. In particular, the higher K-groups of $C$ can be recovered as Grothendieck groups:
		\[
			K^{\cont}_n(C) \cong K^{\cont}_0(\Shv(\R^n;C)).
		\]
	\end{example}

	\subsection*{Idempotent algebras and the compact operators}
	A large supply of dualizable categories comes from \emph{idempotent} maps of ring spectra $A \to B$, i.e.\ maps for which $B \otimes_A B \to B$ is an equivalence. Equivalently, the base change functor $- \otimes_A B\colon \D(A) \to \D(B)$ is a localization (a \emph{homological epimorphism}, see \Cref{def:Homological_Epimorphism}). Its kernel
	\[
		\D(A \to B) := \ker\big(\D(A) \to \D(B)\big)
	\]
	is then dualizable by \Cref{thm:Dualizable_Characterizations}(1), and
	\[
		K^{\cont}(\D(A \to B)) \simeq \fib\big(K(A) \to K(B)\big).
	\]
	In fact, \emph{every} dualizable category is of this form (\Cref{prop:Every_Dualizable_Kernel}): the Calkin sequence above exhibits $C$ as the kernel of $\Ind(C^{\kappa}) \to \Ind(\Calk^{\cont}_{\kappa}(C))$, and for $\kappa$ large enough one can replace $C^{\kappa}$ by a subcategory with a single generator.

	\begin{example}
		\label[example]{ex:Compact_Operators_Intro}
		By a theorem of Wodzicki, if $A$ is a unital C$^*$-algebra and $I \subseteq A$ a closed two-sided ideal, then $A \to A/I$ is idempotent. Applying this to $A = B(\Hh)$ and $I = \Kk$, we get a dualizable category $\D(\Kk) := \D(B(\Hh) \to B(\Hh)/\Kk)$, and
		\[
			K^{\cont}(\D(\Kk)) \simeq \fib\big(K(B(\Hh)) \to K(B(\Hh)/\Kk)\big) = K(\Kk) \simeq \KU.
		\]
		So this is the category we were looking for in \Cref{sec:Small_Categories_Fall_Short}! Since $K(B(\Hh)) = 0$ by an Eilenberg swindle, $K^{\cont}(\D(\Kk))$ is the loop space of the K-theory of the \emph{Calkin algebra} $B(\Hh)/\Kk$; this is the origin of the name ``Calkin category''.

		The category $\D(\Kk)$ also admits an intrinsic description, without choosing $\Hh$: it is the full subcategory of $\Ind(\D^b(\text{separable Hilbert spaces}))$ generated under colimits by sequential colimits $\colim(\Hh_0 \to \Hh_1 \to \Hh_2 \to \cdots)$ along compact operators. Since the tensor product of two compact operators is again compact, the tensor product of Hilbert spaces makes $\D(\Kk)$ into a commutative algebra in $\Pr^{\dual}$. This gives an $\mathbb{E}_{\infty}$-ring structure on $K^{\cont}(\D(\Kk)) \simeq \KU$ which comes from evident categorical structure, but which is quite hard to see from the non-unital ring perspective.
	\end{example}

	\subsection*{Nuclear modules and K-theory of analytic spaces}
	There is a non-archimedean analogue of the previous example. Let $R$ be a commutative ring, $I = (f_1,\dots,f_m) \subseteq R$ a finitely generated ideal, and let $P := (\bigoplus_{\N} R)^{\wedge}_I$ be the $I$-completed countable direct sum, which plays the role of a separable Hilbert space. In place of $B(\Hh)$ we take the endomorphism ring $\End_R(P)$, and in place of the Calkin algebra $B(\Hh)/\Kk$ the $I$-adically completed endomorphism ring
	\[
		\Calk^{\mathrm{top}}_R(P) := \End_{\Calk(R)}(P)^{\wedge}_I
	\]
	of $P$ in the Calkin category of $\Perf(R)$. The fiber of $\End_R(P) \to \Calk^{\mathrm{top}}_R(P)$ is $\big(\prod_{\N}R \otimes_R \bigoplus_{\N}R\big)^{\wedge}_I$, which one may think of as the $I$-adic completion of the endomorphisms of finite rank, just as $\Kk$ is the norm closure of the operators of finite rank. Efimov shows that $\End_R(P) \to \Calk^{\mathrm{top}}_R(P)$ is a homological epimorphism, whose kernel is the category of \emph{nuclear} modules over $R^{\wedge}_I$ in the sense of Clausen--Scholze \cite[Corollary~7.7]{Efimov_Inverse_Limits}:
	\[
		\Nuc^{CS}(R^{\wedge}_I) \simeq \D\big(\End_R(P) \to \Calk^{\mathrm{top}}_R(P)\big).
	\]
	Alternatively, $\Nuc^{CS}(R^{\wedge}_I)$ can be obtained from the category of solid $R^{\wedge}_I$-modules by a ``rigidification'' procedure, see \Cref{sec:Nuclear_Modules}. Efimov also introduced a variant $\Nuc(R^{\wedge}_I)$, which contains $\Nuc^{CS}(R^{\wedge}_I)$ as a full subcategory with the same compact objects \cite[Corollary~7.6]{Efimov_Inverse_Limits}. In this introduction we will not distinguish between the two, and for $R$ itself $I$-adically complete we simply write $\Nuc(R)$. Efimov computed the continuous K-theory:

	\begin{theorem}[{Efimov, \cite[Corollary~6.3]{Efimov_Inverse_Limits}}]
		\label[theorem]{thm:Kcont_Nuclear}
		Let $R$ be a commutative ring, and $I = (f_1, \dots, f_m) \subseteq R$. Then
		\[
			K^{\cont}(\Nuc(R^{\wedge}_I)) = \lim_{n \to \infty} K(R \quot (f_1^n, \dots,f^n_m)).
		\]
	\end{theorem}
	The same formula holds for $\Nuc^{CS}(R^{\wedge}_I)$ \cite[Corollary~7.10]{Efimov_Inverse_Limits}, and if $R$ is Noetherian, the right-hand side is $\lim_n K(R/I^n)$.

	These categories are exactly what is needed to fix the problem with analytic spaces from \Cref{sec:Small_Categories_Fall_Short}. If $\operatorname{Spa}(A,A^+) \subseteq \operatorname{Spa}(R,R)$ is a rational open subset, then $A$ is an \emph{idempotent} commutative algebra in $\Nuc(R)$, and $\Nuc(A) \simeq \Mod_A(\Nuc(R))$. Tensor products of such algebras correspond to intersections of rational opens, and covers by rational opens give pullback squares of algebras. So from the point of view of nuclear modules, rational open subsets behave exactly like Zariski open subsets of schemes, and the Thomason--Trobaugh argument goes through verbatim. This leads to the following result of Andreychev:

	\begin{theorem}[{Andreychev, \cite[Satz~4.2 and Satz~5.12]{Andreychev_Thesis}}]
		\label[theorem]{thm:Nuclear_Analytic_Descent}
		Let $(A,A^+)$ be a complete sheafy analytic Huber pair. The assignments $U \mapsto \Nuc(\mathcal{O}(U))$ and $U \mapsto K^{\cont}(\Nuc(\mathcal{O}(U)))$, for $U \subseteq \operatorname{Spa}(A,A^+)$ rational open, satisfy descent for covers by rational open subsets.
	\end{theorem}

	As an application, one recovers a theorem of Kerz--Strunk--Tamme \cite{Kerz_Strunk_Tamme_Blowups} in a conceptual way: if $f\colon X \to Y$ is a proper map of Noetherian schemes which is an isomorphism outside a closed subscheme $Z \subseteq Y$, then the square
	\[
		\begin{tikzcd}
			K(Y) \rar \dar["f^*"'] & \lim_n K(Z_n) \dar["f^*"] \\
			K(X) \rar & \lim_n K((f^{-1}Z)_n)
		\end{tikzcd}
	\]
	is a pullback, where $Z_n$ and $(f^{-1}Z)_n$ denote the $n$-th infinitesimal thickenings. Indeed, using the localization sequences for K-theory with supports, one may replace $K(Y)$ and $K(X)$ by $K_Z(Y)$ and $K_{f^{-1}Z}(X)$. The theory of nuclear modules then identifies the cofibers of the horizontal maps with the K-theory of the generic fibers $(Y^{\wedge}_Z)_{\eta}$ and $(X^{\wedge}_{f^{-1}Z})_{\eta}$ of the formal completions, and these generic fibers agree because $f$ is an isomorphism away from $Z$.

	\subsection*{Outlook: archimedean analytic geometry}
	% Source: Clausen, Lecture 3 (19.07.2023)
	The same ideas apply in archimedean functional analysis. There is a canonical rigid category $\D^{\mathrm{an}}(\C)$ (obtained for instance by rigidifying the category of liquid $\C$-modules, independently of the choice of $p$), and Clausen--Scholze have announced results relating the continuous K-theory of complex analytic spaces built from it to topological K-theory and to (Hodge-filtered) de Rham cohomology. We will not discuss this further in these notes.

	\section{Overview of the notes}
	\label[section]{sec:Overview}

	The remaining chapters develop the theory sketched above in detail.
	\begin{itemize}
		\item \Cref{chap:Localization_Sequences} recalls the necessary background on stable categories: Verdier quotients, short exact sequences, localization sequences in $\PrLst$, the Thomason--Neeman theorem, semiorthogonal decompositions, and localizing invariants.
		\item \Cref{chap:Compactly_Assembled_Dualizable} introduces compactly assembled and dualizable categories and proves Lurie's characterization (\Cref{thm:Lurie_Dualizable_Characterization}), as well as the fact that dualizable categories are $\omega_1$-presentable.
		\item \Cref{chap:Examples} discusses examples of dualizable categories, including Keller's example, categories of the form $\D(A \to B)$, and nuclear modules.
		\item \Cref{chap:Compact_Morphisms_AB6} develops the internal point of view through compact morphisms, and relates dualizability to the axiom (AB6).
		\item \Cref{chap:Calkin_Limits_Colimits} studies Calkin categories and constructs colimits and limits in $\Pr^{\dual}$.
		\item \Cref{chap:Flatness} proves Efimov's characterization of dualizable categories as the flat objects of $\PrLst$.
		\item \Cref{chap:Localizing_Invariants_Dualizable} constructs the continuous extension of localizing invariants and proves \Cref{thm:Efimov_Continuous_Extension}.
		\item \Cref{chap:Sheaves} computes the continuous K-theory of categories of sheaves on continuous posets and locally compact Hausdorff spaces, proving \Cref{thm:Kcont_Sheaves}.
		\item \Cref{chap:Aleph_1_Presentability} shows that $\Pr^{\dual}$ is itself $\aleph_1$-presentable.
	\end{itemize}
	Background on filtered categories, presentability, idempotent completeness, mates, and tensor products of presentable categories is collected in \Cref{app:Categorical_Background}.


	% Theorem numbers for Efimov_Large_Categories refer to arXiv:2405.12169v4.
	\chapter{Stable categories and localization sequences}
	\label[chapter]{chap:Localization_Sequences}

	In this chapter we collect the basic theory of localization sequences, both for small stable categories and for presentable stable categories, and introduce localizing invariants. Throughout, we denote by $\Cat_{\st}$ the category of small stable categories and exact functors, and by $\Cat_{\perf} \subseteq \Cat_{\st}$ the full subcategory of idempotent complete ones (this category is denoted $\Cat^{\perf}$ in \cite{BGT2013Universal, Efimov_Large_Categories}). For a stable category $B$ we write $\hom_B(x,y) \in \Sp$ for the mapping spectra. A full subcategory of a stable category is called \emph{stable} if it contains $0$ and is closed under fibers and cofibers, and \emph{thick} if it is moreover closed under retracts.

	\section{Verdier sequences}
	\label[section]{sec:Verdier_Sequences}

	% Source: written from the literature (Nikolaus--Scholze, Hermitian K-theory II)

	The basic construction of this chapter is the quotient of a small stable category $B$ by a stable subcategory $A$, obtained by forcing the objects of $A$ to become zero. In the setting of stable categories this is simply a localization in the sense of Dwyer--Kan, and the following theorem shows that it behaves just like Verdier's classical quotient of triangulated categories: it is again stable, it has the expected universal property, and its mapping spectra are computed by a colimit over $A$. We then introduce the corresponding notion of exact sequence, the Verdier sequences.

	\begin{theorem}[{\cite[Theorem~I.3.3]{NikolausScholze2018Cyclic}}]
		\label[theorem]{thm:Verdier_Quotient}
		Let $B$ be a small stable category and $A \subseteq B$ a full stable subcategory. Let $W_A$ be the class of morphisms of $B$ whose cofiber lies in $A$, and let $q\colon B \to B/A := B[W_A^{-1}]$ be the localization. Then:
		\begin{enumerate}
			\item The category $B/A$ is stable and $q$ is exact.
			\item For every stable category $E$, restriction along $q$ induces an equivalence between $\Fun^{\ex}(B/A,E)$ and the full subcategory of $\Fun^{\ex}(B,E)$ spanned by the functors that vanish on $A$.
			\item For $x,y \in B$, the category $A_{/y}$ is filtered and
			\[
				\Map_{B/A}(qx,qy) \simeq \colim_{(a \to y) \in A_{/y}} \Map_B\big(x,\cofib(a \to y)\big).
			\]
		\end{enumerate}
		The category $B/A$ is called the \emph{Verdier quotient} of $B$ by $A$.
	\end{theorem}

	Note that $q$ is essentially surjective, being a localization. Moreover, part (2) says that $B \to B/A$ is the cofiber of $A \hookrightarrow B$ in $\Cat_{\st}$. Indeed, for a stable category $E$, exact functors out of the pushout $0 \sqcup_A B$ in $\Cat_{\st}$ are given by
	\[
		\Fun^{\ex}(B,E) \times_{\Fun^{\ex}(A,E)} \{0\},
	\]
	the category of exact functors $F\colon B \to E$ together with an equivalence $F|_A \simeq 0$. Since $0$ is a zero object of $\Fun^{\ex}(A,E)$, the space of such equivalences is either empty or contractible, so this is the full subcategory of $\Fun^{\ex}(B,E)$ spanned by the functors vanishing on $A$.

	\begin{lemma}
		\label[lemma]{lem:Kernel_Verdier_Quotient}
		In the situation of \Cref{thm:Verdier_Quotient}, the kernel of $q\colon B \to B/A$ is the thick closure of $A$ in $B$, i.e.\ the full subcategory of retracts of objects of $A$. In particular, $B/A \simeq B/\bar{A}$, where $\bar{A}$ is the thick closure of $A$.
	\end{lemma}
	\prooffor{lem:Kernel_Verdier_Quotient}
	\begin{proof}
		Since $q$ is exact, its kernel is a thick subcategory of $B$, and it contains $A$ since $0 \to a$ lies in $W_A$ for $a \in A$. Conversely, let $x \in B$ with $q(x) \simeq 0$. Then $\id_{qx}$ is nullhomotopic, so by \Cref{thm:Verdier_Quotient}(3) (applied with $y = x$) there exists $a \to x$ with $a \in A$ such that the canonical map $x \to c := \cofib(a \to x)$ is nullhomotopic. Since $a \simeq \fib(x \to c)$ and the fiber of a nullhomotopic map $x \to c$ is $x \oplus \Omega c$, it follows that $x$ is a retract of $a$. This shows that the kernel of $q$ consists of retracts of objects of $A$, and every retract of an object of $A$ lies in the thick subcategory $\ker(q)$. The last statement follows from the universal property \Cref{thm:Verdier_Quotient}(2), since an exact functor vanishes on $A$ if and only if it vanishes on $\bar{A}$.
	\end{proof}

	\begin{remark}
		\label[remark]{rem:Nullhomotopies_Property}
		Let $F\colon B \to C$ be a functor into a pointed category $C$. Since the constant functor $0$ with value a zero object is a terminal object of $\Fun(B,C)$, the space of natural transformations $F \to 0$ is contractible. Hence the space of equivalences $F \simeq 0$ is either empty or contractible, and it is non-empty if and only if $F$ sends every object to a zero object. So for a composable pair $A \xrightarrow{i} B \xrightarrow{p} C$ of functors between pointed categories, a nullhomotopy $pi \simeq 0$ is unique if it exists, and its existence is a \emph{property} of the pair $(i,p)$. We will therefore write ``$pi \simeq 0$'' as a condition throughout, without specifying a nullhomotopy.
	\end{remark}

	\begin{definition}
		\label[definition]{def:Verdier_Sequence}
		A sequence $A \xrightarrow{i} B \xrightarrow{p} C$ in $\Cat_{\st}$ with $pi \simeq 0$ is called a \emph{Verdier sequence} if it is both a fiber sequence and a cofiber sequence in $\Cat_{\st}$. Explicitly, this means that $i$ is fully faithful with essential image $\ker(p)$, and that the induced functor $B/A \to C$ is an equivalence. In this case $p$ is called a \emph{Verdier projection}.
	\end{definition}

	By \Cref{lem:Kernel_Verdier_Quotient}, for every full stable subcategory $A \subseteq B$ the sequence $\bar{A} \hookrightarrow B \to B/A$ is a Verdier sequence. See \cite[Lemma~A.1.8 and Proposition~A.1.9]{Hermitian_II} for a detailed discussion of Verdier sequences.

	\section{Idempotent completion and short exact sequences}
	\label[section]{sec:Karoubi_Sequences}

	% Source: written from the literature (Lurie HA, Efimov, Hermitian K-theory II); Mathew, Lecture 1 (MPIM Bonn, 09.09.2024)

	Every small stable category $A$ admits an idempotent completion $A \hookrightarrow A^{\idem}$, which is again stable \cite[Corollary~1.1.3.7]{lurie2016HA}; it is fully faithful and every object of $A^{\idem}$ is a retract of an object of $A$, and $A \mapsto A^{\idem}$ is left adjoint to the inclusion $\Cat_{\perf} \subseteq \Cat_{\st}$ (see \Cref{sec:Idempotent_Completeness}). In particular, cofibers in $\Cat_{\perf}$ are computed as $(B/A)^{\idem}$.

	\begin{definition}
		\label[definition]{def:Short_Exact_Sequence}
		A sequence $A \xrightarrow{i} B \xrightarrow{p} C$ in $\Cat_{\perf}$ with $pi \simeq 0$ is called a \emph{short exact sequence} (or \emph{Karoubi sequence}) if it is both a fiber sequence and a cofiber sequence in $\Cat_{\perf}$. Explicitly, this means that $i$ is fully faithful with essential image $\ker(p)$, and that the induced functor $(B/A)^{\idem} \to C$ is an equivalence. In this case $p$ is called a \emph{Karoubi projection}.
	\end{definition}

	We follow \cite{Efimov_Large_Categories} in calling these short exact sequences; the term Karoubi sequence is used in \cite{Hermitian_II}.

	\begin{lemma}
		\label[lemma]{lem:Characterization_Short_Exact}
		A sequence $A \xrightarrow{i} B \xrightarrow{p} C$ in $\Cat_{\perf}$ with $pi \simeq 0$ is short exact if and only if $i$ is fully faithful with essential image $\ker(p)$, the induced functor $B/A \to C$ is fully faithful, and every object of $C$ is a retract of an object in the image of $p$.
	\end{lemma}
	\prooffor{lem:Characterization_Short_Exact}
	\begin{proof}
		The functor $\bar{p}\colon B/A \to C$ factors as $B/A \hookrightarrow (B/A)^{\idem} \to C$. If $\bar{p}$ is fully faithful, then so is $(B/A)^{\idem} \to C$, since mapping spaces between retracts of objects are retracts of mapping spaces between those objects. If moreover every object of $C$ is a retract of an object $p(b) = \bar{p}(qb)$, then $(B/A)^{\idem} \to C$ is essentially surjective since $(B/A)^{\idem}$ is idempotent complete. Conversely, if $(B/A)^{\idem} \to C$ is an equivalence, then $\bar{p}$ is fully faithful and every object of $C$ is a retract of an object of $B/A$, and hence of an object $p(b)$, as $q$ is essentially surjective.
	\end{proof}

	\begin{remark}
		\label[remark]{rem:Verdier_Vs_Karoubi}
		Every Verdier sequence between idempotent complete categories is short exact. The converse fails whenever $A \subseteq B$ is a thick subcategory such that the Verdier quotient $B/A$ is not idempotent complete: then $A \to B \to (B/A)^{\idem}$ is short exact but not a Verdier sequence.
	\end{remark}

	% Source: Mathew, Lecture 1 (MPIM Bonn, 09.09.2024)
	\begin{remark}
		\label[remark]{rem:Mapping_Spectra_Quotient}
		Let $A \subseteq B$ be a full stable subcategory with quotient functor $q\colon B \to B/A$. Applying \Cref{thm:Verdier_Quotient}(3) to all shifts of $x$, and using that the colimit is filtered and hence commutes with cofibers, we obtain the following formula for the mapping spectra:
		\begin{align*}
			\hom_{B/A}(qx,qy) &\simeq \colim_{(a \to y) \in A_{/y}} \hom_B\big(x,\cofib(a \to y)\big) \\
			&\simeq \cofib\Big(\colim_{(a \to y) \in A_{/y}} \hom_B(x,a) \to \hom_B(x,y)\Big).
		\end{align*}
		Since $B/A \to (B/A)^{\idem}$ is fully faithful, the same formula computes mapping spectra in the idempotent complete quotient.
	\end{remark}

	\begin{remark}
		\label[remark]{rem:Lifting_Objects_Quotient}
		In a short exact sequence $A \to B \xrightarrow{p} C$, the functor $p$ need not be essentially surjective: in general, objects of $C$ are only retracts of objects in the image of $p$. It is possible to say precisely which objects lift. By Thomason's classification of dense subcategories \cite{Thomason_Classification}, the essential image of the dense subcategory $B/A \subseteq C$ consists precisely of those objects $c \in C$ whose class $[c] \in K_0(C)$ lies in the image of $K_0(B) \to K_0(C)$. For example, $c \oplus \Sigma c$ always lifts to $B$, since $[c \oplus \Sigma c] = [c] - [c] = 0$.
	\end{remark}

	\begin{example}[{Thomason--Trobaugh, \cite{Thomason_Trobaugh}}]
		\label[example]{ex:Thomason_Trobaugh_Sequence}
		Let $X$ be a qcqs scheme and $U \subseteq X$ a quasi-compact open subscheme with closed complement $Z$. Then
		\[
			\Perf_Z(X) \hookrightarrow \Perf(X) \to \Perf(U)
		\]
		is a short exact sequence, where $\Perf_Z(X)$ denotes the perfect complexes on $X$ whose restriction to $U$ vanishes. In general, the restriction functor $\Perf(X) \to \Perf(U)$ is not essentially surjective; by \Cref{rem:Lifting_Objects_Quotient}, a perfect complex on $U$ extends to $X$ if and only if its class in $K_0(U)$ extends.
	\end{example}

	\section{Localization sequences of presentable stable categories}
	\label[section]{sec:Localization_Sequences_PrL}

	% Source: written from the literature (Lurie HTT, Efimov, Thomason--Neeman); Mathew, Lecture 1 (MPIM Bonn, 09.09.2024); Talk 1 (Marc Hoyois, 13.10.2026)

	We now turn to large categories. Recall that $\PrL_{\st}$ denotes the category of presentable stable categories and colimit-preserving functors. A colimit-preserving functor $F\colon D \to E$ between presentable categories has a right adjoint $F^R$ (\Cref{thm:Adjoint_Functor_Theorem}); recall that $F$ is called \emph{strongly continuous} if $F^R$ again preserves colimits. We denote by $\Pr^{\mathrm{LL}}_{\st} \subseteq \PrL_{\st}$ the subcategory with the same objects and strongly continuous functors.

	\begin{definition}
		\label[definition]{def:Localization_Sequence_PrL}
		A sequence $K \xrightarrow{i} D \xrightarrow{p} E$ in $\PrL_{\st}$ with $pi \simeq 0$ is called a \emph{localization sequence} if $i$ is fully faithful with essential image $\ker(p)$ and $p$ is a \emph{localization}, i.e.\ its right adjoint $p^R$ is fully faithful. It is called \emph{strongly continuous} if moreover $i$ and $p$ are strongly continuous.
	\end{definition}

	For a full subcategory $K \subseteq D$ we write $K^{\perp} = \{x \in D \mid \Map_D(k,x) \simeq * \text{ for all } k \in K\}$ for its right orthogonal.

	\begin{proposition}
		\label[proposition]{prop:Localization_Sequences_PrL}
		\leavevmode
		\begin{enumerate}
			\item Let $D$ be a presentable stable category and $K \subseteq D$ a full stable subcategory which is closed under colimits and presentable. Then the inclusion $i\colon K \hookrightarrow D$ has a right adjoint $i^R$, the full subcategory $K^{\perp} \subseteq D$ is presentable, its inclusion admits a left adjoint $L\colon D \to K^{\perp}$ given by
			\[
				L(x) = \cofib(ii^R x \to x),
			\]
			and $K \xrightarrow{i} D \xrightarrow{L} K^{\perp}$ is a localization sequence. It is both a fiber sequence and a cofiber sequence in $\PrL_{\st}$.
			\item Conversely, let $K \xrightarrow{i} D \xrightarrow{p} E$ be a localization sequence. Then $p^R$ induces an equivalence $E \simeq K^{\perp}$ under which $p$ corresponds to $L$, and there is a natural cofiber sequence
			\[
				ii^R \to \id_D \to p^Rp.
			\]
		\end{enumerate}
	\end{proposition}
	\prooffor{prop:Localization_Sequences_PrL}
	\begin{proof}
		(1) The inclusion $i$ preserves colimits, so it has a right adjoint by \Cref{thm:Adjoint_Functor_Theorem}. Define $L(x) = \cofib(ii^Rx \to x)$, using the counit. For $k \in K$, the map $\hom_D(k,ii^Rx) \to \hom_D(k,x)$ is an equivalence since $i$ is fully faithful, so $\hom_D(k,L(x)) \simeq 0$; that is, $L(x) \in K^{\perp}$. For $y \in K^{\perp}$, applying $\hom_D(-,y)$ to the cofiber sequence $ii^Rx \to x \to L(x)$ shows that $\hom_D(L(x),y) \to \hom_D(x,y)$ is an equivalence, since $\hom_D(ii^Rx,y) \simeq 0$. Hence $L\colon D \to K^{\perp}$ is left adjoint to the inclusion $K^{\perp} \subseteq D$, with unit $x \to L(x)$. The functor $L$ is accessible since $i^R$ is (being a right adjoint between presentable categories), so $K^{\perp}$ is presentable \cite[Section~5.5.4]{lurie2009HTT}. Since its right adjoint is fully faithful, $L$ is a localization, and $L(x) \simeq 0$ if and only if $ii^Rx \to x$ is an equivalence, if and only if $x \in K$. This shows that $K \to D \to K^{\perp}$ is a localization sequence. It is a fiber sequence in $\PrL_{\st}$ since limits in $\PrL$ are computed in $\widehat{\Cat}$. For the cofiber sequence, let $F\colon D \to E$ be a colimit-preserving functor to a presentable stable category $E$. Since $L$ is a localization, restriction along $L$ identifies $\Fun^{\colim}(K^{\perp},E)$ with the full subcategory of $\Fun^{\colim}(D,E)$ of functors inverting the maps $x \to x'$ with $L(x) \simeq L(x')$. Since $L$ is exact, these are precisely the maps whose fiber lies in $\ker(L) = K$, and an exact functor inverts all such maps if and only if it vanishes on $K$. This is precisely the universal property of the cofiber of $i$ in $\PrL_{\st}$.

		(2) Since $p^R$ is fully faithful and $\Map_D(k,p^Re) \simeq \Map_E(pk,e) \simeq *$ for $k \in K$, the functor $p^R$ lands in $K^{\perp}$. Conversely, let $x \in K^{\perp}$ and let $f \to x \to p^Rp(x)$ be the fiber sequence of the unit. As $pp^R \simeq \id$, we have $p(f) \simeq 0$, so $f \in K$. Then $\hom_D(f,x) \simeq 0$ and $\hom_D(f,p^Rpx) \simeq \hom_E(pf,px) \simeq 0$, so $\hom_D(f,f) \simeq 0$ and hence $f \simeq 0$. So $x \simeq p^Rp(x)$, which shows that $p^R$ induces an equivalence $E \simeq K^{\perp}$. Under this equivalence, $p$ is identified with the left adjoint $L$ of the inclusion $K^{\perp} \subseteq D$, and the unit $x \to p^Rp(x) = L(x)$ is the cofiber of $ii^Rx \to x$ by (1).
	\end{proof}

	\begin{lemma}
		\label[lemma]{lem:Strong_Continuity_Localization_Sequence}
		Let $K \xrightarrow{i} D \xrightarrow{p} E$ be a localization sequence. Then $i^R$ preserves colimits if and only if $p^R$ preserves colimits. In particular, the sequence is strongly continuous as soon as $i$ or $p$ is strongly continuous.
	\end{lemma}
	\prooffor{lem:Strong_Continuity_Localization_Sequence}
	\begin{proof}
		Since $i$ is fully faithful and preserves colimits, $i^R$ preserves colimits if and only if $ii^R$ does: indeed, for a diagram $(x_j)$ the map $i(\colim_j i^Rx_j) \simeq \colim_j ii^Rx_j \to ii^R(\colim_j x_j)$ is the image under $i$ of the comparison map for $i^R$. Similarly, $p^R$ preserves colimits if and only if $p^Rp$ does: if $p^Rp$ preserves colimits and $(e_j)$ is a diagram in $E$, then $\colim_j e_j \simeq p(\colim_j p^Re_j)$, so $p^R(\colim_j e_j) \simeq p^Rp(\colim_j p^Re_j) \simeq \colim_j p^Rpp^Re_j \simeq \colim_j p^Re_j$. Finally, by \Cref{prop:Localization_Sequences_PrL}(2) there is a cofiber sequence of exact functors $ii^R \to \id \to p^Rp$, so $ii^R$ preserves colimits if and only if $p^Rp$ does.
	\end{proof}

	Recall that a presentable stable category $E$ is \emph{compactly generated} if $E \simeq \Ind(E^{\omega})$, equivalently if there is a set $G$ of compact objects of $E$ such that an object $e \in E$ with $\hom_E(g,e) \simeq 0$ for all $g \in G$ is zero.

	\begin{lemma}
		\label[lemma]{lem:Compact_Generators_Thick_Closure}
		Let $E$ be a presentable stable category and $G \subseteq E^{\omega}$ a set of compact objects such that $\hom_E(g,e) \simeq 0$ for all $g \in G$ implies $e \simeq 0$. Then $E^{\omega}$ is the thick closure of $G$, and $E \simeq \Ind(E^{\omega})$.
	\end{lemma}
	\prooffor{lem:Compact_Generators_Thick_Closure}
	\begin{proof}
		Let $T \subseteq E^{\omega}$ be the thick closure of $G$; it is essentially small. The functor $\Ind(T) \to E$ extending the inclusion is fully faithful since $T$ consists of compact objects, and its essential image is closed under colimits since $T$ is closed under finite colimits \cite[Section~5.3.5]{lurie2009HTT}. Its right adjoint is conservative: if $\hom_E(t,e) \simeq 0$ for all $t \in T$ then $e \simeq 0$ by assumption. Hence $\Ind(T) \simeq E$. Since $T$ is idempotent complete, the compact objects of $\Ind(T)$ are precisely the objects of $T$ (\Cref{thm:Accessible_Recollections}(1)), so $E^{\omega} = T$.
	\end{proof}

	\begin{lemma}
		\label[lemma]{lem:Compact_Objects_Localization_Sequence}
		Let $K \xrightarrow{i} D \xrightarrow{p} E$ be a strongly continuous localization sequence. Then $i$ and $p$ preserve compact objects, and $K^{\omega} = K \cap D^{\omega}$. If $D$ is compactly generated, then so is $E$, and $E^{\omega}$ is the thick closure of $p(D^{\omega})$.
	\end{lemma}
	\prooffor{lem:Compact_Objects_Localization_Sequence}
	\begin{proof}
		A colimit-preserving functor whose right adjoint preserves filtered colimits preserves compact objects. If $k \in K$ is compact in $D$, then it is compact in $K$ since $i$ is fully faithful and preserves colimits. If $D$ is compactly generated and $e \in E$ satisfies $\hom_E(p(d),e) \simeq \hom_D(d,p^Re) \simeq 0$ for all $d \in D^{\omega}$, then $p^Re \simeq 0$ and hence $e \simeq pp^Re \simeq 0$. The claim follows from \Cref{lem:Compact_Generators_Thick_Closure}.
	\end{proof}

	\begin{definition}[{\cite[Definition~1.9]{Efimov_Large_Categories}}]
		\label[definition]{def:Homological_Epimorphism}
		A functor $f\colon A \to B$ in $\Cat_{\perf}$ is called a \emph{homological epimorphism} if $\Ind(f)\colon \Ind(A) \to \Ind(B)$ is a localization, i.e.\ its right adjoint is fully faithful.
	\end{definition}

	\begin{theorem}[Thomason--Neeman]
		\label[theorem]{thm:Thomason_Neeman}
		\leavevmode
		\begin{enumerate}
			\item If $A \xrightarrow{i} B \xrightarrow{p} C$ is a short exact sequence in $\Cat_{\perf}$, then $\Ind(A) \to \Ind(B) \to \Ind(C)$ is a strongly continuous localization sequence. In particular, every Karoubi projection is a homological epimorphism.
			\item Conversely, if $K \xrightarrow{i} D \xrightarrow{p} E$ is a strongly continuous localization sequence in which $K$ and $D$ are compactly generated, then $E$ is compactly generated and $K^{\omega} \to D^{\omega} \to E^{\omega}$ is a short exact sequence.
		\end{enumerate}
	\end{theorem}
	\prooffor{thm:Thomason_Neeman}
	\begin{proof}
		(1) The functor $\Ind(i)$ is fully faithful, as mapping spaces between ind-objects are computed by the formula of \Cref{constr:Ind_Large}. Its essential image $K \subseteq \Ind(B)$ is closed under colimits and presentable, so by \Cref{prop:Localization_Sequences_PrL}(1) we obtain a localization sequence $K \to \Ind(B) \xrightarrow{L} K^{\perp}$ which is a cofiber sequence in $\PrL_{\st}$. For every presentable stable category $E$, we have natural equivalences
		\[
			\begin{aligned}
				\Fun^{\colim}(K^{\perp},E) &\simeq \Fun^{\colim}_{K \mapsto 0}(\Ind(B),E) \simeq \Fun^{\ex}_{A \mapsto 0}(B,E) \\
				&\simeq \Fun^{\ex}(B/A,E) \simeq \Fun^{\ex}((B/A)^{\idem},E) \simeq \Fun^{\colim}(\Ind(C),E),
			\end{aligned}
		\]
		using \Cref{thm:Ind_Kappa_Universal_Property}(3), \Cref{thm:Verdier_Quotient}(2), the fact that $E$ is idempotent complete, and $(B/A)^{\idem} \simeq C$. Under these equivalences, precomposition with $L$ corresponds to precomposition with $\Ind(p)$. By the Yoneda lemma there is an equivalence $K^{\perp} \simeq \Ind(C)$ under which $L$ corresponds to $\Ind(p)$. Hence $\Ind(A) \to \Ind(B) \to \Ind(C)$ is a localization sequence. Since $\Ind(i)$ preserves compact objects, its right adjoint preserves filtered colimits, hence all colimits; so the sequence is strongly continuous by \Cref{lem:Strong_Continuity_Localization_Sequence}.

		(2) Put $A = K^{\omega}$ and $B = D^{\omega}$. By \Cref{lem:Compact_Objects_Localization_Sequence}, $i$ restricts to a fully faithful functor $A \to B$ whose image $A = K \cap B$ is a thick subcategory of $B$. By \Cref{lem:Kernel_Verdier_Quotient}, the sequence $A \to B \to (B/A)^{\idem}$ is short exact, so by (1) we obtain a localization sequence $\Ind(A) \to \Ind(B) \to \Ind((B/A)^{\idem})$. As $K \simeq \Ind(A)$ and $D \simeq \Ind(B)$, both $E$ and $\Ind((B/A)^{\idem})$ are identified with $K^{\perp}$ compatibly with the functors from $D$, by \Cref{prop:Localization_Sequences_PrL}(2). Hence $E \simeq \Ind((B/A)^{\idem})$ is compactly generated with $E^{\omega} \simeq (B/A)^{\idem}$, and the claim follows.
	\end{proof}

	Thus, by \Cref{thm:Thomason_Neeman}, $\Ind$ turns short exact sequences of small categories into strongly continuous localization sequences of compactly generated categories. Conversely, a strongly continuous localization sequence $K \to D \to E$ with $D$ compactly generated comes from a short exact sequence of compact objects if and only if $K$ is compactly generated; strong continuity alone does not guarantee this. In the language of \Cref{def:Homological_Epimorphism}: a functor $f\colon A \to B$ in $\Cat_{\perf}$ is a homological epimorphism if and only if $\ker(\Ind(f)) \to \Ind(A) \to \Ind(B)$ is a strongly continuous localization sequence, and it is a Karoubi projection if and only if moreover $\ker(\Ind(f))$ is compactly generated. In general, the kernel of a homological epimorphism need not be compactly generated, see \Cref{ex:Keller} below; it is, however, always a dualizable category (\Cref{chap:Compactly_Assembled_Dualizable}).

	% Source: Mathew, Lecture 1 (MPIM Bonn, 09.09.2024)
	\Cref{thm:Thomason_Neeman} gives a description of idempotent complete quotients in terms of large categories. Recall from \Cref{cor:Ind_As_Exact_Functors} that $\Ind(B) \simeq \Fun^{\ex}(B\catop,\Sp)$ for $B \in \Cat_{\perf}$.

	\begin{corollary}
		\label[corollary]{cor:Karoubi_Quotient_Via_Ind}
		Let $i\colon A \hookrightarrow B$ be a fully faithful functor in $\Cat_{\perf}$. Under the identification $\Ind(-) \simeq \Fun^{\ex}((-)\catop,\Sp)$, the right adjoint $\Ind(i)^R\colon \Ind(B) \to \Ind(A)$ is given by restriction along $i$, and
		\[
			(B/A)^{\idem} \simeq \ker\big(\Ind(i)^R\big)^{\omega},
		\]
		where $\ker(\Ind(i)^R) \subseteq \Fun^{\ex}(B\catop,\Sp)$ is the full subcategory of exact functors vanishing on $A$.
	\end{corollary}

	For the description of exact presheaves on $B/A$ by this vanishing condition, see \href{https://animarium.org/tag/1155/}{[SCT, Corollary~19.6.8]}.

	\prooffor{cor:Karoubi_Quotient_Via_Ind}
	\begin{proof}
		For $a \in A$ and $F \in \Ind(B)$ we have $\hom_{\Ind(A)}(a,\Ind(i)^RF) \simeq \hom_{\Ind(B)}(ia,F) \simeq F(ia)$ by the Yoneda lemma, so $\Ind(i)^R(F) \simeq F \circ i\catop$. Since $A$ is idempotent complete and $i$ is fully faithful, the essential image of $i$ is a thick subcategory, so $A \to B \to (B/A)^{\idem}$ is short exact by \Cref{lem:Kernel_Verdier_Quotient}. By \Cref{thm:Thomason_Neeman}(1), $\Ind(A) \to \Ind(B) \to \Ind((B/A)^{\idem})$ is a localization sequence, so by \Cref{prop:Localization_Sequences_PrL}(2) we have $\Ind((B/A)^{\idem}) \simeq \Ind(A)^{\perp}$. Finally, $F \in \Ind(A)^{\perp}$ if and only if $\hom(a,\Ind(i)^RF) \simeq \hom(ia,F) \simeq 0$ for all $a \in A$, if and only if $\Ind(i)^RF \simeq 0$. Taking compact objects gives the claim.
	\end{proof}

	% Source: Talk 1 (Marc Hoyois, 13.10.2026)
	\begin{remark}
		\label[remark]{rem:Ind_Fully_Faithful_Localization}
		The functor $\Ind$ induces an equivalence of categories
		\[
			\Ind\colon \Cat_{\perf} \rightleftarrows \PrL_{\omega,\st}\noloc (-)^{\omega},
		\]
		where $\PrL_{\omega,\st}$ denotes the category of compactly generated presentable stable categories and colimit-preserving functors that preserve compact objects (equivalently, strongly continuous functors), see \Cref{thm:Accessible_Recollections} and \Cref{lem:Compact_Generators_Thick_Closure}. For a functor $f\colon A \to B$ in $\Cat_{\perf}$:
		\begin{itemize}
			\item $f$ is fully faithful if and only if $\Ind(f)$ is fully faithful (by the formula for mapping spaces in \Cref{constr:Ind_Large}), if and only if $\Ind(f)$ is a monomorphism \cite[Proposition~1.13]{Efimov_Large_Categories};
			\item If $f$ is a Karoubi projection, then $\Ind(f)$ is a localization (\Cref{thm:Thomason_Neeman}), or equivalently $\Ind(f)$ is an epimorphism \cite[Proposition~1.13]{Efimov_Large_Categories}. The converse does not hold, see \Cref{ex:Keller}.
		\end{itemize}
	\end{remark}

	% Source: Talk 1 (Marc Hoyois, 13.10.2026)
	\begin{example}[{Keller, \cite{Keller_Smashing}}]
		\label[example]{ex:Keller}
		Let $A$ be a ring and let $0 \neq I \subseteq \mathrm{Jac}(A)$ be a Tor-unital ideal. Then $A/I \otimes^L_A A/I \simeq A/I$, so the restriction functor $\Mod_{A/I} \to \Mod_A$ is fully faithful and $f\colon \Perf_A \to \Perf_{A/I}$ is a homological epimorphism. The category
		\[
			\ker(\Mod_A \to \Mod_{A/I})
		\]
		contains no nonzero compact objects. Note that $\Mod_A \simeq \Ind(\Perf_A)$ and $\Mod_{A/I} \simeq \Ind(\Perf_{A/I})$. So the kernel of $f$ is zero, whereas the large kernel contains the nonzero module $I$, since $(A/I) \otimes_A^L I \simeq 0$. This gives a counterexample to the converse of the second statement of \Cref{rem:Ind_Fully_Faithful_Localization}: $\Ind(f)$ is a localization, but $f$ is not a Karoubi projection.

		For example, consider $A = k[t,t^{1/2},t^{1/4}, \dots]_{(t,t^{1/2}, t^{1/4}, \dots)}$, for $k$ a field. Let $\mm$ be the maximal ideal. Then the map $A \to A/\mm$ is an example.
	\end{example}

	\section{Semiorthogonal decompositions}
	\label[section]{sec:Semiorthogonal_Decompositions}

	% Source: written from the literature (Efimov, Section 2.11); Mathew, Lecture 1 (MPIM Bonn, 09.09.2024)

	The simplest localization sequences are the split ones. In this case the middle category $E$ is glued from two subcategories $C$ and $D$ between which there are no maps in one direction, and every object of $E$ is uniquely an extension of an object of $C$ by an object of $D$. Such a splitting is called a semiorthogonal decomposition. We follow Efimov's convention: in $E = \lra{C_1,\dots,C_n}$ there are no nonzero maps from $C_i$ to $C_j$ for $i > j$. We first treat the presentable setting, which is the one relevant for dualizable categories, and then explain the analogous notion for small stable categories and its relation to short exact sequences.

	\begin{definition}[{\cite[Definition~2.74]{Efimov_Large_Categories}}]
		\label[definition]{def:Semiorthogonal_Decomposition}
		A \emph{semiorthogonal decomposition} of a presentable stable category $E$ is a collection of full presentable stable subcategories $C_1,\dots,C_n \subseteq E$ such that
		\begin{enumerate}
			\item the inclusion functors $C_i \hookrightarrow E$ are strongly continuous;
			\item for $i > j$, $x \in C_i$ and $y \in C_j$ we have $\hom_E(x,y) \simeq 0$;
			\item the smallest full stable subcategory of $E$ containing $C_1,\dots,C_n$ is $E$ itself.
		\end{enumerate}
		In this case we write $E = \lra{C_1,\dots,C_n}$.
	\end{definition}

	\begin{proposition}
		\label[proposition]{prop:Two_Term_SOD}
		Let $E = \lra{C,D}$ be a semiorthogonal decomposition, with inclusions $i_C\colon C \hookrightarrow E$ and $i_D\colon D \hookrightarrow E$. Then:
		\begin{enumerate}
			\item Every $x \in E$ fits into a cofiber sequence $d \to x \to c$ with $d \in D$ and $c \in C$, which is unique up to unique equivalence and functorial in $x$.
			\item The functor $i_C$ admits a left adjoint $i_C^L$ and $i_D$ admits a right adjoint $i_D^R$, given by $x \mapsto c$ and $x \mapsto d$ respectively; there is a natural cofiber sequence $i_Di_D^R \to \id_E \to i_Ci_C^L$, and $i_C^Li_D \simeq 0$.
			\item The sequence $D \xrightarrow{i_D} E \xrightarrow{i_C^L} C$ is a strongly continuous localization sequence, split by $i_C$.
		\end{enumerate}
	\end{proposition}
	\prooffor{prop:Two_Term_SOD}
	\begin{proof}
		Note first that $\hom_E(d,c) \simeq 0$ for $d \in D$ and $c \in C$. Suppose $d \to x \to c$ is a cofiber sequence with $d \in D$ and $c \in C$. For $d' \in D$ we then have $\hom_E(d',d) \xrightarrow{\simeq} \hom_E(d',x)$, since $\hom_E(d',c) \simeq 0$; so $d \to x$ exhibits $d$ as the value at $x$ of a right adjoint to $i_D$. Dually, for $c' \in C$ we have $\hom_E(c,c') \xrightarrow{\simeq} \hom_E(x,c')$, so $x \to c$ exhibits $c$ as the value at $x$ of a left adjoint of $i_C$. In particular, such a sequence is unique up to unique equivalence and functorial in $x$, wherever it exists.

		Let $E' \subseteq E$ be the full subcategory of those $x$ that admit such a cofiber sequence. It contains $C$ (take $0 \to c \to c$) and $D$ (take $d \to d \to 0$) and is closed under shifts. It is closed under cofibers: a map $x \to x'$ in $E'$ induces, by the universal properties above, a map of cofiber sequences, and passing to cofibers gives a cofiber sequence $\cofib(d \to d') \to \cofib(x \to x') \to \cofib(c \to c')$ with first term in $D$ and last term in $C$, as these are stable subcategories. Hence $E'$ is a stable subcategory containing $C$ and $D$, so $E' = E$ by assumption (3). This proves (1), and by \cite[Proposition~5.2.4.2]{lurie2009HTT} also the existence of the adjoints in (2), together with the cofiber sequence $i_Di_D^R \to \id \to i_Ci_C^L$. Applying it to $x = i_D(d)$ gives $i_C^Li_D(d) \simeq 0$.

		For (3), $i_C^L$ is a localization since its right adjoint $i_C$ is fully faithful, and $i_C^L(x) \simeq 0$ if and only if $i_Di_D^R x \to x$ is an equivalence, if and only if $x$ lies in the essential image of $D$. The right adjoints $i_D^R$ and $i_C$ preserve colimits by assumption (1) and since $C \subseteq E$ is closed under colimits, so the sequence is strongly continuous. Finally, $i_C^L i_C \simeq \id_C$ as $i_C$ is fully faithful, so $i_C$ splits $i_C^L$.
	\end{proof}

	% Source: Talk 6 (Jacopo Mazzei, 17.11.2026); cf. Efimov, Remark 2.75
	Conversely, a semiorthogonal decomposition can be recognized from the adjoints of the inclusion functors. In particular, the projections onto the two pieces are themselves strongly continuous.

	\begin{proposition}[{\cite[Remark~2.75]{Efimov_Large_Categories}}]
		\label[proposition]{prop:SOD_Characterization}
		Let $i_1\colon C \to E$ and $i_2\colon D \to E$ be strongly continuous functors between presentable stable categories. Then $i_1$ and $i_2$ are fully faithful and exhibit $E = \lra{i_1(C), i_2(D)}$ as a semiorthogonal decomposition if and only if the following conditions hold:
		\begin{enumerate}
			\item $i_1$ admits a left adjoint $i_1^L$ and $i_2$ admits a right adjoint $i_2^R$;
			\item the counit $i_1^Li_1 \to \id_C$ and the unit $\id_D \to i_2^Ri_2$ are equivalences;
			\item $i_1^Li_2 \simeq 0$, or equivalently $i_2^Ri_1 \simeq 0$;
			\item the sequence $i_2i_2^R \to \id_E \to i_1i_1^L$, formed by the counit and the unit, is a cofiber sequence.
		\end{enumerate}
		In this case, the functors $i_1^L$ and $i_2^R$ are strongly continuous.
	\end{proposition}

	The composite in (4) admits a unique nullhomotopy, as $\Map(i_2i_2^R,i_1i_1^L) \simeq \Map(i_1^Li_2i_2^R,i_1^L) \simeq *$ by (3); so the condition in (4) makes sense. The equivalence of the two formulations of (3) follows from $\hom_D(d,i_2^Ri_1c) \simeq \hom_E(i_2d,i_1c) \simeq \hom_C(i_1^Li_2d,c)$.

	% Source: Talk 6 (Jacopo Mazzei, 17.11.2026); strong continuity of the projections added from Efimov, Remark 2.75
	\prooffor{prop:SOD_Characterization}
	\begin{proof}
		Suppose first that $i_1$ and $i_2$ are fully faithful and $E = \lra{i_1(C),i_2(D)}$. Then (1), (3) and (4) are \Cref{prop:Two_Term_SOD}(2), applied to the essential images, and (2) holds as $i_1$ and $i_2$ are fully faithful.

		Conversely, assume (1)--(4). By (2), the functors $i_1$ and $i_2$ are fully faithful. For $c \in C$ and $d \in D$ we have $\hom_E(i_2d,i_1c) \simeq \hom_C(i_1^Li_2d,c) \simeq 0$ by (3), so the semiorthogonality condition holds. For $e \in E$, the cofiber sequence $i_2i_2^Re \to e \to i_1i_1^Le$ of (4) shows that $e$ lies in the stable subcategory generated by the essential images of $i_1$ and $i_2$. Finally, the inclusions are strongly continuous by assumption, and their essential images are closed under colimits, as $i_1$ and $i_2$ preserve colimits.

		It remains to show that $i_1^L$ and $i_2^R$ are strongly continuous. The right adjoint of $i_1^L$ is $i_1$, which preserves colimits. The right adjoint of $i_1i_1^L$ is $i_1i_1^R$, which preserves colimits since $i_1$ is strongly continuous. Passing to right adjoints is an equivalence between the stable category of colimit-preserving endofunctors of $E$ and the opposite of the stable category of accessible limit-preserving endofunctors, so it carries the fiber sequence $i_2i_2^R \to \id_E \to i_1i_1^L$ to a cofiber sequence $i_1i_1^R \to \id_E \to (i_2i_2^R)^R$. Hence $(i_2i_2^R)^R$ preserves colimits. Since $i_2$ is fully faithful, we have
		\[
			\hom_D(i_2^Re,d) \simeq \hom_E(i_2i_2^Re,i_2d) \simeq \hom_E(e,(i_2i_2^R)^R i_2 d),
		\]
		so $(i_2i_2^R)^R \circ i_2$ is right adjoint to $i_2^R$ and preserves colimits.
	\end{proof}

	% Source: Mathew, Lecture 1 (MPIM Bonn, 09.09.2024)
	\begin{remark}
		\label[remark]{rem:SOD_Small}
		\Cref{def:Semiorthogonal_Decomposition} makes sense for any stable category $E$ and full stable subcategories $C_1,\dots,C_n \subseteq E$, omitting condition (1). No presentability or size assumption is needed. Parts (1) and (2) of \Cref{prop:Two_Term_SOD} hold in this setting with the same proof, and the sequence $D \xrightarrow{i_D} E \xrightarrow{i_C^L} C$ is a Verdier sequence: $i_C^L$ is a localization with kernel $D$, and it inverts precisely the maps whose cofiber lies in $D$. If $E$ is idempotent complete, then so are $C$ and $D$, being thick subcategories of $E$ (as kernels of the exact functors $i_D^R$ and $i_C^L$), and the sequence is short exact. (In the talk, the definition was given simultaneously for presentable stable categories, for dualizable categories and for small stable categories.)
	\end{remark}

	Conversely, short exact sequences admitting adjoints give rise to semiorthogonal decompositions.

	\begin{proposition}
		\label[proposition]{prop:SOD_From_Short_Exact}
		Let $A \xrightarrow{i} B \xrightarrow{p} C$ be a short exact sequence in $\Cat_{\perf}$. Then $i$ admits a right adjoint $i^R$ if and only if $p$ admits a right adjoint $p^R$. In this case:
		\begin{enumerate}
			\item The functor $p^R$ is fully faithful, with essential image the right orthogonal $A^{\perp} = \{x \in B \mid \hom_B(ia,x) \simeq 0 \text{ for all } a \in A\}$;
			\item There is a semiorthogonal decomposition $B = \lra{p^R(C),\, i(A)}$;
			\item The sequence $C \xrightarrow{p^R} B \xrightarrow{i^R} A$ is again short exact.
		\end{enumerate}
	\end{proposition}
	\prooffor{prop:SOD_From_Short_Exact}
	\begin{proof}
		Assume first that $i^R$ exists, and for $b \in B$ let $L(b) := \cofib(ii^Rb \to b)$. As in the proof of \Cref{prop:Localization_Sequences_PrL}(1), $L(b) \in A^{\perp}$ and $L\colon B \to A^{\perp}$ is left adjoint to the inclusion. Moreover, $L(b) \simeq 0$ if and only if $b \simeq ii^Rb$ lies in $A$, so $L$ inverts precisely the maps whose cofiber lies in $A$. Since $L$ is a reflective localization, it exhibits $A^{\perp}$ as the localization of $B$ at these maps \cite[Proposition~5.2.7.12]{lurie2009HTT}, so $B/A \simeq A^{\perp}$. As $A^{\perp}$ is closed under retracts in the idempotent complete category $B$, it is idempotent complete, so $C \simeq (B/A)^{\idem} \simeq A^{\perp}$, with $p$ corresponding to $L$. Hence $p$ admits the fully faithful right adjoint $A^{\perp} \hookrightarrow B$.

		Conversely, assume that $p^R$ exists. For $a \in A$ and $c \in C$ we have $\hom_B(ia,p^Rc) \simeq \hom_C(pia,c) \simeq 0$, so $p^R$ lands in $A^{\perp}$. Every $x \in A^{\perp}$ is local with respect to the maps $s\colon b' \to b$ with $\cofib(s) \in A$, i.e.\ $s^*\colon \hom_B(b,x) \to \hom_B(b',x)$ is an equivalence. Hence $\hom_{B/A}(qb,qx) \simeq \hom_B(b,x)$ for all $b \in B$ and $x \in A^{\perp}$ (for instance by \Cref{rem:Mapping_Spectra_Quotient}, applied to $B\catop$). Applying this to $x = p^Rc$, we find that the composite
		\[
			\hom_B(b,p^Rc) \to \hom_C(pb,pp^Rc) \to \hom_C(pb,c),
		\]
		whose second map is induced by the counit, is the adjunction equivalence. Its first map is an equivalence as well: since the functor $B/A \to C$ induced by $p$ is fully faithful, it is identified with the equivalence $\hom_B(b,p^Rc) \simeq \hom_{B/A}(qb,qp^Rc)$. Hence so is the second map, i.e.\ the counit $pp^Rc \to c$ induces an equivalence on $\hom_C(pb,-)$ for all $b \in B$. Since every object of $C$ is a retract of an object of the form $pb$, the counit is an equivalence, i.e.\ $p^R$ is fully faithful. For $b \in B$, let $f \to b \to p^Rpb$ be the fiber sequence of the unit. Then $p(f) \simeq 0$, so $f$ lies in $A$, and for $a \in A$ we get $\hom_B(ia,f) \simeq \hom_B(ia,b)$ since $\hom_B(ia,p^Rpb) \simeq 0$. Hence $b \mapsto f$ is right adjoint to $i$.

		It remains to prove (1)--(3). For (1), it remains to show that every $x \in A^{\perp}$ lies in the image of $p^R$. The fiber $f$ of the unit $x \to p^Rpx$ lies in $A$, so $\hom_B(f,x) \simeq 0$ and $\hom_B(f,p^Rpx) \simeq 0$, whence $\hom_B(f,f) \simeq 0$ and $f \simeq 0$. For (2), we have $\hom_B(ia,p^Rc) \simeq 0$, and every $b \in B$ fits into the cofiber sequence $ii^Rb \to b \to p^Rpb$. For (3), the kernel of $i^R$ is $A^{\perp} = p^R(C)$ by (1), and $i^R$ is a coreflective localization since its left adjoint $i$ is fully faithful; applying the first paragraph to opposite categories, $i^R$ exhibits $A$ as the Verdier quotient $B/p^R(C)$.
	\end{proof}

	\section{Calkin categories}
	\label[section]{sec:Calkin_Small}

	% Source: Mathew, Lecture 1 (MPIM Bonn, 09.09.2024)
	For $A \in \Cat_{\perf}$ and a regular cardinal $\kappa$, we write $\Ind(A)^{\kappa}$ for the full subcategory of $\kappa$-compact objects of $\Ind(A)$. It is essentially small and stable, and it contains $A = \Ind(A)^{\omega}$ as a thick subcategory.

	\begin{definition}[{\cite[Section~2.9]{Efimov_Large_Categories}}]
		\label[definition]{def:Calkin_Small}
		Let $A \in \Cat_{\perf}$ and let $\kappa$ be an uncountable regular cardinal. The \emph{$\kappa$-Calkin category} of $A$ is
		\[
			\Calk_{\kappa}(A) := \big(\Ind(A)^{\kappa}/A\big)^{\idem},
		\]
		so that there is a short exact sequence $A \hookrightarrow \Ind(A)^{\kappa} \to \Calk_{\kappa}(A)$.
	\end{definition}

	\begin{remark}
		\label[remark]{rem:Calkin_Mapping_Spectra}
		The mapping spectra in $\Calk_{\kappa}(A)$ can be computed using \Cref{rem:Mapping_Spectra_Quotient}. Let $x = \quotcolim_{i \in I} x_i$ and $y = \quotcolim_{j \in J} y_j$ be objects of $\Ind(A)^{\kappa}$, with $I$ and $J$ filtered. The functor $J \to A_{/y}$ sending $j$ to $y_j \to y$ is cofinal, and $\hom_{\Ind(A)}(x,y_j) \simeq \lim_i \hom_A(x_i,y_j)$. Hence
		\[
			\hom_{\Calk_{\kappa}(A)}(x,y) \simeq \cofib\Big(\colim_{j}\lim_{i}\hom_A(x_i,y_j) \to \lim_{i}\colim_{j}\hom_A(x_i,y_j)\Big).
		\]
		So the Calkin category measures the failure of the limit over $i$ to commute with the filtered colimit over $j$.

		To justify the cofinality assertion, fix $a \to y$ with $a \in A$. The comma category $(a \to y)/J$ is the category of elements of the functor
		\[
			j \mapsto \fib_{a \to y}\big(\Map_A(a,y_j) \to \Map_{\Ind(A)}(a,y)\big).
		\]
		Its groupoid completion is the colimit of this functor. Since $a$ is compact and filtered colimits of animae commute with finite limits, this colimit is the fiber of the identity of $\Map_{\Ind(A)}(a,y)$ over $a \to y$, hence is contractible. The cofinality criterion therefore applies.
	\end{remark}

	\begin{example}
		\label[example]{ex:Calkin_Perf}
		Let $R$ be an $\bbE_{\infty}$-ring and $A = \Perf_R$, so that $\Ind(A) \simeq \Mod_R$. For $M \in \Mod_R$ and a perfect $R$-module $P$ we have $\hom_R(M,P) \simeq \hom_R(M,R) \otimes_R P$, since $P$ is dualizable. Writing $M' \in \Mod_R$ as a filtered colimit of perfect modules, \Cref{rem:Mapping_Spectra_Quotient} gives
		\[
			\hom_{\Calk_{\kappa}(\Perf_R)}(M,M') \simeq \cofib\big(\hom_R(M,R) \otimes_R M' \to \hom_R(M,M')\big)
		\]
		for $M,M' \in \Mod_R^{\kappa}$. For instance, if $R = k$ is a field and $V$ is a $k$-vector space of dimension $< \kappa$, the map $V^{\vee} \otimes_k V \to \End_k(V)$ is injective with image the finite rank endomorphisms, so
		\[
			\End_{\Calk_{\kappa}(\Perf_k)}(V) \simeq \End_k(V)/\{\text{finite rank endomorphisms}\}.
		\]
		This is the algebraic analogue of the Calkin algebra $B(\Hh)/\Kk(\Hh)$ of bounded operators modulo compact operators on a Hilbert space, which explains the name.
	\end{example}

	\section{Localizing invariants}
	\label[section]{sec:Localizing_Invariants}

	% Source: Mathew, Lecture 1 (MPIM Bonn, 09.09.2024); Talk 1 (Marc Hoyois, 13.10.2026)
	We start with some recollections on the functor $K\colon \Cat_{\st} \to \Sp$.

	\begin{definition}
		\label[definition]{def:K_0}
		The \emph{Grothendieck group} $K_0(A)$ of a small stable category $A$ is the free abelian group on the equivalence classes $[x]$ of objects $x \in A$, modulo the relations $[x] = [x'] + [x'']$ for every cofiber sequence $x' \to x \to x''$ in $A$.
	\end{definition}

	Nonconnective algebraic K-theory is a spectrum-valued refinement $K\colon \Cat_{\st} \to \Sp$ of this construction, satisfying $\pi_0K(A) \cong K_0(A)$ for $A$ idempotent complete. Its fundamental property is that it is a localizing invariant in the following sense.

	\begin{definition}
		\label[definition]{def:Localizing_Invariant}
		Let $\Ee$ be a stable category. A functor $F\colon \Cat_{\perf} \to \Ee$ is called a \emph{localizing invariant} if $F(0) \simeq 0$ and $F$ sends short exact sequences (\Cref{def:Short_Exact_Sequence}) to fiber sequences. We extend $F$ to $\Cat_{\st}$ by setting $F(A) := F(A^{\idem})$. We denote by $\Fun_{\loc}(\Cat_{\st},\Ee)$ the full subcategory of localizing invariants.
	\end{definition}

	By \Cref{thm:Thomason_Neeman}, a sequence $A \to B \to C$ in $\Cat_{\perf}$ is short exact if and only if $\Ind(A) \to \Ind(B) \to \Ind(C)$ is a strongly continuous localization sequence in $\PrL_{\st}$.

	\begin{lemma}
		\label[lemma]{lem:Localizing_Invariants_Verdier}
		Let $\Ee$ be stable and let $F\colon \Cat_{\st} \to \Ee$ satisfy $F(0) \simeq 0$. Consider the following conditions:
		\begin{enumerate}
			\item For every Verdier sequence $A \hookrightarrow B \to B/A$ in the sense of \Cref{def:Verdier_Sequence}, $F$ takes this sequence to a fiber sequence.
			\item Each sequence $A \hookrightarrow B \to C$ with $A \to B$ fully faithful and inducing an equivalence $(B/A)^{\idem} \iso C^{\idem}$ is sent to a fiber sequence under $F$.
		\end{enumerate}
		Then (2) implies (1), and (2) holds if and only if $F$ satisfies (1) and inverts the maps $A \to A^{\idem}$. In this case, $F$ is (the extension of) a localizing invariant in the sense of \Cref{def:Localizing_Invariant}.
	\end{lemma}
	\prooffor{lem:Localizing_Invariants_Verdier}
	\begin{proof}
		The sequences in (1) are among those of (2). Suppose that $F$ satisfies (1) and inverts idempotent completions. For a sequence as in (2), the fully faithful functor $A^{\idem} \hookrightarrow B^{\idem}$ has thick image, so
		\[
			A^{\idem} \hookrightarrow B^{\idem} \to B^{\idem}/A^{\idem}
		\]
		is a Verdier sequence. Moreover, $(B^{\idem}/A^{\idem})^{\idem} \simeq (B/A)^{\idem} \simeq C^{\idem}$, by the universal property of the idempotent complete quotient. Applying (1) and idempotent invariance therefore gives (2). Conversely, (2) applied to $0 \to A \to A^{\idem}$ shows that $F(A) \to F(A^{\idem})$ is an isomorphism.
	\end{proof}

	\begin{remark}
		\label[remark]{rem:Verdier_Localizing_Not_Karoubi}
		Condition (1) of \Cref{lem:Localizing_Invariants_Verdier} alone does not imply that $F$ inverts $A \to A^{\idem}$. For instance, connective algebraic K-theory satisfies (1) (see the proof of \cite[Corollary~4.4.15]{Hermitian_II}), but $K_0(A) \to K_0(A^{\idem})$ is in general not surjective. This is why localizing invariants are defined on idempotent complete categories.
	\end{remark}

	% Source: Mathew, Lecture 1 (MPIM Bonn, 09.09.2024)
	\begin{remark}
		\label[remark]{rem:Finitary_Localizing_Invariants}
		One frequently additionally requires that $\Ee$ is presentable and that $F$ preserves filtered colimits (or at least $\kappa$-filtered colimits for some regular cardinal $\kappa$). Algebraic K-theory and $\THH$ are examples of localizing invariants preserving filtered colimits. Topological cyclic homology $\mathrm{TC}$ is localizing, but does not in general preserve filtered colimits.
	\end{remark}

	\begin{example}
		\label[example]{ex:Localizing_Additive}
		Let $B = \lra{C,D}$ be a semiorthogonal decomposition in $\Cat_{\perf}$ (\Cref{rem:SOD_Small}). Then every localizing invariant $F$ sends the short exact sequence $D \to B \to C$ to a fiber sequence $F(D) \to F(B) \to F(C)$, which is split by $F(i_C)$. Hence
		\[
			F(B) \simeq F(C) \oplus F(D).
		\]
		Functors $F\colon \Cat_{\perf} \to \Ee$ with this property are called \emph{additive}.
	\end{example}

	\begin{example}
		\label[example]{ex:Localizing_Calkin}
		Let $A \in \Cat_{\perf}$ and let $F$ be a localizing invariant. The short exact sequence $A \to \Ind(A)^{\omega_1} \to \Calk_{\omega_1}(A)$ of \Cref{def:Calkin_Small} gives a fiber sequence
		\[
			F(A) \to F(\Ind(A)^{\omega_1}) \to F(\Calk_{\omega_1}(A)).
		\]
		Since $\Ind(A)^{\omega_1}$ admits countable coproducts, the middle term vanishes by an Eilenberg swindle: for $T = \bigoplus_{\N}\colon \Ind(A)^{\omega_1} \to \Ind(A)^{\omega_1}$ we have $\id \oplus T \simeq T$, and additivity gives $F(\id) + F(T) = F(T)$, so $F(\id) = 0$. Hence
		\[
			F(A) \simeq \Omega F(\Calk_{\omega_1}(A)).
		\]
		This is the formula that will be used to extend $F$ to dualizable categories, see \Cref{def:Continuous_Extension}.
	\end{example}

	The most fundamental localizing invariant is nonconnective algebraic K-theory, as made precise by the following universal property.

	\begin{theorem}
		\label[theorem]{thm:K_Universal_Property}
		Nonconnective algebraic K-theory $K\colon \Cat_{\perf} \to \Sp$ is the initial localizing invariant equipped with a natural transformation $(-)^{\simeq} \to \Omega^{\infty}K(-)$.
	\end{theorem}

	\begin{remark}
		\label[remark]{rem:K_Universal_Property_Literature}
		In the literature, \Cref{thm:K_Universal_Property} is usually stated for localizing invariants preserving filtered colimits, see \cite{BGT2013Universal} and also \cite{Barwick_Higher_K_Theory}. The version without this condition is explicitly stated in \cite[page~115, footnote~19]{Ramzi_Thesis}. That footnote does not supply a proof of the additional generality.
	\end{remark}

	\begin{definition}[{\cite{BGT2013Universal}}]
		\label[definition]{def:Mot_Loc}
		The category of \emph{localizing motives} $\mathrm{Mot}_{\loc}$ is the presentable stable category equipped with a localizing invariant
		\[
			U_{\loc}\colon \Cat_{\perf} \to \mathrm{Mot}_{\loc}
		\]
		preserving filtered colimits, characterized as follows: for every presentable stable category $\Ee$, precomposition with $U_{\loc}$ induces an equivalence between colimit-preserving functors $\mathrm{Mot}_{\loc}\to\Ee$ and localizing invariants $\Cat_{\perf}\to\Ee$ preserving filtered colimits. We fix a universe for the small categories in $\Cat_{\perf}$ and use a larger universe for functor categories; these size conventions are suppressed in the notation.
	\end{definition}

	Although $\mathrm{Mot}_{\loc}$ is defined by a universal property, it can be constructed explicitly as a localization of a category of spectral presheaves on a suitable small subcategory of $\Cat_{\perf}$. This is the \emph{finitary} category of motives. It is not asserted to classify all the localizing invariants of \Cref{def:Localizing_Invariant}.

	\begin{theorem}[{\cite[Theorem~9.8]{BGT2013Universal}}]
		\label[theorem]{thm:K_Corepresentable}
		For all $A \in \Cat_{\perf}$ there is a natural equivalence
		\[
			\hom_{\mathrm{Mot}_{\loc}}\big(U_{\loc}(\Sp^{\omega}),U_{\loc}(A)\big) \simeq K(A).
		\]
	\end{theorem}

	% Source: Talk 1 (Marc Hoyois, 13.10.2026)
	Marc formulated the universal property of K-theory as follows.

	\begin{theorem}[{Theorem/definition, cf.\ \cite{BGT2013Universal}}]
		\label[theorem]{thm:K_Unit_Day_Convolution}
		$K$ is the unit for the Day convolution operad structure on $\Fun_{\loc}(\Cat_{\perf},\Sp)$, where localizing invariants are not required to preserve filtered colimits.
	\end{theorem}

	Let us explain the statement. The category $\Fun(\Cat_{\perf},\Sp)$ carries the Day convolution symmetric monoidal structure, induced by the tensor product on $\Cat_{\perf}$ and the smash product on $\Sp$. Its unit is the functor $\Sigma^{\infty}_+\Map_{\Cat_{\perf}}(\Sp^{\omega},-) \simeq \Sigma^{\infty}_+(-)^{\simeq}$. The full subcategory $\Fun_{\loc}(\Cat_{\perf},\Sp)$ of localizing invariants is not obviously closed under Day convolution, but it inherits a structure of an operad: a multimorphism $(F_1,\dots,F_n) \to G$ is a natural transformation $F_1(A_1) \otimes \cdots \otimes F_n(A_n) \to G(A_1 \otimes \cdots \otimes A_n)$. The claim is that $K$ is a unit for this operad: precomposing with the natural transformation $\Sigma^{\infty}_+(-)^{\simeq} \to K$ of \Cref{thm:K_Universal_Property} in one of the inputs induces an equivalence on spaces of multimorphisms whose other inputs and target are localizing.

	% Source: written from the literature (universal property as in Ramzi's thesis); proof drafted with generative AI (Claude)
	\prooffor{thm:K_Unit_Day_Convolution}
	\begin{proof}[Proof sketch]
		By \Cref{thm:K_Universal_Property}, precomposition with $\Sigma^{\infty}_+(-)^{\simeq} \to K$ induces an equivalence
		\[
			\Map(K,H) \simeq \Map\big(\Sigma^{\infty}_+(-)^{\simeq},H\big)
		\]
		for every localizing invariant $H$. Now let $F_2,\dots,F_n$ and $G$ be localizing invariants. A multimorphism $(F_1,F_2,\dots,F_n) \to G$ is the same as a natural transformation from $F_1$ to the functor
		\[
			H(A) := \int_{B_2,\dots,B_n} \hom\big(F_2(B_2) \otimes \cdots \otimes F_n(B_n),\, G(A \otimes B_2 \otimes \cdots \otimes B_n)\big),
		\]
		where $\int$ denotes the end expressing naturality in the remaining variables. The functor $H$ is localizing: for fixed $B_2,\dots,B_n$, the functor $A \mapsto A \otimes B_2 \otimes \cdots \otimes B_n$ preserves short exact sequences, $G$ is localizing, and $\hom(X,-)$ and limits preserve fiber sequences. Applying the above equivalence to $H$ proves the claim for the first input, and by symmetry for every input.
	\end{proof}

	This argument uses the universal property of \Cref{thm:K_Universal_Property} without the assumption that localizing invariants preserve filtered colimits, not the universal property of $\mathrm{Mot}_{\loc}$. Note also that we do not claim that the category of localizing invariants is presentable.

	\begin{theorem}[{Tamme, before Efimov, \cite[Theorem~18]{Tamme_Excision}}]
		\label[theorem]{thm:Tamme_Excision}
		Consider a square $Q$ in $\Cat_{\st}$ of the form
		\[
			\begin{tikzcd}
				A \rar["f"] \dar & B \dar \\
				C \rar["g"'] & D.
			\end{tikzcd}
		\]
		Assume that $\Ind(Q)$ is cartesian, and that $\Ind(f)$ and $\Ind(g)$ are localizations. Then every localizing invariant $F$ sends $Q$ to a cartesian square.
	\end{theorem}

	% Source: proof strategy written from the literature (Efimov, Propositions 2.86 and 8.11)
	We will prove this theorem in \Cref{chap:Localizing_Invariants_Dualizable} for accessible localizing invariants with values in an accessible stable category with $\lambda$-filtered colimits for some $\lambda$, which includes K-theory. The idea is as follows. The square $\Ind(Q)$ is a pullback of compactly generated categories along the localization $\Ind(g)$, so by \Cref{prop:Dualizable_Pullbacks}(2) it is also a pullback in $\Pr^{\dual}$, and $\Ind(f)$ is again a localization with $\ker(\Ind(f)) \simeq \ker(\Ind(g))$. The continuous extension $F^{\cont}$ of \Cref{thm:Efimov_Continuous_Extension} sends the two resulting short exact sequences of dualizable categories to fiber sequences, and $F^{\cont}(\Ind(A)) \simeq F(A)$. Hence $F(Q)$ is cartesian \cite[Proposition~8.11]{Efimov_Large_Categories}. In particular, only $\Ind(g)$ needs to be a localization.

	\begin{warning}
		\label[warning]{warn:Excision_Requires_Ind_Localizations}
		One cannot simply apply $F$ to the common kernel of $f$ and $g$. The hypotheses say that $\Ind(f)$ and $\Ind(g)$ are localizations; they do not say that $f$ and $g$ are Karoubi projections. Their small kernels therefore need not give short exact sequences to which $F$ applies.
	\end{warning}

	\begin{remark}
		\label[remark]{rem:Tamme_Lax_Pullback_Proof}
		Tamme's original proof predates continuous K-theory and applies to localizing invariants with values in any stable category, without accessibility assumptions. It uses the lax pullback $P = B \times_D^{\to} C$ of triples $(b,c,\varphi\colon h(b) \to g(c))$, where $h\colon B \to D$ is the right vertical functor. There is a split exact sequence $C \to P \to B$ \cite[Proposition~10]{Tamme_Excision}, and under the hypotheses of \Cref{thm:Tamme_Excision} the functor $(b,c,\varphi) \mapsto \cofib(\varphi)$ fits into an exact sequence $A \to P \to D$ \cite[Theorem~15]{Tamme_Excision}. Comparing the two sequences gives the claim. See \cite{Land_Tamme_Pullbacks} for a generalization to arbitrary pullbacks of ring spectra.
	\end{remark}

	\begin{example}[{Milnor patching, \cite{Milnor_Intro_K_Theory}}]
		\label[example]{ex:Milnor_Patching}
		Consider a pullback square of $\bbE_1$-rings:
		\[
			\begin{tikzcd}
				A \rar \dar & A' \dar \\
				B \rar & B'
			\end{tikzcd}
		\]
		and assume that $A' \otimes_A A' \iso A'$. Then the induced square
		\[
			\begin{tikzcd}
				\Mod_A \rar \dar & \Mod_{A'} \dar \\
				\Mod_B \rar & \Mod_{B'}
			\end{tikzcd}
		\]
		is also a pullback square, and $\Mod_B \to \Mod_{B'}$ is a localization \cite[Lemma~23 and Theorem~26]{Tamme_Excision}. In particular, we deduce:
	\end{example}

	\begin{theorem}[{\cite[Theorem~28]{Tamme_Excision}; see also \cite{Land_Tamme_Pullbacks}}]
		\label[theorem]{thm:Milnor_Patching_K_Theory}
		In the above situation, there is a pullback square
		\[
			\begin{tikzcd}
				K(A) \rar \dar & K(A') \dar \\
				K(B) \rar & K(B').
			\end{tikzcd}
		\]
	\end{theorem}

	\begin{example}
		\label[example]{ex:Milnor_Square_Of_Rings}
		Given a map of rings $A \to B$ that sends an ideal $I \subseteq A$ isomorphically to the induced ideal $IB \subseteq B$, the induced square
		\[
			\begin{tikzcd}
				A \rar \dar & A/I \dar \\
				B \rar & B/IB
			\end{tikzcd}
		\]
		is a pullback of $\bbE_1$-rings.
	\end{example}


	% Theorem numbers for Efimov_Large_Categories refer to arXiv:2405.12169v4.
	\chapter{Compactly assembled and dualizable categories}
	\label[chapter]{chap:Compactly_Assembled_Dualizable}

	% Source: drafted with generative AI (Claude), based on the discussion of compactly assembled categories in Talk 2 (Marc Hoyois, 20.10.2026)
	In \Cref{chap:Introduction_And_Overview} we saw that compact generation is too restrictive for many purposes: for instance, the category $\Shv(\R^n;\Sp)$, for $n \geq 1$, has no nonzero compact objects at all. The goal of this chapter is to introduce the class of categories that replaces the compactly generated ones: the \emph{dualizable} categories, or, in the unstable setting, the \emph{compactly assembled} categories.

	To motivate the definition, let us first ask what makes compact generation so useful. Let $C$ be a category with filtered colimits, and let us call an ind-object $X \in \Ind(C)$ together with an equivalence $\colim X \simeq x$ a \emph{presentation} of $x \in C$. Every object $x$ has a trivial presentation, namely $y(x)$; by the adjunction $\colim \dashv y$, every presentation of $x$ maps uniquely to $y(x)$, compatibly with the identifications of the colimits with $x$. So $y(x)$ is the \emph{terminal}, or coarsest, presentation of $x$. If $C = \Ind(C_0)$ is compactly generated, then $x$ has a second canonical presentation
	\[
		\hat{y}(x) = \quotcolim_{(c \to x) \in (C_0)_{/x}} c,
	\]
	and its usefulness comes from the fact that it maps into every other presentation: since every map $c \to \colim_j z_j$ from a compact object factors through some stage $z_j$, we have
	\[
		\Map_C\big(x,\colim_j z_j\big) \simeq \Map_{\Ind(C)}\big(\hat{y}(x), \quotcolim_j z_j\big)
	\]
	for every filtered diagram $(z_j)_j$ in $C$. In other words, $\hat{y}(x)$ is the \emph{initial}, or finest, presentation of $x$. This formula says that maps out of $x$ into filtered colimits are computed by the formal object $\hat{y}(x)$, even if $x$ itself is not compact.

	A category is \emph{compactly assembled} if every object admits such a finest presentation, i.e.\ if $\colim\colon \Ind(C) \to C$ admits a left adjoint $\hat{y}$. Unlike in the compactly generated case, the terms of the ind-object $\hat{y}(x)$ are no longer required to be compact. A presentable stable category is compactly assembled if and only if it is dualizable (\Cref{thm:Lurie_Dualizable_Characterization}), and this is the class of categories to which K-theory will be extended.

	\section{Compactly assembled categories}
	\label[section]{sec:Compactly_Assembled}

	We now make this precise. In this section $C$ is a category with filtered colimits, not necessarily stable; the stable case is the subject of \Cref{sec:Dualizable_Categories}. Besides the definition, the key notion is that of a \emph{compact morphism}, which in compactly assembled categories takes over the role played by compact objects in compactly generated ones. We also show that the compactly assembled categories are precisely the retracts of compactly generated ones (\Cref{thm:Compactly_Assembled_Retract}), and introduce the corresponding class of strongly continuous functors.

	% Source: Talk 2 (Marc Hoyois, 20.10.2026)
	\begin{definition}
		\label[definition]{def:Compactly_Assembled}
		A category $C$ is called \emph{compactly assembled} if it is accessible, admits filtered colimits, and the colimit functor $\colim\colon \Ind(C) \to C$ admits a left adjoint $\hat{y}\colon C \hookrightarrow \Ind(C)$.
	\end{definition}

	Note that $\hat{y}$ is automatically fully faithful, since its double right adjoint $y\colon C \hookrightarrow \Ind(C)$ is fully faithful (\Cref{lem:Adjoint_Triples}).

	% Source: drafted with generative AI (Claude)
	Since $y$ is fully faithful, the counit $\colim y \to \id$ is invertible. Inverting it gives a map $\id \to \colim y$, which under the adjunction $\hat{y} \dashv \colim$ corresponds to a natural transformation
	\[
		\hat{y} \to y.
	\]
	By construction, applying $\colim$ to $\hat{y}(x) \to y(x)$ gives (up to the identifications $\colim\hat{y}(x) \simeq x \simeq \colim y(x)$) the identity of $x$.

	\begin{lemma}
		\label[lemma]{lem:Compact_Iff_Yhat_Is_Y}
		Let $C$ be compactly assembled and $x \in C$. Then $x$ is compact if and only if the map $\hat{y}(x) \to y(x)$ is an equivalence.
	\end{lemma}
	\prooffor{lem:Compact_Iff_Yhat_Is_Y}
	\begin{proof}
		Consider the natural transformation
		\[
			\Map_{\Ind(C)}(y(x),-) \to \Map_C(x,\colim(-))
		\]
		of functors $\Ind(C) \to \An$ given by applying $\colim$ and using $\colim y(x) \simeq x$. Composing with the adjunction equivalence $\Map_C(x,\colim(-)) \simeq \Map_{\Ind(C)}(\hat{y}(x),-)$, it corresponds by the Yoneda lemma to the map $\hat{y}(x) \to y(x)$, by the description of this map given above. Hence $\hat{y}(x) \to y(x)$ is an equivalence if and only if this transformation is an equivalence. Evaluated at $\quotcolim_j z_j$, it is the canonical map
		\[
			\colim_j \Map_C(x,z_j) \simeq \Map_{\Ind(C)}\big(y(x),\quotcolim_j z_j\big) \to \Map_C\big(x,\colim_j z_j\big),
		\]
		where the first equivalence holds as $y(x)$ is compact in $\Ind(C)$. This is an equivalence for all filtered diagrams $(z_j)$ if and only if $x$ is compact.
	\end{proof}

	% Source: drafted with generative AI (Claude); cf. Clausen, Lecture 2 (18.07.2023) and Mathew, Lecture 3 (MPIM Bonn, 10.09.2024)
	Compact objects of $C$ are those $x$ for which maps out of $x$ into a filtered colimit factor through some stage of the colimit. The following relative version of this notion is the basic tool for studying compactly assembled categories.

	\begin{definition}
		\label[definition]{def:Compact_Morphism}
		Let $C$ be a category with filtered colimits. A morphism $f\colon x \to x'$ in $C$ is called \emph{compact} if the natural transformation
		\[
			f^*\colon \Map_C\big(x',\colim(-)\big) \to \Map_C\big(x,\colim(-)\big)
		\]
		of functors $\Ind(C) \to \An$ factors through the natural transformation
		\[
			\Map_{\Ind(C)}\big(y(x),-\big) \to \Map_C\big(x,\colim(-)\big)
		\]
		given by applying $\colim$ and using $\colim y(x) \simeq x$.
	\end{definition}

	Evaluated at an ind-object $Z = \quotcolim_j z_j$, the second transformation is the canonical map $\colim_j \Map_C(x,z_j) \to \Map_C(x,\colim_j z_j)$. So informally, $f$ is compact if for every map $g\colon x' \to \colim_j z_j$ the composite $gf$ factors through some stage $z_j$, in a way which is natural in the filtered diagram $(z_j)_j$. By definition, $x$ is compact if and only if the second transformation is an equivalence; in that case every morphism $x \to x'$ and every morphism $x'' \to x$ is compact.

	\begin{remark}
		\label[remark]{rem:Compact_Morphisms_Ideal}
		Compact morphisms form a two-sided ideal: if $f\colon x \to x'$ is compact and $k\colon w \to x$ and $h\colon x' \to x''$ are arbitrary, then $hfk$ is compact. Indeed, given a factorization $\theta\colon \Map_C(x',\colim(-)) \to \Map_{\Ind(C)}(y(x),-)$ of $f^*$, the composite $y(k)^* \circ \theta \circ h^*$ is a factorization of $(hfk)^* = k^*f^*h^*$.
	\end{remark}

	\begin{lemma}
		\label[lemma]{lem:Compact_Morphisms_Yhat}
		Let $C$ be compactly assembled and let $f\colon x \to x'$ be a morphism in $C$.
		\begin{enumerate}
			\item The morphism $f$ is compact if and only if $y(f)\colon y(x) \to y(x')$ factors through the map $\hat{y}(x') \to y(x')$.
			\item The identity $\id_x$ is a compact morphism if and only if $x$ is a compact object.
		\end{enumerate}
	\end{lemma}
	\prooffor{lem:Compact_Morphisms_Yhat}
	\begin{proof}
		(1) Write $\varepsilon\colon \hat{y}(x') \to y(x')$, and recall that $\colim(\varepsilon)$ is the identity of $x'$. By the adjunction $\hat{y} \dashv \colim$, we have $\Map_C(x',\colim(-)) \simeq \Map_{\Ind(C)}(\hat{y}(x'),-)$, so by the Yoneda lemma, natural transformations $\theta\colon \Map_C(x',\colim(-)) \to \Map_{\Ind(C)}(y(x),-)$ correspond to maps $g\colon y(x) \to \hat{y}(x')$, via $\theta = g^*$. The composite of $g^*$ with the transformation $\Map_{\Ind(C)}(y(x),-) \to \Map_C(x,\colim(-))$ is precomposition with $\colim(g)\colon x \to x'$. Hence $f$ is compact if and only if there is a map $g\colon y(x) \to \hat{y}(x')$ with $\colim(g) \simeq f$. If $\varepsilon g \simeq y(f)$, then $\colim(g) \simeq \colim(\varepsilon g) \simeq f$. Conversely, if $\colim(g) \simeq f$, then $\varepsilon g$ and $y(f)$ are maps $y(x) \to y(x')$ with the same image under $\colim$; since $\colim\colon \Map_{\Ind(C)}(y(x),y(x')) \to \Map_C(x,x')$ is inverse to the equivalence induced by the fully faithful functor $y$, we get $\varepsilon g \simeq y(f)$.

		(2) By (1), $\id_x$ is compact if and only if $\varepsilon\colon \hat{y}(x) \to y(x)$ admits a section $g$. In that case, $g\varepsilon\colon \hat{y}(x) \to \hat{y}(x)$ corresponds under the adjunction to $\colim(g)\colon x \to x$, which is the identity since $\colim(\varepsilon g) = \id_x$ and $\colim(\varepsilon) = \id_x$. Hence $g\varepsilon \simeq \id$, so $\varepsilon$ is an equivalence and $x$ is compact by \Cref{lem:Compact_Iff_Yhat_Is_Y}. The converse is clear.
	\end{proof}

	\begin{example}
		\label[example]{ex:Compactly_Generated_Compactly_Assembled}
		If $C$ is accessible and compactly generated, then $C$ is compactly assembled. Indeed, if $C = \Ind(C_0)$ with $C_0$ small, then
		\[
			\hat{y} = \Ind(y_{C_0})\colon \Ind(C_0) \to \Ind(\Ind(C_0)),
		\]
		where $y_{C_0}\colon C_0 \hookrightarrow \Ind(C_0) = C$ is the Yoneda embedding. When $C$ is cocomplete, we may take $C_0=C^{\omega}$ by \href{https://animarium.org/tag/01BJ/}{[SHT, Proposition~23.3.7]}.
	\end{example}

	% Source: written from the literature (Johnstone--Joyal; Gierz et al.)
	The following example explains the terminology, and will be generalized in \Cref{chap:Sheaves}.

	\begin{example}[{Continuous posets, \cite{Johnstone_Joyal_Continuous_Categories, Gierz_Continuous_Lattices}}]
		\label[example]{ex:Continuous_Posets}
		Let $P$ be a partially ordered set admitting directed joins, regarded as a category. Then $\Ind(P)$ is the poset of \emph{ideals} of $P$, i.e.\ non-empty directed downward closed subsets $I \subseteq P$, ordered by inclusion, and $\colim\colon \Ind(P) \to P$ sends $I$ to its join $\bigvee I$. The functor $y$ sends $x$ to the principal ideal ${\downarrow}x$, which is the largest ideal with join $\leq x$. A left adjoint $\hat{y}$, if it exists, sends $x$ to the smallest ideal $I$ with $x \leq \bigvee I$.

		Say that $y$ is \emph{way below} $x$, written $y \ll x$, if $y \in I$ for every ideal $I$ with $x \leq \bigvee I$; taking $I = {\downarrow}x$ shows that $y \ll x$ implies $y \leq x$. Then $\hat{y}$ exists if and only if for every $x \in P$ the set ${\downarrow\!\!\downarrow}x = \{y \in P \mid y \ll x\}$ is directed with join $x$, in which case $\hat{y}(x) = {\downarrow\!\!\downarrow}x$; we prove this in \Cref{prop:Continuous_Poset_Characterization}. Posets with this property are classically called \emph{continuous}. An element $x$ is compact if and only if $x \ll x$, in accordance with \Cref{lem:Compact_Iff_Yhat_Is_Y}. For example:
		\begin{enumerate}
			\item In the interval $[0,1]$, we have $y \ll x$ if and only if $y < x$ or $y = 0$. So $[0,1]$ is continuous, and its only compact element is $0$.
			\item Let $X$ be a locally compact Hausdorff space and let $\mathrm{Open}(X)$ be the poset of open subsets. Then $U \ll V$ if and only if there exists a compact subset $K \subseteq X$ with $U \subseteq K \subseteq V$, and $\mathrm{Open}(X)$ is continuous. The compact elements are the compact open subsets; for $X = \R^n$ the only one is $\emptyset$.
			\item The power set of a set $S$ is continuous, and in fact \emph{algebraic}, i.e.\ compactly generated: the compact elements are the finite subsets, and every subset is the directed union of its finite subsets.
		\end{enumerate}
		Note how example (2) foreshadows the fact that $\Shv(\R^n;\Sp)$ is dualizable but has no nonzero compact objects. Continuous posets are discussed in more detail, with proofs, in \Cref{sec:Continuous_Posets}.
	\end{example}

	% Source: Talk 2 (Marc Hoyois, 20.10.2026)
	The following theorem gives an alternative characterization of compactly assembled categories, which shows that they are precisely the retracts of compactly generated ones.

	\begin{theorem}[{Lurie, \cite[Definition~21.1.2.1 and Theorem~21.1.2.10]{SAG}}]
		\label[theorem]{thm:Compactly_Assembled_Retract}
		Let $C$ be a category with filtered colimits. Then the following are equivalent:
		\begin{enumerate}
			\item $C$ is compactly assembled;
			\item There exists a small category $D$ such that $C$ is a retract of $\Ind(D)$ as a category with filtered colimits.
		\end{enumerate}
	\end{theorem}
	\prooffor{thm:Compactly_Assembled_Retract}
	\begin{proof}
		(1) $\implies$ (2): We have an adjunction between $\hat{y}\colon C \hookrightarrow \Ind(C)$ and $\colim\colon \Ind(C) \to C$. Since $\hat{y}$ is fully faithful, this exhibits $C$ as a retract of $\Ind(C)$, with both functors preserving filtered colimits. To ensure we may pick $D$ small, write $C = \Ind_{\kappa}(C^{\kappa})$, with $C^{\kappa}$ small (\Cref{thm:Accessible_Recollections}). Since $\Ind(C) = \bigcup_{\lambda}\Ind(C^{\lambda})$ (\Cref{constr:Ind_Large}) and $C^{\kappa}$ is small, we may find some regular cardinal $\lambda$ such that $\hat{y}(C^{\kappa}) \subseteq \Ind(C^{\lambda})$. Since the inclusion $\Ind(C^{\lambda}) \subseteq \Ind(C)$ preserves filtered colimits, it follows that in fact $\hat{y}(C) \subseteq \Ind(C^{\lambda})$. So it follows that $C$ is already a retract of $\Ind(C^{\lambda})$, giving the claim.

		(2) $\implies$ (1): Assume we have $u\colon C \hookrightarrow \Ind(D)$ with retraction $v\colon \Ind(D) \to C$, and assume $u$ and $v$ preserve filtered colimits. It follows that $C$ is accessible: we have
		\[
			C = \lim\big( \dots \xrightarrow{uv} \Ind(D) \xrightarrow{uv} \Ind(D)\big),
		\]
		and the subcategory $\mathrm{Acc} \subseteq \widehat{\Cat}$ of accessible categories and accessible functors is closed under limits in $\widehat{\Cat}$ (\Cref{thm:Limits_Of_Accessible_Categories}).

		To show the existence of $\hat{y}$, consider the following commutative diagram:
		\[
			\begin{tikzcd}
				\Ind(C) \rar["\colim_C"] \dar[shift right, "\Ind(u)"'] & C \dar[shift right, "u"'] \\
				\Ind(\Ind(D)) \rar["\colim_{\Ind(D)}"'] \uar[shift right, "\Ind(v)"'] & \Ind(D). \uar[shift right, "v"']
			\end{tikzcd}
		\]
		It follows that $\colim_C$ is a retract of $\colim_{\Ind(D)}$. So we are done in light of \Cref{lem:Retract_Left_Adjoint} in the appendix.
	\end{proof}

	\begin{definition}
		\label[definition]{def:Strongly_Continuous_Compactly_Assembled}
		Let $C$ and $D$ be compactly assembled. A functor $f\colon C \to D$ is called \emph{strongly continuous} if it preserves filtered colimits and commutes with $\hat{y}$, meaning that the mate transformation
		\[
			\hat{y}_D f \to \Ind(f) \hat{y}_C
		\]
		of the commutative square
		\[
			\begin{tikzcd}
				\Ind(C) \rar["\colim_C"] \dar["\Ind(f)"'] & C \dar["f"] \\
				\Ind(D) \rar["\colim_D"'] & D
			\end{tikzcd}
		\]
		is a natural isomorphism.
	\end{definition}

	\begin{remark}
		\label[remark]{rem:Strongly_Continuous_Right_Adjoint}
		Assume $f$ admits a right adjoint $g\colon D \to C$. Then the condition that $f$ is strongly continuous is equivalent to $g$ preserving filtered colimits, see \Cref{prop:Strongly_Continuous_Via_Right_Adjoint}.
	\end{remark}

	\section{Dualizable categories}
	\label[section]{sec:Dualizable_Categories}

	% Source: Talk 2 (Marc Hoyois, 20.10.2026); Mathew, Lecture 2 (MPIM Bonn, 09.09.2024)
	We now specialize to presentable stable categories. Recall from \Cref{sec:PrL_Tensor_Products} the tensor product on $\PrL$, and that compactly generated categories are dualizable in $\PrL_{\st}$ by \Cref{cor:Compactly_Generated_Dualizable}.

	\begin{definition}
		\label[definition]{def:Dualizable_Category}
		A \emph{dualizable category} is a dualizable object of $\PrL_{\st}$. We denote by $\Pr^{\dual} \subseteq \PrL_{\st}$ the subcategory whose objects are the dualizable categories and whose morphisms are the internal left adjoints, i.e.\ those functors in $\PrL_{\st}$ whose right adjoint also lies in $\PrL_{\st}$, i.e.\ preserves colimits.
	\end{definition}

	Explicitly, a presentable stable category $C$ is dualizable if there exist a presentable stable category $C^{\vee}$ and colimit-preserving functors
	\[
		\ev\colon C \otimes C^{\vee} \to \Sp, \qquad \coev\colon \Sp \to C^{\vee} \otimes C,
	\]
	satisfying the triangle identities: the composites
	\[
		C \xrightarrow{C \otimes \coev} C \otimes C^{\vee} \otimes C \xrightarrow{\ev \otimes C} C \qquad\text{and}\qquad C^{\vee} \xrightarrow{\coev \otimes C^{\vee}} C^{\vee} \otimes C \otimes C^{\vee} \xrightarrow{C^{\vee} \otimes \ev} C^{\vee}
	\]
	are equivalent to the identity. Since $\PrL_{\st}$ is closed, the dual is then necessarily given by
	\[
		C^{\vee} \simeq \Fun^{\colim}(C,\Sp),
	\]
	the internal hom from $C$ to the unit, with $\ev$ the evaluation functor.

	\begin{example}
		\label[example]{ex:Dual_Of_Modules}
		Let $A$ be an $\bbE_1$-ring and let $\Mod_A$ denote the category of right $A$-modules, so that $\Mod_A \simeq \Ind(\Perf_A)$. Then $\Mod_A$ is dualizable with dual $\Mod_{A\catop}$, the category of left $A$-modules. The evaluation and coevaluation are given by
		\begin{alignat*}{2}
			\ev&\colon \Mod_A \otimes \Mod_{A\catop} \to \Sp, &\qquad (M,N) &\mapsto M \otimes_A N, \\
			\coev&\colon \Sp \to \Mod_{A\catop} \otimes \Mod_A \simeq \mathrm{BMod}_{A,A}, &\qquad \S &\mapsto A,
		\end{alignat*}
		where the coevaluation sends the sphere spectrum to $A$ regarded as an $A$-bimodule. This is a special case of \Cref{cor:Compactly_Generated_Dualizable}, using the equivalence $\Perf_A\catop \simeq \Perf_{A\catop}$, $P \mapsto \hom_A(P,A)$.
	\end{example}

	\begin{remark}
		\label[remark]{rem:Strongly_Continuous_Stable}
		Note that a right adjoint of a functor between stable categories automatically preserves finite colimits, so it preserves all colimits if and only if it preserves filtered colimits. So the morphisms in $\Pr^{\dual}$ are precisely the strongly continuous functors.
	\end{remark}

	\begin{remark}
		\label[remark]{rem:Compactly_Generated_In_Prdual}
		By \Cref{cor:Compactly_Generated_Dualizable}, every compactly generated category $C = \Ind(C_0)$ is dualizable. Moreover, every colimit-preserving functor preserving compact objects is strongly continuous, so a morphism in $\Pr^{\dual}$ by \Cref{rem:Strongly_Continuous_Stable}. So we get a fully faithful inclusion
		\[
			\PrL_{\omega,\st} \subseteq \Pr^{\dual}.
		\]
	\end{remark}

	\begin{remark}
		\label[remark]{rem:Prdual_Formal_Properties}
		Dualizable objects of a symmetric monoidal category are closed under tensor products, and the tensor product of two strongly continuous functors is again strongly continuous: since $- \otimes D$ is a functor of $(\infty,2)$-categories, it preserves adjunctions, so the right adjoint of $f \otimes D$ is $f^R \otimes D$, which again preserves colimits. Hence the symmetric monoidal structure on $\PrL_{\st}$ restricts to $\Pr^{\dual}$. Moreover, the inclusion $\Pr^{\dual} \hookrightarrow \PrL_{\st}$ preserves colimits (\Cref{prop:Prdual_Colimits}), and $\Pr^{\dual}$ admits all limits, although these are \emph{not} computed in $\PrL_{\st}$; see \Cref{thm:General_Limits_Dualizable}.
	\end{remark}

	% Source: Talk 2 (Marc Hoyois, 20.10.2026)
	The following theorem is the main result of this chapter: it shows that dualizable categories are precisely the compactly assembled presentable stable categories.

	\begin{theorem}[{Lurie, \cite[Proposition~D.7.3.1]{SAG}}]
		\label[theorem]{thm:Lurie_Dualizable_Characterization}
		Let $C \in \PrL_{\st}$. Then the following are equivalent:
		\begin{enumerate}
			\item $C$ is dualizable;
			\item $C$ is compactly assembled;
			\item There exists a localization $D \to E$ in $\PrL_{\omega,\st}$ such that $C = \ker(D \to E)$;
			\item $C$ has the left lifting property with respect to any localization in $\PrL_{\st}$: for every localization $D \to E$ and any $C \to E$, there exists a lift $C \to D$:
			\[
				\begin{tikzcd}
					& D \dar \\
					C \rar \arrow[ur, dashed] & E.
				\end{tikzcd}
			\]
		\end{enumerate}
	\end{theorem}
	\prooffor{thm:Lurie_Dualizable_Characterization}
	\begin{proof}
		(2) $\implies$ (1): Every compactly assembled category is a retract of a compactly generated one. \Cref{rem:Compactly_Generated_In_Prdual} shows that compactly generated categories are dualizable, and since dualizable objects are closed under retracts we are done.

		(1) $\implies$ (4): Assume $C$ is dualizable, and consider a cofiber sequence
		\[
			K \hookrightarrow D \to E
		\]
		in $\PrL_{\st}$. Tensoring with $C^{\vee}$ preserves colimits, so gives another cofiber sequence
		\[
			C^{\vee} \otimes K \to C^{\vee} \otimes D \to C^{\vee} \otimes E.
		\]
		In particular, the functor $C^{\vee} \otimes D \to C^{\vee} \otimes E$ is again a localization, hence essentially surjective. It follows that the functor $\Sp \to C^{\vee} \otimes E$ classifying $C \to E$ lifts to a functor $\Sp \to C^{\vee} \otimes D$ classifying a lift $C \to D$.

		(4) $\implies$ (2): Let $C_0 \subseteq C$ be a small dense stable subcategory, and consider the adjunction
		\[
			y\colon C \rightleftarrows \Ind(C_0) \noloc \colim.
		\]
		Since $y$ is fully faithful, it follows that $\colim\colon \Ind(C_0) \to C$ is a localization, and so by (4) this admits a section in $\PrL_{\st}$. It follows that $C$ is a retract of $\Ind(C_0)$, hence it is compactly assembled.

		(3) $\implies$ (2): Consider a fiber sequence
		\[
			C \hookrightarrow D \to E.
		\]
		Since $D \to E$ preserves compact objects, its right adjoint preserves colimits. By \Cref{lem:Strong_Continuity_Localization_Sequence}, the kernel inclusion $C \to D$ also has a colimit-preserving right adjoint $D \to C$. This right adjoint is a retraction, so $C$ is a retract of $D$ in $\PrL_{\st}$.

		(2) $\implies$ (3): As before, we may find some $\lambda$ such that $\colim\colon \Ind(C^{\lambda}) \to C$ admits a left adjoint $C \to \Ind(C^{\lambda})$. Then consider the cofiber sequence
		\[
			C \xhookrightarrow{\hat{y}} \Ind(C^{\lambda}) \to \cofib(\hat{y}) \simeq \fib(\colim).
		\]
		The localization functor $\Ind(C^{\lambda}) \to \cofib(\hat{y})$ admits a right adjoint in $\PrL_{\st}$, and so it preserves compact objects. It in particular follows that $\cofib(\hat{y})$ is compactly generated, i.e.\ $\cofib(\hat{y}) \simeq \Ind(\Calk^{\cont}_{\lambda}(C))$ where $\Calk^{\cont}_{\lambda}(C) := \cofib(\hat{y})^{\omega}$ is the continuous $\lambda$-Calkin category. So this gives (3).
	\end{proof}

	% Source: Mathew, Lecture 2 (MPIM Bonn, 09.09.2024)
	There are several further useful characterizations. Recall that a presentable category $C$ is \emph{$\kappa$-compactly generated} if $\Ind_{\kappa}(C^{\kappa}) \simeq C$.

	\begin{proposition}
		\label[proposition]{prop:Dualizable_Further_Characterizations}
		Let $C \in \PrL_{\st}$. Then the following are equivalent:
		\begin{enumerate}
			\item $C$ is dualizable;
			\item The functor $\colim\colon \Ind(C) \to C$ preserves small products;
			\item For every regular cardinal $\kappa$ such that $C$ is $\kappa$-compactly generated, the functor $\colim\colon \Ind(C^{\kappa}) \to C$ admits a left adjoint;
			\item There exist a regular cardinal $\kappa$ such that $C$ is $\kappa$-compactly generated and the functor $\colim\colon \Ind(C^{\kappa}) \to C$ admits a left adjoint;
			\item There exists a fully faithful strongly continuous functor $C \hookrightarrow D$ into a compactly generated category $D$.
		\end{enumerate}
	\end{proposition}

	Using the description of products in $\Ind(C)$ (\Cref{lem:Products_In_Ind}), condition (2) can be reformulated as Grothendieck's axiom (AB6): filtered colimits distribute over products in $C$, see \Cref{prop:Dualizable_AB6}.

	\prooffor{prop:Dualizable_Further_Characterizations}
	\begin{proof}
		(1) $\implies$ (2): By \Cref{thm:Lurie_Dualizable_Characterization}, $\colim$ admits a left adjoint $\hat{y}$, hence preserves all limits.

		(2) $\implies$ (3): Let $\kappa$ be such that $C$ is $\kappa$-compactly generated, let $j\colon C^{\kappa} \hookrightarrow C$ be the inclusion and $F = \Ind(j)\colon \Ind(C^{\kappa}) \hookrightarrow \Ind(C)$. The functor $F$ is fully faithful. Its right adjoint is restriction of presheaves along $j$: for $X \in \Ind(C)$, the presheaf $d \mapsto \Map_{\Ind(C)}(y(d),X)$ on $C^{\kappa}$ preserves finite limits, so belongs to $\Ind(C^{\kappa})$. Here we use the description of the Ind-completion of a category with finite colimits as finite-limit-preserving presheaves. This constructs $F^R$ without a presentability assumption on the large category $\Ind(C)$. Since representables are compact, this restriction functor preserves filtered colimits. We claim that
		\[
			\colim_C \simeq \colim_{\kappa} \circ F^R,
		\]
		where $\colim_{\kappa}\colon \Ind(C^{\kappa}) \to C$ denotes the colimit functor. Both sides preserve filtered colimits, so it suffices to compare them on objects $y(c)$ with $c \in C$, naturally in $c$. For $d \in C^{\kappa}$ we have $\Map(y(d),F^Ry(c)) \simeq \Map_C(d,c)$, so $F^Ry(c)$ is the ind-object $\quotcolim_{(d \to c) \in (C^{\kappa})_{/c}} d$; note that $(C^{\kappa})_{/c}$ is $\kappa$-filtered with colimit $c$, as $C \simeq \Ind_{\kappa}(C^{\kappa})$. Hence $\colim_{\kappa}F^Ry(c) \simeq c \simeq \colim_C y(c)$, naturally in $c$.

		Now let $(X_s)_{s \in S}$ be a family of objects of $\Ind(C^{\kappa})$. Since $F^R$ preserves products and $F^RF \simeq \id$, we have $\prod_s X_s \simeq F^R(\prod_s FX_s)$, and hence
		\[
			\colim_{\kappa}\Big(\prod_s X_s\Big) \simeq \colim_C\Big(\prod_s FX_s\Big) \simeq \prod_s \colim_C(FX_s) \simeq \prod_s \colim_{\kappa}(X_s),
		\]
		using (2). So $\colim_{\kappa}$ preserves products. It is moreover an exact functor between stable categories, so it preserves finite limits, and hence all small limits. Since it also preserves filtered colimits, it admits a left adjoint by the adjoint functor theorem (\Cref{thm:Adjoint_Functor_Theorem}).

		(3) $\implies$ (4): Every presentable category is $\kappa$-compactly generated for some $\kappa$.

		(4) $\implies$ (5): Let $\hat{y}\colon C \to \Ind(C^{\kappa})$ be the left adjoint of $\colim_{\kappa}$, and let $F$ be as in the proof of (2) $\implies$ (3). Since $\colim_C \circ F$ preserves filtered colimits and restricts to the inclusion on $C^{\kappa}$, we have $\colim_C \circ F \simeq \colim_{\kappa}$, and hence
		\[
			\Map_{\Ind(C^{\kappa})}(X,F^Ry(c)) \simeq \Map_{\Ind(C)}(FX,y(c)) \simeq \Map_C(\colim_C FX,c) \simeq \Map_C(\colim_{\kappa}X,c).
		\]
		So $\colim_{\kappa}$ admits the right adjoint $F^R \circ y$, which is fully faithful since $\colim_{\kappa}F^Ry(c) \simeq c$ as shown above. Hence $\hat{y}$ is fully faithful by \Cref{lem:Adjoint_Triples}. Its right adjoint $\colim_{\kappa}$ preserves colimits, so $\hat{y}$ is strongly continuous, and $\Ind(C^{\kappa})$ is compactly generated.

		(5) $\implies$ (1): Let $i\colon C \hookrightarrow D$ be as in (5). Then $i^R i \simeq \id_C$ and $i^R$ preserves colimits, so $C$ is a retract of $D$ in $\PrL_{\st}$. Since $D$ is dualizable (\Cref{cor:Compactly_Generated_Dualizable}) and dualizable objects are closed under retracts, $C$ is dualizable.
	\end{proof}

	\begin{remark}
		\label[remark]{rem:Further_Characterizations_Later}
		Further characterizations of dualizable categories will be given later in these notes:
		\begin{itemize}
			\item $C$ is equivalent to $\ker(\Mod_A \to \Mod_B)$ for a homological epimorphism of $\bbE_1$-rings $A \to B$ (\Cref{chap:Examples});
			\item $C$ is generated under colimits by \emph{compactly exhaustible} objects, i.e.\ sequential colimits along compact morphisms (\Cref{prop:Dualizable_Compactly_Exhaustible});
			\item $C$ satisfies (AB6), i.e.\ filtered colimits distribute over products in $C$; this is an elementwise reformulation of condition (2) of \Cref{prop:Dualizable_Further_Characterizations} (\Cref{prop:Dualizable_AB6});
			\item $C$ is \emph{flat}, i.e.\ $C \otimes -$ preserves fully faithful functors in $\PrL_{\st}$ (\Cref{chap:Flatness}).
		\end{itemize}
	\end{remark}

	Localization sequences are the main tool for computing continuous K-theory, so we need to understand how dualizability behaves in them. One direction is easy: dualizability passes from the middle term of a strongly continuous localization sequence to the two outer terms. The converse fails; the gluing criterion and a counterexample are developed in \Cref{sec:Extensions_And_Gluing}, after the compact-morphism criteria needed for their proofs.

	% Source: Clausen, Lecture 1 (17.07.2023)
	\begin{proposition}[Localization sequences of dualizable categories]
		\label[proposition]{prop:Localization_Sequence_Dualizable}
		Let $K \xrightarrow{i} D \xrightarrow{p} E$ be a strongly continuous localization sequence in $\PrL_{\st}$ (\Cref{def:Localization_Sequence_PrL}). If $D$ is dualizable, then so are $K$ and $E$, and the sequence is a sequence in $\Pr^{\dual}$.
	\end{proposition}
	\prooffor{prop:Localization_Sequence_Dualizable}
	\begin{proof}
		Since $i$ and $p^R$ are fully faithful, we have $i^Ri \simeq \id_K$ and $pp^R \simeq \id_E$. By strong continuity, $i^R$ and $p^R$ preserve colimits, so $K$ and $E$ are retracts of $D$ in $\PrL_{\st}$, hence dualizable. The functors $i$ and $p$ are strongly continuous by assumption.
	\end{proof}

	\section{\texorpdfstring{$\omega_1$}{omega\_1}-presentability}
	\label[section]{sec:Omega_1_Presentability}

	% Source: Talk 2 (Marc Hoyois, 20.10.2026)
	A striking consequence of compact assembly is that it bounds the size of the generators: every compactly assembled category is generated by its $\omega_1$-compact objects, which moreover admit an explicit description. Recall the natural transformation $\hat{y} \to y$ from \Cref{sec:Compactly_Assembled}.

	\begin{theorem}[{Efimov, \cite[Corollary~2.13 and Proposition~2.16]{Efimov_Large_Categories}}]
		\label[theorem]{thm:Compactly_Assembled_Omega_1}
		Let $C$ be compactly assembled.
		\begin{enumerate}
			\item $C$ is $\omega_1$-compactly generated.
			\item An object $x$ of $C$ is $\omega_1$-compact if and only if it can be written as
			\[
				x = \colim\big( \dots \to x_n \to x_{n+1} \to \dots\big)
			\]
			where each map $y(x_n) \to y(x_{n+1})$ factors as
			\[
				y(x_n) \to \hat{y}(x_{n+1}) \to y(x_{n+1}).
			\]
		\end{enumerate}
	\end{theorem}
	% Source: proof revised using countable directed subposets; cf. Efimov, Corollary 2.13 and Proposition 2.16.
	\prooffor{thm:Compactly_Assembled_Omega_1}[main-proof]
	\begin{proof}
		We first prove (2). For $\impliedby$, write $x = \colim_n x_n$. Then $\hat{y}(x) = \colim_n \hat{y}(x_n) = \colim_n y(x_n)$, where the last equality holds by considering
		\[
			\begin{tikzcd}
				\dots \rar & y(x_n) \rar \arrow[dr] & y(x_{n+1}) \rar & \dots \\
				\dots \rar & \hat{y}(x_n) \rar \uar & \hat{y}(x_{n+1}) \rar \uar & \dots
			\end{tikzcd}
		\]
		filled with the given maps $y(x_n) \to \hat{y}(x_{n+1})$. The top triangle commutes by construction. The fact that we can also make the bottom triangle commute in a way that makes the total square agree with the canonical one follows from the adjunction $\hat{y} \dashv \colim$: after applying the adjunction, this square looks like
		\[
			\begin{tikzcd}
				x_n \rar \dar[equals] & x_{n+1} \dar[equals] \\
				x_n \rar & x_{n+1}.
			\end{tikzcd}
		\]
		Now, each $y(x_n)$ is compact in $\Ind(C)$, and in particular $\omega_1$-compact, so also the colimit $\hat{y}(x) = \colim_n y(x_n)$ is $\omega_1$-compact. But since $\hat{y}$ is fully faithful and preserves filtered colimits, this implies that $x$ is already $\omega_1$-compact in $C$.

		For $\implies$, consider $x \in C$.

		\emph{Observation:} Write $\hat{y}(x) = \colim_{i \in I} y(x_i)$ for a filtered category $I$. By applying $\colim\colon \Ind(C) \to C$, we get $x = \colim_{i \in I} x_i$. Since $\hat{y}$ preserves colimits, we get $\hat{y}(x) = \colim_{i \in I} \hat{y}(x_i)$. Since each $y(x_i)$ is compact, it follows that for each $i \in I$ there exists a morphism $i \to j$ in $I$ such that the map $y(x_i) \to y(x_j)$ factors through $\hat{y}(x_j) \to y(x_j)$.

		Now, assume $x$ is $\omega_1$-compact. Then also $\hat{y}(x)$ is $\omega_1$-compact in $\Ind(C)$: by adjunction this follows because its right adjoint $\colim$ preserves $\omega_1$-filtered colimits. Thus we may write $\hat{y}(x) = \colim_n y(x_n)$ by \Cref{prop:Omega_1_Compact_In_Ind} (applied to $\Ind(C^{\lambda}) \subseteq \Ind(C)$ for $\lambda$ large enough that $\hat{y}(x) \in \Ind(C^{\lambda})$, see \Cref{constr:Ind_Large}). By the previous observation we see that for each $n$ there exists some $k \geq n$ such that the map $y(x_n) \to y(x_k)$ factors through $\hat{y}(x_k)$. By reindexing, we may assume $k = n+1$, giving the claim.

		Finally, we prove (1). As in the proof of \Cref{thm:Compactly_Assembled_Retract}, choose a regular cardinal $\lambda$ such that $\hat{y}(C) \subseteq \Ind(C^{\lambda})$, with $C^{\lambda}$ essentially small. For $x \in C$, choose a presentation $\hat{y}(x) = \colim_{i\in I}y(x_i)$ with $I$ a small filtered category and $x_i \in C^{\lambda}$. Since every filtered category admits a cofinal functor from a filtered poset (\Cref{thm:Filtered_Posets_Cofinal}), we may assume that $I$ is a filtered poset. By the observation above, choose for each $i$ an element $s(i) \geq i$ such that $y(x_i) \to y(x_{s(i)})$ factors through $\hat{y}(x_{s(i)})$.

		Let $\Sigma_s$ be the poset of nonempty countable directed subsets $Q \subseteq I$ closed under $s$, ordered by inclusion. Every countable subset of $I$ is contained in such a $Q$: start with that subset and one element of $I$, and repeatedly add the chosen successors and upper bounds for every pair. The union after countably many steps has the required properties. Applying this to countable unions shows that $\Sigma_s$ is $\omega_1$-filtered and cofinal in the poset of all countable directed subsets. By \Cref{prop:Colimits_Over_Countable_Subposets},
		\[
			x \simeq \colim_{Q\in\Sigma_s}\colim_{i\in Q}x_i.
		\]
		For each $Q$, enumerate its elements and choose a cofinal sequence $i_0\leq i_1\leq\cdots$ in $Q$ with $i_{n+1}\geq s(i_n)$, by successively taking upper bounds with the next enumerated element. The transition $y(x_{i_n})\to y(x_{i_{n+1}})$ factors through $\hat{y}(x_{i_{n+1}})$, using the chosen factorization at $s(i_n)$ and naturality of $\hat{y}\to y$. Thus $\colim_{i\in Q}x_i$ is $\omega_1$-compact by (2).

		This writes every object as an $\omega_1$-filtered colimit of $\omega_1$-compact objects. These objects may be chosen from the essentially small collection of colimits of sequences in $C^{\lambda}$. Every $\omega_1$-compact object is a retract of one of them, by factoring its identity through this presentation. Since $C$ is locally small, the collection of these retracts is essentially small as well. Consequently $C$ is $\omega_1$-compactly generated.
	\end{proof}

	% Source: Mathew, Lecture 2 (MPIM Bonn, 09.09.2024)
	For dualizable categories, there is an alternative and perhaps more intuitive argument for part (1) of \Cref{thm:Compactly_Assembled_Omega_1}, using the presentation as a kernel.

	\prooffor{thm:Compactly_Assembled_Omega_1}[kernel-proof]
	\begin{proof}[Alternative proof of \Cref{thm:Compactly_Assembled_Omega_1}(1) for $C$ dualizable]
		By \Cref{thm:Lurie_Dualizable_Characterization}, we may write $C = \ker(L\colon D \to E)$ for a localization $L$ between compactly generated categories which preserves compact objects. Since $C \subseteq D$ is closed under colimits, every object of $C$ which is $\omega_1$-compact in $D$ is $\omega_1$-compact in $C$.

		Let $x \in C$, and write $x = \colim_{p \in P} d_p$ in $D$ with $P$ a filtered partially ordered set and $d_p \in D^{\omega}$ (\Cref{thm:Filtered_Posets_Cofinal}). Then $\colim_p L(d_p) \simeq L(x) \simeq 0$, and each $L(d_p)$ is compact in $E$. Hence for every $p \in P$ there is some $p' \geq p$ such that the map $L(d_p) \to L(d_{p'})$ is nullhomotopic. Let $\Sigma' \subseteq \Sigma$ be the set of countable filtered subsets $Q \subseteq P$ such that for every $q \in Q$ there is some $q' \in Q$ with $q' \geq q$ and $L(d_q) \to L(d_{q'})$ nullhomotopic. Arguing as in the proof of \Cref{prop:Colimits_Over_Countable_Subposets}, but alternately adding upper bounds and such witnesses $q'$, one sees that every countable collection of elements of $\Sigma$ has an upper bound in $\Sigma'$. Hence $\Sigma'$ is $\omega_1$-filtered and cofinal in $\Sigma$, so by \Cref{prop:Colimits_Over_Countable_Subposets} and \Cref{lem:Cofinality_Posets} we get
		\[
			x \simeq \colim_{Q \in \Sigma'} x_Q, \qquad x_Q := \colim_{q \in Q} d_q.
		\]
		For $Q \in \Sigma'$ we claim that $L(x_Q) \simeq \colim_{q \in Q} L(d_q)$ vanishes, so that $x_Q \in C$. Since $E$ is compactly generated and its compact objects are closed under shifts, it suffices to show that $\pi_0\Map_E(c,L(x_Q)) = 0$ for every compact object $c \in E$. As $c$ is compact and $Q$ is filtered, $\pi_0\Map_E(c,L(x_Q)) \cong \colim_{q \in Q} \pi_0\Map_E(c,L(d_q))$. An element of this colimit is represented by some $f\colon c \to L(d_q)$, and choosing $q' \geq q$ in $Q$ with $L(d_q) \to L(d_{q'})$ nullhomotopic shows that its image in $\pi_0\Map_E(c,L(d_{q'}))$ vanishes. Hence the colimit is zero. Moreover, $x_Q$ is a countable colimit of compact objects, hence $\omega_1$-compact in $D$ (\Cref{cor:Kappa_Small_Colimits_Of_Kappa_Compacts}), and hence in $C$. So every object of $C$ is an $\omega_1$-filtered colimit of $\omega_1$-compact objects.
	\end{proof}

	% Source: Clausen, Lecture 2 (18.07.2023)
	\begin{remark}
		\label[remark]{rem:Omega_1_Compact_Compact_Morphisms}
		By \Cref{lem:Compact_Morphisms_Yhat}, the condition on the transition maps in \Cref{thm:Compactly_Assembled_Omega_1}(2) says precisely that the maps $x_n \to x_{n+1}$ are compact morphisms. So the $\omega_1$-compact objects of a compactly assembled category are precisely the sequential colimits along compact morphisms, and every object is an $\omega_1$-filtered colimit of such. In \Cref{chap:Compact_Morphisms_AB6} we will see that, conversely, a presentable stable category is dualizable if and only if it is generated under colimits by such sequential colimits. While compact \emph{objects} are the right tool to study compactly generated categories, compact \emph{morphisms} are the right tool to study dualizable ones.
	\end{remark}

	% Source: Mathew, Lecture 3 (MPIM Bonn, 10.09.2024); Talk 3 (Max Glaser, 27.10.2026); proof drafted with generative AI (Claude), cf. Efimov's characterization of strongly continuous functors via compact maps
	As a first illustration of this principle, strongly continuous functors between dualizable categories can be detected using compact morphisms, just like functors between compactly generated categories are strongly continuous if and only if they preserve compact objects.

	\begin{proposition}
		\label[proposition]{prop:Strongly_Continuous_Compact_Morphisms}
		Let $F\colon C \to D$ be a colimit-preserving functor between dualizable categories. Then $F$ is strongly continuous if and only if $F$ preserves compact morphisms.
	\end{proposition}
	\prooffor{prop:Strongly_Continuous_Compact_Morphisms}
	\begin{proof}
		Write $\varepsilon\colon \hat{y} \to y$ for the natural transformation from \Cref{sec:Compactly_Assembled}, and let $\beta\colon \hat{y}_DF \to \Ind(F)\hat{y}_C$ be the mate from \Cref{def:Strongly_Continuous_Compactly_Assembled}. By \Cref{prop:Strongly_Continuous_Via_Right_Adjoint}, $F$ is strongly continuous if and only if $\beta$ is an equivalence. We will repeatedly use the following observation: a map $\hat{y}_D(z) \to Y$ in $\Ind(D)$ is determined, via the adjunction $\hat{y}_D \dashv \colim$, by its image under $\colim$ (composed with the identification $z \simeq \colim\hat{y}_D(z)$). Applying $\colim$ to $\varepsilon_{Fx}$, to $\Ind(F)(\varepsilon_x) \circ \beta_x$, and to $\beta_x$ yields the identity of $Fx$ in each case (up to the canonical identifications), by construction of $\varepsilon$ and of the mate $\beta$. Hence
		\begin{equation}
			\label{eq:Mate_Compatible_Epsilon}
			\varepsilon_{Fx} \simeq \Ind(F)(\varepsilon_x) \circ \beta_x\colon \hat{y}_D(Fx) \to y(Fx).
		\end{equation}

		Suppose first that $F$ is strongly continuous, and let $f\colon x \to x'$ be a compact morphism in $C$. By \Cref{lem:Compact_Morphisms_Yhat}, there is $g\colon y(x) \to \hat{y}_C(x')$ with $\varepsilon_{x'}g \simeq y(f)$. Then $g' := \beta_{x'}^{-1} \circ \Ind(F)(g)\colon y(Fx) \to \hat{y}_D(Fx')$ satisfies
		\[
			\varepsilon_{Fx'} \circ g' \simeq \Ind(F)(\varepsilon_{x'}) \circ \Ind(F)(g) \simeq \Ind(F)(y(f)) = y(Ff)
		\]
		by \eqref{eq:Mate_Compatible_Epsilon}, so $Ff$ is compact by \Cref{lem:Compact_Morphisms_Yhat}.

		Conversely, suppose that $F$ preserves compact morphisms. Both functors $\hat{y}_DF$ and $\Ind(F)\hat{y}_C$ preserve colimits: $\hat{y}_C$ and $\hat{y}_D$ are left adjoints, and $\Ind(F)$ is an exact functor between stable categories which preserves filtered colimits. Since $C$ is generated under colimits by its $\omega_1$-compact objects (\Cref{thm:Compactly_Assembled_Omega_1}(1)), it suffices to show that $\beta_x$ is an equivalence for $x$ $\omega_1$-compact. By \Cref{thm:Compactly_Assembled_Omega_1}(2) and \Cref{lem:Compact_Morphisms_Yhat}, we may write $x \simeq \colim(x_0 \to x_1 \to \cdots)$ with compact transition maps, and the proof of ``$\impliedby$'' in \Cref{thm:Compactly_Assembled_Omega_1}(2) shows that the canonical map
		\[
			\hat{y}_C(x) \simeq \colim_n \hat{y}_C(x_n) \xrightarrow{\colim_n \varepsilon_{x_n}} \colim_n y(x_n)
		\]
		is an equivalence in $\Ind(C)$. Since $F$ preserves colimits and compact morphisms, $Fx \simeq \colim_n Fx_n$ with compact transition maps, so likewise $\hat{y}_D(Fx) \to \colim_n y(Fx_n)$ is an equivalence. Now consider the composite
		\[
			\hat{y}_D(Fx) \xrightarrow{\beta_x} \Ind(F)(\hat{y}_C(x)) \xrightarrow{\simeq} \Ind(F)\big(\colim_n y(x_n)\big) \simeq \colim_n y(Fx_n).
		\]
		Applying $\colim$ to it gives the identity of $Fx$ (up to the canonical identifications), so by the observation above it agrees with the canonical equivalence $\hat{y}_D(Fx) \to \colim_n y(Fx_n)$. Hence $\beta_x$ is an equivalence.
	\end{proof}


		\chapter{Compact morphisms and the axiom (AB6)}
	\label[chapter]{chap:Compact_Morphisms_AB6}

	% Source: Talk 4 (Jiapeng Liu, 03.11.2026)
	In compactly generated categories, compact objects control everything: they generate the category, and maps out of them into filtered colimits factor through a finite stage. In dualizable categories there may be no nonzero compact objects at all, and their role is taken over by \emph{compact morphisms}, which we introduced in \Cref{sec:Compactly_Assembled}. We develop their factorization properties and use them to characterize dualizability.

	The argument proceeds from individual maps to whole categories: the lifting criteria lead to interpolation, compact exhaustions give generators, and (AB6) packages the result as a commutation property of products and filtered colimits. The final section uses these criteria to prove automatic dualizability.

	Throughout, for a category $C$ with filtered colimits we write $\colim\colon \Ind(C) \to C$ for the colimit functor and $\varepsilon\colon \hat{y} \to y$ for the natural transformation of \Cref{sec:Compactly_Assembled}, whenever $C$ is compactly assembled. (In the talk, as in \cite{Krause_Nikolaus_Putzstuck_Sheaves}, the colimit functor was denoted $k$.) The talk mostly worked with dualizable categories. Several of the results below hold more generally for compactly assembled categories, and we state them in this generality whenever the proof is the same; see \cite[Sections~2.2--2.4]{Krause_Nikolaus_Putzstuck_Sheaves} for a systematic treatment of the non-stable case.

	\section{Weakly compact morphisms}
	\label[section]{sec:Weakly_Compact_Morphisms}

	Our definition of compact morphisms (\Cref{def:Compact_Morphism}) asks for a factorization which is natural in the filtered diagram. In the talk, as in \cite[Definition~2.37]{Efimov_Large_Categories}, compact morphisms were instead defined by asking for such factorizations one map at a time. We now show that the two notions agree in compactly assembled categories, so that both definitions may be used interchangeably in these notes.

	\begin{definition}[{\cite[Definition~2.2.3]{Krause_Nikolaus_Putzstuck_Sheaves}}]
		\label[definition]{def:Weakly_Compact_Morphism}
		Let $C$ be a category with filtered colimits. A morphism $f\colon x \to x'$ in $C$ is called \emph{weakly compact} if for every filtered diagram $(z_i)_{i \in I}$ in $C$ and every morphism $g\colon x' \to \colim_i z_i$, the composite $gf\colon x \to \colim_i z_i$ factors through the canonical map $z_i \to \colim_i z_i$ for some $i \in I$.
	\end{definition}

	\begin{remark}
		\label[remark]{rem:Weakly_Compact_Lifting}
		Since $y(x)$ is compact in $\Ind(C)$, maps $y(x) \to \quotcolim_i z_i$ are the same as maps $x \to z_i$ for some $i$, up to the usual identifications. Hence $f$ is weakly compact if and only if for every $Z \in \Ind(C)$, every diagram of solid arrows
		\[
			\begin{tikzcd}
				y(x) \rar[dashed] \dar["y(f)"'] & Z \dar["\eta_Z"] \\
				y(x') \rar & y(\colim Z)
			\end{tikzcd}
		\]
		admits a dashed arrow making the square commute. Here $\eta_Z$ is the unit of the adjunction $\colim \dashv y$ of \Cref{constr:Ind_Large}, and a map $y(x') \to y(\colim Z)$ is the same as a map $x' \to \colim Z$, since $y$ is fully faithful.

		A compact morphism in the sense of \Cref{def:Compact_Morphism} is weakly compact: the factorization of $f^*$ through $\Map_{\Ind(C)}(y(x),-)$ provides, in particular, a lift of each point. In the terminology of \cite[Definition~2.2.4]{Krause_Nikolaus_Putzstuck_Sheaves}, compact morphisms in our sense are called \emph{strongly compact}.
	\end{remark}

	\begin{example}
		\label[example]{ex:Weakly_Compact_Basic}
		\leavevmode
		\begin{enumerate}
			\item If $x$ is compact, then $\id_x$ is weakly compact. The converse holds for stable $C$, see \Cref{cor:Compact_Iff_Identity_Weakly_Compact} below.
			\item As in \Cref{rem:Compact_Morphisms_Ideal}, weakly compact morphisms form a two-sided ideal: if $f\colon x \to x'$ is weakly compact and $k\colon w \to x$, $h\colon x' \to x''$ are arbitrary, then for every $g\colon x'' \to \colim_i z_i$ the composite $(gh)f$ factors through some $z_i$, and hence so does $ghfk$.
			\item In particular, every morphism which factors through a compact object is weakly compact.
		\end{enumerate}
	\end{example}

	\begin{proposition}[{\cite[Proposition~2.39]{Efimov_Large_Categories}, \cite[Proposition~2.4.1]{Krause_Nikolaus_Putzstuck_Sheaves}}]
		\label[proposition]{prop:Compact_Morphism_Characterizations}
		Let $C$ be compactly assembled and let $f\colon x \to x'$ be a morphism in $C$. Then the following are equivalent:
		\begin{enumerate}
			\item $f$ is compact (in the sense of \Cref{def:Compact_Morphism});
			\item $f$ is weakly compact;
			\item $y(f)\colon y(x) \to y(x')$ factors through $\varepsilon_{x'}\colon \hat{y}(x') \to y(x')$;
			\item $\hat{y}(f)\colon \hat{y}(x) \to \hat{y}(x')$ factors through $\varepsilon_x\colon \hat{y}(x) \to y(x)$;
			\item $\hat{y}(f)$ is a weakly compact morphism in $\Ind(C)$.
		\end{enumerate}
	\end{proposition}
	\prooffor{prop:Compact_Morphism_Characterizations}
	\begin{proof}
		We use repeatedly that maps $\hat{y}(w) \to Z$ in $\Ind(C)$ correspond, under the adjunction $\hat{y} \dashv \colim$, to maps $w \to \colim Z$; in particular, two maps $\hat{y}(w) \to Z$ are homotopic as soon as their images under $\colim$ are, where we identify $\colim\hat{y}(w) \simeq w$. Similarly, two maps $y(w) \to y(w')$ are homotopic as soon as their images under $\colim$ are, since $y$ is fully faithful. Recall also that $\colim(\varepsilon_w) \simeq \id_w$.

		(1) $\iff$ (3) is \Cref{lem:Compact_Morphisms_Yhat}(1), and (1) $\implies$ (2) was observed in \Cref{rem:Weakly_Compact_Lifting}.

		(2) $\implies$ (3): Write $\hat{y}(x') \simeq \quotcolim_i z_i$. Then $\colim_i z_i \simeq \colim\hat{y}(x') \simeq x'$. Applying weak compactness to the identity $x' \to \colim_i z_i$, we find a map $h\colon x \to z_i$ for some $i$ such that the composite $x \to z_i \to \colim_i z_i \simeq x'$ is homotopic to $f$. Let $g\colon y(x) \to \hat{y}(x')$ be the composite of $y(h)$ with the canonical map $y(z_i) \to \quotcolim_i z_i$. Then $\colim(\varepsilon_{x'} g) \simeq \colim(g) \simeq f$, and hence $\varepsilon_{x'}g \simeq y(f)$.

		(3) $\implies$ (4): Let $g\colon y(x) \to \hat{y}(x')$ with $\varepsilon_{x'}g \simeq y(f)$. Then $g\varepsilon_x$ and $\hat{y}(f)$ are maps $\hat{y}(x) \to \hat{y}(x')$ with $\colim(g\varepsilon_x) \simeq \colim(g) \simeq f \simeq \colim(\hat{y}(f))$, so $\hat{y}(f) \simeq g\varepsilon_x$.

		(4) $\implies$ (5): The object $y(x)$ is compact in $\Ind(C)$, so this follows from \Cref{ex:Weakly_Compact_Basic}(3).

		(5) $\implies$ (2): Let $g\colon x' \to \colim_i z_i$ be a map in $C$, and let $Z = \quotcolim_i z_i = \colim_i y(z_i)$. By adjunction, $g$ corresponds to a map $\tilde{g}\colon \hat{y}(x') \to Z$ with $\colim(\tilde{g}) \simeq g$. By (5), the composite $\tilde{g}\hat{y}(f)\colon \hat{y}(x) \to \colim_i y(z_i)$ factors through some $y(z_i)$, say as $\hat{y}(x) \xrightarrow{u} y(z_i) \to Z$. Applying $\colim$, we find that $gf \simeq \colim(\tilde{g}\hat{y}(f))$ factors as $x \xrightarrow{\colim(u)} z_i \to \colim_i z_i$.
	\end{proof}

	\begin{remark}
		\label[remark]{rem:Weak_Compactness_In_Small_Ind_Categories}
		In (5), weak compactness in $\Ind(C)$ may be replaced by weak compactness in $\Ind(C^{\kappa})$, for any $\kappa$ such that $\hat{y}$ takes values in $\Ind(C^{\kappa})$, by the same argument. In particular, if $C$ is dualizable we may take $\kappa = \omega_1$ (\Cref{thm:Compactly_Assembled_Omega_1}), and then $\Ind(C^{\omega_1})$ is compactly generated.
	\end{remark}

	From now on we simply speak of \emph{compact morphisms} in compactly assembled categories, and use the characterizations of \Cref{prop:Compact_Morphism_Characterizations} freely. The following observation is fundamental; it says that compact morphisms can be ``interpolated''.

	\begin{corollary}[{\cite[Corollary~2.40]{Efimov_Large_Categories}, \cite[Proposition~2.4.4]{Krause_Nikolaus_Putzstuck_Sheaves}}]
		\label[corollary]{cor:Compact_Morphism_Factorization}
		Let $C$ be compactly assembled. Then every compact morphism $f\colon x \to x'$ factors as a composite $x \xrightarrow{f_1} w \xrightarrow{f_2} x'$ of two compact morphisms.
	\end{corollary}
	\prooffor{cor:Compact_Morphism_Factorization}
	\begin{proof}
		By \Cref{prop:Compact_Morphism_Characterizations}, there is $g\colon y(x) \to \hat{y}(x')$ with $\varepsilon_{x'}g \simeq y(f)$. Write $\hat{y}(x') \simeq \quotcolim_i w_i$, so that $x' \simeq \colim_i w_i$. Since $\hat{y}$ preserves colimits, also $\hat{y}(x') \simeq \colim_i \hat{y}(w_i)$ in $\Ind(C)$. As $y(x)$ is compact in $\Ind(C)$, the map $g$ factors as
		\[
			y(x) \xrightarrow{h} \hat{y}(w_i) \xrightarrow{\hat{y}(f_2)} \hat{y}(x')
		\]
		for some $i$, where $f_2\colon w_i \to x'$ is the canonical map. Put $w := w_i$ and $f_1 := \colim(h)\colon x \to w$. Then $\varepsilon_w h$ and $y(f_1)$ have the same image under $\colim$, so $\varepsilon_w h \simeq y(f_1)$, and $f_1$ is compact by \Cref{prop:Compact_Morphism_Characterizations}(3). The map $y(w) \to \quotcolim_j w_j = \hat{y}(x')$ is a lift of $y(f_2)$ along $\varepsilon_{x'}$, since both become $f_2$ after applying $\colim$; so $f_2$ is compact as well. Finally, $f \simeq \colim(g) \simeq \colim(\hat{y}(f_2))\colim(h) \simeq f_2f_1$.
	\end{proof}

	% Source: written from the literature (Krause--Nikolaus--P\"utzst\"uck, Proposition 2.4.4 and Proposition 2.4.6)
	By iterating \Cref{cor:Compact_Morphism_Factorization}, a single compact morphism can be ``interpolated'' by a dense family of compact morphisms. Since this is also needed for other classes of morphisms (for instance trace-class morphisms, see \Cref{prop:Nuclear_Rigid}), we formulate it abstractly. We write $\Q_{[0,1]} = \Q \cap [0,1]$ and $\Q_{\geq 0} = \Q \cap [0,\infty)$, regarded as partially ordered sets.

	\begin{lemma}
		\label[lemma]{lem:Interpolation}
		Let $D$ be a category, let $S$ be a class of morphisms of $D$ which is closed under homotopy and composition, and let $P$ be a class of objects of $D$. Assume that every morphism in $S$ is homotopic to a composite $x \to p \to x'$ of two morphisms in $S$ with $p \in P$.
		\begin{enumerate}
			\item For every morphism $s\colon x \to x'$ in $S$ there is a functor $F\colon \Q_{[0,1]} \to D$ with $F(0) = x$, $F(1) = x'$ and $F(0 \leq 1) \simeq s$, such that $F(a \leq b) \in S$ for all $a < b$ and $F(a) \in P$ for all $0 < a < 1$.
			\item For every sequence $x_0 \to x_1 \to x_2 \to \cdots$ of morphisms in $S$ there is a functor $F\colon \Q_{\geq 0} \to D$ whose restriction to $\N \subseteq \Q_{\geq 0}$ is equivalent to the given sequence, such that $F(a \leq b) \in S$ for all $a < b$ and $F(a) \in P$ for all $a \notin \N$. If $D$ admits sequential colimits, then $\colim_{\Q_{\geq 0}} F \simeq \colim_n x_n$.
		\end{enumerate}
	\end{lemma}
	\prooffor{lem:Interpolation}
	\begin{proof}
		(1) The set $T$ of dyadic rationals in $[0,1]$ is a countable dense totally ordered set with a smallest and a largest element, so it is order-isomorphic to $\Q_{[0,1]}$ by Cantor's back-and-forth argument; we may therefore replace $\Q_{[0,1]}$ by $T$, with $0$ and $1$ preserved. Let $T_k = \{i/2^k \mid 0 \leq i \leq 2^k\}$, so that $T = \bigcup_k T_k$ and hence $N(T) = \bigcup_k N(T_k)$. Thus $\Fun(T,D) \simeq \lim_k \Fun(T_k,D)$, and on cores this is a limit of a tower of animae. So it suffices to construct functors $F_k\colon T_k \to D$, with $F_k(a \leq b) \in S$ for $a < b$ and $F_k(a) \in P$ for $0 < a < 1$, together with equivalences $F_{k+1}|_{T_k} \simeq F_k$. Indeed, a point of the limit of a tower of animae is specified by points at each stage and paths between their successive images; no separate higher compatibility choices are required. As in \Cref{rem:Sequential_Diagrams}, a functor out of a finite totally ordered set amounts to a composable sequence of morphisms, and two such functors are equivalent as soon as they agree on objects and their consecutive morphisms are homotopic. Let $F_0$ be given by $s$. Given $F_k$, factor each morphism $F_k(a \leq a')$ between consecutive elements $a < a'$ of $T_k$ as a composite of two morphisms in $S$ through an object of $P$, and let $F_{k+1}$ be the resulting functor on $T_{k+1}$. Then $F_{k+1}|_{T_k}$ and $F_k$ agree on objects, and their consecutive morphisms are homotopic, so $F_{k+1}|_{T_k} \simeq F_k$. Since $S$ is closed under composition, every $F_{k+1}(a \leq b)$ with $a < b$ is in $S$.

		(2) Apply the construction of (1) simultaneously to all morphisms $x_n \to x_{n+1}$: on the dyadic rationals in $[0,\infty)$, which are the union of the subsets $2^{-k}\N$, we refine the given sequence on $\N = 2^{0}\N$ by factoring consecutive morphisms as before. The set of dyadic rationals in $[0,\infty)$ is order-isomorphic to $\Q_{\geq 0}$ by an isomorphism fixing $\N$ (apply the back-and-forth argument to each interval $[n,n+1]$). Finally, $\N \subseteq \Q_{\geq 0}$ is cofinal, so the colimits agree by \Cref{lem:Cofinality_Posets}.
	\end{proof}

	\begin{proposition}[{\cite[Propositions~2.4.4 and~2.4.6]{Krause_Nikolaus_Putzstuck_Sheaves}}]
		\label[proposition]{prop:Compact_Interpolation}
		Let $C$ be compactly assembled.
		\begin{enumerate}
			\item Every compact morphism $f\colon x \to x'$ factors as $x \to w \to x'$ with both morphisms compact and $w$ an $\omega_1$-compact object.
			\item Every compact morphism $f\colon x \to x'$ extends to a functor $F\colon \Q_{[0,1]} \to C$ with $F(0 \leq 1) \simeq f$, such that $F(a \leq b)$ is compact for all $a < b$ and $F(a)$ is $\omega_1$-compact for $0 < a < 1$.
			\item An object of $C$ is $\omega_1$-compact if and only if it is the colimit of a functor $F\colon \Q_{\geq 0} \to C$, or equivalently of a functor $F\colon \Q \to C$, such that $F(a \leq b)$ is compact for all $a < b$. Moreover, $F$ can be chosen to take values in $\omega_1$-compact objects.
		\end{enumerate}
	\end{proposition}
	\prooffor{prop:Compact_Interpolation}
	\begin{proof}
		(1) By \Cref{cor:Compact_Morphism_Factorization}, applied twice, $f$ factors as $x \to x_0 \to x_0' \to x'$ with all three maps compact. By \Cref{thm:Compactly_Assembled_Omega_1}(1), $x_0'$ is a filtered colimit of $\omega_1$-compact objects, so the compact (hence weakly compact, by \Cref{prop:Compact_Morphism_Characterizations}) morphism $x_0 \to x_0'$ factors as $x_0 \to w \to x_0'$ with $w$ $\omega_1$-compact. Since compact morphisms form an ideal, $x \to x_0 \to w$ and $w \to x_0' \to x'$ are compact.

		(2) Apply \Cref{lem:Interpolation}(1) to the class $S$ of compact morphisms and the class $P$ of $\omega_1$-compact objects, using (1).

		(3) By \Cref{thm:Compactly_Assembled_Omega_1}(2), \Cref{rem:Omega_1_Compact_Compact_Morphisms} and \Cref{prop:Compact_Morphism_Characterizations}, an object $x$ is $\omega_1$-compact if and only if $x \simeq \colim_n x_n$ for a sequence of compact morphisms $x_n \to x_{n+1}$. By (1), each $x_n \to x_{n+1}$ factors as $x_n \to w_n \to x_{n+1}$ through an $\omega_1$-compact $w_n$, with both maps compact. The sequence $w_0 \to w_1 \to \cdots$ of composites $w_n \to x_{n+1} \to w_{n+1}$ again consists of compact morphisms between $\omega_1$-compact objects, and has colimit $x$ by cofinality. Applying \Cref{lem:Interpolation}(2) to it (with $P$ the $\omega_1$-compact objects) gives a functor $F\colon \Q_{\geq 0} \to C$ with the required properties and colimit $x$. Since $\Q_{> 0} \subseteq \Q_{\geq 0}$ is cofinal and $\Q_{> 0}$ is order-isomorphic to $\Q$ (both are countable dense totally ordered sets without endpoints), we may equivalently use $\Q$. Conversely, the colimit of such a functor is the colimit of its restriction to the cofinal subset $\N$, which is $\omega_1$-compact by \Cref{thm:Compactly_Assembled_Omega_1}(2).
	\end{proof}

	For stable categories there is a useful reformulation of compactness, which only involves ind-objects with vanishing colimit. This is the definition of compact morphisms used by Clausen and Mathew.

	\begin{proposition}
		\label[proposition]{prop:Compact_Morphisms_Stable}
		Let $C$ be a stable category with filtered colimits (e.g.\ $C \in \PrL_{\st}$), and let $f\colon x \to x'$ be a morphism in $C$. Then $f$ is weakly compact if and only if for every $Z \in \Ind(C)$ with $\colim Z \simeq 0$, the map
		\[
			f^*\colon \pi_0\Map_{\Ind(C)}(y(x'),Z) \to \pi_0\Map_{\Ind(C)}(y(x),Z)
		\]
		is zero. Explicitly, this means: for every filtered diagram $(z_i)_{i \in I}$ with $\colim_i z_i \simeq 0$ and every map $g\colon x' \to z_i$, there is $i' \geq i$ such that the composite $x \xrightarrow{f} x' \xrightarrow{g} z_i \to z_{i'}$ is nullhomotopic.
	\end{proposition}
	\prooffor{prop:Compact_Morphisms_Stable}
	\begin{proof}
		Note that $\Ind(C)$ is stable, and that filtered colimits in $C$ as well as the colimit functor $\Ind(C) \to C$ are exact.

		Suppose that $f$ is weakly compact, and let $Z = \quotcolim_{i \in I} z_i$ with $\colim_i z_i \simeq 0$; by \Cref{thm:Filtered_Posets_Cofinal} we may assume that $I$ is a filtered partially ordered set. A map $y(x') \to Z$ is represented by some $g\colon x' \to z_{i_0}$. For $i \geq i_0$ write $g_i\colon x' \to z_i$ for the composite with the transition map, and $w_i := \fib(g_i)$. As $\{i \geq i_0\} \subseteq I$ is cofinal and filtered colimits are exact, the canonical map
		\[
			\colim_{i \geq i_0} w_i \to \fib\big(x' \to \colim_{i \geq i_0} z_i\big) \simeq x'
		\]
		is an equivalence. Applying weak compactness of $f$ to the inverse of this equivalence, we find $i \geq i_0$ and a map $b\colon x \to w_i$ such that $f$ is homotopic to the composite of $b$ with $w_i \to x'$. Hence $g_if$ factors through $g_i \circ (w_i \to x') \simeq 0$, so $g_if \simeq 0$, and the image of our map under $f^*$ vanishes.

		Conversely, suppose the condition holds, and let $g\colon x' \to \colim_i z_i$. Put $Z = \quotcolim_i z_i$, let $\eta_Z\colon Z \to y(\colim Z)$ be the unit of $\colim \dashv y$, and let $W = \cofib(\eta_Z)$. Then $\colim W \simeq \cofib(\colim(\eta_Z)) \simeq 0$, since $\colim(\eta_Z)$ is an equivalence. In the exact sequence of abelian groups
		\[
			\pi_0\Map(y(x),Z) \xrightarrow{(\eta_Z)_*} \pi_0\Map(y(x),y(\colim Z)) \to \pi_0\Map(y(x),W),
		\]
		the class of $y(gf) = y(g)y(f)$ maps to the image under $f^*$ of the class of the composite $y(x') \xrightarrow{y(g)} y(\colim Z) \to W$, which vanishes by assumption. Hence $y(gf) \simeq \eta_Z\psi$ for some $\psi\colon y(x) \to Z$. Since $y(x)$ is compact, $\psi$ is represented by a map $\psi_0\colon x \to z_i$ for some $i$, and $\eta_Z\psi$ is then $y$ of the composite $x \xrightarrow{\psi_0} z_i \to \colim_i z_i$. As $y$ is fully faithful, $gf$ is homotopic to this composite.
	\end{proof}

	\begin{remark}
		\label[remark]{rem:Compact_Morphisms_Other_Definitions}
		In the literature one also finds weaker, ``elementwise'' definitions of compact morphisms. For instance, Clausen and Mathew call a morphism $f\colon x \to x'$ in a presentable stable category compact if for every filtered diagram $(z_j)_j$ with $\colim_j z_j \simeq 0$, every map $x' \to z_j$ becomes nullhomotopic after precomposition with $f$ and postcomposition with $z_j \to z_{j'}$ for some $j' \geq j$. In compactly assembled categories this agrees with \Cref{def:Compact_Morphism}, by \Cref{prop:Compact_Morphism_Characterizations,prop:Compact_Morphisms_Stable}.
	\end{remark}

	\begin{remark}
		\label[remark]{rem:Stable_Compact_Morphism_Definitions}
		The vanishing-colimit test in \Cref{prop:Compact_Morphisms_Stable} characterizes weak compactness in any stable category with filtered colimits. Compact assembly is used only to identify this with the natural factorization condition of \Cref{def:Compact_Morphism}, via \Cref{prop:Compact_Morphism_Characterizations}.
	\end{remark}

	\begin{corollary}[{cf.\ \cite[Remark~2.38]{Efimov_Large_Categories}}]
		\label[corollary]{cor:Compact_Iff_Identity_Weakly_Compact}
		Let $C$ be a stable category with filtered colimits. Then $x \in C$ is compact if and only if $\id_x$ is weakly compact.
	\end{corollary}
	\prooffor{cor:Compact_Iff_Identity_Weakly_Compact}
	\begin{proof}
		Suppose that $\id_x$ is weakly compact. By \Cref{prop:Compact_Morphisms_Stable}, we have $\pi_0\Map_{\Ind(C)}(y(x),W) = 0$ for every $W \in \Ind(C)$ with $\colim W \simeq 0$. Applying this to all shifts of $W$, we get $\hom_{\Ind(C)}(y(x),W) \simeq 0$. For $Z = \quotcolim_i z_i$, applying this to the fiber $W$ of the unit $\eta_Z\colon Z \to y(\colim Z)$ shows that
		\begin{align*}
			\colim_i \hom_C(x,z_i) &\simeq \hom_{\Ind(C)}(y(x),Z) \\
			&\to \hom_{\Ind(C)}(y(x),y(\colim Z)) \simeq \hom_C(x,\colim_i z_i)
		\end{align*}
		is an equivalence, i.e.\ $x$ is compact. The converse is \Cref{ex:Weakly_Compact_Basic}(1).
	\end{proof}

	\section{Compactly exhaustible objects}
	\label[section]{sec:Compactly_Exhaustible}

	In \Cref{sec:Omega_1_Presentability} we saw that the $\omega_1$-compact objects of a compactly assembled category are precisely the sequential colimits along compact morphisms (\Cref{thm:Compactly_Assembled_Omega_1} and \Cref{rem:Omega_1_Compact_Compact_Morphisms}). We now show that, for presentable stable categories, the existence of enough such objects characterizes dualizability.

	\begin{definition}
		\label[definition]{def:Compactly_Exhaustible}
		Let $C$ be a category with filtered colimits. A \emph{compact exhaustion} of an object $x \in C$ is a sequence $x_0 \to x_1 \to x_2 \to \cdots$ of weakly compact morphisms together with an equivalence $\colim_n x_n \simeq x$. An object is called \emph{compactly exhaustible} if it admits a compact exhaustion.
	\end{definition}

	\begin{proposition}[{\cite[Theorem~2.2.15]{Krause_Nikolaus_Putzstuck_Sheaves}}]
		\label[proposition]{prop:Dualizable_Compactly_Exhaustible}
		Let $C \in \PrL_{\st}$. Then $C$ is dualizable if and only if $C$ is generated under colimits by compactly exhaustible objects.
	\end{proposition}
	\prooffor{prop:Dualizable_Compactly_Exhaustible}
	\begin{proof}
		If $C$ is dualizable, then $C$ is generated under colimits by its $\omega_1$-compact objects, which are compactly exhaustible by \Cref{thm:Compactly_Assembled_Omega_1} together with \Cref{prop:Compact_Morphism_Characterizations}.

		Conversely, assume that $C$ is generated under colimits by compactly exhaustible objects. By \Cref{thm:Lurie_Dualizable_Characterization}, it suffices to show that $\colim\colon \Ind(C) \to C$ admits a left adjoint, i.e.\ that for every $x \in C$ the functor $\Map_C(x,\colim(-))\colon \Ind(C) \to \An$ is corepresentable. The category $\Ind(C)$ is stable and admits small colimits (any small diagram lies in some $\Ind(C^{\lambda})$, and the inclusions $\Ind(C^{\lambda}) \subseteq \Ind(C^{\mu})$ are exact and preserve filtered colimits, see \Cref{constr:Ind_Large}). Hence the class of $x$ for which $\Map_C(x,\colim(-))$ is corepresentable is closed under small colimits: if $x \simeq \colim_a x_a$ and $\Map_C(x_a,\colim(-)) \simeq \Map_{\Ind(C)}(X_a,-)$ naturally, then
		\begin{align*}
			\Map_C(x,\colim(-)) &\simeq \lim_a \Map_C(x_a,\colim(-)) \\
			&\simeq \lim_a \Map_{\Ind(C)}(X_a,-) \simeq \Map_{\Ind(C)}(\colim_a X_a,-).
		\end{align*}
		So it suffices to treat a compactly exhaustible object $x \simeq \colim_n x_n$, with $x_n \to x_{n+1}$ weakly compact. Put $X := \colim_n y(x_n) \in \Ind(C)$, the colimit being taken in $\Ind(C)$. Then $\colim X \simeq \colim_n x_n \simeq x$, as $\colim$ is a left adjoint and $\colim \circ y \simeq \id$. We claim that the map
		\[
			\Map_{\Ind(C)}(X,Z) \to \Map_C(\colim X, \colim Z) \simeq \Map_C(x,\colim Z)
		\]
		induced by $\colim$ is an equivalence, naturally in $Z \in \Ind(C)$. Under the adjunction $\colim \dashv y$, it is identified with the map $(\eta_Z)_*\colon \Map(X,Z) \to \Map(X,y(\colim Z))$. Since everything is stable, it suffices to show that $\hom_{\Ind(C)}(X,W) \simeq 0$ for $W = \fib(\eta_Z)$, which satisfies $\colim W \simeq 0$ (as $\colim(\eta_Z)$ is an equivalence).

		Put $h_n := \hom_{\Ind(C)}(y(x_n),W)$. We have $\hom(X,W) \simeq \lim_n h_n$, so there are Milnor short exact sequences
		\[
			0 \to {\lim_n}^1 \pi_{m+1}h_n \to \pi_m\hom(X,W) \to \lim_n \pi_m h_n \to 0.
		\]
		Now $\pi_m h_n \simeq \pi_0\Map(y(x_n),\Omega^mW)$, and $\colim \Omega^mW \simeq 0$. Since $x_n \to x_{n+1}$ is weakly compact, \Cref{prop:Compact_Morphisms_Stable} shows that all transition maps in the tower $(\pi_m h_n)_n$ are zero, for every $m \in \Z$. A tower of abelian groups with zero transition maps has vanishing $\lim$ and $\lim^1$, so $\pi_m\hom(X,W) = 0$ for all $m$, as desired.
	\end{proof}

	% Source: written from the literature (Krause--Nikolaus--P\"utzst\"uck)
	\begin{remark}
		\label[remark]{rem:Compact_Exhaustions_In_Presentable_Categories}
		In \Cref{prop:Dualizable_Compactly_Exhaustible}, stability is used to reduce to ind-objects with vanishing colimit. For general presentable categories, the analogous statement holds provided filtered colimits in $C$ are exact, with ``dualizable'' replaced by ``compactly assembled'' \cite[Theorem~2.2.15]{Krause_Nikolaus_Putzstuck_Sheaves}.
	\end{remark}

	\section{A universal property of \texorpdfstring{$\hat{y}$}{y-hat}}
	\label[section]{sec:Universal_Property_Yhat}

	For an object $x$ of a compactly assembled category $C$, the functor $\Map_{\Ind(C)}(y(x),\hat{y}(-))$ is a ``filtered-colimit-preserving approximation'' of $\Map_C(x,-)$: if $x$ is compact, the two agree, and in general the former is the best approximation of the latter by a functor that preserves filtered colimits. To make this precise, for categories $C$ and $D$ with filtered colimits we write $\Fun^{\mathrm{filt}}(C,D) \subseteq \Fun(C,D)$ for the full subcategory of functors preserving filtered colimits.

	\begin{proposition}[{\cite[Proposition~2.4.8]{Krause_Nikolaus_Putzstuck_Sheaves}}]
		\label[proposition]{prop:Filtered_Coreflection}
		Let $C$ be compactly assembled, and let $D$ be a category with filtered colimits. Then the inclusion $\Fun^{\mathrm{filt}}(C,D) \hookrightarrow \Fun(C,D)$ admits a right adjoint, given by
		\[
			F \mapsto \colim_D \circ \Ind(F) \circ \hat{y},
		\]
		where $\Ind(F)\colon \Ind(C) \to \Ind(D)$ is the extension of $F$ preserving filtered colimits.
	\end{proposition}
	\prooffor{prop:Filtered_Coreflection}
	\begin{proof}
		Write $\colim = \colim_C$. The adjunction $\hat{y} \dashv \colim$ induces an adjunction
		\[
			\colim^*\colon \Fun(C,D) \rightleftarrows \Fun(\Ind(C),D)\noloc \hat{y}^*
		\]
		given by precomposition: a transformation $G\colim \to H$ corresponds to the composite $G \to G\colim\hat{y} \to H\hat{y}$, using the unit $\id \to \colim\hat{y}$. Since $\colim$ and $\hat{y}$ preserve filtered colimits, this restricts to an adjunction between $\Fun^{\mathrm{filt}}(C,D)$ and $\Fun^{\mathrm{filt}}(\Ind(C),D)$. By the universal property of $\Ind(C)$ (the analogue of \Cref{thm:Ind_Kappa_Universal_Property}(2) for large $C$, see \Cref{constr:Ind_Large}), restriction along $y$ is an equivalence $\Fun^{\mathrm{filt}}(\Ind(C),D) \simeq \Fun(C,D)$, with inverse $F \mapsto \colim_D \circ \Ind(F)$. Under this equivalence, the left adjoint $\colim^*$ becomes $G \mapsto G\colim y \simeq G$, i.e.\ the inclusion, and the right adjoint $\hat{y}^*$ becomes $F \mapsto \colim_D \circ \Ind(F) \circ \hat{y}$.
	\end{proof}

	\begin{corollary}
		\label[corollary]{cor:Yhat_Universal_Property}
		Let $C$ be compactly assembled and $x \in C$. Then the natural transformation
		\[
			\Map_{\Ind(C)}\big(y(x),\hat{y}(-)\big) \xrightarrow{(\varepsilon)_*} \Map_{\Ind(C)}\big(y(x),y(-)\big) \simeq \Map_C(x,-)
		\]
		exhibits $\Map_{\Ind(C)}(y(x),\hat{y}(-))$ as the terminal object of $\Fun^{\mathrm{filt}}(C,\An)_{/\Map_C(x,-)}$. If $C$ is moreover stable, the same holds for the mapping spectra $\hom_{\Ind(C)}(y(x),\hat{y}(-)) \to \hom_C(x,-)$ in $\Fun^{\mathrm{filt}}(C,\Sp)_{/\hom_C(x,-)}$.
	\end{corollary}
	\prooffor{cor:Yhat_Universal_Property}
	\begin{proof}
		For a right adjoint $R$ of a fully faithful inclusion $i$, the counit $iR(F) \to F$ is a terminal object of the slice over $F$, since $\Map(G,R(F)) \simeq \Map(iG,F)$. Apply this to \Cref{prop:Filtered_Coreflection} with $D = \An$ and $F = \Map_C(x,-)$. For $Z = \quotcolim_j z_j \in \Ind(C)$ we have
		\[
			\colim_{\An}\Ind(F)(Z) \simeq \colim_j \Map_C(x,z_j) \simeq \Map_{\Ind(C)}(y(x),Z),
		\]
		as $y(x)$ is compact in $\Ind(C)$. Hence the right adjoint sends $F$ to $\Map_{\Ind(C)}(y(x),\hat{y}(-))$. To identify the counit, write $\tilde{F} = \colim_{\An} \circ \Ind(F)$. By construction, the counit at $F$ is obtained by restricting the counit $\tilde{F}\hat{y}\colim \to \tilde{F}$ of the precomposition adjunction along $y$, and the latter is $\tilde{F}$ applied to the counit $\hat{y}\colim \to \id$ of $\hat{y} \dashv \colim$. At $y(c)$, this counit is the map $\hat{y}(c) \simeq \hat{y}(\colim y(c)) \to y(c)$, which is $\varepsilon_c$ by definition of $\varepsilon$. So the counit is $\tilde{F}(\varepsilon)$, i.e.\ the map induced by $\varepsilon$. The stable case is identical, with $D = \Sp$.
	\end{proof}

	\section{The axiom (AB6)}
	\label[section]{sec:AB6}

	We finally discuss the characterization of dualizable categories by the axiom (AB6), due to Clausen and Scholze. It expresses the fact that infinite products commute with filtered colimits, in the strongest possible sense. Concretely, (AB6) is the elementwise form of condition (2) of \Cref{prop:Dualizable_Further_Characterizations}, that $\colim\colon \Ind(C) \to C$ preserves products; the translation between the two uses the description of products in $\Ind(C)$ given in \Cref{lem:Products_In_Ind}.

	\begin{definition}
		\label[definition]{def:AB6}
		Let $C$ be a category with small products and filtered colimits. We say that $C$ satisfies \emph{(AB6)} if for every set $S$ and every family of functors $F_s\colon J_s \to C$, $s \in S$, with each $J_s$ filtered, the canonical map
		\[
			\colim_{(j_s)_s \in \prod_s J_s} \prod_{s \in S} F_s(j_s) \to \prod_{s \in S} \colim_{j \in J_s} F_s(j)
		\]
		is an equivalence. Note that $\prod_s J_s$ is again filtered.
	\end{definition}

	\begin{example}
		\label[example]{ex:AB6_Sets_And_Animae}
		The category of sets and the category $\An$ of animae satisfy (AB6) \cite[Section~2.7]{Efimov_Large_Categories}.
	\end{example}
	% Source: written from the literature (Efimov, Section 2.7); proof added with AI assistance.
	\prooffor{ex:AB6_Sets_And_Animae}
	\begin{proof}
		For sets, surjectivity follows by choosing, for each $s$, a stage $j_s$ representing the given element of $\colim_{J_s}F_s$. For injectivity, two representatives which agree in each colimit agree at a later stage in each coordinate. Choosing these stages gives equality in the colimit over $\prod_s J_s$. Here we use the axiom of choice.

		For animae, we may replace each $J_s$ by a cofinal directed poset, by \Cref{thm:Filtered_Posets_Cofinal}. The product of these cofinal functors is cofinal. Indeed, if $u\colon P\to J$ is such a functor and $j\in J$, the comma category $P\times_J J_{j/}$ is filtered: any finite diagram in it has an underlying cone in $P$, and the finitely many compatibility data in the animae $\Map_J(j,u(p))$ can be extended at a later stage, since their filtered colimit is contractible by cofinality. Here we use that filtered colimits of animae commute with finite limits. Thus the comma categories for the product functor are products of nonempty filtered categories, which are again nonempty and filtered (using choice), and hence contractible. The set calculation proves that the comparison map induces a bijection on $\pi_0$, since $\pi_0$ commutes with products and filtered colimits of animae.

		Fix a point in the source, represented by points $x_s \in F_s(j_s)$. Restrict to the cofinal product of the tails $J_{s,\geq j_s}$, so that all diagrams are pointed by the images of $x_s$. For $n \geq 1$, homotopy groups commute with products and filtered colimits of pointed animae. Applying the set calculation to these homotopy groups therefore gives
		\[
			\colim_{(k_s)_s \geq (j_s)_s}\prod_s \pi_n(F_s(k_s),x_s)
			\cong \prod_s \colim_{k_s \geq j_s}\pi_n(F_s(k_s),x_s).
		\]
		Thus the comparison map induces isomorphisms on all homotopy groups at every basepoint, and is an isomorphism of animae.
	\end{proof}

	% Source: written when preparing these notes, with AI assistance
	\begin{remark}
		\label[remark]{rem:AB6_Versus_Naive_Commutation}
		Despite its informal description, (AB6) does \emph{not} imply that products commute with filtered colimits in the naive sense: for a set $S$, a filtered category $J$ and a functor $F\colon S \times J \to C$, the canonical map
		\[
			\colim_{j \in J} \prod_{s \in S} F(s,j) \to \prod_{s \in S} \colim_{j \in J} F(s,j)
		\]
		is in general not an equivalence. Taking $J_s = J$ for all $s$, this map factors as
		\[
			\colim_{j \in J} \prod_{s \in S} F(s,j) \to \colim_{(j_s)_s \in J^S} \prod_{s \in S} F(s,j_s) \to \prod_{s \in S} \colim_{j \in J} F(s,j),
		\]
		where the first map is induced by the diagonal $J \to J^S$ and the second map is the one from \Cref{def:AB6}. (AB6) says that the second map is an equivalence, but the diagonal is in general not cofinal: a family $(j_s)_s$ need not have a common upper bound in $J$.

		Indeed, the naive statement already fails in $\mathrm{Set}$. Let $S = J = \N$ and $F(s,j) = \{0,\dots,j\}$, with the inclusions as transition maps. Then $\prod_s \colim_j F(s,j) = \N^{\N}$, whereas $\colim_j \prod_s F(s,j)$ is the set of bounded sequences. In the colimit over $J^S$, on the other hand, a sequence $(n_s)_s$ is reached at the stage $(n_s)_s$. Similarly, in $\Sp$, which satisfies (AB6) by \Cref{prop:Dualizable_AB6}, the functor $F(s,j) = \bigoplus_{k \leq j} \Z$ gives the map $\bigoplus_{\N}\prod_{\N}\Z \to \prod_{\N}\bigoplus_{\N}\Z$, which is not surjective on $\pi_0$: the family whose $s$-th component is the $s$-th basis vector does not have uniformly bounded support.

		If $J$ is $\kappa$-filtered for a regular cardinal $\kappa > |S|$, then the diagonal $J \to J^S$ is cofinal, so in this case (AB6) does imply the naive commutation.
	\end{remark}

	(AB6) in $C$ is equivalent to the statement that the colimit functor $\colim\colon \Ind(C) \to C$ preserves products. This rests on the following description of products in $\Ind(C)$.

	\begin{lemma}[{\cite[Proposition~2.50]{Efimov_Large_Categories}}]
		\label[lemma]{lem:Products_In_Ind}
		Let $C$ be a locally small category with small products. Then $\Ind(C)$ has small products, and for filtered diagrams $(x_j)_{j \in J_s}$ in $C$, $s \in S$, we have
		\[
			\prod_{s \in S} \quotcolim_{j \in J_s} x_j \simeq \quotcolim_{(j_s)_s \in \prod_s J_s} \prod_{s \in S} x_{j_s}.
		\]
	\end{lemma}
	\prooffor{lem:Products_In_Ind}
	\begin{proof}
		We show that the right-hand side has the universal property of the product. Since every object of $\Ind(C)$ is a filtered colimit of representables, it suffices to test against $y(c)$, $c \in C$. As $y(c)$ is compact, and using (AB6) in $\An$, we get natural equivalences
		\begin{align*}
			\Map\Big(y(c),\quotcolim_{(j_s)_s} \prod_s x_{j_s}\Big) &\simeq \colim_{(j_s)_s} \prod_s \Map_C(c,x_{j_s}) \\
			&\simeq \prod_s \colim_{j \in J_s} \Map_C(c,x_j) \\
			&\simeq \prod_s \Map\Big(y(c),\quotcolim_{j \in J_s} x_j\Big). \qedhere
		\end{align*}
	\end{proof}

	\begin{proposition}[{Clausen--Scholze, see \cite[Proposition~2.51]{Efimov_Large_Categories}}]
		\label[proposition]{prop:Dualizable_AB6}
		Let $C \in \PrL_{\st}$. Then $C$ is dualizable if and only if $C$ satisfies (AB6).
	\end{proposition}
	\prooffor{prop:Dualizable_AB6}
	\begin{proof}
		By \Cref{lem:Products_In_Ind}, and since $\colim$ sends $\quotcolim_{(j_s)_s}\prod_s x_{j_s}$ to $\colim_{(j_s)_s}\prod_s x_{j_s}$, the axiom (AB6) for $C$ says exactly that $\colim\colon \Ind(C) \to C$ preserves products of families of ind-objects. Every family of objects of $\Ind(C)$ is of this form. So (AB6) holds if and only if $\colim$ preserves small products, which is equivalent to dualizability of $C$ by \Cref{prop:Dualizable_Further_Characterizations}.
	\end{proof}

	\begin{remark}
		\label[remark]{rem:AB6_And_Small_Ind_Completions}
		In the talk, the implication ``(AB6) $\implies$ dualizable'' was argued directly with $\colim\colon \Ind(C^{\kappa}) \to C$. Note that products in $\Ind(C^{\kappa})$ are not computed by the formula of \Cref{lem:Products_In_Ind}, since a product of $\kappa$-compact objects need not be $\kappa$-compact. This is why the argument passes through the large category $\Ind(C)$, as in the proof of (2) $\implies$ (3) in \Cref{prop:Dualizable_Further_Characterizations}.
	\end{remark}

	\begin{remark}
		\label[remark]{rem:AB6_And_Compact_Assembly}
		Since $\Sp$ is compactly generated, hence dualizable, it satisfies (AB6). More generally, by \cite[Remark~2.52]{Efimov_Large_Categories} a presentable, not necessarily stable, category is compactly assembled if and only if it satisfies (AB6) and filtered colimits in it commute with finite limits.
	\end{remark}

	\section{Automatic dualizability and images}
	\label[section]{sec:Automatic_Dualizability}

	% Source: Talk 5 (Anton Zakrewski, 10.11.2026); Zakrewski's notes for the talk; Efimov, Section 2.8
	As an application of the two characterizations of dualizability in terms of compact exhaustions and (AB6), we show that dualizability propagates along strongly continuous functors whose images generate the target. As a consequence, a strongly continuous functor out of a dualizable category has a well-behaved image. These results will be used in \Cref{chap:Calkin_Limits_Colimits} to study colimits and limits of dualizable categories. We begin with a criterion for strong continuity.

	\begin{lemma}[{\cite[Proposition~1.8]{Efimov_Large_Categories}}]
		\label[lemma]{lem:Automatic_Strong_Continuity}
		Let $G\colon D \to E$ be a colimit-preserving functor between presentable stable categories, and let $F_i\colon C_i \to D$, $i \in I$, be a family of strongly continuous functors from presentable stable categories $C_i$. Assume that $D$ is generated under colimits by the essential images of the $F_i$, and that each composite $G \circ F_i$ is strongly continuous. Then $G$ is strongly continuous.
	\end{lemma}
	\prooffor{lem:Automatic_Strong_Continuity}
	\begin{proof}
		The right adjoints $F_i^R$ are jointly conservative: if $F_i^R(x) \simeq 0$ for all $i$, then the full subcategory of objects $d \in D$ with $\hom_D(d,x) \simeq 0$ is closed under colimits and contains the essential images of the $F_i$, so it is all of $D$, and in particular $x \simeq 0$. Now let $(e_a)_a$ be a diagram in $E$. Since the functors $F_i^R$ and $F_i^RG^R \simeq (GF_i)^R$ preserve colimits, applying $F_i^R$ to the comparison map $\colim_a G^R(e_a) \to G^R(\colim_a e_a)$ gives the comparison map for $(GF_i)^R$, which is an equivalence. Hence the comparison map is an equivalence, i.e.\ $G^R$ preserves colimits.
	\end{proof}

	\begin{proposition}[{\cite[Proposition~2.53]{Efimov_Large_Categories}}]
		\label[proposition]{prop:Generation_Dualizable}
		Let $D$ be a presentable stable category, and let $F_i\colon C_i \to D$, $i \in I$, be a family of strongly continuous functors from dualizable categories $C_i$. If $D$ is generated under colimits by the essential images of the $F_i$, then $D$ is dualizable.
	\end{proposition}
	\prooffor{prop:Generation_Dualizable}
	\begin{proof}
		As in the proof of \Cref{lem:Automatic_Strong_Continuity}, the right adjoints $F_i^R\colon D \to C_i$ are jointly conservative. They preserve small products, being right adjoints, and filtered colimits, by strong continuity. Let $S$ be a set and let $(x_{s,j})_{j \in J_s}$, $s \in S$, be filtered diagrams in $D$. Applying $F_i^R$ to the comparison map
		\[
			\colim_{(j_s)_s \in \prod_s J_s} \prod_{s \in S} x_{s,j_s} \to \prod_{s \in S} \colim_{j \in J_s} x_{s,j}
		\]
		yields the corresponding comparison map for the diagrams $(F_i^R x_{s,j})_j$ in $C_i$, which is an equivalence by \Cref{prop:Dualizable_AB6}, as $C_i$ is dualizable. Hence the comparison map in $D$ is an equivalence. So $D$ satisfies (AB6), and is dualizable by \Cref{prop:Dualizable_AB6}.
	\end{proof}

	% Source: written from the literature (Efimov, Proposition 2.76; Krause--Nikolaus--P\"utzst\"uck, Theorem 2.2.15)
	\prooffor{prop:Generation_Dualizable}[compact-exhaustion]
	\begin{proof}[Second proof]
		A strongly continuous functor $F\colon C \to D$ preserves weakly compact morphisms. Indeed, let $f\colon x \to x'$ be weakly compact in $C$, let $(z_j)_j$ be a filtered diagram in $D$, and let $g\colon F(x') \to \colim_j z_j$. The adjoint map $x' \to F^R(\colim_j z_j) \simeq \colim_j F^R(z_j)$ becomes, after precomposition with $f$, a map factoring through some $F^R(z_j)$. Passing back to adjoints, $gF(f)$ factors through $z_j$. Hence $F$ also preserves compact exhaustions (\Cref{def:Compactly_Exhaustible}).

		By \Cref{prop:Dualizable_Compactly_Exhaustible}, each $C_i$ is generated under colimits by compactly exhaustible objects. Their images under $F_i$ are compactly exhaustible, and they generate $D$ under colimits, since the localizing subcategory they generate contains the essential images of the $F_i$. So $D$ is dualizable by \Cref{prop:Dualizable_Compactly_Exhaustible}.
	\end{proof}

	% Source: Talk 5 (Anton Zakrewski, 10.11.2026); Efimov, proof of Proposition 2.53
	\begin{remark}
		\label[remark]{rem:Products_Dualizable}
		In \cite{Efimov_Large_Categories}, the proof of \Cref{prop:Generation_Dualizable} is first reduced to the case of a single functor, using the induced functor $\prod_i C_i \to D$. For this, note that the product $\prod_i C_i$, computed in $\widehat{\Cat}$, is also the coproduct of the $C_i$ in $\PrL_{\st}$, as colimits in $\PrL$ are computed as limits along the right adjoints (\Cref{sec:PrL_Tensor_Products}). If each $C_i$ is dualizable, then so is $\prod_i C_i$: products and filtered colimits in $\prod_i C_i$ are computed componentwise, so (AB6) holds in $\prod_i C_i$ since it holds in each $C_i$. Alternatively, if $C_i$ is a retract of a compactly generated category $\widetilde{C}_i$ in $\PrL_{\st}$, then $\prod_i C_i$ is a retract of $\prod_i\widetilde{C}_i$, which is compactly generated by the objects with a single nonzero component, which is compact.
	\end{remark}

	% Source: Talk 5 (Anton Zakrewski, 10.11.2026); Efimov, Definition 2.54 and Corollary 2.55
	\begin{definition}[{\cite[Definition~2.54]{Efimov_Large_Categories}}]
		\label[definition]{def:Image_Functor}
		Let $F\colon C \to D$ be a colimit-preserving functor between presentable stable categories. The \emph{image} $\Im(F) \subseteq D$ is the smallest full stable subcategory of $D$ which contains the essential image of $F$ and is closed under colimits, i.e.\ the localizing subcategory generated by the essential image of $F$.
	\end{definition}

	If $C$ is generated under colimits by a set $S$ of objects, then $\Im(F)$ is generated under colimits by the set $F(S)$, and hence is presentable \cite[Section~1.3]{Efimov_Large_Categories}. By \Cref{prop:Localization_Sequences_PrL}(1), the cofiber of $F$ in $\PrL_{\st}$ is $D \to \Im(F)^{\perp}$, whose kernel is $\Im(F)$. So $\Im(F)$ is equivalently the fiber of the cofiber of $F$ in $\PrL_{\st}$.

	\begin{corollary}[{\cite[Corollary~2.55]{Efimov_Large_Categories}}]
		\label[corollary]{cor:Image_Dualizable}
		Let $F\colon C \to D$ be a strongly continuous functor between presentable stable categories, and assume that $C$ is dualizable. Then $\Im(F)$ is dualizable, and the inclusion $j\colon \Im(F) \hookrightarrow D$ is strongly continuous.
	\end{corollary}
	\prooffor{cor:Image_Dualizable}
	\begin{proof}
		Write $F = j \circ F'$ with $F'\colon C \to \Im(F)$. Since $j$ is fully faithful, the right adjoint of $F'$ is $F^R \circ j$, which preserves colimits. So $F'$ is strongly continuous, and its essential image generates $\Im(F)$ under colimits. Hence $\Im(F)$ is dualizable by \Cref{prop:Generation_Dualizable}. Since $j \circ F' = F$ is strongly continuous, \Cref{lem:Automatic_Strong_Continuity} shows that $j$ is strongly continuous.
	\end{proof}

	The two criteria have different uses. Compact exhaustions describe objects and allow constructions by successive factorizations; (AB6) detects dualizability through jointly conservative right adjoints. The image and generation results above make the latter method available without choosing an exhaustion for each object.

	% Theorem numbers for Efimov_Large_Categories refer to arXiv:2405.12169v4.
	\chapter{Examples of dualizable categories}
	\label[chapter]{chap:Examples}

	% Source: Talk 3 (Max Glaser, 27.10.2026)
	In this chapter we discuss several sources of dualizable (and, more generally, compactly assembled) categories:
	\begin{itemize}
		\item Kernels of homological epimorphisms of ring spectra, and in particular almost ring theory (\Cref{sec:Homological_Epimorphisms,sec:Almost_Ring_Theory}).
		\item Continuous posets, as a class of non-stable compactly assembled categories (\Cref{sec:Continuous_Posets}).
		\item Gluing constructions and extensions, including a nondualizable extension of compactly generated categories (\Cref{sec:Extensions_And_Gluing}).
	\end{itemize}

	\section{Kernels of homological epimorphisms}
	\label[section]{sec:Homological_Epimorphisms}

	% Source: Clausen, Lecture 2 (18.07.2023); Mathew, Lecture 3 (MPIM Bonn, 10.09.2024)
	Let $A \to B$ be a map of $\bbE_1$-rings, and consider the base change functor $B \otimes_A -\colon \D(A) \to \D(B)$ between the categories of left modules, with right adjoint the restriction functor $\res\colon \D(B) \to \D(A)$. Recall from \Cref{def:Homological_Epimorphism} that $A \to B$ (or rather $\Perf(A) \to \Perf(B)$) is a \emph{homological epimorphism} if $B \otimes_A -$ is a localization, i.e.\ if $\res$ is fully faithful. Since $\res$ preserves colimits, this is equivalent to the counit $B \otimes_A \res(N) \to N$ being an equivalence for all $N \in \D(B)$, and since both sides preserve colimits in $N$ and $\D(B)$ is generated by $B$, this is in turn equivalent to the multiplication map
	\[
		B \otimes_A B \to B
	\]
	being an equivalence. In the language of Clausen, $A \to B$ is then called \emph{idempotent}.

	\begin{proposition}
		\label[proposition]{prop:Kernel_Homological_Epi}
		Let $A \to B$ be a homological epimorphism of $\bbE_1$-rings, and put
		\[
			\D(A \to B) := \ker\big(B \otimes_A -\colon \D(A) \to \D(B)\big).
		\]
		Then $\D(A \to B)$ is dualizable, and the sequence $\D(A \to B) \hookrightarrow \D(A) \to \D(B)$ is a strongly continuous localization sequence. In particular,
		\[
			K^{\cont}(\D(A \to B)) \simeq \fib\big(K(A) \to K(B)\big).
		\]
	\end{proposition}
	\prooffor{prop:Kernel_Homological_Epi}
	\begin{proof}
		The functor $B \otimes_A -$ is a localization by assumption, and it is strongly continuous since its right adjoint $\res$ preserves colimits. Hence $\D(A \to B) \hookrightarrow \D(A) \to \D(B)$ is a strongly continuous localization sequence (\Cref{lem:Strong_Continuity_Localization_Sequence}), and $\D(A \to B)$ is dualizable by \Cref{prop:Localization_Sequence_Dualizable}. The final statement follows from the definition of $K^{\cont}$, see \Cref{sec:Continuous_K_Theory_Intro}.
	\end{proof}

	\begin{proposition}
		\label[proposition]{prop:Every_Dualizable_Kernel}
		Every dualizable category is of the form $\D(E \to E')$ for some homological epimorphism $E \to E'$ of $\bbE_1$-rings.
	\end{proposition}
	\prooffor{prop:Every_Dualizable_Kernel}
	\begin{proof}
		Let $C$ be dualizable. Choose a set $S$ of representatives of the isomorphism classes of $C^{\omega_1}$, a regular cardinal $\kappa > |S|$, and put $G := \bigoplus_{x \in S} x$. As a $\kappa$-small sum of $\kappa$-compact objects, $G$ is $\kappa$-compact. Let $A \subseteq C^{\kappa}$ be the thick subcategory generated by $G$; it contains $C^{\omega_1}$, since every object of $C^{\omega_1}$ is a retract of $G$. The functor $\hat{y}\colon C \to \Ind(C^{\kappa})$ takes values in $\Ind(C^{\omega_1}) \subseteq \Ind(A)$ (\Cref{thm:Compactly_Assembled_Omega_1}, see also \cite[Corollary~2.14]{Efimov_Large_Categories}), and hence is a fully faithful left adjoint of $\colim\colon \Ind(A) \to C$. Its essential image $\hat{y}(C)$ is closed under colimits, with right orthogonal $\ker(\colim\colon \Ind(A) \to C)$. Since the latter is closed under colimits in $\Ind(A)$, the localization
		\[
			Q\colon \Ind(A) \to \Ind(A)/\hat{y}(C) \simeq \ker(\colim)
		\]
		has a colimit-preserving right adjoint and thus preserves compact objects. As $Q$ is essentially surjective and $\Ind(A)$ is compactly generated by $G$, the category $\ker(\colim)$ is compactly generated by the single object $Q(G)$. Hence $\ker(\colim) \simeq \Ind(B)$ for $B$ the thick subcategory generated by $Q(G)$, and $Q$ is $\Ind$ of a homological epimorphism $A \to B$ with kernel $\hat{y}(C) \simeq C$; this is the Calkin sequence of \Cref{thm:Lurie_Dualizable_Characterization} with $C^{\kappa}$ replaced by $A$. Finally, by Morita theory \cite[Theorem~7.1.2.1]{lurie2016HA} we have $\Ind(A) \simeq \D(E)$ for $E := \End(G)\catop$ and $\Ind(B) \simeq \D(E')$ for $E' := \End(Q(G))\catop$, under which $Q$ corresponds to $E' \otimes_E -$, where the ring map $E \to E'$ is induced by $Q$ on endomorphisms. Indeed, $Q(G)$ has the resulting $E$-action, so the two colimit-preserving functors agree on the generator together with its endomorphism action, as required by Morita theory. So $C \simeq \D(E \to E')$.
	\end{proof}
	% In Clausen's lectures this was attributed to Efimov; we have not found it stated in this form in Efimov's papers. The proof above was written with AI assistance and should be checked.

	There is a convenient criterion for a map $A \to B$ with fiber $I$ to be a homological epimorphism. Recall that a non-unital ring spectrum $I$ is called \emph{h-unital} if the multiplication $I \otimes_{\S\oplus I} I \to I$ is an isomorphism, where $\S\oplus I$ denotes the unitalization. Equivalently, the nonunital relative tensor product, defined by the semisimplicial bar construction with terms $I^{\otimes(n+2)}$, maps isomorphically to $I$; see \cite[Section~2.10 and Lemma~3.5.5]{Krause_Nikolaus_Putzstuck_Sheaves}.

	\begin{theorem}[{\cite[Section~2.10, especially Corollary~2.10.6]{Krause_Nikolaus_Putzstuck_Sheaves}}]
		\label[theorem]{thm:H_Unital}
		For a non-unital ring spectrum $I$, the following are equivalent:
		\begin{enumerate}
			\item $I$ is h-unital;
			\item the map $\S \oplus I \to \S$ is a homological epimorphism;
			\item for every fiber sequence $I \to A \to B$ in which $A \to B$ is a map of unital ring spectra, the map $A \to B$ is a homological epimorphism.
		\end{enumerate}
	\end{theorem}

	Conversely, if $A \to B$ is a homological epimorphism of connective ring spectra and $\pi_0(A) \to \pi_0(B)$ is surjective, then $\fib(A \to B)$ is h-unital \cite[Proposition~2.10.7]{Krause_Nikolaus_Putzstuck_Sheaves}. The surjectivity hypothesis cannot be omitted: the fiber of $k[t]\to k[t,t^{-1}]$ is not h-unital \cite[Example~2.10.8]{Krause_Nikolaus_Putzstuck_Sheaves}.

	\begin{example}
		\label[example]{ex:Homological_Epis}
		The following maps are homological epimorphisms, and hence give rise to dualizable categories $\D(A \to B)$:
		\begin{enumerate}
			\item A localization $A \to A[S^{-1}]$ of a commutative ring at a multiplicative subset. More generally, if $A$ is a commutative Noetherian ring, the kernel of every homological epimorphism out of $A$ is compactly generated: the associated localization of $\D(A)$ is smashing, so this follows from \cite[Theorem~3.3 and Corollary~3.4]{Neeman_Chromatic_Tower}.
			\item A surjection $A \twoheadrightarrow B$ of perfect $\F_p$-algebras: the higher $\mathrm{Tor}^A$-groups between perfect $A$-algebras vanish by \cite[Lemma~3.16]{Bhatt_Scholze_Witt_Grassmannian}, so $B \otimes^{\bL}_A B \simeq B$. Similarly, if $A$ is an integral perfectoid ring and $B$ is a perfect $\F_p$-algebra, any surjection $A \twoheadrightarrow B$ is a homological epimorphism. Indeed, the derived $p$-completed tensor product of perfectoid rings is discrete by the proof of \cite[Lemma~10.3]{Bhatt_Scholze_Prisms}. Here $B \otimes^{\bL}_A B$ is already derived $p$-complete, since it is a $B$-module, and its degree-zero homology is $B$.
			\item (Wodzicki) If $A$ is a unital C$^*$-algebra and $I \subseteq A$ a closed two-sided $*$-ideal, then $A \to A/I$ is a homological epimorphism. For $A = B(\Hh)$ and $I = \Kk(\Hh)$, this gives the category $\D(\Kk)$ from \Cref{ex:Compact_Operators_Intro}.
			\item (Almost ring theory) If $A$ is a ring and $I \subseteq A$ is a two-sided ideal with $I^2 = I$ which is flat as a left or right $A$-module, then $A \to A/I$ is a homological epimorphism. We discuss this in detail in \Cref{sec:Almost_Ring_Theory}.
		\end{enumerate}
	\end{example}
	% Source: written from the literature (Neeman, Theorem 3.3; Bhatt--Scholze, Witt vectors, Lemma 3.16, and Prisms, proof of Lemma 10.3).

	\section{Almost ring theory}
	\label[section]{sec:Almost_Ring_Theory}

	% Source: Talk 3 (Max Glaser, 27.10.2026)
	Almost ring theory (Faltings, Gabber--Ramero \cite{Gabber_Ramero_Almost}) provides one of the most natural sources of dualizable categories which are not compactly generated. Throughout this section, let $A$ be an (associative, unital) ring and let $I \subseteq A$ be a two-sided ideal such that
	\begin{itemize}
		\item $I^2 = I$, and
		\item $I$ is flat as a right $A$-module.
	\end{itemize}
	We work with left $A$-modules, and write $A\text{-}\Mod$ for the abelian category of left $A$-modules and $\D(A)$ for its derived category.

	The forgetful functor $(A/I)\text{-}\Mod \hookrightarrow A\text{-}\Mod$ is fully faithful, with essential image the $A$-modules $M$ with $IM = 0$. Since $I^2 = I$, this subcategory is closed under extensions: if $0 \to M' \to M \to M'' \to 0$ is exact with $IM' = 0 = IM''$, then $IM = I^2M \subseteq I M' = 0$. As it is clearly closed under subobjects and quotients, it is a \emph{Serre subcategory}.

	\begin{definition}
		\label[definition]{def:Almost_Modules}
		The category of \emph{$(A,I)$-almost modules} is the Serre quotient
		\[
			(A,I)^a\text{-}\Mod := A\text{-}\Mod/(A/I)\text{-}\Mod.
		\]
	\end{definition}

	\begin{example}
		\label[example]{ex:Almost_Modules}
		\leavevmode
		\begin{enumerate}
			\item If $I = A$, then $(A/I)\text{-}\Mod = 0$ and $(A,A)^a\text{-}\Mod = A\text{-}\Mod$.
			\item The motivating examples come from non-discrete valuation rings: let $A$ be a valuation ring of rank one whose value group is not discrete, e.g.\ $A = \Oo_{\C_p}$, and let $I = \mm$ be its maximal ideal. Then $\mm^2 = \mm$, since every element of $\mm$ is a product of two elements of $\mm$, and $\mm$ is flat since it is torsion-free over a valuation ring. An $A$-module is called \emph{almost zero} if it is killed by $\mm$; for instance, the residue field $A/\mm$ becomes zero in the almost category. This is the setting of Faltings' almost purity theorem in $p$-adic Hodge theory.
		\end{enumerate}
	\end{example}

	\subsection*{The derived picture}
	We have an adjoint triple
	\[
		\begin{tikzcd}[column sep=large]
			\D(A) \arrow[r, shift left=2.5, "L"] \arrow[r, shift right=2.5, "R"'] & \D(A/I) \arrow[l, hook', "\res" description]
		\end{tikzcd}
	\]
	where $\res$ denotes restriction of scalars, $L = A/I \otimes^{\bL}_A -$ and $R = \bR\Hom_A(A/I,-)$.

	\begin{lemma}
		\label[lemma]{lem:Almost_Restriction_FF}
		The functor $\res\colon \D(A/I) \to \D(A)$ is fully faithful. In particular, $L$ is a localization and $A \to A/I$ is a homological epimorphism.
	\end{lemma}
	\prooffor{lem:Almost_Restriction_FF}
	\begin{proof}
		As explained at the beginning of \Cref{sec:Homological_Epimorphisms}, it suffices to show that the multiplication map $A/I \otimes^{\bL}_A A/I \to A/I$ is an equivalence. Using the short exact sequence $0 \to I \to A \to A/I \to 0$ of right $A$-modules, we find
		\begin{align*}
			\mathrm{Tor}^A_0(A/I,A/I) &= A/I, \\
			\mathrm{Tor}^A_1(A/I,A/I) &\cong I/I^2 = 0, \\
			\mathrm{Tor}^A_k(A/I,A/I) &\cong \mathrm{Tor}^A_{k-1}(I,A/I) = 0 \quad \text{for } k \geq 2,
		\end{align*}
		where the last vanishing uses flatness of $I$ as a right $A$-module.
	\end{proof}

	\begin{lemma}
		\label[lemma]{lem:Almost_Essential_Image}
		The essential image of $\res\colon \D(A/I) \hookrightarrow \D(A)$ is the full subcategory
		\[
			\D_{(A/I)\text{-}\Mod}(A) := \{X \in \D(A) \mid IH_n(X) = 0 \text{ for all } n \in \Z\}
		\]
		of complexes whose homology lies in the Serre subcategory $(A/I)\text{-}\Mod$. Moreover, the kernel of $L$ consists of those $X \in \D(A)$ with $X \simeq I \otimes^{\bL}_A X$, and $I \otimes^{\bL}_A X \in \ker(L)$ for every $X \in \D(A)$.
	\end{lemma}
	\prooffor{lem:Almost_Essential_Image}
	\begin{proof}
		Since $\res$ is fully faithful with left adjoint $L$, its essential image consists of those $X$ for which the unit $X \to \res L(X)$ is an equivalence. Tensoring the short exact sequence $0 \to I \to A \to A/I \to 0$ of $A$-bimodules with $X$ gives a cofiber sequence
		\[
			I \otimes^{\bL}_A X \to X \to \res L(X),
		\]
		so $X$ lies in the essential image if and only if $I \otimes^{\bL}_A X \simeq 0$. Since $I$ is flat as a right module, $H_n(I \otimes^{\bL}_A X) \cong I \otimes_A H_n(X)$. Finally, for a left $A$-module $M$, we have $I \otimes_A M = 0$ if and only if $IM = 0$: the multiplication map $I \otimes_A M \to IM$ is surjective, and conversely if $IM = 0$ then $M$ is an $A/I$-module and $I \otimes_A M \cong (I \otimes_A A/I) \otimes_{A/I} M \cong I/I^2 \otimes_{A/I} M = 0$.

		For the second statement, $L(I \otimes^{\bL}_A X) \simeq (A/I \otimes^{\bL}_A I) \otimes^{\bL}_A X$, and $A/I \otimes^{\bL}_A I \simeq \fib(A/I \to A/I \otimes^{\bL}_A A/I) \simeq 0$ by \Cref{lem:Almost_Restriction_FF}. The characterization of $\ker(L)$ follows from the cofiber sequence above.
	\end{proof}

	\begin{proposition}
		\label[proposition]{prop:Almost_Derived_Category}
		There is an equivalence
		\[
			\D((A,I)^a\text{-}\Mod) \simeq \ker\big(L\colon \D(A) \to \D(A/I)\big),
		\]
		and $\ker(L) = \D(A \to A/I)$ is dualizable. More precisely, $\D((A,I)^a\text{-}\Mod) \hookrightarrow \D(A) \xrightarrow{L} \D(A/I)$ is a strongly continuous localization sequence, and $\D((A,I)^a\text{-}\Mod)$ is also the kernel of $L$ computed in $\Pr^{\dual}$.
	\end{proposition}
	\prooffor{prop:Almost_Derived_Category}
	\begin{proof}
		Since $(A/I)\text{-}\Mod \subseteq A\text{-}\Mod$ is a Serre subcategory of a Grothendieck abelian category which is closed under colimits, the derived category of the Serre quotient is the Verdier quotient
		\[
			\D((A,I)^a\text{-}\Mod) \simeq \D(A)/\D_{(A/I)\text{-}\Mod}(A).
		\]
		Indeed, the quotient functor $A\text{-}\Mod \to (A,I)^a\text{-}\Mod$ is exact and has a fully faithful right adjoint, and for such functors this identification is proved in \cite[Section~7.7]{Krause_Localization}; see also \cite[Example~2.26]{Efimov_Large_Categories}. It thus remains to identify the Verdier quotient $\D(A)/\D_{(A/I)\text{-}\Mod}(A)$ with $\ker(L)$. Consider the functor $q := I \otimes^{\bL}_A -\colon \D(A) \to \ker(L)$, which lands in $\ker(L)$ by \Cref{lem:Almost_Essential_Image}. For $K \in \ker(L)$ and $X \in \D(A)$, applying $\Map(K,-)$ to the cofiber sequence $I \otimes^{\bL}_A X \to X \to \res L(X)$ and using $\Map(K,\res L X) \simeq \Map(LK,LX) \simeq *$ shows that $\Map(K,qX) \iso \Map(K,X)$. So $q$ is right adjoint to the fully faithful inclusion $\ker(L) \hookrightarrow \D(A)$, and hence $q$ exhibits $\ker(L)$ as the localization of $\D(A)$ at the morphisms inverted by $q$ \cite[Proposition~5.2.7.12]{lurie2009HTT}. Since $q$ is exact, these are the morphisms whose cofiber lies in $\ker(q)$, which is $\D_{(A/I)\text{-}\Mod}(A)$ by \Cref{lem:Almost_Essential_Image}. Hence $\ker(L) \simeq \D(A)/\D_{(A/I)\text{-}\Mod}(A)$.

		By \Cref{prop:Kernel_Homological_Epi} and \Cref{lem:Almost_Restriction_FF}, $\ker(L)$ is dualizable and its inclusion into $\D(A)$ is strongly continuous; explicitly, its right adjoint is $X \mapsto I \otimes^{\bL}_A X$ by \Cref{lem:Almost_Essential_Image}. To see that $\ker(L)$ is also the kernel in $\Pr^{\dual}$, let $E$ be dualizable and $F\colon E \to \D(A)$ strongly continuous with $LF \simeq 0$. Then $F$ factors uniquely through a colimit-preserving functor $F'\colon E \to \ker(L)$, and $F'$ is strongly continuous since its right adjoint is the composite of the colimit-preserving functors $\ker(L) \hookrightarrow \D(A) \xrightarrow{F^R} E$.
	\end{proof}

	\begin{remark}
		\label[remark]{rem:Almost_General}
		In the talk, the last part of \Cref{prop:Almost_Derived_Category} was deduced from the general fact that the inclusion $\Pr^{\dual} \hookrightarrow \PrL_{\st}$ preserves pullbacks along localizations (\Cref{prop:Dualizable_Pullbacks}(2)). Efimov \cite[Example~2.26]{Efimov_Large_Categories} shows that $\D((A,I)^a\text{-}\Mod)$ is dualizable even without any flatness assumption on $I$: in general, $\D(A) \to \D_{(A/I)\text{-}\Mod}(A)$ is still a strongly continuous localization, identified with $B \otimes_A -$ for a connective DG ring $B$ with $H_0(B) = A/I$, and $B = A/I$ if $I$ is flat as a right or left $A$-module.
	\end{remark}

	The following is a special case of Keller's example (\Cref{ex:Keller}), and shows that almost module categories are typically very far from being compactly generated.

	\begin{proposition}
		\label[proposition]{prop:Almost_No_Compacts}
		If $I \subseteq \mathrm{Jac}(A)$ is contained in the Jacobson radical of $A$, then $\D((A,I)^a\text{-}\Mod)$ has no nonzero compact objects.
	\end{proposition}
	\prooffor{prop:Almost_No_Compacts}
	\begin{proof}
		By \Cref{prop:Almost_Derived_Category}, the inclusion $\ker(L) \hookrightarrow \D(A)$ is strongly continuous and hence preserves compact objects. So it suffices to show: if $P \in \Perf(A)$ satisfies $A/I \otimes^{\bL}_A P \simeq 0$, then $P \simeq 0$.

		Suppose $P \neq 0$ and let $n$ be minimal with $H_n(P) \neq 0$. We may represent $P$ by a bounded complex of finitely generated projective modules concentrated in homological degrees $\geq n$: starting from any bounded complex $P_\bullet$ of finitely generated projectives, if $a < n$ is the lowest degree with $P_a \neq 0$, then $P_{a+1} \to P_a$ is surjective as $H_a(P) = 0$, hence split, and we may replace $P_{a+1}$ by its kernel (again finitely generated projective) and remove $P_a$. In this representation, $H_n(P) = \coker(P_{n+1} \to P_n)$ is a finitely generated $A$-module, and by right exactness of the tensor product
		\[
			0 = H_n(A/I \otimes^{\bL}_A P) \cong A/I \otimes_A H_n(P) \cong H_n(P)/IH_n(P).
		\]
		So $H_n(P) = IH_n(P) \subseteq \mathrm{Jac}(A)H_n(P)$, and Nakayama's lemma implies $H_n(P) = 0$, a contradiction.
	\end{proof}

	\begin{example}
		\label[example]{ex:Almost_Modules_Over_Cp}
		In the situation of \Cref{ex:Almost_Modules}(2), $A = \Oo_{\C_p}$ is local with maximal ideal $\mm = \mathrm{Jac}(A)$. Hence $\D((\Oo_{\C_p},\mm)^a\text{-}\Mod)$ is a dualizable category without nonzero compact objects.
	\end{example}

	\subsection*{Hilbert-space models for operator ideals}

	% Source: Clausen, Lectures 2--3 (18.--19.07.2023); the proof below was added when preparing these notes, with AI assistance.
	\begin{proposition}
		\label[proposition]{prop:Operator_Ideal_Models}
		Let $A = B(\Hh)$ for a separable infinite-dimensional complex Hilbert space, and let $J$ be either $\Kk(\Hh)$ or $S_{0+} = \bigcap_{p>0} S_p(\Hh)$. Let $E$ be the additive category of separable complex Hilbert spaces and bounded linear maps, with the split exact structure, and put $D = \Ind(\D^b(E))$. Then $A \to A/J$ is a homological epimorphism. Under Morita theory, its kernel identifies with the full stable subcategory $C_J \subseteq D$ generated under colimits by sequences of Hilbert spaces whose transition maps belong to $J$. Here membership in $J$ means compactness, respectively membership in every Schatten class, also for operators between different Hilbert spaces.

		The completed Hilbert tensor product restricts to a symmetric monoidal structure on $C_J$ whose unit and multiplication functors are strongly continuous. In particular, $C_J \in \CAlg(\Pr^{\dual})$.
	\end{proposition}
	\prooffor{prop:Operator_Ideal_Models}
	\begin{proof}
		We use right $A$-modules in this proof. Every separable Hilbert space is a retract of $\Hh$, and every bounded idempotent has closed image. Hence $V \mapsto \Hom_E(\Hh,V)$ identifies $E$ with the finitely generated projective right $A$-modules, and $D$ with their unbounded derived category. To recover our left-module convention, choose a conjugation $c$ on $\Hh$: the complex-linear transpose $a \mapsto ca^*c$ identifies $A\catop$ with $A$ and preserves both ideals. The intrinsic category $C_J$ itself does not require this choice.

		We first record a factorization fact: given countably many $a_n \in J$, there is a positive injective $d \in J$ with $a_n \in dA$ for every $n$. Include a fixed positive injective operator in $J$ among the $a_n$, and choose $\lambda_n > 0$ sufficiently small that
		\[
			s = \sum_n \lambda_n a_na_n^*
		\]
		converges in operator norm and, when $J = S_{0+}$, in every Schatten quasi-norm $\|{-}\|_{1/m}$, $m \geq 1$. Such a choice is possible by imposing the first $n$ quasi-norm bounds on the $n$th summand, with summable bounds in the corresponding quasi-norm metrics. Then $s$ is positive and injective and belongs to $J$, as does $d=s^{1/2}$. Since $a_na_n^* \leq \lambda_n^{-1}s$, the map $dx \mapsto a_n^*x$ is bounded on the dense image of $d$. Extending it to $\Hh$ and taking adjoints gives $a_n = db_n$ with $b_n \in A$.

		Consequently, $J$ is the directed union of the right ideals $dA$ with $d \in J$ positive and injective. Each $dA$ is a free right module of rank one, so $J$ is flat on the right. Also $J^2=J$: polar decomposition writes $a=u|a|^{1/2}|a|^{1/2}$ with both factors in $J$. Both ideals are closed under positive square roots; for $S_{0+}$ this follows directly from the definition by singular values. The Tor calculation in \Cref{lem:Almost_Restriction_FF} now shows that $A \to A/J$ is a homological epimorphism.

		For a positive injective $d \in J$, put
		\[
			T_d = \bigcup_{n\geq 0} d^{1/2^n}A \subseteq J.
		\]
		Under the identifications $d^{1/2^n}A \cong A$, this is the sequential colimit of free modules with transition maps given by left multiplication by $d^{1/2^{n+1}} \in J$. The family of submodules $T_d$ is directed: apply the factorization fact to all the roots of two given operators to obtain $e$ with both submodules contained in $eA \subseteq T_e$. It covers $J$, so $J$ is a filtered colimit of such telescopes.

		Let $L = - \otimes_A^{\bL} A/J$, and let $C = \ker(L)$. Its inclusion $i\colon C \to D$ has the colimit-preserving right adjoint $r(M)=M\otimes_A^{\bL}J$. This follows by tensoring the fiber sequence $J \to A \to A/J$ with $M$ and using that $L$ is a localization. Since $D$ is generated as a stable category under colimits by $A$, the category $C$ is generated in the same sense by $r(A)=J$. Every $J$-transition telescope lies in $C$, because applying $L$ gives a sequence with zero transition maps, whose colimit is zero. Conversely, the preceding expression for $J$ as a filtered colimit of the $T_d$ puts $J$ in the stable subcategory generated by these telescopes. Thus $C=C_J$.

		The completed Hilbert tensor product extends to a colimit-preserving multiplication $\mu_D\colon D\otimes D\to D$. It preserves compact objects: it sends pairs of Hilbert spaces to Hilbert spaces, and compact objects are obtained from these by finite stable operations and retracts. Thus $\mu_D^R$ preserves colimits. Tensor products of two $J$-transition telescopes again have $J$-transitions, using the cofinal diagonal in $\N\times\N$. For $S_{0+}$, the relevant assertion follows from
		\[
			\sum_{i,j}(s_i(a)s_j(b))^p
			=\left(\sum_i s_i(a)^p\right)\left(\sum_j s_j(b)^p\right)<\infty
		\]
		for every $p>0$; for compact operators it follows by finite-rank approximation. Hence the tensor product restricts to $C_J$. The unit $\C$ belongs to $C_J$, since its identity has finite rank, and is compact there because it is compact in $D$ and the inclusion preserves colimits.

		Finally, $i\otimes i$ has right adjoint $r\otimes r$: the adjunction $i\dashv r$ is an adjunction between colimit-preserving functors, so it can be tensored. The restricted multiplication has right adjoint
		\[
			\mu_{C_J}^R = (r\otimes r)\mu_D^R i,
		\]
		as follows by applying these adjunctions to mapping spectra. All three functors preserve colimits. The unit functor is strongly continuous by compactness of $\C$, and $C_J$ is dualizable by \Cref{prop:Kernel_Homological_Epi}. This proves the last assertion.
	\end{proof}

	\section{Continuous posets}
	\label[section]{sec:Continuous_Posets}

	% Source: Talk 3 (Max Glaser, 27.10.2026); proofs completed with generative AI (Claude)
	We now study a class of non-stable compactly assembled categories in detail, which was already briefly discussed in \Cref{ex:Continuous_Posets}. They will reappear in \Cref{chap:Sheaves}, where we study sheaves on continuous posets. See \cite{Gierz_Continuous_Lattices} for a comprehensive treatment.

	Throughout, let $P$ be a partially ordered set, regarded as a category.

	\begin{lemma}
		\label[lemma]{lem:Ind_Of_Poset}
		The category $\Ind(P)$ is (equivalent to) a partially ordered set. Explicitly, every object of $\Ind(P)$ is of the form $\quotcolim_{s \in S} s$ for a directed subset $S \subseteq P$, and for directed subsets $S,T \subseteq P$ we have
		\begin{align*}
			\Map_{\Ind(P)}\Big(\quotcolim_{s \in S} s, \quotcolim_{t \in T} t\Big) &= \lim_{s \in S\catop}\colim_{t \in T}\Map_P(s,t) \\
			&\simeq \begin{cases} * & \text{if for all } s \in S \text{ there is } t \in T \text{ with } s \leq t, \\ \emptyset & \text{otherwise.}\end{cases}
		\end{align*}
		Hence $\Ind(P)$ is equivalent to the poset of \emph{ideals} of $P$, i.e.\ non-empty, directed, downward closed subsets $J \subseteq P$, ordered by inclusion, via $S \mapsto {\downarrow}S := \{p \in P \mid p \leq s \text{ for some } s \in S\}$.
	\end{lemma}
	\prooffor{lem:Ind_Of_Poset}
	\begin{proof}
		By \Cref{thm:Filtered_Posets_Cofinal}, every object of $\Ind(P)$ is a colimit of representables indexed by a filtered poset $I$, i.e.\ of the form $\quotcolim_{i \in I} p_i$ for a functor $p\colon I \to P$. The image $S = p(I)$ is a directed subset, and the map $I \to S$ is cofinal by \Cref{lem:Cofinality_Posets} (for $s \in S$, the set $\{i \mid s \leq p_i\}$ is non-empty and upward closed in the filtered poset $I$, hence filtered). This proves the first claim. The formula for mapping spaces follows from \Cref{constr:Ind_Large}: each $\Map_P(s,t)$ is either empty or contractible, filtered colimits of such spaces are again empty or contractible (and contractible if and only if one of the terms is), and a limit of such spaces is contractible if and only if all terms are. In particular, all mapping spaces in $\Ind(P)$ are empty or contractible, so $\Ind(P)$ is a poset. Finally, the condition in the formula only depends on ${\downarrow}S$ and ${\downarrow}T$: it says precisely that ${\downarrow}S \subseteq {\downarrow}T$.
	\end{proof}

	Suppose now that $P$ admits directed suprema. Then the colimit functor $\colim\colon \Ind(P) \to P$ sends a directed subset $S$ to $\sup S$, and the Yoneda embedding sends $b$ to ${\downarrow}b$.

	\begin{lemma}
		\label[lemma]{lem:Posets_Accessible}
		Let $P$ be a small partially ordered set admitting directed suprema. Then $P$ is accessible.
	\end{lemma}
	\prooffor{lem:Posets_Accessible}
	\begin{proof}
		Let $\kappa$ be a regular cardinal with $\kappa > |P|$. Let $D \subseteq P$ be a $\kappa$-directed subset. Since $|D| < \kappa$, the set $D$ itself has an upper bound in $D$ (\Cref{lem:Kappa_Filtered_Posets}), i.e.\ $D$ has a largest element. It follows that every $\kappa$-filtered diagram in $P$ has a cofinal constant subdiagram, so that $\kappa$-filtered colimits in $P$ are ``trivial'': every object of $P$ is $\kappa$-compact, and the Yoneda embedding $P \to \Ind_{\kappa}(P)$ is an equivalence, as every $\kappa$-filtered colimit of representable presheaves on $P$ is representable. Hence $P \simeq \Ind_{\kappa}(P^{\kappa})$ is $\kappa$-accessible.
	\end{proof}

	\begin{definition}
		\label[definition]{def:Continuous_Poset}
		A partially ordered set $P$ is called \emph{continuous} if it admits directed suprema and the functor $\colim = \sup\colon \Ind(P) \to P$ admits a left adjoint $\hat{y}\colon P \to \Ind(P)$.
	\end{definition}

	By \Cref{lem:Posets_Accessible}, a small poset is continuous if and only if it is compactly assembled. Unwinding the definitions using \Cref{lem:Ind_Of_Poset}, the existence of the left adjoint $\hat{y}$ means: for every $b \in P$ there is a directed subset $\hat{y}(b) \subseteq P$ such that for every directed subset $T \subseteq P$,
	\begin{equation}
		\label{eq:Yhat_Poset}
		b \leq \sup T \iff \text{for all } s \in \hat{y}(b) \text{ there exists } t \in T \text{ with } s \leq t.
	\end{equation}

	\begin{definition}
		\label[definition]{def:Way_Below}
		Let $P$ be a poset admitting directed suprema, and let $a,b \in P$. We say that $a$ is \emph{way below} $b$, written $a \ll b$, if for every directed subset $T \subseteq P$ with $b \leq \sup T$ there exists $t \in T$ with $a \leq t$. We write
		\[
			P_{\ll b} := \{a \in P \mid a \ll b\}.
		\]
	\end{definition}

	\begin{remark}
		\label[remark]{rem:Way_Below_Basics}
		\leavevmode
		\begin{enumerate}
			\item If $a \ll b$, then $a \leq b$: apply the definition to $T = \{b\}$.
			\item If $a' \leq a \ll b \leq b'$, then $a' \ll b'$.
			\item An element $b$ is compact in $P$ (i.e.\ $\Map_P(b,-)$ preserves filtered colimits) if and only if $b \ll b$.
			\item If $P$ is continuous, then $\hat{y}(b) \subseteq P_{\ll b}$ for every $b$. Indeed, let $s \in \hat{y}(b)$ and let $T$ be directed with $b \leq \sup T$. By \eqref{eq:Yhat_Poset}, there exists $t \in T$ with $s \leq t$, so $s \ll b$.
		\end{enumerate}
	\end{remark}

	\begin{proposition}
		\label[proposition]{prop:Continuous_Poset_Characterization}
		Let $P$ be a poset admitting directed suprema. Then $P$ is continuous if and only if for every $b \in P$ the subset $P_{\ll b}$ is directed and $\sup P_{\ll b} = b$. In this case, $\hat{y}(b) = P_{\ll b}$, viewed as an object of $\Ind(P)$.
	\end{proposition}
	\prooffor{prop:Continuous_Poset_Characterization}
	\begin{proof}
		Suppose that each $P_{\ll b}$ is directed with supremum $b$. We claim that $b \mapsto P_{\ll b}$ satisfies \eqref{eq:Yhat_Poset}, and hence defines a left adjoint to $\sup$ (a monotone map, by \Cref{rem:Way_Below_Basics}(2)). If $b \leq \sup T$, then by definition every $a \ll b$ satisfies $a \leq t$ for some $t \in T$. Conversely, if every $a \in P_{\ll b}$ is bounded by some element of $T$, then $b = \sup P_{\ll b} \leq \sup T$.

		Conversely, suppose that $P$ is continuous. We claim that
		\[
			P_{\ll b} = \{a \in P \mid a \leq s \text{ for some } s \in \hat{y}(b)\} = {\downarrow}\hat{y}(b).
		\]
		For ``$\supseteq$'', note that $\hat{y}(b) \subseteq P_{\ll b}$ by \Cref{rem:Way_Below_Basics}(4), and $P_{\ll b}$ is downward closed by \Cref{rem:Way_Below_Basics}(2). For ``$\subseteq$'', note that $\hat{y}$ is fully faithful (\Cref{lem:Adjoint_Triples}), so the counit $\sup\hat{y}(b) \leq b$ is an equality, i.e.\ $b = \sup \hat{y}(b)$. Applying the definition of $a \ll b$ to the directed set $T = \hat{y}(b)$ gives $s \in \hat{y}(b)$ with $a \leq s$. Since $\hat{y}(b)$ is directed, so is ${\downarrow}\hat{y}(b) = P_{\ll b}$, and $\sup P_{\ll b} = \sup \hat{y}(b) = b$. Finally, $P_{\ll b}$ and $\hat{y}(b)$ define the same object of $\Ind(P)$ by \Cref{lem:Ind_Of_Poset}.
	\end{proof}

	\begin{example}
		\label[example]{ex:Continuous_Poset_Reals}
		The poset $P = \R \cup \{\infty\}$ (with its usual order) admits directed suprema: the supremum of a non-empty subset is its usual supremum if the subset is bounded, and $\infty$ otherwise. We claim that $a \ll b$ if and only if $a < b$. If $a < b$ and $T$ is directed with $b \leq \sup T$, then $\sup T > a$, so some $t \in T$ satisfies $t > a$. Conversely, suppose $a \ll b$; then $a \leq b$ by \Cref{rem:Way_Below_Basics}(1). If $a = b$, consider the directed set $T = \{b - 1/n \mid n \geq 1\}$ (for $b \in \R$) resp.\ $T = \N$ (for $b = \infty$): it has supremum $b$, but no element $t \in T$ satisfies $b \leq t$. Hence $a < b$. It follows that $P_{\ll b} = \{a \mid a < b\}$ is directed with supremum $b$, so $P$ is continuous by \Cref{prop:Continuous_Poset_Characterization}, with $\hat{y}(b) = \quotcolim_{a < b} a$. Note that $P$ has no compact elements.
	\end{example}

	\begin{example}
		\label[example]{ex:Continuous_Poset_Opens}
		Let $X$ be a locally compact Hausdorff space and let $\Open(X)$ be the poset of open subsets of $X$, ordered by inclusion. Directed suprema in $\Open(X)$ are unions. We claim that
		\[
			U \ll V \iff \overline{U} \text{ is compact and } \overline{U} \subseteq V,
		\]
		in which case we write $U \Subset V$. Suppose first that $\overline{U}$ is compact and $\overline{U} \subseteq V$, and let $(W_i)_{i \in I}$ be a directed family of open subsets with $V \subseteq \bigcup_i W_i$. Then the $W_i$ cover the compact set $\overline{U}$, so finitely many of them do, and since the family is directed, a single $W_i$ contains $\overline{U} \supseteq U$. Conversely, let $U \ll V$. Since $X$ is locally compact Hausdorff, every point $x \in V$ has an open neighbourhood $W$ with $\overline{W}$ compact and $\overline{W} \subseteq V$. Hence $V$ is the union of the family $\{W \in \Open(X) \mid W \Subset V\}$, which is directed since $W_1 \cup W_2 \Subset V$ if $W_1,W_2 \Subset V$. By assumption, $U \subseteq W$ for some $W \Subset V$, and then $\overline{U} \subseteq \overline{W}$ is a closed subset of a compact set, hence compact, and $\overline{U} \subseteq \overline{W} \subseteq V$.

		The argument shows moreover that $\Open(X)_{\ll V}$ is directed with union $V$, so $\Open(X)$ is a continuous poset by \Cref{prop:Continuous_Poset_Characterization}, with $\hat{y}(V) = \quotcolim_{U \Subset V} U$. Its compact elements are the open subsets $U$ with $U \Subset U$, i.e.\ the compact open subsets of $X$.
	\end{example}

	\begin{remark}
		\label[remark]{rem:Continuous_Posets_Sheaves}
		If $X$ is a locally compact Hausdorff space, then the morphism $h_U \to h_V$ of representable sheaves on $X$ is a compact morphism in $\Shv(X;\Sp)$ whenever $U \Subset V$ (Mathew, Lecture 3). Since every open $V$ is the union of a sequence $U_0 \Subset U_1 \Subset \cdots$ if $X$ is second countable, this suggests (and, suitably generalized, proves) that $\Shv(X;\Sp)$ is generated by sequential colimits along compact morphisms and hence dualizable, cf.\ \Cref{rem:Omega_1_Compact_Compact_Morphisms}. We will return to this in \Cref{chap:Sheaves}.
	\end{remark}

	\section{Extensions and semiorthogonal gluing}
	\label[section]{sec:Extensions_And_Gluing}

	We first show that semiorthogonal decompositions and gluings behave well with respect to dualizability. We then explain why a general extension behaves differently. The Tate-object classification at the end of this section is a more specialized application and may be omitted on a first reading.

	% Source: Talk 6 (Jacopo Mazzei, 17.11.2026); Efimov, Propositions 2.76 and 2.77
	\begin{proposition}[{\cite[Proposition~2.76]{Efimov_Large_Categories}}]
		\label[proposition]{prop:SOD_Dualizable}
		Let $E = \lra{C,D}$ be a semiorthogonal decomposition of presentable stable categories in the sense of \Cref{def:Semiorthogonal_Decomposition}. Then $E$ is dualizable if and only if $C$ and $D$ are dualizable. In this case $E^{\omega} = \lra{C^{\omega},D^{\omega}}$.
	\end{proposition}
	\prooffor{prop:SOD_Dualizable}
	\begin{proof}
		The projections $i_C^L$ and $i_D^R$ preserve colimits, so $C$ and $D$ are retracts of $E$ in $\PrL_{\st}$. Thus dualizability of $E$ implies that of $C$ and $D$.

		Conversely, if $C$ and $D$ are dualizable, then $E$ is dualizable by \Cref{prop:Generation_Dualizable}, since the strongly continuous inclusions $i_C$ and $i_D$ jointly generate $E$.

		The inclusions preserve compact objects. The projection $i_C^L$ does too, since its right adjoint $i_C$ preserves colimits. For compact $e \in E$, the cofiber sequence $i_Di_D^Re \to e \to i_Ci_C^Le$ therefore shows that $i_Di_D^Re$ is compact. Since $i_D$ is fully faithful and preserves filtered colimits, $i_D^Re$ is compact in $D$. This restricts the semiorthogonal decomposition to compact objects.
	\end{proof}

	\begin{construction}[Gluing along a functor]
		\label[construction]{constr:Semiorthogonal_Gluing}
		Let $F\colon D \to C$ be an accessible exact functor between presentable stable categories. Let $E$ be the category of triples $(x,y,u)$ with $x \in C$, $y \in D$, and $u\colon x \to F(y)$. The same comma category occurs in the construction of a recollement in \href{https://animarium.org/tag/026K/}{\emph{Lectures on Topos Theory}, Construction~5.48}. This is a presentable stable category \cite[Proposition~1.2]{Efimov_Large_Categories}; colimits are computed on $x$ and $y$, using the natural maps $\colim F(y_j) \to F(\colim y_j)$.

		The fully faithful colimit-preserving functors
		\[
			i_C(c) = (c,0,0), \qquad i_D(d) = (0,d,0)
		\]
		have right adjoints $i_C^R(x,y,u) = \fib(u)$ and $i_D^R(x,y,u) = y$. Moreover, $i_C^L(x,y,u) = x$ is left adjoint to $i_C$, and every object fits into a cofiber sequence
		\[
			i_D(y) \to (x,y,u) \to i_C(x).
		\]
		We have $\hom_E(i_D(d),i_C(c)) \simeq 0$. The functor $i_D$ is always strongly continuous, since its right adjoint $(x,y,u) \mapsto y$ preserves colimits.

		In \cite[Definition~1.1]{Efimov_Large_Categories}, the inclusion of $C$ is instead defined as $c \mapsto (c[-1],0,0)$. Its right adjoint is then $(x,y,u) \mapsto \cofib(u)$, and its composite with $i_D$ is $F$ itself, while in our convention $i_C^Ri_D \simeq \Omega F$. This shift plays no role in what follows.
	\end{construction}

	% Source: Talk 6 (Jacopo Mazzei, 17.11.2026); Efimov, Proposition 2.77
	\begin{proposition}[{\cite[Proposition~2.77]{Efimov_Large_Categories}}]
		\label[proposition]{prop:Gluing_Dualizable}
		Let $F\colon D \to C$ be an accessible exact functor between dualizable categories, and let $E$ be the gluing of \Cref{constr:Semiorthogonal_Gluing}. Then the following are equivalent:
		\begin{enumerate}
			\item $F$ preserves colimits;
			\item $E$ is dualizable, and $i_C$ and $i_D$ are strongly continuous.
		\end{enumerate}
		In this case, $i_C$ and $i_D$ exhibit a semiorthogonal decomposition $E = \lra{C,D}$.
	\end{proposition}
	\prooffor{prop:Gluing_Dualizable}
	\begin{proof}
		(2) $\implies$ (1): If $i_C$ is strongly continuous, then $i_C^R$ preserves colimits, and hence so does $i_C^Ri_D \simeq \Omega F$. So $F$ preserves colimits. This does not use the dualizability of $E$.

		(1) $\implies$ (2): If $F$ preserves colimits, then $i_C^R(x,y,u) = \fib(u\colon x \to F(y))$ preserves colimits, so $i_C$ is strongly continuous. The functor $i_D$ is always strongly continuous. The functors $i_C$ and $i_D$ are fully faithful, we have $\hom_E(i_D(d),i_C(c)) \simeq 0$, and every object of $E$ is an extension of an object of $i_C(C)$ by one of $i_D(D)$. So we obtain a semiorthogonal decomposition $E = \lra{C,D}$ in the sense of \Cref{def:Semiorthogonal_Decomposition}, and $E$ is dualizable by \Cref{prop:SOD_Dualizable}.
	\end{proof}

	The strong continuity condition in \Cref{prop:Gluing_Dualizable} is essential: a general strongly continuous localization sequence with dualizable outer terms need not have dualizable middle term (\Cref{ex:Extension_Not_Dualizable}).

	% Source: Talk 6 (Jacopo Mazzei, 17.11.2026); Efimov, Proposition 2.78; proof completed with generative AI (Claude)
	Semiorthogonal decompositions are inherited by categories of strongly continuous functors. For presentable stable categories $T$ and $E$, we write $\Fun^{\mathrm{LL}}(T,E) \subseteq \Fun^{\colim}(T,E)$ for the full subcategory of strongly continuous functors. It is a stable subcategory: shifts of strongly continuous functors are strongly continuous, and the right adjoint of the cofiber of a map $G \to G'$ of colimit-preserving functors is the fiber of the induced map $G'^R \to G^R$ of right adjoints, which preserves colimits if $G^R$ and $G'^R$ do. In the following, we use the notion of a semiorthogonal decomposition of a (not necessarily presentable) stable category from \Cref{rem:SOD_Small}.

	\begin{proposition}[{\cite[Proposition~2.78]{Efimov_Large_Categories}}]
		\label[proposition]{prop:SOD_Functor_Categories}
		Let $E = \lra{i_1(C),i_2(D)}$ be a semiorthogonal decomposition of a presentable stable category, let $\pi_1 := i_1^L\colon E \to C$ and $\pi_2 := i_2^R\colon E \to D$ be the projections, and let $T$ be a presentable stable category.
		\begin{enumerate}
			\item A colimit-preserving functor $G\colon T \to E$ is strongly continuous if and only if $\pi_1G$ and $\pi_2G$ are strongly continuous. Postcomposition with $i_1$ and $i_2$ gives a semiorthogonal decomposition
			\[
				\Fun^{\mathrm{LL}}(T,E) = \lra{\Fun^{\mathrm{LL}}(T,C),\Fun^{\mathrm{LL}}(T,D)}.
			\]
			\item A colimit-preserving functor $G\colon E \to T$ is strongly continuous if and only if $Gi_1$ and $Gi_2$ are strongly continuous. Precomposition with $\pi_1$ and $\pi_2$ gives a semiorthogonal decomposition
			\[
				\Fun^{\mathrm{LL}}(E,T) = \lra{\Fun^{\mathrm{LL}}(C,T),\Fun^{\mathrm{LL}}(D,T)}.
			\]
		\end{enumerate}
	\end{proposition}
	\prooffor{prop:SOD_Functor_Categories}
	\begin{proof}
		Recall from \Cref{prop:SOD_Characterization} that $i_1$, $i_2$, $\pi_1$ and $\pi_2$ are strongly continuous, that $\pi_1i_1 \simeq \id$, $\pi_2i_2 \simeq \id$ and $\pi_1i_2 \simeq 0$, and that there is a cofiber sequence $i_2\pi_2 \to \id_E \to i_1\pi_1$. In particular, composition with any of these four functors preserves strong continuity.

		(1) If $G$ is strongly continuous, then so are $\pi_1G$ and $\pi_2G$. Conversely, if $\pi_1G$ and $\pi_2G$ are strongly continuous, then so are $i_1\pi_1G$ and $i_2\pi_2G$, and hence so is $G$, by the cofiber sequence $i_2\pi_2G \to G \to i_1\pi_1G$ and the stability of $\Fun^{\mathrm{LL}}(T,E)$. Let $j_1 := i_1 \circ -$ and $j_2 := i_2 \circ -$. The pointwise adjunctions $\pi_1 \dashv i_1$ and $i_2 \dashv \pi_2$ show that $\pi_1 \circ -$ is left adjoint to $j_1$ and that $\pi_2 \circ -$ is right adjoint to $j_2$. Hence $j_1$ and $j_2$ are fully faithful, since $\pi_1i_1 \simeq \id$ and $\pi_2i_2 \simeq \id$, and
		\[
			\Map(j_2(H'),j_1(H)) \simeq \Map(\pi_1i_2H',H) \simeq *
		\]
		for $H \in \Fun^{\mathrm{LL}}(T,C)$ and $H' \in \Fun^{\mathrm{LL}}(T,D)$. Finally, every $G$ is an extension $i_2\pi_2G \to G \to i_1\pi_1G$ of an object in the image of $j_1$ by one in the image of $j_2$.

		(2) The criterion for strong continuity follows from \Cref{lem:Automatic_Strong_Continuity}, applied to the functors $i_1$ and $i_2$, whose essential images generate $E$; the converse holds since $i_1$ and $i_2$ are strongly continuous. Let $j_1' := - \circ \pi_1$ and $j_2' := - \circ \pi_2$. For an adjunction $L \dashv R$, precomposition with $R$ is left adjoint to precomposition with $L$. Hence $j_1'$ has the left adjoint $- \circ i_1$, and $j_2'$ has the right adjoint $- \circ i_2$. Again $j_1'$ and $j_2'$ are fully faithful, as $\pi_1i_1 \simeq \id$ and $\pi_2i_2 \simeq \id$. Moreover,
		\[
			\Map(j_2'(H'),j_1'(H)) \simeq \Map(H',H\pi_1i_2) \simeq *
		\]
		for $H \in \Fun^{\mathrm{LL}}(C,T)$ and $H' \in \Fun^{\mathrm{LL}}(D,T)$. Finally, precomposing $G$ with the cofiber sequence $i_2\pi_2 \to \id_E \to i_1\pi_1$ gives a cofiber sequence $j_2'(Gi_2) \to G \to j_1'(Gi_1)$.
	\end{proof}

	\begin{warning}
		\label[warning]{warn:Extensions_Not_Dualizable}
		The converse of \Cref{prop:Localization_Sequence_Dualizable} fails: if $K$ and $E$ are dualizable, then $D$ need not be dualizable, even if $K$ and $E$ are compactly generated. An explicit counterexample is given in \Cref{ex:Extension_Not_Dualizable} below. However, dualizable categories are closed under pullbacks along strongly continuous localizations, see \Cref{prop:Dualizable_Pullbacks}(2).
	\end{warning}

	% Source: Efimov, Lecture 4 (MPIM Bonn, 10.09.2024)
	Let us take a closer look at extensions of compactly generated categories. Consider a strongly continuous localization sequence
	\[
		\Ind(A) \xrightarrow{F} D \xrightarrow{G} \Ind(B)
	\]
	in $\PrL_{\st}$, with $A,B \in \Cat_{\perf}$. Since $\hom_D(Fa,G^Rb) \simeq \hom_{\Ind(B)}(GFa,b) \simeq 0$, the category $D$ admits a semiorthogonal decomposition
	\[
		D = \lra{G^R(\Ind(B)),F(\Ind(A))}
	\]
	in the ordinary sense of semiorthogonality and generation. It need not satisfy \Cref{def:Semiorthogonal_Decomposition}: although $G^R$ preserves colimits, its right adjoint need not do so. Thus \Cref{prop:SOD_Dualizable} does not apply to an arbitrary extension. Since $G$ is strongly continuous, $G^R$ admits a right adjoint $G^{RR}$, and the extension $D$ is encoded by the accessible functor
	\[
		G^{RR} \circ F\colon \Ind(A) \to \Ind(B).
	\]
	An accessible exact functor $\Phi\colon \Ind(A) \to \Ind(B)$ is in turn determined by the functor $\Psi\colon B \to \Pro(\Ind(A))$ sending $b \in B$ to the pro-object corepresenting the accessible functor $x \mapsto \hom_{\Ind(B)}(b,\Phi(x))$.

	Following Hennion \cite{Hennion_Tate}, we denote by $\mathrm{Tate}^{\mathrm{el}}(A) \subseteq \Pro(\Ind(A))$ the stable subcategory generated by $\Ind(A)$ and $\Pro(A)$ (the \emph{elementary Tate objects}), and by $\mathrm{Tate}(A) \subseteq \Pro(\Ind(A))$ its idempotent completion. By \cite[Theorem~3]{Hennion_Tate}, an object of $\Pro(\Ind(A))$ lies in $\mathrm{Tate}^{\mathrm{el}}(A)$ if and only if it is an extension $u \to t \to v$ with $u \in \Pro(A)$ and $v \in \Ind(A)$.

	\begin{proposition}[{Efimov \cite[Proposition~6.3]{Efimov_Large_Categories}}]
		\label[proposition]{prop:Extension_Dualizable_Tate}
		In the above situation, let $\Psi\colon B \to \Pro(\Ind(A))$ be the functor corresponding to $\Phi = G^{RR} \circ F$. Then the following are equivalent:
		\begin{enumerate}
			\item $D$ is dualizable;
			\item $D$ is compactly generated;
			\item The essential image of $\Psi$ is contained in $\mathrm{Tate}(A)$.
		\end{enumerate}
	\end{proposition}
	\prooffor{prop:Extension_Dualizable_Tate}
	\begin{proof}
		We first show that $D$ is compactly generated if and only if for every $b \in B$ there is a compact object $d \in D^{\omega}$ with $G(d) \simeq b \oplus \Sigma b$. If $D$ is compactly generated, this follows from \Cref{thm:Thomason_Neeman}(2) and \Cref{rem:Lifting_Objects_Quotient}. Conversely, suppose such lifts $d_b$ exist. Since $F$ is strongly continuous, the objects $F(a)$, $a \in A$, are compact, and we claim that together with the $d_b$ they generate $D$. Indeed, if $x \in D$ is right orthogonal to all of them, then $F^R(x) = 0$, so $x \simeq G^RG(x)$, and $0 \simeq \hom_D(d_b, G^RG(x)) \simeq \hom_{\Ind(B)}(b \oplus \Sigma b, G(x))$ for all $b$, so $G(x) = 0$ and hence $x = 0$.

		(1) $\Rightarrow$ (2): If $D$ is dualizable, then the required lifts exist by \cite[Proposition~2.10]{Efimov_Large_Categories} (which reduces to the compactly generated case by embedding $D$ strongly continuously into a compactly generated category). (2) $\Rightarrow$ (1) is clear.

		(2) $\Leftrightarrow$ (3): Let $d \in D$, and put $z := G(d)$ and $w := F^R(d)$, so that we have a fiber sequence $F(w) \to d \to G^R(z)$. Applying $\hom_D(-,F(w'))$ for $w' \in \Ind(A)$ and using $\hom_D(G^R(z),F(w')) \simeq \hom_{\Ind(B)}(z,\Phi(w'))$, we obtain a fiber sequence of functors $\Ind(A) \to \Sp$
		\[
			\hom_{\Ind(B)}(z,\Phi(-)) \to \hom_D(d,F(-)) \to \hom_{\Ind(A)}(w,-).
		\]
		Every filtered colimit in $D$ is an extension of a filtered colimit in $F(\Ind(A))$ by one in $G^R(\Ind(B))$, as $F^R$ and $G$ preserve colimits. Since $\hom_D(d,G^R(-)) \simeq \hom_{\Ind(B)}(z,-)$, it follows that $d$ is compact if and only if $z$ is compact and $\hom_D(d,F(-))$ preserves filtered colimits. Now identify exact accessible functors $\Ind(A) \to \Sp$ with $\Pro(\Ind(A))\catop$; under this identification, the functors preserving filtered colimits correspond to $\Pro(A)$. The first term of the fiber sequence corresponds to $\Psi(z)$ by definition, and the last one to $w \in \Ind(A)$. So $d$ is compact if and only if $z \in B$ and there is a cofiber sequence $u \to \Psi(z) \to \Sigma w$ in $\Pro(\Ind(A))$ with $u \in \Pro(A)$. Conversely, given $z \in B$ and a cofiber sequence $u \to \Psi(z) \to v$ with $u \in \Pro(A)$ and $v \in \Ind(A)$, the map $\Psi(z) \to v$ corresponds to a map $z \to \Phi(v)$, i.e.\ to a map $G^R(z) \to F(v)$, and its fiber $d$ is a compact object with $G(d) \simeq z$ by the same computation. Using Hennion's theorem, we conclude that $z \in B$ lifts to $D^{\omega}$ if and only if $\Psi(z) \in \mathrm{Tate}^{\mathrm{el}}(A)$.

		Combining this with the first paragraph, $D$ is compactly generated if and only if $\Psi(b) \oplus \Sigma\Psi(b) \in \mathrm{Tate}^{\mathrm{el}}(A)$ for all $b \in B$. This implies $\Psi(b) \in \mathrm{Tate}(A)$. Conversely, if $\Psi(b) \in \mathrm{Tate}(A)$, then $\Psi(b) \oplus \Sigma\Psi(b) \in \mathrm{Tate}^{\mathrm{el}}(A)$ by Thomason's classification of dense subcategories, as in \Cref{rem:Lifting_Objects_Quotient}.
	\end{proof}

	\begin{remark}
		\label[remark]{rem:Tate_Classification_Of_Extensions}
		Applied to compactly generated $D$, the proof shows that short exact sequences $A \to D_0 \to B$ in $\Cat_{\perf}$ correspond to exact functors $B \to \mathrm{Tate}(A)$; this classification first appeared in \cite[Theorem~A.5.1 and Corollary~A.5.10]{Hermitian_II}, see also \cite[Theorem~6.4]{Efimov_Large_Categories}.
	\end{remark}

	\begin{example}[Efimov]
		\label[example]{ex:Extension_Not_Dualizable}
		Let $k$ be a field, and consider the category
		\[
			D := \Big\{(V,W,\phi) \;\Big|\; V,W \in \D(k),\ \phi\colon \bigoplus_{\N}V \to \bigoplus_{\N}W\Big\}.
		\]
		The functors $W \mapsto (0,W,0)$ and $(V,W,\phi) \mapsto V$ exhibit $D$ as an extension of $\D(k)$ by $\D(k)$, both of which are compactly generated. Nevertheless, $D$ is not dualizable.

		\emph{Proof} (following \cite[Proposition~6.8]{Efimov_Large_Categories})\emph{.} Here $\Phi(W) = \prod_{\N}\bigoplus_{\N}W$, since $\hom(\bigoplus_{\N}V,\bigoplus_{\N}W) \simeq \hom(V,\prod_{\N}\bigoplus_{\N}W)$. By \Cref{prop:Extension_Dualizable_Tate}, it suffices to show that $\Psi(k) \notin \mathrm{Tate}(\Perf(k))$, where $\Psi(k)$ is the pro-object corepresenting $W \mapsto \hom(k,\Phi(W)) = \prod_{\N}\bigoplus_{\N}W$. Using (AB6) in $\D(k)$, we may rewrite this functor as
		\[
			\prod_{\N}\bigoplus_{\N}W \simeq \colim_{f\colon \N \to \N}\ \prod_{n \in \N}\bigoplus_{i=0}^{f(n)-1}W \simeq \colim_{f\colon \N \to \N}\hom(X_f,W), \qquad X_f := \bigoplus_{n \in \N}k^{f(n)},
		\]
		where $\N^{\N}$ is ordered pointwise, and for $f \leq g$ the transition map is induced by the projection $X_g \to X_f$. Hence $\Psi(k) \simeq \text{``}\!\lim\nolimits_f\!\text{} X_f$, a pro-system in which all transition maps are split epimorphisms.

		Let $q\colon \D(k) \to \Calk(\Perf(k))$ be the projection to $\Calk(\Perf(k)) := \bigcup_{\kappa}\Calk_{\kappa}(\Perf(k))$ (the transition functors are fully faithful by \Cref{rem:Calkin_Mapping_Spectra}); this is an idempotent complete stable category, and $q(x) = 0$ if and only if $x \in \Perf(k)$. Suppose that $\Psi(k) \in \mathrm{Tate}(\Perf(k))$. Then $\Psi(k) \oplus \Sigma\Psi(k) \in \mathrm{Tate}^{\mathrm{el}}(\Perf(k))$, so there is a cofiber sequence $u \to \Psi(k) \oplus \Sigma\Psi(k) \to v$ with $u \in \Pro(\Perf(k))$ and $v \in \D(k)$. Applying $\Pro(q)$ kills $u$, so $\Pro(q)(\Psi(k)) = \text{``}\!\lim\nolimits_f\!\text{} q(X_f)$ is a retract of the constant pro-object $q(v)$, and hence is itself isomorphic to a constant pro-object $c$.

		We now use the following elementary observation: if a pro-system $(Y_f)_f$ indexed by a directed poset, with split epimorphisms as transition maps $t_{gf}\colon Y_g \to Y_f$, is isomorphic to a constant pro-object $c$, then there is $f_0$ such that $t_{gf}$ is an isomorphism for all $f_0 \leq f \leq g$. Indeed, the isomorphism is given by maps $p_f\colon Y_f \to c$ for $f \geq f_0$ (compatible with the $t_{gf}$) and a compatible family $s_f\colon c \to Y_f$, such that $p_fs_f = \id_c$ for $f \geq f_0$, and for every $f$ there is $g \geq f,f_0$ with $s_fp_g = t_{gf}$. If $\sigma$ is a section of $t_{gf}$, then $s_f(p_g\sigma) = \id$. So for $f \geq f_0$, the map $s_f$ is both a split mono and a split epi, hence an isomorphism, and then $t_{gf} = s_fs_g^{-1}$ is an isomorphism for $f_0 \leq f \leq g$.

		Applied to $Y_f = q(X_f)$, this gives a contradiction: for every $f$, the fiber of $X_{f+1} \to X_f$ is $\bigoplus_{\N}k \notin \Perf(k)$, so $q(X_{f+1}) \to q(X_f)$ is not an isomorphism.
	\end{example}


		\chapter{Rigidity and nuclear objects}
	\label[chapter]{chap:Rigidity_And_Nuclear_Objects}

	This chapter develops the symmetric monoidal counterpart of compact assembly. Trace-class morphisms take the role of compact morphisms, and nuclear objects are built from compact objects along trace-class maps. We first discuss rigidity and self-duality, then construct the rigidification, study nuclear objects and examine the analytic examples. The interpolation results of \Cref{chap:Compact_Morphisms_AB6} will be used in the rigidification argument.

	\section{Rigid symmetric monoidal categories}
	\label[section]{sec:Rigid_Categories}

	% Source: Talk 3 (Max Glaser, 27.10.2026); Clausen, Lecture 3 (19.07.2023); literature (Gaitsgory--Rozenblyum, Ramzi, Efimov), proofs partly completed with generative AI (Claude)
	A rich source of dualizable categories are \emph{rigid} symmetric monoidal categories, introduced by Gaitsgory--Rozenblyum \cite[Chapter~1, Section~9]{GaitsgoryRozenbluym2017StudyDAG}. We have already met them in \Cref{def:Rigid_Category}; we now give a more detailed treatment.

	\begin{definition}[{\cite[Chapter~1, Definition~9.1.2]{GaitsgoryRozenbluym2017StudyDAG}}]
		\label[definition]{def:Rigid_Category_Detailed}
		Let $C \in \CAlg(\PrL_{\st})$ be a presentably symmetric monoidal stable category, with unit $\unit_C$ and multiplication $\mu\colon C \otimes C \to C$. Then $C$ is called \emph{rigid} if
		\begin{enumerate}
			\item the unit $\unit_C$ is compact in $C$, or equivalently, the unit functor $u\colon \Sp \to C$, $\S \mapsto \unit_C$, is strongly continuous;
			\item the right adjoint $\mu^R\colon C \to C \otimes C$ preserves colimits, and its canonical lax $C \otimes C$-linear structure is strict; equivalently, $\mu$ admits a right adjoint in $\Mod_{C \otimes C}(\PrL_{\st})$.
		\end{enumerate}
	\end{definition}

	\begin{remark}
		\label[remark]{rem:Equivalent_Rigidity_Definitions}
		The equivalence of the two formulations in (2) is a general fact about adjunctions internal to module categories: a right adjoint of a $C \otimes C$-linear functor inherits a lax $C \otimes C$-linear structure, and it is a right adjoint in $\Mod_{C \otimes C}(\PrL_{\st})$ if and only if it preserves colimits and this lax structure is strict. \Cref{def:Rigid_Category_Detailed} is the definition of Gaitsgory--Rozenblyum, and it agrees with Ramzi's definition \cite[Definitions~4.5 and~4.36]{Ramzi_Locally_Rigid} (Ramzi additionally requires $C$ to be dualizable, which is automatic by \Cref{prop:Rigid_Self_Dual_Detailed}). It is also equivalent to \Cref{def:Rigid_Category}, which asks that for every finite set $I$ the multiplication $\mu_I\colon C^{\otimes I} \to C$ admits a right adjoint in $\Mod_{C^{\otimes I}}(\PrL_{\st})$. Indeed, for $I = \emptyset$ the map $\mu_{\emptyset}$ is the unit functor $u\colon \Sp \to C$, and $\Mod_{\Sp}(\PrL_{\st}) = \PrL_{\st}$, so this case is condition (1); the case $|I| = 1$ is vacuous, and the case $|I| = 2$ is condition (2). Conversely, if (2) holds, then by induction $\mu_{I \sqcup \{*\}} = \mu \circ (\mu_I \otimes \id_C)$ is a composite of two functors admitting right adjoints in $\Mod_{C^{\otimes I} \otimes C}(\PrL_{\st})$: for $\mu_I \otimes \id_C$ this follows by tensoring the adjunction for $\mu_I$ with $C$, and for $\mu$ by restricting the $C \otimes C$-linear adjunction along the algebra map $\mu_I \otimes \id_C\colon C^{\otimes I} \otimes C \to C \otimes C$.
	\end{remark}

	\begin{proposition}[{\cite[Chapter~1, Section~9.2.1]{GaitsgoryRozenbluym2017StudyDAG}}]
		\label[proposition]{prop:Rigid_Self_Dual_Detailed}
		Let $C$ be rigid. Then $C$ is self-dual in $\PrL_{\st}$, with evaluation and coevaluation given by
		\[
			\ev\colon C \otimes C \xrightarrow{\mu} C \xrightarrow{\Gamma := \hom_C(\unit_C,-)} \Sp, \qquad \coev\colon \Sp \xrightarrow{u} C \xrightarrow{\mu^R} C \otimes C.
		\]
		In particular, $C$ is dualizable.
	\end{proposition}
	\prooffor{prop:Rigid_Self_Dual_Detailed}
	\begin{proof}
		Note first that both functors preserve colimits: $\Gamma$ preserves colimits since $\unit_C$ is compact (and $C$ is stable), and $\mu^R$ preserves colimits by assumption. Moreover, $\Gamma$ is right adjoint to $u$. Write $\eta := \mu^R(\unit_C) \in C \otimes C$, so that $\coev$ is the colimit-preserving functor $\Sp \to C \otimes C$ sending $\S$ to $\eta$. We verify the first triangle identity, i.e.\ that the composite
		\[
			C \simeq \Sp \otimes C \xrightarrow{\coev \otimes C} C \otimes C \otimes C \xrightarrow{C \otimes \ev} C \otimes \Sp \simeq C
		\]
		is equivalent to the identity; the second one follows by symmetry. Since $C \otimes \ev = (C \otimes \Gamma) \circ (C \otimes \mu)$, the composite sends $c \in C$ to
		\[
			(C \otimes \Gamma)\big((C \otimes \mu)(\eta \boxtimes c)\big) \simeq (C \otimes \Gamma)\big(\eta \cdot (\unit_C \boxtimes c)\big),
		\]
		where $\cdot$ denotes the multiplication in $C \otimes C$. Now we use the linearity of $\mu^R$: for $a \in C \otimes C$ and $x \in C$ we have $\mu^R(\mu(a) \otimes x) \simeq a \cdot \mu^R(x)$. Taking $a = \unit_C \boxtimes c$ and $x = \unit_C$ gives
		\[
			\eta \cdot (\unit_C \boxtimes c) = (\unit_C \boxtimes c) \cdot \mu^R(\unit_C) \simeq \mu^R(c).
		\]
		Hence the composite is $(C \otimes \Gamma) \circ \mu^R\colon C \to C$, naturally in $c$ (as all functors involved preserve colimits, it suffices to identify the functors, which the above does naturally in $c$). Since $C \otimes - $ is a functor of $(\infty,2)$-categories, $C \otimes \Gamma$ is right adjoint to $C \otimes u\colon C \to C \otimes C$, $x \mapsto x \boxtimes \unit_C$. Hence $(C \otimes \Gamma) \circ \mu^R$ is right adjoint to $\mu \circ (C \otimes u) \simeq \id_C$, and is therefore equivalent to the identity.
	\end{proof}

	The linearity of the right adjoint to multiplication is what makes the triangle identity work. Trace-class morphisms express the corresponding duality data for individual morphisms.

	Next, we introduce trace-class morphisms, following Clausen--Scholze \cite{Scholze_Analytic_Geometry}.

	\begin{definition}
		\label[definition]{def:Trace_Class}
		Let $C$ be a symmetric monoidal category. A morphism $f\colon x \to y$ in $C$ is called \emph{trace-class} if there exist an object $z \in C$ and maps $\alpha\colon \unit_C \to y \otimes z$ and $\beta\colon z \otimes x \to \unit_C$ such that $f$ is homotopic to the composite
		\[
			x \simeq \unit_C \otimes x \xrightarrow{\alpha \otimes x} y \otimes z \otimes x \xrightarrow{y \otimes \beta} y \otimes \unit_C \simeq y.
		\]
	\end{definition}

	\begin{lemma}
		\label[lemma]{lem:Identity_Trace_Class}
		Let $C$ be a symmetric monoidal category whose homotopy category is idempotent complete; for instance, $C$ stable and idempotent complete \cite[Lemma~1.2.4.6]{lurie2016HA}. Then an object $x \in C$ is dualizable if and only if $\id_x$ is trace-class.
	\end{lemma}
	\prooffor{lem:Identity_Trace_Class}
	\begin{proof}
		If $x$ is dualizable with dual $x^{\vee}$, then the coevaluation and evaluation are a witness for $\id_x$. Conversely, let $(z,\alpha,\beta)$ be a witness for $\id_x$, so that $\alpha\colon \unit \to x \otimes z$ and $\beta\colon z \otimes x \to \unit$ satisfy the first triangle identity $(x \otimes \beta)(\alpha \otimes x) \simeq \id_x$, but not necessarily the second one. Consider
		\[
			e := (\beta \otimes z)(z \otimes \alpha)\colon z \to z \otimes x \otimes z \to z.
		\]
		Inserting the first triangle identity in the middle of $e \circ e$ shows that $e \circ e \simeq e$, and similarly $\beta(e \otimes x) \simeq \beta$ and $(x \otimes e)\alpha \simeq \alpha$. As the homotopy category is idempotent complete, $e$ splits: $e \simeq sr$ with $r\colon z \to z'$, $s\colon z' \to z$ and $rs \simeq \id_{z'}$. Put $\alpha' := (x \otimes r)\alpha$ and $\beta' := \beta(s \otimes x)$. Then
		\[
			(x \otimes \beta')(\alpha' \otimes x) \simeq (x \otimes \beta)(x \otimes e \otimes x)(\alpha \otimes x) \simeq (x \otimes \beta)(\alpha \otimes x) \simeq \id_x
		\]
		and $(\beta' \otimes z')(z' \otimes \alpha') \simeq r e s \simeq rsrs \simeq \id_{z'}$. So $(\alpha',\beta')$ exhibits $z'$ as a dual of $x$; since dualizability can be checked in the homotopy category, $x$ is dualizable.
	\end{proof}

	\begin{remark}
		\label[remark]{rem:Trace_Class_Basics}
		\leavevmode
		\begin{enumerate}
			\item If $C$ is closed, write $x^{\vee} := \iHom(x,\unit_C)$ for the \emph{predual} of $x$. Then $f$ is trace-class if and only if the map $\unit_C \to \iHom(x,y)$ classifying $f$ factors through the canonical map $y \otimes x^{\vee} \to \iHom(x,y)$: given $(z,\alpha,\beta)$ as above, $\beta$ is adjoint to a map $z \to x^{\vee}$, and composing $\alpha$ with it gives the desired factorization; conversely, given a factorization $\unit_C \to y \otimes x^{\vee}$, take $z = x^{\vee}$ and $\beta$ the evaluation map.
			\item Trace-class morphisms form a two-sided ideal: if $f$ is trace-class with witness $(z,\alpha,\beta)$, then $hfk$ is trace-class with witness $(z,(h \otimes z)\alpha,\beta(z \otimes k))$.
			\item If $C$ is stable and idempotent complete (for instance presentable and stable), then the identity of $x$ is trace-class if and only if $x$ is dualizable, see \Cref{lem:Identity_Trace_Class}.
			\item A symmetric monoidal functor preserves trace-class morphisms, since it preserves the data $(z,\alpha,\beta)$.
		\end{enumerate}
	\end{remark}

	\begin{lemma}
		\label[lemma]{lem:Trace_Class_Compact}
		Let $C \in \CAlg(\PrL_{\st})$ be such that $\unit_C$ is compact. Then every trace-class morphism in $C$ is compact in the sense of \Cref{def:Compact_Morphism}.
	\end{lemma}
	\prooffor{lem:Trace_Class_Compact}
	\begin{proof}
		Let $f\colon x \to y$ be trace-class with witness $(z,\alpha,\beta)$. We construct a factorization of $f^*\colon \Map_C(y,\colim(-)) \to \Map_C(x,\colim(-))$ through $\Map_{\Ind(C)}(y(x),-)$, naturally in $W = \quotcolim_j w_j \in \Ind(C)$. Given $g\colon y \to \colim_j w_j$, consider
		\[
			(g \otimes z) \circ \alpha\colon \unit_C \to (\colim_j w_j) \otimes z \simeq \colim_j (w_j \otimes z).
		\]
		Since $\unit_C$ is compact, this defines a point of $\colim_j \Map_C(\unit_C,w_j \otimes z)$, and postcomposing with the natural map
		\[
			\colim_j \Map_C(\unit_C,w_j \otimes z) \to \colim_j \Map_C(x,w_j), \qquad \gamma \mapsto (w_j \otimes \beta) \circ (\gamma \otimes x),
		\]
		we obtain a point of $\colim_j \Map_C(x,w_j) = \Map_{\Ind(C)}(y(x),W)$. All these maps are natural in $W$, and the image of this point in $\Map_C(x,\colim_j w_j)$ is $(\colim_j w_j \otimes \beta) \circ ((g \otimes z)\alpha \otimes x) \simeq g \circ (y \otimes \beta) \circ (\alpha \otimes x) \simeq gf$.
	\end{proof}

	The converse holds for rigid categories; this is the main structural result about rigid categories. It was formulated in Clausen's lectures as the \emph{definition} of rigidity for dualizable symmetric monoidal categories.

	\begin{theorem}[{Ramzi \cite[Corollaries~2.6, 4.55 and~4.60, Example~2.9]{Ramzi_Locally_Rigid}; see also \cite[Propositions~1.6 and~1.7]{Efimov_Inverse_Limits}}]
		\label[theorem]{thm:Rigid_Trace_Class}
		Let $C \in \CAlg(\PrL_{\st})$. Then the following are equivalent:
		\begin{enumerate}
			\item $C$ is rigid;
			\item $\unit_C$ is compact, and $C$ is generated under colimits by objects of the form $\colim(x_0 \to x_1 \to x_2 \to \cdots)$ in which each map $x_n \to x_{n+1}$ is trace-class;
			\item $C$ is dualizable, and a morphism in $C$ is compact if and only if it is trace-class.
		\end{enumerate}
	\end{theorem}

	We will not prove this theorem. Note however that some implications are formal given what we know: (1) $\implies$ the first half of (3) is \Cref{prop:Rigid_Self_Dual_Detailed}, and (3) $\implies$ (2) follows from \Cref{lem:Compact_Morphisms_Yhat}(2), since $\id_{\unit_C}$ is trace-class and hence compact, and \Cref{rem:Omega_1_Compact_Compact_Morphisms}, which shows that a dualizable category is generated by sequential colimits along compact morphisms. The hard part is that in a rigid category every compact morphism is trace-class \cite[Corollary~4.55, Example~2.9 and Corollary~2.6]{Ramzi_Locally_Rigid}, and that (2) implies (1) \cite[Corollary~4.60]{Ramzi_Locally_Rigid}. Note that the dualizability hypothesis in (3) was not stated in the talk.

	\begin{corollary}
		\label[corollary]{cor:Rigid_Functors_Strongly_Continuous}
		Let $F\colon C \to D$ be a colimit-preserving symmetric monoidal functor between rigid categories. Then $F$ is strongly continuous. Consequently, the category $\mathrm{Rig}$ of rigid symmetric monoidal categories and colimit-preserving symmetric monoidal functors is a full subcategory of $\CAlg(\Pr^{\dual})$.
	\end{corollary}
	\prooffor{cor:Rigid_Functors_Strongly_Continuous}
	\begin{proof}
		By \Cref{thm:Rigid_Trace_Class}, the compact morphisms in $C$ are trace-class. Since $F$ is symmetric monoidal, it preserves trace-class morphisms (\Cref{rem:Trace_Class_Basics}(4)), and trace-class morphisms in $D$ are compact by \Cref{lem:Trace_Class_Compact}. So $F$ preserves compact morphisms, and hence is strongly continuous by \Cref{prop:Strongly_Continuous_Compact_Morphisms}.
	\end{proof}

	\begin{remark}
		\label[remark]{rem:Strong_Continuity_From_Rigid_Source}
		The argument shows more generally that a colimit-preserving symmetric monoidal functor $F\colon C \to D$ with $C$ rigid and $D$ dualizable with compact unit is strongly continuous. In fact, the dualizability hypothesis on $D$ can be removed \cite[Proposition~1.9]{Efimov_Inverse_Limits}.
	\end{remark}

	\begin{example}
		\label[example]{ex:Rigid_Compactly_Generated}
		Let $C \in \CAlg(\PrL_{\st})$ be compactly generated. Then $C$ is rigid if and only if $\unit_C$ is compact and every compact object of $C$ is dualizable, if and only if the compact objects of $C$ are precisely the dualizable objects \cite[Chapter~1, Lemma~9.1.5 and Corollary~9.1.7]{GaitsgoryRozenbluym2017StudyDAG}. Indeed, if $\unit_C$ is compact, then dualizable objects are compact (as $\Map(x,-) \simeq \Map(\unit_C,x^{\vee} \otimes -)$), and conversely if all compact objects are dualizable, then $C = \Ind(C^{\omega})$ is generated by objects whose identity is trace-class, so (2) of \Cref{thm:Rigid_Trace_Class} holds. Examples include $\D(R)$ for a commutative ring $R$, more generally $\D_{\mathrm{qc}}(X)$ for a qcqs scheme $X$, and $\Sp$ itself.
	\end{example}

	\begin{example}
		\label[example]{ex:Shv_Compact_Hausdorff_Rigid}
		Let $X$ be a compact Hausdorff space. Then $\Shv(X;\Sp)$ is rigid; this can be seen as a manifestation of Verdier duality, see \cite[Example~1.4]{Efimov_Inverse_Limits}. For instance, the unit (the constant sheaf) is compact since global sections $\Gamma(X;-)$ commute with filtered colimits for compact $X$ \cite[Corollary~7.3.4.12]{lurie2009HTT}. In particular, $\Shv(X;\Sp)$ is dualizable and self-dual. For every presentable stable category $D$ we have $\Shv(X;D) \simeq \Shv(X;\Sp) \otimes D$ \cite[Proposition~2.8.7]{Krause_Nikolaus_Putzstuck_Sheaves}, so $\Shv(X;D)$ is dualizable whenever $D$ is. We will study these categories in detail in \Cref{chap:Sheaves}.
	\end{example}

	\begin{example}
		\label[example]{ex:Locally_Rigid_Sheaf_Categories}
		If $X$ is merely locally compact Hausdorff and not compact, then $\Shv(X;\Sp)$ is not rigid, since the unit is not compact. It is still dualizable, and even \emph{locally rigid} in the sense of Ramzi \cite[Definition~4.5 and Example~4.10]{Ramzi_Locally_Rigid} (see also \cite[Definition~1.36 and Example~1.38]{Efimov_Inverse_Limits} and \cite[Example~4.2.9, Corollary~4.4.13]{Krause_Nikolaus_Putzstuck_Sheaves}): it satisfies condition (2) of \Cref{def:Rigid_Category_Detailed}. Indeed, for $j\colon X \hookrightarrow X^+$ the one-point compactification, $j_!$ identifies $\Shv(X;\Sp)$ with a smashing ideal in the rigid category $\Shv(X^+;\Sp)$.
	\end{example}

	\section{Rigidification}
	\label[section]{sec:Rigidification}

	% Source: Clausen, Lecture 3 (19.07.2023); written from the literature (Ramzi, Construction 4.79 and Theorem 4.81; Efimov, Proposition 1.23 and Remark 1.24).
	A symmetric monoidal functor preserves trace-class maps, and in a rigid category the compact maps are precisely the trace-class maps. This suggests looking for a rigid category over $C$ through which every symmetric monoidal colimit-preserving functor from a rigid category factors.

	\begin{definition}[Rigidification]
		\label[definition]{def:Rigidification}
		Let $C$ be a presentably symmetric monoidal stable category. Its \emph{rigidification} is a rigid category $C^{\rig}$ with a symmetric monoidal colimit-preserving functor $C^{\rig}\to C$ such that composition induces an equivalence
		\[
			\Map_{\mathrm{Rig}}(R,C^{\rig})\simeq\Map_{\CAlg(\PrL_{\st})}(R,C)
		\]
		for every rigid category $R$.
	\end{definition}

	To construct it, choose an uncountable regular cardinal $\kappa$ such that $C\in\CAlg(\PrL_{\st,\kappa})$. Write $y\colon C^\kappa\to\Ind(C^\kappa)$ for the Yoneda embedding. Given a diagram $F\colon\Q_{\geq 0}\to C^\kappa$ whose transition maps $F(a)\to F(b)$ are trace-class in $C$ for all $a<b$, form its telescope
	\[
		\colim_{a\in\Q_{\geq 0}} y(F(a))\quad\in\Ind(C^\kappa).
	\]
	The full subcategory generated under colimits by these telescopes is $C^{\rig}$. The functor to $C$ takes the displayed formal colimit to the actual colimit of $F$ \cite[Construction~4.79 and Theorem~4.81]{Ramzi_Locally_Rigid}.

	Here are the main steps of Ramzi's proof. The tensor product of two telescopes is again a telescope, since the diagonal in $\Q_{\geq 0}\times\Q_{\geq 0}$ is cofinal and tensor products preserve trace-class maps. The constant telescope on the unit gives the compact object $y(\unit_C)$. To prove rigidity, Ramzi studies maps between telescopes that factor through a stage of the target. A cut at a later rational index gives an intermediate telescope through which such a map factors again. Ramzi shows that these maps detect nonzero objects and applies his rigidity criterion \cite[Corollary~4.59 and proof of Theorem~4.81]{Ramzi_Locally_Rigid}.

	For the universal property, let $R$ be rigid. Ramzi lifts a symmetric monoidal colimit-preserving functor $R\to C$ along $\Ind(C^\kappa)\to C$ using the lifting theorem for dualizable sources. The cardinal condition for this lift is automatic: compact maps in $R$ are trace-class, so their images in $C$ are trace-class. Since $\unit_C$ is $\kappa$-compact, these images are $\kappa$-compact maps, which implies that the functor preserves $\kappa$-compact objects. He then uses rational trace-class presentations in $R$ to show that the lift lands in the subcategory generated by the displayed telescopes. The equivalence of mapping animae in \Cref{def:Rigidification} follows from the lifting equivalence. In particular the construction is independent of the choice of $\kappa$ \cite[Theorem~4.81]{Ramzi_Locally_Rigid}.

	When $\unit_C$ is compact, the functor $C^{\rig}\to C$ is fully faithful. Indeed, trace-class maps in $C$ are then compact by \Cref{lem:Trace_Class_Compact}. Efimov observes that telescopes of such maps lie in the image of the partially defined left adjoint to realization $\Ind(C^\kappa)\to C$. On this image realization is inverse to the left adjoint. If moreover $C$ is compactly generated, each compact map factors through a compact object. In this case Efimov identifies $C^{\rig}$ inside $C$ as the subcategory generated under colimits by the actual colimits of rational trace-class diagrams in $C^\omega$ \cite[Proposition~1.23 and Remark~1.24(i)]{Efimov_Inverse_Limits}.

	\begin{remark}
		\label[remark]{rem:Rigidification_Via_Idempotent_Ideals}
		Clausen gives the same construction in his third lecture \citep{Clausen_Efimov_K_Theory_Lecture_3} as a special case of a construction for a presentable stable category $C$ and a two-sided ideal $S$ of maps. He assumes that $S$ contains the zero maps, is closed under finite sums and shifts in both directions, and satisfies the following cofiber condition. Given composable maps of arrows $u\to v\to w$ in $C$, if both maps on sources and both maps on targets lie in $S$, then $\cofib(u)\to\cofib(w)$ lies in $S$.

		Clausen forms $(C,S)^{\dual}$ by taking formal rational telescopes with transitions in $S$ and closing under colimits. For a dualizable category $D$, strongly continuous functors $D\to(C,S)^{\dual}$ correspond to colimit-preserving functors $D\to C$ that take compact maps into $S$. For symmetric monoidal $C$ he takes $S$ to be the trace-class maps to obtain $C^{\rig}$. Rational indexing matters here. If every map in $S$ factors as two maps in $S$, then \Cref{lem:Interpolation} lets us refine any sequence of maps in $S$ to a rational diagram with the same colimit. Only this last simplification uses the factorization hypothesis.
	\end{remark}

	\section{Nuclear objects}
	\label[section]{sec:Nuclear_Modules}

	% Source: Talk 3 (Max Glaser, 27.10.2026); literature (Clausen--Scholze, Efimov, Ramzi, Andreychev), drafted with generative AI (Claude)
	In condensed mathematics one constructs categories of nuclear modules inside larger categories of complete modules. We first study the abstract construction of Clausen and Scholze \cite{Scholze_Analytic_Geometry}. Nuclear objects can be built from compact objects along trace-class maps. We will use this description to determine when the nuclear objects form a rigid category.

	Throughout this section, let $C \in \CAlg(\PrL_{\st})$ be compactly generated with compact unit. We write $x^{\vee} = \iHom(x,\unit_C)$ for the predual of $x$.

	\begin{definition}[{\cite[Definition~1.25]{Efimov_Inverse_Limits}; cf.\ \cite[Definition~8.5]{Clausen_Scholze_Complex_Geometry} and \cite[Definitions~13.10 and~13.12]{Scholze_Analytic_Geometry}}]
		\label[definition]{def:Nuclear}
		An object $M \in C$ is called \emph{nuclear} if every morphism $P \to M$ from a compact object $P \in C^{\omega}$ is trace-class. We denote by $\Nuc(C) \subseteq C$ the full subcategory of nuclear objects. An object $M$ is called \emph{basic nuclear} if $M \simeq \colim(P_0 \to P_1 \to P_2 \to \cdots)$ with each $P_n$ compact and each map $P_n \to P_{n+1}$ trace-class.
	\end{definition}

	Note that a compact object $P$ is nuclear if and only if $\id_P$ is trace-class, i.e.\ if and only if $P$ is dualizable (\Cref{rem:Trace_Class_Basics}(3)). So the nuclear objects should be thought of as the objects which are ``locally dualizable''.

	For ordinary vector spaces, a tensor $v\otimes\ell\in x\otimes P^\vee$ gives the linear map $p\mapsto\ell(p)v$. The image of $x\otimes P^\vee\to\Hom(P,x)$ consists of the finite-rank maps. A finite-rank map $f\colon x\to y$ lets us go part of the way back: write $f(v)=\sum_i\ell_i(v)y_i$. Then for \emph{any} $b\colon P\to x$, the composite $fb$ is represented by the tensor $\sum_i y_i\otimes(\ell_i\circ b)$.

	The same construction works in $C$. The comparison $x\otimes P^\vee\to\iHom(P,x)$ is adjoint to evaluation on $P$. A trace-class presentation of $f$ plays the role of the expression $\sum_i y_i\otimes\ell_i$. Even when the comparison at $x$ has no inverse, a trace-class transition $x\to y$ lets us go diagonally from maps at $x$ to tensors at $y$.

	\begin{lemma}
		\label[lemma]{lem:Trace_Class_Interleaving}
		Let $f\colon x\to y$ be trace-class and let $P\in C$. The naturality square for the tensor-to-internal-Hom comparison admits a diagonal making both triangles commute up to homotopy:
		\[
			\begin{tikzcd}[column sep=large]
				x\otimes P^{\vee} \rar["f\otimes P^{\vee}"] \dar & y\otimes P^{\vee} \dar \\
				\iHom(P,x) \rar["{\iHom(P,f)}"'] \arrow[ur,dashed] & \iHom(P,y)
			\end{tikzcd}
		\]
	\end{lemma}
	\prooffor{lem:Trace_Class_Interleaving}
	\begin{proof}
		Choose a tensor $t\colon\unit_C\to y\otimes x^\vee$ representing $f$ as in \Cref{rem:Trace_Class_Basics}(1). For the diagonal, tensor with $t$ and then compose a functional on $x$ with a map from $P$ to $x$:
		\[
			\iHom(P,x)\xrightarrow{t\otimes\id}
			y\otimes x^\vee\otimes\iHom(P,x)\longrightarrow y\otimes P^\vee.
		\]
		The last arrow uses internal composition $x^\vee\otimes\iHom(P,x)\to P^\vee$, since $x^\vee=\iHom(x,\unit_C)$.

		For the lower triangle, evaluate on $P$. By associativity of composition, composing a functional with a map from $P$ and then evaluating on $P$ is the same as first evaluating the map into $x$ and then applying the functional. Contracting $t$ against this $x$ gives $f$, so the composite is postcomposition with $f$.

		For the upper triangle, start with $x\otimes P^\vee$. The composite $x^\vee\otimes x\otimes P^\vee\to x^\vee\otimes\iHom(P,x)\to P^\vee$ uses the evaluation $x^\vee\otimes x\to\unit_C$ and leaves the factor $P^\vee$ untouched. The composite around this triangle is therefore the contraction of $t$ against $x$, tensored with $\id_{P^\vee}$. This is $f\otimes P^\vee$.
	\end{proof}

	\begin{lemma}
		\label[lemma]{lem:Nuclear_Colimit_Interleaving}
		Let $M \simeq \colim_{i \in I} M_i$ for a directed poset $I$ such that for every $i$ there is $j \geq i$ with $M_i \to M_j$ trace-class. Let $P, Q \in C$ be such that $Q$ and $Q \otimes P$ are compact. Then the map $\hom_C(Q, M \otimes P^{\vee}) \to \hom_C(Q \otimes P, M)$ induced by $M \otimes P^{\vee} \to \iHom(P,M)$ is an equivalence of spectra.
	\end{lemma}
	\prooffor{lem:Nuclear_Colimit_Interleaving}
	\begin{proof}
		Apply $\pi_0\Map_C(Q,-)$ to the square of \Cref{lem:Trace_Class_Interleaving}. By adjunction the bottom row now consists of homotopy classes of maps out of $Q\otimes P$. Since $Q$ and $Q\otimes P$ are compact and tensoring with $P^\vee$ preserves colimits, the comparison at $M$ is, on $\pi_0$, the filtered colimit of the vertical comparisons at the stages $M_i$.

		For surjectivity, represent a class by a map $Q\otimes P\to M_i$. Choose a trace-class transition $M_i\to M_j$. The diagonal gives a class in $\pi_0\Map_C(Q,M_j\otimes P^\vee)$, and the lower triangle says it represents the required map at stage $j$.

		For injectivity, take two tensor classes with the same image in the colimit. Pass to a common stage where their images in the bottom row agree, and then choose a trace-class transition out of that stage. Applying the diagonal to the common image gives the same tensor class for both. By the upper triangle this is exactly the image of each original class at the later stage.

		This proves bijectivity on $\pi_0$. Replacing $Q$ by $\Sigma^n Q$ gives the same conclusion on $\pi_n$ for every $n\in\Z$, hence the equivalence of spectra.
	\end{proof}

	\begin{proposition}[Clausen--Scholze; see {\cite[Theorem~8.6]{Clausen_Scholze_Complex_Geometry} and \cite[Propositions~1.27 and~1.28]{Efimov_Inverse_Limits}}]
		\label[proposition]{prop:Nuclear_Basics}
		\leavevmode
		\begin{enumerate}
			\item For $M \in C$, the following are equivalent:
			\begin{enumerate}
				\item[(a)] $M$ is nuclear;
				\item[(b)] for every compact $P$, the map $\hom_C(\unit_C,P^{\vee} \otimes M) \to \hom_C(P,M)$ is an equivalence of spectra (equivalently, since we may replace $P$ by its shifts, the corresponding map of mapping spaces is an equivalence for every compact $P$);
				\item[(d)] $M \simeq \colim_{i \in I} M_i$ for a directed poset $I$ and compact objects $M_i$, such that for every $i$ there is $j \geq i$ with $M_i \to M_j$ trace-class.
			\end{enumerate}
			Moreover, these conditions are implied by
			\begin{enumerate}
				\item[(c)] for every compact $P$, the map $P^{\vee} \otimes M \to \iHom(P,M)$ is an equivalence,
			\end{enumerate}
			and if the tensor product of any two compact objects of $C$ is compact, then (c) is equivalent to (a), (b) and (d).
			\item The subcategory $\Nuc(C) \subseteq C$ is closed under colimits. It is $\omega_1$-presentable, and its $\omega_1$-compact objects are precisely the basic nuclear objects.
		\end{enumerate}
	\end{proposition}
	The main step is to turn the pointwise trace-class condition into a presentation whose transition maps are eventually trace-class. The interleaving lemma then upgrades surjectivity on $\pi_0$ to the isomorphism of mapping spectra in (b). For (2), we decompose these presentations into countable pieces, as in the proof of $\omega_1$-presentability for compactly assembled categories.
	\prooffor{prop:Nuclear_Basics}
	\begin{proof}
		We first prove (1). (b) $\Rightarrow$ (a): On $\pi_0$, surjectivity of the comparison in (b) says that every map $P\to M$ is represented by a tensor $\unit_C\to M\otimes P^\vee$. This is precisely the trace-class condition.

		(a) $\Rightarrow$ (d): Since $C$ is compactly generated, $M$ is a filtered colimit of compact objects, and by \Cref{thm:Filtered_Posets_Cofinal} we may write $M \simeq \colim_{i \in I} M_i$ with $I$ a directed poset and $M_i$ compact; denote by $t_{ij}$ the transition maps and by $\iota_i\colon M_i \to M$ the canonical maps. Fix $i$. Since $\iota_i$ is trace-class, it is represented by a tensor $w \in \pi_0\Map_C(\unit_C, M \otimes M_i^{\vee})$. As $\unit_C$ is compact and $M \otimes M_i^{\vee} \simeq \colim_j M_j \otimes M_i^{\vee}$, the class $w$ is the image of some $w' \in \pi_0\Map_C(\unit_C, M_j \otimes M_i^{\vee})$ with $j \geq i$. Let $g\colon M_i\to M_j$ be the map represented by $w'$. It is trace-class, and $\iota_jg\simeq\iota_i$ because the image of $w'$ at $M$ represents $\iota_i$. As $M_i$ is compact, $\pi_0\Map_C(M_i,M) \simeq \colim_k \pi_0\Map_C(M_i,M_k)$, so there is $k \geq j$ with $t_{jk}g \simeq t_{ik}$. Since trace-class morphisms form an ideal (\Cref{rem:Trace_Class_Basics}(2)), $t_{ik}$ is trace-class.

		(d) $\Rightarrow$ (b): This is \Cref{lem:Nuclear_Colimit_Interleaving} with $Q = \unit_C$.

		(c) $\Rightarrow$ (b): Apply $\hom_C(\unit_C,-)$ and use $\hom_C(\unit_C,\iHom(P,M)) \simeq \hom_C(P,M)$.

		(d) $\Rightarrow$ (c), if compact objects are closed under tensor products: For compact $Q$, the map $\hom_C(Q,P^{\vee} \otimes M) \to \hom_C(Q,\iHom(P,M)) \simeq \hom_C(Q \otimes P,M)$ is an equivalence by \Cref{lem:Nuclear_Colimit_Interleaving}. Since $C$ is compactly generated, the compact objects detect equivalences, so $P^{\vee} \otimes M \to \iHom(P,M)$ is an equivalence.

		We now prove (2). By (1), $M$ is nuclear if and only if the maps in (b) are equivalences. Since $\unit_C$ and $P$ are compact, the functors $\hom_C(\unit_C,P^{\vee} \otimes -)$ and $\hom_C(P,-)$ preserve filtered colimits, and being exact they preserve all small colimits. Hence $\Nuc(C)$ is closed under colimits in $C$; in particular, colimits in $\Nuc(C)$ are computed in $C$.

		Basic nuclear objects are nuclear by (d) (with $I = \N$). They are countable colimits of compact objects, hence $\omega_1$-compact in $C$ (\Cref{cor:Kappa_Small_Colimits_Of_Kappa_Compacts}), and hence also in $\Nuc(C)$.

		Next, every nuclear $M$ is an $\omega_1$-filtered colimit of basic nuclear objects. Write $M \simeq \colim_{i \in I} M_i$ as in (d), and choose for every $i$ an element $\sigma(i) \geq i$ such that $M_i \to M_{\sigma(i)}$ is trace-class. Let $\Sigma_{\sigma}$ be the set of countable directed subsets $J \subseteq I$ with $\sigma(J) \subseteq J$, ordered by inclusion. The proof of \Cref{prop:Colimits_Over_Countable_Subposets} applies verbatim to $\Sigma_{\sigma}$ instead of $\Sigma$ (adding in the inductive construction also $\sigma(u)$ for $u \in U_k$), and shows that $\Sigma_{\sigma}$ is $\omega_1$-filtered and that $M \simeq \colim_{J \in \Sigma_{\sigma}} M_J$ with $M_J = \colim_{i \in J} M_i$. Each $M_J$ is basic nuclear: enumerate $J = \{q_0,q_1,\dots\}$, and choose inductively $j_0 = q_0$ and $j_{n+1} \in J$ with $j_{n+1} \geq \sigma(j_n), q_{n+1}$. Then $n \mapsto j_n$ is cofinal (\Cref{lem:Cofinality_Posets}), so $M_J \simeq \colim_n M_{j_n}$, and each $M_{j_n} \to M_{j_{n+1}}$ factors through the trace-class map $M_{j_n} \to M_{\sigma(j_n)}$, hence is trace-class.

		The basic nuclear objects form an essentially small collection, since $C^{\omega}$ is essentially small. Thus $\Nuc(C)$ admits small colimits and is generated under $\omega_1$-filtered colimits by an essentially small collection of $\omega_1$-compact objects, so it is $\omega_1$-accessible by \cite[Proposition~5.4.2.2]{lurie2009HTT}, and hence $\omega_1$-presentable.

		Finally, let $X \in \Nuc(C)$ be $\omega_1$-compact. Writing $X \simeq \colim_{J \in \Sigma_{\sigma}} X_J$ as above, the identity of $X$ factors through some $X_J$, so $X$ is a retract of a basic nuclear object $Y = \colim_n a_n$, where the $a_n$ are compact and the transition maps $a_n \to a_{n+1}$ are trace-class. Now we argue as in the proof of \Cref{prop:Omega_1_Compact_In_Ind}, applied in $C \simeq \Ind(C^{\omega})$: by \Cref{lem:Retract_As_Telescope} and \Cref{lem:Strictification_Sequential_Maps}, we find a strictly increasing $f\colon \N \to \N$ with $f(n) \geq n$ and a natural transformation $\eta\colon a_{\bullet} \to a_{f(\bullet)}$ such that $X \simeq \colim_n \Phi(n,n)$, where $\Phi(k,n) = a_{f^k(n)}$. The transition map $\Phi(n,n) \to \Phi(n+1,n+1)$ is the composite of the transition map $a_{f^n(n)} \to a_{f^n(n+1)}$ of $a_{\bullet}$, which is trace-class as $f^n(n+1) > f^n(n)$, with $\eta_{f^n(n+1)}$. So it is trace-class, and $X$ is basic nuclear.
	\end{proof}

	\begin{remark}
		\label[remark]{rem:Nuclear_Definition_Pi0}
		\Cref{prop:Nuclear_Basics} shows that the condition in \Cref{def:Nuclear}, which only asks that the maps $\pi_0\Map_C(\unit_C,P^{\vee} \otimes M) \to \pi_0\Map_C(P,M)$ be surjective, already implies that $\hom_C(\unit_C,P^{\vee} \otimes M) \to \hom_C(P,M)$ is an equivalence. Clausen--Scholze \cite[Definition~8.5]{Clausen_Scholze_Complex_Geometry} instead define nuclearity by the latter condition (and \cite[Definition~13.10]{Scholze_Analytic_Geometry} by an analogous condition for analytic rings), whereas \Cref{def:Nuclear} follows \cite[Definition~1.25]{Efimov_Inverse_Limits}. The two definitions agree; the injectivity comes from the fact that, by (a) $\Rightarrow$ (d), $M$ is a filtered colimit of compact objects along trace-class maps, for which the two sides become interleaved (\Cref{lem:Nuclear_Colimit_Interleaving}).
	\end{remark}

	In the setting of this section, \Cref{sec:Rigidification} realizes $C^{\rig}$ inside $C$ by rationally indexed trace-class diagrams of compact objects. The following criterion lets us recognize the same objects using sequences.

	\begin{proposition}[{\cite[Proposition~1.29]{Efimov_Inverse_Limits}}]
		\label[proposition]{prop:Nuclear_Rigid}
		If every trace-class map between compact objects of $C$ factors through a compact object by two trace-class maps, then $C^{\rig}=\Nuc(C)$ as full subcategories of $C$.
	\end{proposition}
	\prooffor{prop:Nuclear_Rigid}
	\begin{proof}[Sketch]
		Each of the rationally indexed generators of $C^{\rig}$ from \Cref{sec:Rigidification} is nuclear by \Cref{prop:Nuclear_Basics}(1)(d). Since nuclear objects are closed under colimits, $C^{\rig}\subseteq\Nuc(C)$. Conversely, let $P_0 \to P_1 \to \cdots$ be a sequence of trace-class maps between compact objects. By the assumption and \Cref{lem:Interpolation}(2), applied to $D = C^{\omega}$ and the class of trace-class morphisms (which is closed under composition by \Cref{rem:Trace_Class_Basics}(2)), it extends to a functor $F\colon \Q_{\geq 0} \to C^{\omega}$ with trace-class transition maps and the same colimit. Restricting $F$ to the cofinal subset $\Q_{>0} \cong \Q$, we see that the basic nuclear objects lie in $C^{\rig}$. By \Cref{prop:Nuclear_Basics}(2), these generate $\Nuc(C)$.
	\end{proof}

	\section{Examples from condensed mathematics}
	\label[section]{sec:Condensed_Nuclear_Examples}

	The solid integers give a first example in which we can identify the nuclear objects explicitly. Passing to an adically complete ring leads to two categories of nuclear modules, obtained by rigidifying different ambient categories. We then look at individual nuclear modules over $\Z_p$ and at the corresponding questions over $\R$ and $\C$. In the solid and liquid examples, the compact objects are closed under tensor products, so the internal-Hom criterion of \Cref{prop:Nuclear_Basics}(1)(c) applies. We refer to \cite{Scholze_Analytic_Geometry} for the solid and liquid theory, to \cite{Clausen_Scholze_Complex_Geometry} for the gaseous theory, and to \cite{Andreychev_Thesis} and \cite[Section~7.2]{Efimov_Inverse_Limits} for the comparison with Efimov's theory.

	\begin{example}[Solid modules and nuclear modules over formal schemes]
		\label[example]{ex:Nuclear_Solid}
		Let $\D(\mathrm{Solid}_{\Z})$ be the derived category of light solid abelian groups, as in Efimov's convention in \cite[Section~7.2]{Efimov_Inverse_Limits}. This is a compactly generated symmetric monoidal category, with compact projective generator $\prod_{\N}\Z$ and compact unit $\Z$. It is far from rigid: the compact object $\prod_{\N}\Z$ is not dualizable (its predual is $\bigoplus_{\N}\Z$). In fact, there is a symmetric monoidal equivalence
		\[
			\D(\mathrm{Solid}_{\Z}) \simeq \Ind\big(\D^b(\mathrm{Proj}^{\omega_1}\text{-}\Z)\catop\big),
		\]
		where $\D^b(\mathrm{Proj}^{\omega_1}\text{-}\Z) \subseteq \D(\Z)$ is the stable subcategory generated by $\bigoplus_{\N}\Z$, under which $\prod_{\N}\Z$ corresponds to $(\bigoplus_{\N}\Z)\catop$ \cite[Section~7.2]{Efimov_Inverse_Limits}.

		In this example the nuclear objects admit a familiar description. A complex of light solid abelian groups is nuclear precisely when it comes from an ordinary complex of abelian groups via the discrete embedding \cite[Lecture~VIII, Exercise~1]{Clausen_Scholze_Complex_Geometry}. Thus $\Nuc(\D(\mathrm{Solid}_{\Z}))\simeq\D(\Z)$.

		More generally, let $R$ be a commutative ring and $I = (a_1,\dots,a_m) \subseteq R$ a finitely generated ideal, and let $R^{\wedge}_I$ denote its derived $I$-adic completion. We use the induced solid analytic structure, with the notation of \cite[Section~7.2]{Efimov_Inverse_Limits}. All completions and tensor products in this example are derived; for Noetherian $R$, the derived completion of $R$ agrees with its ordinary completion. Its category of modules $\D(\mathrm{Solid}_{(R^{\wedge}_I,\Z)}) = \Mod_{R^{\wedge}_I}(\D(\mathrm{Solid}_{\Z}))$ is compactly generated by $R^{\wedge}_I \otimes^{\square} \prod_{\N}\Z$, with compact unit. The category of \emph{nuclear solid modules} is
		\[
			\Nuc^{CS}(R^{\wedge}_I) := \Nuc\big(\D(\mathrm{Solid}_{(R^{\wedge}_I,\Z)})\big).
		\]
		Efimov identifies this category with the rigidification of the ambient solid module category:
		\[
			\Nuc^{CS}(R^{\wedge}_I) \simeq \D(\mathrm{Solid}_{(R^{\wedge}_I,\Z)})^{\rig}.
		\]

		The key technical input is that the preduals of compact objects are nuclear, which goes back to Andreychev \cite[Lemma~3.8 and Korollar~3.14]{Andreychev_Thesis}; see \cite[Propositions~7.4 and~7.5]{Efimov_Inverse_Limits} (Efimov deduces the rigidity from this via \cite[Propositions~1.33 and~1.34]{Efimov_Inverse_Limits} rather than via \Cref{prop:Nuclear_Rigid}). In particular:
		\begin{itemize}
			\item $\Nuc^{CS}(R^{\wedge}_I)$ is rigid, hence dualizable, but in general not compactly generated.
			\item Its compact objects are the perfect complexes over $R^{\wedge}_I$.
			\item Its continuous K-theory is $K^{\cont}(\Nuc^{CS}(R^{\wedge}_I)) \simeq \lim_n K(R \quot (a_1^n,\dots,a_m^n))$, which is $\lim_n K(R/I^n)$ if $R$ is Noetherian (Efimov \cite[Corollary~7.10]{Efimov_Inverse_Limits}, as conjectured by Clausen--Scholze), cf.\ \Cref{thm:Kcont_Nuclear}.
		\end{itemize}

		To compare the two constructions, keep the ring fixed and change the ambient category. Rigidifying the solid module category gives $\Nuc^{CS}(R^{\wedge}_I)$ as above. Rigidifying the category $\D_{I\text{-}\mathrm{compl}}(R)$ of derived $I$-complete $R$-modules gives Efimov’s variant $\Nuc(R^{\wedge}_I)$ \cite[Definition~3.28 and Proposition~4.3]{Efimov_Inverse_Limits}, which can be described as $\lim^{\dual}_n \D(R \quot (a_1^n,\dots,a_m^n))$, the limit computed in $\Pr^{\dual}$ \cite[Proposition~5.23]{Efimov_Inverse_Limits}. It receives a fully faithful strongly continuous functor from $\Nuc^{CS}(R^{\wedge}_I)$ \cite[Corollary~7.6]{Efimov_Inverse_Limits}, which induces an equivalence on compact objects, both sides having compact objects $\Perf(R^{\wedge}_I)$ \cite[Proposition~3.29]{Efimov_Inverse_Limits}. The ambient category of derived complete modules can itself fail to be rigid. For $R=\Z$ and $I=(p)$ its unit $\Z_p$ is not compact.

		% Source: written from the literature (Efimov, Localizing invariants of inverse limits, Corollary 7.20(i) and the discussion after Remark 7.21).
		Compact generation fails already for $\Nuc^{CS}(\Z_p)$. Indeed, suppose otherwise. Its compact objects are $\Perf(\Z_p)$, so the canonical $\Z$-linear functor $\D(\Z_p) \to \Nuc^{CS}(\Z_p)$ would be an equivalence. But Efimov computes
		\[
			\mathrm{HH}^{\cont}(\Nuc^{CS}(\Z_p)/\Z) \simeq \Z_p
		\]
		\cite[Corollary~7.20(i)]{Efimov_Inverse_Limits}, whereas
		\[
			\mathrm{HH}_1(\Z_p/\Z) \cong \Omega^1_{\Z_p/\Z} \cong \Omega^1_{\Q_p/\Q} \neq 0.
		\]
		Here $\mathrm{HH}^{\cont}$ denotes the continuous extension of relative Hochschild homology, which agrees with ordinary Hochschild homology on compactly generated categories. This gives the required contradiction; see also the discussion after \cite[Remark~7.21]{Efimov_Inverse_Limits}.

		% Source: written from the literature (Efimov, Corollary 7.7); moved from the introductory overview.
		There is also an operator-algebra presentation, analogous to the compact-operator example of \Cref{ex:Compact_Operators_Intro}. For the same pair $(R,I)$, take $P := (\bigoplus_{\N} R)^{\wedge}_I$ be the $I$-completed countable direct sum, which plays the role of a separable Hilbert space. In place of $B(\Hh)$ we take the endomorphism ring $\End_R(P)$, and in place of the Calkin algebra $B(\Hh)/\Kk$ the $I$-adically completed endomorphism ring
		\[
			\Calk^{\mathrm{top}}_R(P) := \End_{\Calk(R)}(P)^{\wedge}_I
		\]
		of $P$ in the Calkin category of $\Perf(R)$. The fiber of $\End_R(P) \to \Calk^{\mathrm{top}}_R(P)$ is $\big(\prod_{\N}R \otimes_R \bigoplus_{\N}R\big)^{\wedge}_I$, which one may think of as the $I$-adic completion of the endomorphisms of finite rank, just as $\Kk$ is the norm closure of the operators of finite rank. Efimov shows that $\End_R(P) \to \Calk^{\mathrm{top}}_R(P)$ is a homological epimorphism, and identifies the kernel of extension of scalars along this map with the category of \emph{nuclear} modules over $R^{\wedge}_I$ in the sense of Clausen--Scholze \cite[Corollary~7.7]{Efimov_Inverse_Limits}:
		\[
			\Nuc^{CS}(R^{\wedge}_I) \simeq \D\big(\End_R(P) \to \Calk^{\mathrm{top}}_R(P)\big).
		\]
	\end{example}

	\begin{example}[Banach spaces as nuclear modules]
		\label[example]{ex:Nuclear_Banach}
		Let $S=\N\cup\{\infty\}$ be the one-point compactification of the discrete set $\N$. A continuous function $S\to\Z_p$ is a convergent sequence together with its limit. Evaluation at $\infty$ splits by the constant functions, giving
		\[
			C(S,\Z_p)\simeq \Z_p\oplus c_0(\N,\Z_p),
		\]
		where $c_0(\N,\Z_p)$ consists of sequences tending to zero. By \cite[Lemma~3.8]{Andreychev_Thesis}, $C(S,\Z_p)$ is nuclear over $\Z_p^{\square}$ for every profinite $S$. Its summand $c_0(\N,\Z_p)$ is therefore nuclear, and inverting $p$ gives the nuclear solid $\Q_p$-vector space $c_0(\N,\Q_p)$.

		The convergence condition makes a difference. The module of all sequences, the compact objects such as $\prod_{\N}\Z_p$, which are of ``profinite'' or ``compact'' type. For instance, for every profinite set $S$ the module $C(S,\Z_p)$ of continuous functions is nuclear over $\Z_p^{\square}$ \cite[Lemma~3.8]{Andreychev_Thesis}. Taking $S = \N \cup \{\infty\}$, it follows that the $p$-adic Banach space $c_0(\N,\Z_p)$ of null sequences, a direct summand of $C(S,\Z_p)$, is nuclear, and inverting $p$ gives nuclear solid $\Q_p$-vector spaces such as $c_0(\N,\Q_p)$. In contrast, the compact object $\prod_{\N}\Z_p$ of $\D(\mathrm{Solid}_{\Z_p})$, is not nuclear: a compact nuclear object is dualizable, and the dualizable objects of $\D(\mathrm{Solid}_{\Z_p})$ are just the perfect complexes of $\Z_p$-modules \cite[Theorem~14.10]{Scholze_Analytic_Geometry}.
	\end{example}

	\begin{example}[Liquid and gaseous modules]
		\label[example]{ex:Nuclear_Liquid_Gaseous}
		Over the real and complex numbers, the solid theory is replaced by the \emph{liquid} theory: for $0 < p \leq 1$ there is a compactly generated symmetric monoidal category $\D(\mathrm{Liq}_p(\R))$ of $p$-liquid $\R$-vector spaces \cite{Scholze_Analytic_Geometry}, and one can form $\Nuc(\R^{p\text{-}\mathrm{liq}})$. The nuclear objects of $\D(\mathrm{Liq}_p)$ do not depend on $p$ \cite[Corollary~8.18]{Clausen_Scholze_Complex_Geometry}.

		According to Clausen \citep{Clausen_Efimov_K_Theory_Lecture_3}, there is an explicit model for the rigidification of the category of liquid $\C$-vector spaces in terms of operator algebras:
		\[
			\D^{\mathrm{liq}}(\C)^{\rig} \simeq \D\big(B(\Hh) \to B(\Hh)/S_{0+}\big).
		\]
		Here $\Hh$ is a separable infinite-dimensional complex Hilbert space and $B(\Hh)$ is its algebra of bounded operators. The ideal $S_{0+}=\bigcap_{q>0} S_q$ consists of compact operators whose singular values are $q$-summable for every $q>0$. As in the solid example, the category on the right is the kernel of extension of scalars along the displayed algebra map.

		Finally, the gaseous theory of Clausen--Scholze \cite{Clausen_Scholze_Complex_Geometry} provides yet another analytic ring structure on $\C$, with a category of nuclear modules $\Nuc(\C^{\mathrm{gas}})$. Rigidification produces rigid, hence dualizable, categories whose continuous K-theory is a candidate for analytic K-theory, cf.\ \Cref{sec:First_Computations}.
	\end{example}

	\begin{remark}
		\label[remark]{rem:Nuclear_Descent}
		On an affinoid adic space the categories of nuclear modules also glue. For a complete sheafy analytic Huber pair $(A,A^+)$, Andreychev proves that $U\mapsto\Nuc(\mathcal O(U))$ satisfies descent for covers by rational opens of $\operatorname{Spa}(A,A^+)$. The same holds after applying continuous K-theory, as recalled in \Cref{thm:Nuclear_Analytic_Descent}.
	\end{remark}


		\chapter{Calkin categories, colimits and limits of dualizable categories}
	\label[chapter]{chap:Calkin_Limits_Colimits}

	% Source: Talk 5 (Anton Zakrewski, 10.11.2026); Talk 6 (Jacopo Mazzei, 17.11.2026)
	In this chapter we study the category $\Pr^{\dual}$ of dualizable categories and strongly continuous functors as a category in its own right. Colimits in $\Pr^{\dual}$ turn out to be computed in $\PrL_{\st}$, and filtered colimits in $\Pr^{\dual}$ are much better behaved than in $\PrL_{\st}$: they preserve fully faithful functors and short exact sequences. Limits in $\Pr^{\dual}$ are more subtle and are in general \emph{not} computed in $\PrL_{\st}$. Kernels and pullbacks exist, but they are given by the largest dualizable subcategories of the corresponding kernels and pullbacks in $\PrL_{\st}$ on which the relevant functors are strongly continuous.

	Along the way we introduce the continuous Calkin category, which is the key ingredient in the definition of localizing invariants of dualizable categories (\Cref{chap:Localizing_Invariants_Dualizable}). The main result about it is that it sends short exact sequences of dualizable categories to short exact sequences of small categories (\Cref{prop:Calkin_Short_Exact}). Its functoriality relies on the description of colimits, which we therefore treat first.

	We follow \cite[Sections~2.9--2.13 and~2.15]{Efimov_Large_Categories}. Recall from \Cref{sec:Localization_Sequences_PrL} that $\Pr^{\mathrm{LL}}_{\st} \subseteq \PrL_{\st}$ denotes the subcategory of presentable stable categories and strongly continuous functors. By \Cref{rem:Compactly_Generated_In_Prdual} and \Cref{rem:Strongly_Continuous_Stable}, we have full subcategories
	\[
		\PrL_{\omega,\st} \subseteq \Pr^{\dual} \subseteq \Pr^{\mathrm{LL}}_{\st}.
	\]

	\section{Colimits of dualizable categories}
	\label[section]{sec:Colimits_Dualizable}

	% Source: Talk 5 (Anton Zakrewski, 10.11.2026); Zakrewski's notes for the talk; Efimov, Propositions 1.7, 1.13 and 2.63 and Corollary 2.64
	Colimits in $\PrL_{\st}$ are computed as limits along the right adjoints (\Cref{sec:PrL_Tensor_Products}). If all functors in the diagram are strongly continuous, then the right adjoints preserve colimits, and this allows us to show that colimits in $\Pr^{\mathrm{LL}}_{\st}$ and in $\Pr^{\dual}$ are computed in $\PrL_{\st}$.

	\begin{proposition}[{\cite[Proposition~1.7]{Efimov_Large_Categories}}]
		\label[proposition]{prop:PrLL_Colimits}
		Let $I \to \Pr^{\mathrm{LL}}_{\st}$, $i \mapsto C_i$, be a small diagram, let $C := \colim_i C_i$ be its colimit in $\PrL_{\st}$, and let $F_i\colon C_i \to C$ be the structure functors.
		\begin{enumerate}
			\item The functors $F_i$ are strongly continuous.
			\item A colimit-preserving functor $G\colon C \to D$ to a presentable stable category $D$ is strongly continuous if and only if each composite $G \circ F_i$ is strongly continuous.
		\end{enumerate}
		In particular, $\Pr^{\mathrm{LL}}_{\st}$ is cocomplete, and the inclusion $\Pr^{\mathrm{LL}}_{\st} \to \PrL_{\st}$ preserves colimits.
	\end{proposition}
	\prooffor{prop:PrLL_Colimits}
	\begin{proof}
		We identify $C$ with the limit $\lim_i C_i$ in $\widehat{\Cat}$, taken along the right adjoints $F_{ij}^R$ of the transition functors $F_{ij}$. Under this identification, $F_i^R\colon C \to C_i$ is the projection. Since the transition functors $F_{ij}^R$ preserve colimits, colimits in $\lim_i C_i$ exist and are computed pointwise \cite[Remark~5.4.7.13]{lurie2009HTT}. In particular, the projections $F_i^R$ preserve colimits, which proves (1). Similarly, the right adjoint $G^R\colon D \to C \simeq \lim_i C_i$ corresponds to the compatible family of functors $F_i^RG^R \simeq (GF_i)^R$. As colimits in the limit are computed pointwise, $G^R$ preserves colimits if and only if each $(GF_i)^R$ does, which proves (2).

		For the final assertion, note that $\Map_{\Pr^{\mathrm{LL}}_{\st}}(C,D) \subseteq \Map_{\PrL_{\st}}(C,D)$ is a union of path components, and similarly for the $C_i$. By (2), the equivalence $\Map_{\PrL_{\st}}(C,D) \simeq \lim_i \Map_{\PrL_{\st}}(C_i,D)$ restricts to an equivalence $\Map_{\Pr^{\mathrm{LL}}_{\st}}(C,D) \simeq \lim_i \Map_{\Pr^{\mathrm{LL}}_{\st}}(C_i,D)$, so that $C$, together with the cocone of (1), is also a colimit in $\Pr^{\mathrm{LL}}_{\st}$.
	\end{proof}

	\begin{proposition}[{\cite[Proposition~2.63]{Efimov_Large_Categories}}]
		\label[proposition]{prop:Prdual_Colimits}
		The full subcategory $\Pr^{\dual} \subseteq \Pr^{\mathrm{LL}}_{\st}$ is closed under small colimits. In particular, $\Pr^{\dual}$ is cocomplete, and the inclusion $\Pr^{\dual} \to \PrL_{\st}$ preserves colimits.
	\end{proposition}
	\prooffor{prop:Prdual_Colimits}
	\begin{proof}
		Let $I \to \Pr^{\dual}$, $i \mapsto C_i$, be a small diagram, and let $C = \colim_i C_i$ in $\PrL_{\st}$. By \Cref{prop:PrLL_Colimits}, it suffices to show that $C$ is dualizable. The structure functors $F_i\colon C_i \to C$ are strongly continuous by \Cref{prop:PrLL_Colimits}(1). Their right adjoints are the projections $\lim_i C_i \to C_i$, which are jointly conservative. Hence the localizing subcategory $L \subseteq C$ generated by the essential images of the $F_i$, which is generated by a set of objects and hence presentable, satisfies $L^{\perp} = 0$. By \Cref{prop:Localization_Sequences_PrL}(1), we get $L = C$. So $C$ is dualizable by \Cref{prop:Generation_Dualizable}.
	\end{proof}

	As a first consequence, we can describe the epimorphisms in these categories.

	\begin{proposition}[{\cite[Propositions~1.13 and~1.14 and Corollary~2.64]{Efimov_Large_Categories}}]
		\label[proposition]{prop:Epimorphisms_Localizations}
		A morphism $F\colon C \to D$ in $\PrL_{\st}$, $\Pr^{\mathrm{LL}}_{\st}$ or $\Pr^{\dual}$ is an epimorphism if and only if it is a localization, i.e.\ $F^R$ is fully faithful.
	\end{proposition}
	\prooffor{prop:Epimorphisms_Localizations}
	\begin{proof}
		A morphism $F\colon C \to D$ in a category with pushouts is an epimorphism if and only if the codiagonal $\nabla\colon D \sqcup_C D \to D$ is an equivalence. By \Cref{prop:PrLL_Colimits,prop:Prdual_Colimits}, pushouts in $\Pr^{\mathrm{LL}}_{\st}$ and $\Pr^{\dual}$ are computed in $\PrL_{\st}$. So it suffices to treat $\PrL_{\st}$.

		Computing the pushout as a pullback along the right adjoints, we have $D \sqcup_C D \simeq D \times_C D$, where both functors $D \to C$ are $F^R$, and the right adjoint of $\nabla$ is the diagonal $\delta\colon D \to D \times_C D$. So $F$ is an epimorphism if and only if $\delta$ is an equivalence. If $F^R$ is fully faithful, then $F^R$ is a monomorphism in $\widehat{\Cat}$, so $\delta$ is an equivalence. Conversely, suppose that $\delta$ is an equivalence, and let $x,y \in D$. Then the map
		\begin{align*}
			\hom_D(x,y) \to \hom_{D \times_C D}(\delta x,\delta y) &\simeq \hom_D(x,y) \times_{\hom_C(F^Rx,F^Ry)} \hom_D(x,y) \\
			&\simeq \hom_D(x,y) \oplus \fib\big(\hom_D(x,y) \to \hom_C(F^Rx,F^Ry)\big)
		\end{align*}
		is the diagonal, and it is an equivalence. Here the splitting of the pullback is induced by the first projection and its section given by the diagonal, so the diagonal corresponds to the inclusion of the first summand. We find that $\hom_D(x,y) \to \hom_C(F^Rx,F^Ry)$ is an equivalence. So $F^R$ is fully faithful.
	\end{proof}

	\begin{remark}
		\label[remark]{rem:Monomorphisms_Fully_Faithful}
		Dually, the monomorphisms in $\PrL_{\st}$, in $\Pr^{\mathrm{LL}}_{\st}$ and in $\Pr^{\dual}$ are precisely the fully faithful functors \cite[Propositions~1.13 and~1.14 and Corollary~2.87]{Efimov_Large_Categories}. For $\Pr^{\dual}$, this requires the description of pullbacks in $\Pr^{\dual}$, see \Cref{cor:Monomorphisms_Prdual} below.
	\end{remark}

	\section{Continuous Calkin categories}
	\label[section]{sec:Continuous_Calkin}

	% Source: Talk 5 (Anton Zakrewski, 10.11.2026); Zakrewski's notes for the talk; Efimov, Section 2.9
	For a small idempotent complete stable category $A$, the Calkin category $\Calk_{\kappa}(A)$ of \Cref{sec:Calkin_Small} fits into a short exact sequence $A \to \Ind(A)^{\kappa} \to \Calk_{\kappa}(A)$. For K-theory, the middle term is invisible: it admits countable coproducts, so its K-theory vanishes by the Eilenberg swindle, and $K(\Calk_{\kappa}(A)) \simeq \Sigma K(A)$. We now extend this construction to dualizable categories, which may have no nonzero compact objects at all. Throughout this section, $\kappa$ denotes an uncountable regular cardinal.

	Let $C$ be dualizable. By \Cref{thm:Compactly_Assembled_Omega_1}(1), $C$ is generated under colimits by $\omega_1$-compact objects, hence by a set of $\kappa$-compact objects, so $C$ is $\kappa$-compactly generated \cite[Lemma~A.21 and Addendum~A.22]{Ramzi_Dualizable_Presentable}. By \Cref{prop:Dualizable_Further_Characterizations}, the functor $\colim\colon \Ind(C^{\kappa}) \to C$ thus admits a fully faithful, strongly continuous left adjoint $\hat{y}\colon C \to \Ind(C^{\kappa})$. As in the proof of \Cref{thm:Lurie_Dualizable_Characterization}, its cofiber in $\PrL_{\st}$ is compactly generated.

	\begin{definition}[{\cite[Definition~2.58]{Efimov_Large_Categories}}]
		\label[definition]{def:Calkin_Continuous}
		Let $C$ be a dualizable category. The \emph{continuous $\kappa$-Calkin category} $\Calk^{\cont}_{\kappa}(C)$ is the small idempotent complete stable category of compact objects of the quotient
		\[
			\Ind(C^{\kappa})/\hat{y}(C) \simeq \ker\big(\colim\colon \Ind(C^{\kappa}) \to C\big),
		\]
		so that we have a strongly continuous localization sequence
		\[
			C \xrightarrow{\hat{y}} \Ind(C^{\kappa}) \to \Ind(\Calk^{\cont}_{\kappa}(C)).
		\]
	\end{definition}

	Here the identification of the quotient with the kernel comes from \Cref{prop:Localization_Sequences_PrL}: the right orthogonal of $\hat{y}(C)$ consists of those $X$ with $\Map(x,\colim X) \simeq \Map(\hat{y}x,X) \simeq *$ for all $x \in C$.

	% Source: written from the literature (Efimov, Proposition 2.59); moved from the small Calkin discussion.
	\begin{remark}
		\label[remark]{rem:Calkin_Small_Vs_Continuous}
		If $C = \Ind(A)$ is compactly generated, the continuous Calkin category $\Calk^{\cont}_{\kappa}(C)$ of \Cref{def:Calkin_Continuous} agrees with $\Calk_{\kappa}(A)$ \cite[Proposition~2.59]{Efimov_Large_Categories}. Indeed, by \Cref{ex:Compactly_Generated_Compactly_Assembled}, the functor $\hat{y}\colon C \to \Ind(C^{\kappa})$ is $\Ind(j)$ for the inclusion $j\colon A \hookrightarrow C^{\kappa} = \Ind(A)^{\kappa}$, so by \Cref{thm:Thomason_Neeman}(1) its cofiber is $\Ind\big((\Ind(A)^{\kappa}/A)^{\idem}\big) = \Ind(\Calk_{\kappa}(A))$.
	\end{remark}

	\begin{construction}
		\label[construction]{constr:Calkin_Functoriality}
		Let $F\colon C \to D$ be a strongly continuous functor between dualizable categories. Then $F$ preserves $\kappa$-compact objects, since $F^R$ preserves colimits, so it restricts to a functor $F^{\kappa}\colon C^{\kappa} \to D^{\kappa}$. The left adjoint of $\colim\colon \Ind(C) \to C$ is the composite of $\hat{y}\colon C \to \Ind(C^{\kappa})$ with the inclusion $\Ind(C^{\kappa}) \subseteq \Ind(C)$: indeed, by the proof of \Cref{prop:Dualizable_Further_Characterizations}, the colimit functor of $\Ind(C)$ factors as $\colim_{\kappa} \circ r$, where $r$ is the right adjoint of the inclusion $\Ind(C^{\kappa}) \subseteq \Ind(C)$. Since $\Ind(F)$ restricts to $\Ind(F^{\kappa})$ on $\Ind(C^{\kappa})$, \Cref{prop:Strongly_Continuous_Via_Right_Adjoint} gives a commutative square
		\[
			\begin{tikzcd}
				C \rar["\hat{y}_C"] \dar["F"'] & \Ind(C^{\kappa}) \dar["\Ind(F^{\kappa})"] \\
				D \rar["\hat{y}_D"'] & \Ind(D^{\kappa})
			\end{tikzcd}
		\]
		of strongly continuous functors. Passing to cofibers in $\Pr^{\mathrm{LL}}_{\st}$, which are computed in $\PrL_{\st}$ by \Cref{prop:PrLL_Colimits}, we obtain a strongly continuous functor $\Ind(\Calk^{\cont}_{\kappa}(C)) \to \Ind(\Calk^{\cont}_{\kappa}(D))$. It preserves compact objects, and hence restricts to an exact functor
		\[
			\Calk^{\cont}_{\kappa}(F)\colon \Calk^{\cont}_{\kappa}(C) \to \Calk^{\cont}_{\kappa}(D).
		\]
		The squares above are coherent in $F$: the colimit functors form a natural transformation $\Ind((-)^{\kappa}) \to \id$, and taking mates along $\hat{y}_C \dashv \colim_C$ gives the natural maps $\hat{y}_D F \to \Ind(F^{\kappa})\hat{y}_C$ used above. These are isomorphisms for strongly continuous $F$, by \Cref{prop:Strongly_Continuous_Via_Right_Adjoint}. The compatibility of mates with composition supplies a natural transformation $\id \to \Ind((-)^{\kappa})$ in $\Pr^{\mathrm{LL}}_{\st}$. Applying the cofiber functor and then taking compact objects therefore defines a functor $\Calk^{\cont}_{\kappa}\colon \Pr^{\dual} \to \Cat_{\perf}$, including its higher coherences.

		Under the identifications of the quotients with kernels, the functor $\Ind(\Calk^{\cont}_{\kappa}(C)) \to \Ind(\Calk^{\cont}_{\kappa}(D))$ is the restriction of $\Ind(F^{\kappa})$. Indeed, the functors $\colim_D \circ \Ind(F^{\kappa})$ and $F \circ \colim_C$ preserve colimits and agree on $C^{\kappa}$, so $\Ind(F^{\kappa})$ maps $\ker(\colim_C)$ into $\ker(\colim_D)$. The left adjoint $q_C$ of the inclusion $\ker(\colim_C) \subseteq \Ind(C^{\kappa})$ is given by $q_C(X) = \cofib(\hat{y}_C\colim_C X \to X)$. Using the square above, we get $\Ind(F^{\kappa})(q_CX) \simeq q_D(\Ind(F^{\kappa})X)$, which identifies the functor induced on the quotients with the restriction of $\Ind(F^{\kappa})$.
	\end{construction}

	% Source: Mathew, Lecture 5 (MPIM Bonn, 11.09.2024); Efimov, Proposition 2.62
	The mapping spectra in the continuous Calkin category have a simple description. For $x \in C^{\kappa}$ we write $\bar{x} \in \Calk^{\cont}_{\kappa}(C)$ for the image of $y(x)$ under the quotient functor $\Ind(C^{\kappa}) \to \Ind(\Calk^{\cont}_{\kappa}(C))$.

	\begin{proposition}[{\cite[Proposition~2.62]{Efimov_Large_Categories}}]
		\label[proposition]{prop:Calkin_Mapping_Spectra}
		Let $C$ be dualizable and $x,x' \in C^{\kappa}$. Then there is a natural equivalence
		\[
			\hom_{\Calk^{\cont}_{\kappa}(C)}(\bar{x},\bar{x}') \simeq \cofib\big(\hom_{\Ind(C)}(y(x),\hat{y}(x')) \to \hom_C(x,x')\big),
		\]
		where the map is induced by $\hat{y}(x') \to y(x')$. Moreover, $\Calk^{\cont}_{\kappa}(C)$ is generated by the objects $\bar{x}$ under finite colimits and retracts.
	\end{proposition}
	\prooffor{prop:Calkin_Mapping_Spectra}
	\begin{proof}
		Identify $\Ind(\Calk^{\cont}_{\kappa}(C))$ with $\ker(\colim) \subseteq \Ind(C^{\kappa})$ as in \Cref{def:Calkin_Continuous}. The quotient functor is then the left adjoint $q$ of the inclusion, with $q(X) \simeq \cofib(\hat{y}\colim X \to X)$, see \Cref{constr:Calkin_Functoriality}. Since $\colim y(x') \simeq x'$, we obtain
		\[
			\hom(\bar{x},\bar{x}') \simeq \hom_{\Ind(C^{\kappa})}\big(y(x),\cofib(\hat{y}(x') \to y(x'))\big) \simeq \cofib\big(\hom(y(x),\hat{y}(x')) \to \hom(y(x),y(x'))\big),
		\]
		and $\hom(y(x),y(x')) \simeq \hom_C(x,x')$. Mapping spectra in $\Ind(C^{\kappa})$ and $\Ind(C)$ agree, as $\Ind(C^{\kappa}) \subseteq \Ind(C)$ is fully faithful. Finally, $q$ is a localization preserving compact objects, and every object of the quotient is the image under $q$ of a filtered colimit of representables. If this image is compact, its identity factors through some $q(y(x))$, so it is a retract of $\bar{x}$.
	\end{proof}

	For example, if $C = \Mod_A$ for an $\bbE_1$-ring $A$, then $\hat{y}(x')$ is the ind-object of perfect modules over $x'$, and since $\hom_A(x,P) \simeq P \otimes_A \hom_A(x,A)$ for perfect $P$, one finds $\hom_{\Ind(C)}(y(x),\hat{y}(x')) \simeq x' \otimes_A \hom_A(x,A)$ \cite[after Proposition~2.62]{Efimov_Large_Categories}. This recovers the formula of \Cref{ex:Calkin_Perf}.

	The main property of the continuous Calkin construction is that it preserves short exact sequences. We first fix terminology.

	\begin{definition}
		\label[definition]{def:Short_Exact_Sequence_Dualizable}
		A \emph{short exact sequence} of dualizable categories is a strongly continuous localization sequence $C \to D \to E$ in $\PrL_{\st}$ (\Cref{def:Localization_Sequence_PrL}) in which $C$, $D$ and $E$ are dualizable. Equivalently, it is a sequence in $\Pr^{\dual}$ which is both a fiber sequence and a cofiber sequence in $\PrL_{\st}$ (\Cref{prop:Localization_Sequences_PrL}).
	\end{definition}

	By \Cref{prop:Localization_Sequence_Dualizable}, it suffices to require that $D$ is dualizable. A short exact sequence of dualizable categories is also both a fiber sequence and a cofiber sequence in $\Pr^{\dual}$, see \Cref{rem:Short_Exact_Fiber_Cofiber} below.

	% Source: Mathew, Lecture 5 (MPIM Bonn, 11.09.2024); Efimov, Proposition 2.73
	Short exact sequences are stable under pushout. This rests on the following observation.

	\begin{proposition}[{\cite[Proposition~2.73]{Efimov_Large_Categories}}]
		\label[proposition]{prop:Pushouts_Fully_Faithful}
		In each of $\PrL_{\st}$, $\Pr^{\mathrm{LL}}_{\st}$ and $\Pr^{\dual}$, the class of fully faithful functors is closed under pushouts.
	\end{proposition}
	\prooffor{prop:Pushouts_Fully_Faithful}
	\begin{proof}
		Pushouts in $\Pr^{\mathrm{LL}}_{\st}$ and $\Pr^{\dual}$ are computed in $\PrL_{\st}$ (\Cref{prop:PrLL_Colimits,prop:Prdual_Colimits}), so it suffices to treat $\PrL_{\st}$. Let $i\colon A \to B$ be fully faithful and $f\colon A \to A'$ arbitrary, and let $j\colon A' \to B' := B \sqcup_A A'$ be the induced functor. Computing the pushout along right adjoints, $B'$ is the pullback $B \times_A A'$ of $i^R$ and $f^R$, whose objects are triples $(b,a',\theta)$ with $\theta\colon i^Rb \simeq f^Ra'$, and $j^R$ is the projection $p$ to $A'$. Consider the functor $L\colon A' \to B'$, $a' \mapsto (if^Ra',a',\theta_{a'})$, where $\theta_{a'}\colon i^Rif^Ra' \simeq f^Ra'$ is the inverse of the unit. For $(b,a'',\theta) \in B'$, we have
		\[
			\Map_{B'}\big(La',(b,a'',\theta)\big) \simeq \Map_B(if^Ra',b) \times_{\Map_A(f^Ra',f^Ra'')} \Map_{A'}(a',a'') \simeq \Map_{A'}(a',a''),
		\]
		since the map $\Map_B(if^Ra',b) \simeq \Map_A(f^Ra',i^Rb) \xrightarrow{\theta} \Map_A(f^Ra',f^Ra'')$ is an equivalence. So $L$ is left adjoint to $p = j^R$, hence $L \simeq j$. As $pL \simeq \id$, the functor $j$ is fully faithful.
	\end{proof}

	\begin{corollary}
		\label[corollary]{cor:Pushout_Short_Exact}
		Let $A \to B \to E$ be a short exact sequence of dualizable categories, and let $A \to A'$ be a strongly continuous functor to a dualizable category $A'$. Then $A' \to B \sqcup_A A' \to E$ is a short exact sequence of dualizable categories, where the pushout is formed in $\Pr^{\dual}$, or equivalently in $\PrL_{\st}$.
	\end{corollary}
	\prooffor{cor:Pushout_Short_Exact}
	\begin{proof}
		Put $B' := B \sqcup_A A'$. It is dualizable by \Cref{prop:Prdual_Colimits}, and the functors $A' \to B' \to E$ are strongly continuous by \Cref{prop:PrLL_Colimits}. Pasting pushout squares shows that $B' \to E$ is the cofiber of $A' \to B'$ in $\PrL_{\st}$, and $A' \to B'$ is fully faithful by \Cref{prop:Pushouts_Fully_Faithful}. Hence $A' \to B' \to E$ is a localization sequence by \Cref{prop:Localization_Sequences_PrL}(1).
	\end{proof}

	\begin{proposition}[{\cite[Proposition~2.61]{Efimov_Large_Categories}}]
		\label[proposition]{prop:Calkin_Short_Exact}
		The functor $\Calk^{\cont}_{\kappa}\colon \Pr^{\dual} \to \Cat_{\perf}$ sends short exact sequences of dualizable categories to short exact sequences.
	\end{proposition}
	\prooffor{prop:Calkin_Short_Exact}
	\begin{proof}
		Let $C \xrightarrow{F} D \xrightarrow{G} E$ be a short exact sequence of dualizable categories. As shown above, $C$ and $D$ are $\kappa$-compactly generated. Since $G$ preserves $\kappa$-compact objects, \cite[Proposition~A.26]{Ramzi_Dualizable_Presentable} shows that $C^{\kappa} \to D^{\kappa} \to E^{\kappa}$ is a short exact sequence. By \Cref{thm:Thomason_Neeman}(1), the sequence $\Ind(C^{\kappa}) \to \Ind(D^{\kappa}) \to \Ind(E^{\kappa})$ is a strongly continuous localization sequence. By \Cref{constr:Calkin_Functoriality}, we obtain a commutative diagram
		\[
			\begin{tikzcd}
				C \rar["\hat{y}"] \dar["F"'] & \Ind(C^{\kappa}) \rar \dar & \Ind(\Calk^{\cont}_{\kappa}(C)) \dar["\Phi"] \\
				D \rar["\hat{y}"] \dar["G"'] & \Ind(D^{\kappa}) \rar \dar & \Ind(\Calk^{\cont}_{\kappa}(D)) \dar \\
				E \rar["\hat{y}"'] & \Ind(E^{\kappa}) \rar & \Ind(\Calk^{\cont}_{\kappa}(E))
			\end{tikzcd}
		\]
		in which the rows and the first two columns are cofiber sequences in $\PrL_{\st}$. Since colimits commute with colimits, the third column is a cofiber sequence in $\PrL_{\st}$ as well.

		We claim that $\Phi$ is fully faithful. By \Cref{constr:Calkin_Functoriality}, $\Phi$ is identified with the restriction of $\Ind(F^{\kappa})$ to the kernels of the colimit functors. Since $F$ is fully faithful, so are $F^{\kappa}$ and $\Ind(F^{\kappa})$, and hence so is $\Phi$.

		Since $\Phi$ is fully faithful and colimit-preserving, its essential image is closed under colimits, and the cofiber sequence of the third column is a localization sequence by \Cref{prop:Localization_Sequences_PrL}(1). All functors involved preserve compact objects, so it is strongly continuous. By \Cref{thm:Thomason_Neeman}(2), the sequence $\Calk^{\cont}_{\kappa}(C) \to \Calk^{\cont}_{\kappa}(D) \to \Calk^{\cont}_{\kappa}(E)$ of compact objects is short exact.
	\end{proof}

	\section{Filtered colimits of dualizable categories}
	\label[section]{sec:Filtered_Colimits_Dualizable}

	% Source: Talk 5 (Anton Zakrewski, 10.11.2026); Zakrewski's notes for the talk; Efimov, Propositions 2.65, 2.67 and 2.70, Corollary 2.68 and Remark 2.69
	Filtered colimits in $\Cat_{\perf}$ are exact: they preserve fully faithful functors and short exact sequences. The analogous statement fails in $\PrL_{\st}$ (\Cref{rem:Filtered_Colimits_Fail_PrL}). We now show that it does hold in $\Pr^{\mathrm{LL}}_{\st}$ and in $\Pr^{\dual}$, whose colimits are computed in $\PrL_{\st}$ by \Cref{sec:Colimits_Dualizable}. As a consequence, the functor $C \mapsto C^{\omega}$ on $\Pr^{\dual}$ preserves filtered colimits. By \Cref{thm:Filtered_Posets_Cofinal}, every filtered colimit can be computed as a colimit over a directed poset, so we may restrict attention to directed posets.

	\begin{lemma}[{\cite[Proposition~2.65]{Efimov_Large_Categories}}]
		\label[lemma]{lem:Filtered_Colimits_PrL}
		Let $I$ be a directed poset, let $I \to \PrL_{\st}$, $i \mapsto C_i$, be a diagram with transition functors $F_{ij}\colon C_i \to C_j$, let $C := \colim_i C_i$ in $\PrL_{\st}$, and let $F_i\colon C_i \to C$ be the structure functors.
		\begin{enumerate}
			\item The counits induce an equivalence $\colim_i F_iF_i^R \simeq \id_C$.
			\item If the functors $F_{ij}$ are strongly continuous, then for all $i,j \in I$ there is a natural equivalence
			\[
				\colim_{k \geq i,j} F_{jk}^RF_{ik} \simeq F_j^RF_i.
			\]
		\end{enumerate}
	\end{lemma}
	\prooffor{lem:Filtered_Colimits_PrL}
	\begin{proof}
		We identify $C$ with $\lim_i C_i$, taken along the functors $F_{ij}^R$, so that $F_i^R$ becomes the projection. Mapping spectra in this limit are the limits of the mapping spectra.

		(1) For $i \leq j$, the counit of $F_{ij} \dashv F_{ij}^R$ induces a map $F_iF_i^R \simeq F_jF_{ij}F_{ij}^RF_j^R \to F_jF_j^R$, so the functors $F_iF_i^R$ form an $I$-indexed diagram with a cocone to $\id_C$, given by the counits. For $x,y \in C$ we get natural equivalences
		\[
			\hom_C\big(\colim_i F_iF_i^Rx,y\big) \simeq \lim_i \hom_C(F_iF_i^Rx,y) \simeq \lim_i \hom_{C_i}(F_i^Rx,F_i^Ry) \simeq \hom_C(x,y),
		\]
		whose composite is induced by the cocone. By the Yoneda lemma, the cocone exhibits $\id_C$ as the colimit.

		(2) By \Cref{prop:PrLL_Colimits}(1), the functors $F_i$ are strongly continuous, so the projections $F_j^R$ preserve colimits, and the transition functors $F_{jj'}^R$ of the limit preserve colimits by assumption. Fix $i$, and for $x \in C_i$ consider the family of objects
		\[
			F_i'(x)_j := \colim_{k \geq i,j} F_{jk}^RF_{ik}(x) \in C_j, \qquad j \in I.
		\]
		For $k \leq \ell$, the transition map in this colimit is
		\[
			F_{jk}^RF_{ik} \longrightarrow F_{jk}^RF_{k\ell}^RF_{k\ell}F_{ik} \simeq F_{j\ell}^RF_{i\ell},
		\]
		induced by the unit of $F_{k\ell} \dashv F_{k\ell}^R$. The units for composites give these diagrams coherently. For $j \leq j'$, we have $F_{jj'}^R F_i'(x)_{j'} \simeq \colim_{k \geq i,j'} F_{jk}^RF_{ik}(x) \simeq F_i'(x)_j$, using that $F_{jj'}^R$ preserves colimits and that $\{k \geq i,j'\} \subseteq \{k \geq i,j\}$ is cofinal. These identifications are induced by restriction to common cofinal tails and the composition isomorphisms for right adjoints. Their compatibility for chains of indices follows from that of the units above and the universal property of colimits. Thus they define a functor $F_i'\colon C_i \to \lim_j C_j \simeq C$, and the claim is equivalent to $F_i' \simeq F_i$. We show that $F_i'$ is left adjoint to $F_i^R$. For $x \in C_i$ and $y = (y_j)_j \in C$, we have
		\begin{align*}
			\hom_C(F_i'x,y) &\simeq \lim_j \hom_{C_j}\Big(\colim_{k \geq i,j} F_{jk}^RF_{ik}x,\, y_j\Big) \\
			&\simeq \lim_{(j,k) \in P\catop} \hom_{C_j}(F_{jk}^RF_{ik}x,y_j) \\
			&\simeq \lim_{j \geq i} \hom_{C_j}(F_{ij}x,y_j) \simeq \lim_{j \geq i} \hom_{C_i}(x,F_{ij}^Ry_j) \\
			&\simeq \hom_{C_i}(x,y_i) = \hom_{C_i}(x,F_i^Ry).
		\end{align*}
		Here $P = \{(j,k) \mid i \leq j \leq k\}$, with the product order. In the second step we first restrict the limit over $j$ to the cofinal subset $I_{\geq i}$. The functor $(j,k) \mapsto \hom_{C_j}(F_{jk}^RF_{ik}x,y_j)$ is contravariant in $(j,k)$. In the third step we restrict along the diagonal $I_{\geq i} \to P$, $j \mapsto (j,j)$, which is cofinal: for $(j,k) \in P$, the poset of $m \geq i$ with $(j,k) \leq (m,m)$ is $I_{\geq k}$, which is directed and hence weakly contractible. Thus the opposite of the diagonal is initial, as required for this limit over $P\catop$. On the diagonal, $F_{jj}^R = \id$. The fourth step is the adjunction $F_{ij} \dashv F_{ij}^R$, and the remaining diagram is constant with value $\hom_{C_i}(x,y_i)$, as $F_{ij}^Ry_j \simeq y_i$. The unit of the resulting adjunction $F_i' \dashv F_i^R$ is the map $x \to F_i'(x)_i$ from the summand $k=i$. To identify the resulting isomorphism with the canonical map, use the cocones
		\[
			F_{jk}^RF_{ik} \longrightarrow F_{jk}^RF_k^RF_kF_{ik} \simeq F_j^RF_i,
		\]
		induced by the unit of $F_k \dashv F_k^R$. They define $F_i' \to F_i$, and at $j=k=i$ carry the unit just described to the unit of $F_i \dashv F_i^R$. By the adjunction, this map is the isomorphism between the two left adjoints. Evaluating at the $j$-th component proves the claim with its canonical cocone.
	\end{proof}

	\begin{proposition}[{\cite[Proposition~2.67]{Efimov_Large_Categories}}]
		\label[proposition]{prop:Filtered_Colimits_Fully_Faithful}
		\leavevmode
		\begin{enumerate}
			\item In $\Pr^{\mathrm{LL}}_{\st}$ and in $\Pr^{\dual}$, the class of fully faithful functors is closed under filtered colimits.
			\item A filtered colimit of strongly continuous localization sequences, formed in $\PrL_{\st}$, is again a strongly continuous localization sequence. In particular, filtered colimits of short exact sequences of dualizable categories are short exact.
		\end{enumerate}
	\end{proposition}
	\prooffor{prop:Filtered_Colimits_Fully_Faithful}
	\begin{proof}
		(1) By \Cref{prop:Prdual_Colimits}, it suffices to treat $\Pr^{\mathrm{LL}}_{\st}$. Let $I$ be a directed poset, let $(C_i)$ and $(D_i)$ be $I$-indexed diagrams in $\Pr^{\mathrm{LL}}_{\st}$ with transition functors $F_{ij}$ and $G_{ij}$, and let $\Phi_i\colon C_i \to D_i$ be a natural transformation whose components are fully faithful and strongly continuous. Let $C = \colim_i C_i$ and $D = \colim_i D_i$, with structure functors $F_i$ and $G_i$, and let $\Phi\colon C \to D$ be the induced functor, which is strongly continuous by \Cref{prop:PrLL_Colimits}. We show that the unit $\id_C \to \Phi^R\Phi$ is an equivalence. Both functors preserve colimits, the essential images of the $F_i$ generate $C$ under colimits, and the functors $F_j^R$ are jointly conservative (see the proof of \Cref{prop:Prdual_Colimits}). So it suffices to show that $F_j^RF_i \to F_j^R\Phi^R\Phi F_i$ is an equivalence for all $i,j \in I$. Using $\Phi F_i \simeq G_i\Phi_i$, \Cref{lem:Filtered_Colimits_PrL}(2), the strong continuity of $\Phi_j$ and the fully faithfulness of $\Phi_k$, we compute
		\begin{align*}
			F_j^R\Phi^R\Phi F_i &\simeq \Phi_j^RG_j^RG_i\Phi_i \simeq \Phi_j^R\Big(\colim_{k \geq i,j} G_{jk}^RG_{ik}\Big)\Phi_i \simeq \colim_{k \geq i,j} \Phi_j^RG_{jk}^RG_{ik}\Phi_i \\
			&\simeq \colim_{k \geq i,j} F_{jk}^R\Phi_k^R\Phi_kF_{ik} \simeq \colim_{k \geq i,j} F_{jk}^RF_{ik} \simeq F_j^RF_i.
		\end{align*}
		More precisely, under the canonical cocones in \Cref{lem:Filtered_Colimits_PrL}(2), the map $F_j^RF_i \to F_j^R\Phi^R\Phi F_i$ is the colimit of the maps
		\[
			F_{jk}^RF_{ik} \longrightarrow F_{jk}^R\Phi_k^R\Phi_kF_{ik}
		\]
		induced by the units of $\Phi_k \dashv \Phi_k^R$. This follows by pasting the adjunction units for the squares $\Phi_kF_{ik} \simeq G_{ik}\Phi_i$ and for the colimit cocones. Each of these maps is an isomorphism because $\Phi_k$ is fully faithful, so their colimit is an isomorphism as required.

		(2) Let $C_i \to D_i \to E_i$ be a filtered diagram of strongly continuous localization sequences, and let $C \to D \to E$ be its colimit in $\PrL_{\st}$. By \Cref{prop:PrLL_Colimits}, it is a sequence of strongly continuous functors. Since colimits commute with colimits, $D \to E$ is the cofiber of $C \to D$ in $\PrL_{\st}$. By (1), $C \to D$ is fully faithful, and it preserves colimits. Hence $C \to D \to E$ is a localization sequence by \Cref{prop:Localization_Sequences_PrL}(1). If the $D_i$ are dualizable, then so is $D$ by \Cref{prop:Prdual_Colimits}.
	\end{proof}

	\begin{corollary}[{\cite[Corollary~2.68]{Efimov_Large_Categories}}]
		\label[corollary]{cor:Filtered_Union_Subcategories}
		Let $C$ be a presentable stable category, and let $(C_i)_{i \in I}$ be a directed family of presentable stable full subcategories $C_i \subseteq C$ whose inclusion functors are strongly continuous. If $C$ is generated under colimits by the $C_i$, then the functor $\colim_i C_i \to C$ is an equivalence, where the colimit is formed in $\PrL_{\st}$.
	\end{corollary}
	\prooffor{cor:Filtered_Union_Subcategories}
	\begin{proof}
		For $i \leq j$, the right adjoint of $C_i \hookrightarrow C_j$ is the restriction of the right adjoint of $C_i \hookrightarrow C$, so the transition functors are strongly continuous. The colimit of the constant diagram with value $C$ over the weakly contractible poset $I$ is $C$. Applying \Cref{prop:Filtered_Colimits_Fully_Faithful}(1) to the inclusions $C_i \hookrightarrow C$ shows that $\colim_i C_i \to C$ is fully faithful. Its essential image is closed under colimits and contains all $C_i$, so it is all of $C$.
	\end{proof}

	\begin{remark}
		\label[remark]{rem:Filtered_Colimits_Fail_PrL}
		\Cref{prop:Filtered_Colimits_Fully_Faithful} and \Cref{cor:Filtered_Union_Subcategories} fail in $\PrL_{\st}$, i.e.\ without the strong continuity assumptions. Efimov \cite[Remark~2.69]{Efimov_Large_Categories} considers $C = \D(\Z)$ and $C_n = \D(\Z[p_k^{-1} \mid k \geq n]) \subseteq C$, where $p_k$ is the $k$-th prime number. The subcategories $C_1 \subseteq C_2 \subseteq \cdots$ generate $C$ under colimits, but $\colim_n C_n \to C$ is not an equivalence, as one shows using \Cref{lem:Filtered_Colimits_PrL}(1).
	\end{remark}

	To compute compact objects in filtered colimits, we use the following observation. Let $C \to \Ind(A) \to \Ind(B)$ be a short exact sequence of dualizable categories, where $A, B \in \Cat_{\perf}$ and the second functor is $\Ind(p)$ for a functor $p\colon A \to B$. Then
	\begin{equation}
		\label{eq:Compacts_Of_Kernel}
		C^{\omega} \simeq \ker(p\colon A \to B).
	\end{equation}
	Indeed, a fully faithful strongly continuous functor preserves compact objects, and it detects them since it preserves filtered colimits and is fully faithful. So $C^{\omega}$ consists of the compact objects of $\Ind(A)$, i.e.\ the objects of $A$, which lie in $C = \ker(\Ind(p))$.

	% Source: written from the literature (Lurie, Higher Algebra, Theorem 1.1.4.4 and Proposition 1.1.4.6); supplementary arguments added when checking Talk 5.
	We recall the categorical facts used below. Limits and filtered colimits of stable categories and exact functors are computed in $\Cat$ \cite[Theorem~1.1.4.4 and Proposition~1.1.4.6]{lurie2016HA}. The same holds for $\Cat_{\perf}$: idempotent completeness is closed under limits; in a filtered colimit, an idempotent in the homotopy category is represented by an endomorphism at some stage, and its idempotence holds at a later stage, where it splits. Here we use the characterization of idempotent completeness of a stable category by idempotent completeness of its homotopy category.

	Filtered colimits in $\Cat$ commute with finite limits, and taking the core preserves filtered colimits. One can check these two facts in the quasicategory model: filtered colimits are computed on underlying simplicial sets; an edge is invertible precisely when it has a two-sided inverse in the homotopy category, a condition witnessed by finite data. Thus taking the maximal Kan complex commutes with filtered colimits. Pullbacks of categories can be computed by triples $(x,y,u)$ with $u$ an isomorphism between the two images; this construction, as well as finite products, commutes with filtered colimits. Consequently filtered colimits in $\Cat_{\perf}$ commute with finite limits as well.

	\begin{proposition}[{\cite[Proposition~2.70]{Efimov_Large_Categories}}]
		\label[proposition]{prop:Compact_Objects_Filtered_Colimits}
		The functor $(-)^{\omega}\colon \Pr^{\dual} \to \Cat_{\perf}$ preserves filtered colimits. In particular, $\Sp$ is a compact object of $\Pr^{\dual}$.
	\end{proposition}
	\prooffor{prop:Compact_Objects_Filtered_Colimits}
	\begin{proof}
		Let $(C_i)_{i \in I}$ be a filtered diagram in $\Pr^{\dual}$, and put $C := \colim_i C_i$, $A := \colim_i C_i^{\omega_1}$ and $B := \colim_i \Calk^{\cont}_{\omega_1}(C_i)$, the latter two colimits taken in $\Cat_{\perf}$. The functor $\Ind\colon \Cat_{\perf} \to \PrL_{\st}$ preserves colimits: for $T \in \PrL_{\st}$, we have $\Fun^{\colim}(\Ind(A'),T) \simeq \Fun^{\ex}(A',T)$ by \Cref{thm:Ind_Kappa_Universal_Property}(3), and $\Fun^{\ex}(-,T)$ sends colimits in $\Cat_{\perf}$ to limits, since $T$ is idempotent complete. By \Cref{prop:Calkin_Short_Exact,prop:Filtered_Colimits_Fully_Faithful}(2), applied to the Calkin sequences of the $C_i$ for $\kappa = \omega_1$, we therefore obtain a short exact sequence
		\[
			C \to \Ind(A) \to \Ind(B).
		\]
		By \eqref{eq:Compacts_Of_Kernel}, applied to this sequence and to the Calkin sequences of the $C_i$, we get
		\[
			C^{\omega} \simeq \ker(A \to B) \simeq \colim_i \ker\big(C_i^{\omega_1} \to \Calk^{\cont}_{\omega_1}(C_i)\big) \simeq \colim_i C_i^{\omega},
		\]
		where the middle equivalence uses the commutation of filtered colimits with finite limits established above.

		Finally, a colimit-preserving functor $\Sp \to C$ is determined by the image $x$ of $\S$, and it is strongly continuous if and only if its right adjoint $\hom_C(x,-)$ preserves colimits, i.e.\ if and only if $x$ is compact. So $\Map_{\Pr^{\dual}}(\Sp,C) \simeq (C^{\omega})^{\simeq}$. Since $(-)^{\simeq}\colon \Cat_{\perf} \to \An$ preserves filtered colimits, $\Sp$ is compact in $\Pr^{\dual}$.
	\end{proof}

	To see why the core functor is conservative on stable categories and exact functors, observe that $\Map_C(x,y)$ is a retract of the anima of automorphisms of $x \oplus y$: send $f$ to the shear $\left(\begin{smallmatrix}1&0\\f&1\end{smallmatrix}\right)$, and send an automorphism $u$ to $p_yui_x$. These maps are natural under exact functors. A functor inducing an isomorphism on cores therefore induces isomorphisms on these automorphism animae and hence on their retracts $\Map_C(x,y)$; it is also essentially surjective, so it is an equivalence.

	Conversely, compactness of $\Sp$ implies that $(-)^{\omega}$ preserves filtered colimits, since $(-)^{\simeq}\colon \Cat_{\perf} \to \An$ is conservative \cite[Proposition~2.70]{Efimov_Large_Categories}.

	\section{Finite limits of dualizable categories}
	\label[section]{sec:Finite_Limits_Dualizable}

	% Source: Talk 6 (Jacopo Mazzei, 17.11.2026); Efimov, Section 2.12
	Limits in $\Pr^{\dual}$ are much more complicated than colimits. In general, the kernel in $\PrL_{\st}$ of a strongly continuous functor between dualizable categories is not dualizable, or its inclusion is not strongly continuous (\Cref{ex:Dualizable_Kernel_DVR}). Instead, the kernel in $\Pr^{\dual}$ is the largest dualizable subcategory of the kernel whose inclusion is strongly continuous. The existence of such a largest subcategory relies on the following refinement of the description of $\omega_1$-compact objects in terms of compact morphisms.

	\begin{proposition}[{\cite[Proposition~2.81]{Efimov_Large_Categories}}]
		\label[proposition]{prop:Q_Exhaustion}
		Let $C$ be a dualizable category and let $x \in C^{\omega_1}$. Then there exists a functor $\Phi\colon (\Q,\leq) \to C$ such that
		\begin{enumerate}
			\item $\colim_{a \in \Q} \Phi(a) \simeq x$;
			\item for every $a \in \Q$, the canonical map $\colim_{b < a} \Phi(b) \to \Phi(a)$ is an equivalence;
			\item for all $a < b$ in $\Q$, the morphism $\Phi(a) \to \Phi(b)$ is compact.
		\end{enumerate}
	\end{proposition}
	\prooffor{prop:Q_Exhaustion}
	\begin{proof}
		By \Cref{prop:Compact_Interpolation}(3), there is a functor $F\colon \Q \to C$ with $\colim_{\Q} F \simeq x$ such that $F(a) \to F(b)$ is compact for all $a < b$. Define $\Phi(a) := \colim_{b < a} F(b)$; formally, $\Phi$ is the left Kan extension of $F \circ \mathrm{pr}_1$ along $\mathrm{pr}_2\colon \{(b,a) \in \Q^2 \mid b < a\} \to \Q$. For $a < a'$, choose $c$ with $a < c < a'$. Then $\Phi(a) \to \Phi(a')$ factors as
		\[
			\Phi(a) \to F(a) \to F(c) \to \Phi(a'),
		\]
		so it is compact, since $F(a) \to F(c)$ is compact and compact morphisms form an ideal (\Cref{rem:Compact_Morphisms_Ideal}). Since $\Q$ is dense, we have $\colim_{b < a}\Phi(b) \simeq \colim_{b < a}\colim_{c < b} F(c) \simeq \colim_{c < a} F(c) = \Phi(a)$, and similarly $\colim_{\Q}\Phi \simeq \colim_{\Q} F \simeq x$.
	\end{proof}

	\begin{lemma}[{\cite[Lemma~2.82]{Efimov_Large_Categories}}]
		\label[lemma]{lem:Largest_Dualizable_Subcategory}
		Let $C$ be a dualizable category, and let $D \subseteq C$ be a full stable subcategory closed under colimits. Then there is a largest full subcategory $E \subseteq D$ such that $E$ is dualizable and the inclusion $E \hookrightarrow C$ is strongly continuous. Moreover, $E$ is generated under colimits by the objects $\colim_{\Q} F$, where $F\colon \Q \to D$ is a functor such that $F(a) \to F(b)$ is compact in $C$ for all $a < b$.
	\end{lemma}
	\prooffor{lem:Largest_Dualizable_Subcategory}
	\begin{proof}
		Let $E' \subseteq D$ be the full stable subcategory generated under colimits by the objects $\colim_{\Q} F$ as in the statement. Each such object is the colimit of the sequence $F(0) \to F(1) \to \cdots$ of compact morphisms, hence is $\omega_1$-compact in $C$ by \Cref{thm:Compactly_Assembled_Omega_1}(2) and \Cref{rem:Omega_1_Compact_Compact_Morphisms}. As $C^{\omega_1}$ is essentially small, $E'$ is generated under colimits by a set of objects, and hence is presentable.

		Let $F\colon \Q \to D$ be as in the statement. Replacing $F$ by $a \mapsto \colim_{b < a} F(b)$ as in the proof of \Cref{prop:Q_Exhaustion}, which does not change $\colim_{\Q} F$ or leave $D$, we may assume that $\colim_{b < a} F(b) \simeq F(a)$ for all $a$. Then each $F(a)$ lies in $E'$, being the colimit of the restriction of $F$ to $\Q_{<a}$, which is order-isomorphic to $\Q$. As in the proof of \Cref{prop:Strongly_Continuous_Compact_Morphisms}, we have
		\[
			\hat{y}_C\big(\colim_{\Q} F\big) \simeq \quotcolim\big(F(0) \to F(1) \to F(2) \to \cdots\big),
		\]
		since the maps $F(n) \to F(n+1)$ are compact. This lies in the essential image of the fully faithful functor $\Ind(E') \to \Ind(C)$. This essential image is closed under colimits, since $\Ind(E') \to \Ind(C)$ is exact and preserves filtered colimits. Since $\hat{y}_C$ preserves colimits, it follows that $\hat{y}_C$ maps $E'$ into $\Ind(E')$. For $e \in E'$ and $Z \in \Ind(E')$, we then have
		\[
			\Map_{\Ind(E')}(\hat{y}_Ce,Z) \simeq \Map_{\Ind(C)}(\hat{y}_Ce,Z) \simeq \Map_C(e,\colim_C Z) \simeq \Map_{E'}(e,\colim_{E'}Z),
		\]
		as the inclusion $E' \hookrightarrow C$ preserves colimits. So $\colim_{E'}\colon \Ind(E') \to E'$ admits the left adjoint $\hat{y}_C|_{E'}$, i.e.\ $E'$ is compactly assembled, hence dualizable by \Cref{thm:Lurie_Dualizable_Characterization}. Moreover, the inclusion $E' \hookrightarrow C$ commutes with $\hat{y}$ by construction, so it is strongly continuous (\Cref{def:Strongly_Continuous_Compactly_Assembled} and \Cref{rem:Strongly_Continuous_Right_Adjoint}).

		Now let $E'' \subseteq D$ be any full subcategory such that $E''$ is dualizable and the inclusion $E'' \hookrightarrow C$ is strongly continuous. Let $x \in E''$ be $\omega_1$-compact. By \Cref{prop:Q_Exhaustion}, there is a functor $F\colon \Q \to E''$ with $\colim_{\Q} F \simeq x$ such that $F(a) \to F(b)$ is compact in $E''$ for $a < b$. These maps are compact in $C$ by \Cref{prop:Strongly_Continuous_Compact_Morphisms}, and $\colim_{\Q} F \simeq x$ in $C$ since the inclusion preserves colimits. Hence $x \in E'$. As $E''$ is generated under colimits by its $\omega_1$-compact objects (\Cref{thm:Compactly_Assembled_Omega_1}(1)), we get $E'' \subseteq E'$. So $E := E'$ is the largest such subcategory.
	\end{proof}

	\begin{remark}
		\label[remark]{rem:Dualizable_Coreflection}
		Efimov also shows that the inclusion $\Pr^{\dual} \hookrightarrow \Pr^{\mathrm{LL}}_{\st}$ admits a right adjoint \cite[Remark~2.57]{Efimov_Large_Categories}: for a presentable stable category $C$, the strictly full dualizable subcategories of $C$ with strongly continuous inclusion form a small complete lattice, whose join is given by images of products in the sense of \Cref{def:Image_Functor} (using \Cref{cor:Image_Dualizable} and \Cref{rem:Products_Dualizable}). Its largest element is the value of the right adjoint at $C$.
	\end{remark}

	\begin{definition}
		\label[definition]{def:Dualizable_Kernel}
		Let $F\colon C \to D$ be a strongly continuous functor between dualizable categories, and let $\ker(F) \subseteq C$ be its kernel in $\PrL_{\st}$. The \emph{dualizable kernel} $\ker^{\dual}(F) \subseteq \ker(F)$ is the largest full subcategory of $\ker(F)$ which is dualizable and whose inclusion into $C$ is strongly continuous; it exists by \Cref{lem:Largest_Dualizable_Subcategory}.
	\end{definition}

	\begin{proposition}[{\cite[Proposition~2.83]{Efimov_Large_Categories}}]
		\label[proposition]{prop:Dualizable_Kernel}
		Let $F\colon C \to D$ be a strongly continuous functor between dualizable categories. Then $\ker^{\dual}(F) \to C \xrightarrow{F} D$ is a fiber sequence in $\Pr^{\dual}$.
	\end{proposition}
	\prooffor{prop:Dualizable_Kernel}
	\begin{proof}
		Write $j\colon \ker^{\dual}(F) \hookrightarrow C$ for the inclusion. Let $T$ be dualizable and let $G\colon T \to C$ be strongly continuous with $FG \simeq 0$. Then $\Im(G) \subseteq \ker(F)$, and by \Cref{cor:Image_Dualizable}, $\Im(G)$ is dualizable and its inclusion into $C$ is strongly continuous. By maximality, $\Im(G) \subseteq \ker^{\dual}(F)$, so $G \simeq jG'$ for a colimit-preserving functor $G'\colon T \to \ker^{\dual}(F)$. Its right adjoint $G^Rj$ preserves colimits, so $G'$ is strongly continuous. Since $j$ is fully faithful, $j \circ -$ identifies $\Map_{\Pr^{\dual}}(T,\ker^{\dual}(F))$ with the union of those path components of $\Map_{\Pr^{\dual}}(T,C)$ which consist of functors $G$ with $FG \simeq 0$. Since a nullhomotopy of $FG$ is unique if it exists (\Cref{rem:Nullhomotopies_Property}), the latter is the fiber of $F \circ -\colon \Map_{\Pr^{\dual}}(T,C) \to \Map_{\Pr^{\dual}}(T,D)$ over the zero functor.
	\end{proof}

	\begin{example}[{\cite[Example~2.84]{Efimov_Large_Categories}}]
		\label[example]{ex:Dualizable_Kernel_DVR}
		Let $R$ be a discrete valuation ring with maximal ideal $\mathfrak{m}$ and fraction field $K$, and consider the base change functor $F\colon \D(R) \to \D(R/\mathfrak{m})$, which is strongly continuous. Its kernel in $\PrL_{\st}$ consists of the $R$-modules on which a uniformizer acts invertibly, i.e.\ $\ker(F) \simeq \D(K)$. However, the inclusion $\D(K) \hookrightarrow \D(R)$ is not strongly continuous, since its right adjoint $\hom_R(K,-)$ does not preserve colimits, $K$ not being a compact $R$-module. The only localizing subcategories of $\D(K)$ are $0$ and $\D(K)$, so $\ker^{\dual}(F) = 0$.
	\end{example}

	Efimov uses a variant of this example to show that sequential colimits in $\Pr^{\dual}$ do not commute with kernels, so that $\Pr^{\dual}$ is not compactly assembled \cite[Corollary~2.85]{Efimov_Large_Categories}.

	% Source: Talk 6 (Jacopo Mazzei, 17.11.2026); Efimov, Proposition 2.86; part (2) and the details of part (1) added from the literature with generative AI (Claude)
	We now describe pullbacks in $\Pr^{\dual}$. Let $F\colon A \to C$ and $G\colon B \to C$ be strongly continuous functors between dualizable categories. Their \emph{lax pullback} $L$ is the category of triples $(a,b,f)$ with $a \in A$, $b \in B$ and $f\colon F(a) \to G(b)$. By adjunction, it is equivalent to the gluing of \Cref{constr:Semiorthogonal_Gluing} along the colimit-preserving functor $F^RG\colon B \to A$. By \Cref{prop:Gluing_Dualizable}, $L$ is dualizable, and we have a semiorthogonal decomposition $L = \lra{i_A(A),i_B(B)}$ with $i_A(a) = (a,0,0)$ and $i_B(b) = (0,b,0)$, and with projections $\pi_A(a,b,f) = a$ and $\pi_B(a,b,f) = b$. Consider the colimit-preserving functor
	\[
		\Phi\colon L \to C, \qquad (a,b,f) \mapsto \cofib\big(f\colon F(a) \to G(b)\big).
	\]
	Since $\Phi i_A \simeq \Sigma F$ and $\Phi i_B \simeq G$ are strongly continuous, $\Phi$ is strongly continuous by \Cref{prop:SOD_Functor_Categories}(2). Its kernel $\ker(\Phi)$ is the full subcategory of triples in which $f$ is an equivalence, i.e.\ $\ker(\Phi) \simeq A \times_C B$ is the pullback in $\PrL_{\st}$, which is computed in $\widehat{\Cat}$.

	\begin{proposition}[{\cite[Proposition~2.86]{Efimov_Large_Categories}}]
		\label[proposition]{prop:Dualizable_Pullbacks}
		Let $F\colon A \to C$ and $G\colon B \to C$ be strongly continuous functors between dualizable categories, and let $A \times_C^{\dual} B := \ker^{\dual}(\Phi)$, with $\Phi$ as above.
		\begin{enumerate}
			\item The category $A \times^{\dual}_C B$ is a pullback of $F$ and $G$ in $\Pr^{\dual}$. The functor $A \times^{\dual}_C B \to A \times_C B$ is fully faithful and strongly continuous, and its essential image is the largest full subcategory $E \subseteq A \times_C B$ such that $E$ is dualizable and the functors $E \to A$ and $E \to B$ are strongly continuous.
			\item If $G$ is a localization, then $A \times^{\dual}_C B \simeq A \times_C B$, and the projection $A \times_C B \to A$ is a localization.
		\end{enumerate}
	\end{proposition}
	\prooffor{prop:Dualizable_Pullbacks}
	\begin{proof}
		(1) Let $T$ be dualizable. By \Cref{prop:Dualizable_Kernel}, $\Map_{\Pr^{\dual}}(T,A \times^{\dual}_C B)$ is the union of the path components of $\Map_{\Pr^{\dual}}(T,L)$ consisting of functors $H$ with $\Phi H \simeq 0$, i.e.\ of strongly continuous functors $T \to L$ landing in $\ker(\Phi) = A \times_C B$. By \Cref{prop:SOD_Functor_Categories}(1), a colimit-preserving functor $H\colon T \to L$ is strongly continuous if and only if $\pi_AH$ and $\pi_BH$ are. So $\Map_{\Pr^{\dual}}(T,A \times^{\dual}_C B)$ is the union of the path components of $\Map_{\PrL_{\st}}(T,A \times_C B)$ consisting of functors whose two components are strongly continuous. As $A \times_C B$ is a pullback in $\PrL_{\st}$, this is
		\[
			\Map_{\Pr^{\dual}}(T,A) \times_{\Map_{\Pr^{\dual}}(T,C)} \Map_{\Pr^{\dual}}(T,B),
		\]
		naturally in $T$. This proves that $A \times^{\dual}_C B$ is a pullback in $\Pr^{\dual}$.

		Write $k\colon A \times^{\dual}_C B \hookrightarrow A \times_C B$ and $j\colon A \times_C B \hookrightarrow L$ for the inclusions. Then $jk$ is strongly continuous and $j$ preserves colimits, so the right adjoint $k^R \simeq (jk)^R j$ preserves colimits, i.e.\ $k$ is strongly continuous. Finally, colimits in $L$ are computed on the components $a$ and $b$, so a functor $E \to L$ preserves colimits if and only if its composites with $\pi_A$ and $\pi_B$ do. Together with \Cref{prop:SOD_Functor_Categories}(1), this shows that for a full subcategory $E \subseteq A \times_C B$, the functors $E \to A$ and $E \to B$ are strongly continuous if and only if the inclusion $E \to L$ is. So the final assertion follows from the definition of $\ker^{\dual}(\Phi)$.

		(2) Suppose that $G^R$ is fully faithful. For $c \in C$ and $(a,b,f) \in L$, we have
		\[
			\Map_L\big((a,b,f),(0,G^Rc,0)\big) \simeq \fib\big(\Map_B(b,G^Rc) \to \Map_C(F(a),GG^Rc)\big) \simeq \Map_C\big(\cofib(f),c\big),
		\]
		using $GG^R \simeq \id$. So $\Phi$ has the right adjoint $\Phi^R(c) = (0,G^Rc,0)$, which preserves colimits and satisfies $\Phi\Phi^R \simeq GG^R \simeq \id$. Hence $\Phi$ is a strongly continuous localization. By \Cref{lem:Strong_Continuity_Localization_Sequence} and \Cref{prop:Localization_Sequence_Dualizable}, its kernel $\ker(\Phi) = A \times_C B$ is dualizable and its inclusion into $L$ is strongly continuous. Hence $\ker^{\dual}(\Phi) = \ker(\Phi)$.

		The functor $R\colon A \to A \times_C B$, $a \mapsto (a,G^RF(a),\theta)$, where $\theta\colon F(a) \simeq GG^RF(a)$ is the inverse of the counit, is right adjoint to the projection: for $(a',b',\theta') \in A \times_C B$, the anima of maps $(a',b',\theta') \to R(a)$ is
		\[
			\Map_A(a',a) \times_{\Map_C(F(a'),F(a))} \Map_B(b',G^RF(a)) \simeq \Map_A(a',a),
		\]
		since $\Map_B(b',G^RF(a)) \simeq \Map_C(G(b'),F(a)) \simeq \Map_C(F(a'),F(a))$ via $\theta'$. The composite of $R$ with the projection is the identity, so $R$ is fully faithful and the projection is a localization.
	\end{proof}

	\begin{corollary}[{\cite[Corollary~2.87]{Efimov_Large_Categories}}]
		\label[corollary]{cor:Monomorphisms_Prdual}
		A strongly continuous functor $F\colon C \to D$ between dualizable categories is a monomorphism in $\Pr^{\dual}$ if and only if it is fully faithful.
	\end{corollary}
	\prooffor{cor:Monomorphisms_Prdual}
	\begin{proof}
		If $F$ is fully faithful, then $F \circ -$ is fully faithful on functor categories, so it induces monomorphisms on mapping spaces. Conversely, if $F$ is a monomorphism, then the diagonal $C \to C \times^{\dual}_D C$ is an equivalence. Composing with the fully faithful functor $C \times^{\dual}_D C \to C \times_D C$ of \Cref{prop:Dualizable_Pullbacks}(1), we find that the diagonal $C \to C \times_D C$ is fully faithful. As in the proof of \Cref{prop:Epimorphisms_Localizations}, it follows that $\hom_C(x,y) \to \hom_D(Fx,Fy)$ is an equivalence for all $x,y \in C$.
	\end{proof}

	\begin{remark}
		\label[remark]{rem:Short_Exact_Fiber_Cofiber}
		A short exact sequence $C \xrightarrow{i} D \xrightarrow{p} E$ of dualizable categories (\Cref{def:Short_Exact_Sequence_Dualizable}) is both a fiber sequence and a cofiber sequence in $\Pr^{\dual}$. It is a cofiber sequence since colimits in $\Pr^{\dual}$ are computed in $\PrL_{\st}$ (\Cref{prop:Prdual_Colimits}). It is a fiber sequence since $\ker^{\dual}(p) = \ker(p) = C$, as the inclusion $C \hookrightarrow D$ is strongly continuous; alternatively, apply \Cref{prop:Dualizable_Pullbacks}(2) to the functors $0 \to E$ and $p$.
	\end{remark}

	\section{General limits of dualizable categories}
	\label[section]{sec:General_Limits_Dualizable}

	% Source: Talk 7 (Christoph Winges, 24.11.2026); Efimov, Proposition 2.89 and Theorem 2.91; proofs completed with generative AI (Claude)
	We now show that $\Pr^{\dual}$ admits all small limits. The idea is to embed each dualizable category $C$ into the compactly generated category $\Ind(C^{\kappa})$ via its Calkin sequence. Limits of compactly generated categories of this form are easy to compute, and the limit in $\Pr^{\dual}$ is then cut out by the vanishing of the images in the Calkin quotients. Throughout this section, $\kappa$ denotes an uncountable regular cardinal. We write $\PrL_{\st,\kappa}$ for the category of $\kappa$-compactly generated stable categories and colimit-preserving functors which preserve $\kappa$-compact objects. Every dualizable category is $\kappa$-compactly generated (see \Cref{sec:Continuous_Calkin}), and strongly continuous functors preserve $\kappa$-compact objects, so $\Pr^{\dual}$ is a subcategory of $\PrL_{\st,\kappa}$.

	\begin{proposition}[{\cite[Proposition~2.89]{Efimov_Large_Categories}}]
		\label[proposition]{prop:Dualizable_Ind_Adjunction}
		The inclusion $\Pr^{\dual} \to \PrL_{\st,\kappa}$ admits a right adjoint $R_{\kappa}$, given by $R_{\kappa}(D) = \Ind(D^{\kappa})$ and $R_{\kappa}(F) = \Ind(F^{\kappa})$.
	\end{proposition}
	\prooffor{prop:Dualizable_Ind_Adjunction}
	\begin{proof}
		The category $\Ind(D^{\kappa})$ is compactly generated, hence dualizable, and $\Ind(F^{\kappa})$ preserves compact objects, hence is strongly continuous. We take as unit the functors $\hat{y}_C\colon C \to \Ind(C^{\kappa})$, which are strongly continuous and natural in $C \in \Pr^{\dual}$ by \Cref{constr:Calkin_Functoriality}. As counit we take the functors $\colim_D\colon \Ind(D^{\kappa}) \to D$, which are natural in $D$. They preserve $\kappa$-compact objects: the $\kappa$-compact objects of $\Ind(D^{\kappa})$ are retracts of $\kappa$-small colimits of objects $y(d)$ with $d \in D^{\kappa}$, which are sent to retracts of $\kappa$-small colimits in $D^{\kappa}$.

		It remains to verify the triangle identities. The composite $C \xrightarrow{\hat{y}} \Ind(C^{\kappa}) \xrightarrow{\colim} C$ is the identity since $\hat{y}$ is fully faithful. For the other one, consider the composite
		\[
			\Ind(D^{\kappa}) \xrightarrow{\hat{y}} \Ind\big(\Ind(D^{\kappa})^{\kappa}\big) \xrightarrow{\Ind(\colim_D^{\kappa})} \Ind(D^{\kappa}).
		\]
		Both functors preserve colimits, so it suffices to consider its restriction to compact objects $y(d)$, $d \in D^{\kappa}$. By \Cref{lem:Compact_Iff_Yhat_Is_Y}, we have $\hat{y}(y(d)) \simeq y(y(d))$, which is sent to $y(\colim_D y(d)) \simeq y(d)$. So the composite is the identity.
	\end{proof}

	Explicitly, a functor $F\colon C \to D$ in $\PrL_{\st,\kappa}$ from a dualizable category $C$ corresponds to the strongly continuous functor $\Ind(F^{\kappa}) \circ \hat{y}_C\colon C \to \Ind(D^{\kappa})$ \cite[Remark~2.90]{Efimov_Large_Categories}. For a diagram $i \mapsto C_i$ in $\Pr^{\dual}$ we write $\lim^{\dual}_i C_i$ for its limit in $\Pr^{\dual}$.

	\begin{theorem}[{\cite[Theorem~2.91]{Efimov_Large_Categories}}]
		\label[theorem]{thm:General_Limits_Dualizable}
		The category $\Pr^{\dual}$ admits all small limits. More precisely, let $i \mapsto C_i$ be a small diagram in $\Pr^{\dual}$. Then
		\[
			\lim^{\dual}_i C_i \simeq \ker^{\dual}\Big(\Ind\big(\lim_i C_i^{\kappa}\big) \to \Ind\big(\lim_i \Calk^{\cont}_{\kappa}(C_i)\big)\Big),
		\]
		where the limits on the right are computed in $\Cat_{\perf}$, and the functor is induced by the quotient functors $p_i\colon C_i^{\kappa} \to \Calk^{\cont}_{\kappa}(C_i)$.
	\end{theorem}
	\prooffor{thm:General_Limits_Dualizable}
	\begin{proof}
		\emph{Step 1: the ambient category.} Passing to $\kappa$-compact objects identifies $\PrL_{\st,\kappa}$ with the category of small idempotent complete stable categories with $\kappa$-small colimits and exact functors preserving them \cite[Propositions~5.5.7.8 and~5.5.7.10]{lurie2009HTT}. Limits in the latter category are computed in $\Cat$ \cite[Remark~5.4.7.13]{lurie2009HTT}, so $\lim_i C_i$ in $\PrL_{\st,\kappa}$ is the $\kappa$-compactly generated category with $\kappa$-compact objects $\lim_i C_i^{\kappa}$. Since right adjoints preserve limits, \Cref{prop:Dualizable_Ind_Adjunction} shows that $M := \Ind(\lim_i C_i^{\kappa})$ is a limit of the diagram $i \mapsto \Ind(C_i^{\kappa})$ in $\Pr^{\dual}$.

		\emph{Step 2: the limit.} Let $T$ be dualizable. A strongly continuous functor $G\colon T \to \Ind(C_i^{\kappa})$ factors through $\hat{y}\colon C_i \hookrightarrow \Ind(C_i^{\kappa})$ if and only if its composite with $\Ind(p_i)$ vanishes, since $C_i$ is the kernel of $\Ind(p_i)$; in this case, the factorization is strongly continuous, as its right adjoint is $G^R \circ \hat{y}$. Full faithfulness of $\hat{y}$ identifies $\Map_{\Pr^{\dual}}(T,C_i)$ with the union of components cut out by this vanishing condition. Passing to the limit and using Step 1, we see that $\lim_i \Map_{\Pr^{\dual}}(T,C_i)$ consists of the components of $\Map_{\Pr^{\dual}}(T,M)$ whose projection to each $\Ind(C_i^\kappa)$ is killed by $\Ind(p_i)$.

		Let $K \subseteq M$ be the full subcategory of objects whose projection to $\Ind(C_i^\kappa)$ is killed by $\Ind(p_i)$ for every $i$. It is a full stable subcategory closed under colimits, so by \Cref{lem:Largest_Dualizable_Subcategory} it has a largest full subcategory $L$ which is dualizable and whose inclusion into $M$ is strongly continuous. A strongly continuous functor $H\colon T \to M$ satisfies the above condition if and only if $\Im(H) \subseteq K$. By \Cref{cor:Image_Dualizable} and the maximality of $L$, this is the case if and only if $\Im(H) \subseteq L$, i.e.\ if and only if $H$ factors through $L$; the factorization is then strongly continuous, as in the proof of \Cref{prop:Dualizable_Kernel}. Hence $\Map_{\Pr^{\dual}}(T,L) \simeq \lim_i \Map_{\Pr^{\dual}}(T,C_i)$, naturally in $T$, and $L$ is a limit of the diagram $i \mapsto C_i$ in $\Pr^{\dual}$.

		\emph{Step 3: the kernel formula.} Let $\Phi\colon M \to \Ind(\lim_i \Calk^{\cont}_{\kappa}(C_i))$ be the functor induced by the $p_i$. We must show that a strongly continuous functor $H\colon T\to M$ from a dualizable category satisfies $\Phi H\simeq 0$ precisely when it satisfies the componentwise vanishing condition of Step 2.

		To see why we can test vanishing componentwise here, consider any small diagram $i\mapsto D_i$ in $\Cat_{\perf}$. Its projections induce a fully faithful functor
		\[
			\Ind(\lim_i D_i)\longrightarrow \lim\nolimits^{\dual}_i\Ind(D_i).
		\]
		Indeed, Step 2 embeds the target fully faithfully in $\Ind(\lim_i \Ind(D_i)^\kappa)$. By the explicit adjunction of \Cref{prop:Dualizable_Ind_Adjunction} and the formula $\hat y=\Ind(y)$ for compactly generated categories (\Cref{lem:Compact_Iff_Yhat_Is_Y}), the composite is $\Ind$ applied to the limit of the Yoneda embeddings $D_i\to\Ind(D_i)^\kappa$. Limits of fully faithful functors are fully faithful, and $\Ind$ preserves full faithfulness.

		Apply this observation to the small Calkin categories $\Calk^{\cont}_{\kappa}(C_i)$. After composing $\Phi H$ with the displayed functor, its $i$th component is the projection of $H$ to $\Ind(C_i^\kappa)$ followed by $\Ind(p_i)$. These components vanish for all $i$ precisely when the functor to the limit vanishes. Full faithfulness then gives the desired criterion for $\Phi H$.

		Now let $E \subseteq M$ be a full subcategory which is dualizable and whose inclusion $j$ into $M$ is strongly continuous. Applying the claim to $H = j$, we see that $E \subseteq K$ if and only if $E \subseteq \ker(\Phi)$. Hence $L = \ker^{\dual}(\Phi)$.
	\end{proof}

	The proof also shows that it suffices to test the vanishing of the composites to the Calkin quotients one $i$ at a time: in Step 2 we only use the limit of the Calkin terms in $\Cat_{\perf}$, not their limit in $\Pr^{\dual}$, which at that point is not yet known to exist.

	% Source: Mathew, Lecture 7 (MPIM Bonn, 12.09.2024)
	Compact objects of a limit of dualizable categories can be computed directly, since $(-)^{\omega}$ is a right adjoint.

	\begin{proposition}
		\label[proposition]{prop:Ind_Compact_Adjunction}
		The functor $\Ind\colon \Cat_{\perf} \to \Pr^{\dual}$ is fully faithful and left adjoint to $(-)^{\omega}\colon \Pr^{\dual} \to \Cat_{\perf}$. Consequently:
		\begin{enumerate}
			\item $\Ind$ preserves colimits and sends short exact sequences in $\Cat_{\perf}$ to short exact sequences of dualizable categories;
			\item for every small diagram $i \mapsto C_i$ in $\Pr^{\dual}$, we have $(\lim^{\dual}_i C_i)^{\omega} \simeq \lim_i C_i^{\omega}$, the limit computed in $\Cat_{\perf}$.
		\end{enumerate}
	\end{proposition}
	\prooffor{prop:Ind_Compact_Adjunction}
	\begin{proof}
		Let $A \in \Cat_{\perf}$ and let $C$ be dualizable. By \Cref{thm:Ind_Kappa_Universal_Property}(3), colimit-preserving functors $\Ind(A) \to C$ correspond to exact functors $A \to C$. Such a functor is strongly continuous if and only if it preserves compact objects, i.e.\ if and only if the corresponding exact functor lands in $C^{\omega}$, since the compact objects of $\Ind(A)$ are the objects of $A$. Hence $\Fun^{\mathrm{LL}}(\Ind(A),C) \simeq \Fun^{\ex}(A,C^{\omega})$, naturally in $A$ and $C$, and passing to cores gives the adjunction. Its unit $A \to \Ind(A)^{\omega}$ is an equivalence, so $\Ind$ is fully faithful. Now (1) follows since left adjoints preserve colimits, together with \Cref{thm:Thomason_Neeman}(1), and (2) follows since right adjoints preserve limits.
	\end{proof}

	% Source: Mathew, Lecture 7 (MPIM Bonn, 12.09.2024)
	% Source: The componentwise description is recorded in Chowdhury, Remark 6.18; cf. Ramzi, Corollary 4.37. Proof added in preparing these notes.
	The proof of \Cref{thm:General_Limits_Dualizable} also gives generators for the limit.

	\begin{corollary}
		\label[corollary]{cor:Dualizable_Limit_Generators}
		In the situation of \Cref{thm:General_Limits_Dualizable}, let $M := \Ind(\lim_i C_i^{\kappa})$ and let $K \subseteq M$ be the full subcategory of objects whose image in $\Ind(\Calk^{\cont}_{\kappa}(C_i))$ vanishes for every $i$. Then $\lim^{\dual}_i C_i \subseteq M$ is generated under colimits by the objects $\colim_{\Q}F$, where $F\colon \Q \to K$ is a functor such that $F(a) \to F(b)$ is compact in $M$ for all $a < b$.

		For $\kappa=\omega_1$, put $A:=\lim_i C_i^{\omega_1}$ and write $y_A\colon A\to M=\Ind(A)$ for the Yoneda embedding. Then $\lim_i^{\dual} C_i$ is generated under colimits by the objects
		\[
			\colim_{r\in\Q} y_A(F(r)),
		\]
		where $F\colon\Q\to A$ ranges over the diagrams for which $F(r)_i\to F(s)_i$ is compact in $C_i$ for every $i$ and every $r<s$ (cf.\ \cite[Definitions~4.20 and~4.22 and Corollary~4.37]{Ramzi_Dualizable_Presentable}).
	\end{corollary}
	\prooffor{cor:Dualizable_Limit_Generators}
	\begin{proof}
		By Step~2 of the proof of \Cref{thm:General_Limits_Dualizable}, $\lim^{\dual}_i C_i$ is the largest full subcategory of $K$ which is dualizable and whose inclusion into $M$ is strongly continuous. Since $M$ is compactly generated, hence dualizable, the first description follows from \Cref{lem:Largest_Dualizable_Subcategory}.

		For the second description, let $F\colon\Q\to A$ have componentwise compact transitions and put $X_t:=\colim_{r<t}y_A(F(r))$ in $M$. By \Cref{lem:Compact_Morphisms_Yhat}, the Yoneda image of each transition in $C_i$ factors through $\hat y_{C_i}$ of its target. Its image in the Calkin quotient is therefore nullhomotopic. Computing $X_t$ along a cofinal increasing sequence in $\Q_{<t}$ shows that its image in each Calkin quotient is zero, so $X_t\in K$. For $t<u$, the map $X_t\to X_u$ factors through the compact object $y_A(F(t))$ of $M$. The first description now shows that $\colim_t X_t\simeq\colim_r y_A(F(r))$ belongs to $\lim_i^{\dual} C_i$.

		Conversely, write $L:=\lim_i^{\dual} C_i$ and take $z\in L^{\omega_1}$. By \Cref{prop:Compact_Interpolation}(3), we can write $z=\colim_{r\in\Q}z_r$ with $z_r\in L^{\omega_1}$ and compact transitions. The projections $L\to C_i$ are strongly continuous, so they induce $P\colon L^{\omega_1}\to A$ and preserve compact morphisms by \Cref{prop:Strongly_Continuous_Compact_Morphisms}. Under the adjunction of \Cref{prop:Dualizable_Ind_Adjunction}, the projection cone corresponds to the inclusion $j\colon L\hookrightarrow M$, so $j\simeq\Ind(P)\hat y_L$. Restricting to the cofinal sequence $\N\subseteq\Q$ and using compactness of the transitions gives $\hat y_L(z)\simeq\colim_r y_{L^{\omega_1}}(z_r)$ as in the proof of \Cref{thm:Compactly_Assembled_Omega_1}(2). Hence $j(z)\simeq\colim_r y_A(P(z_r))$ has the required form. Such $z$ generate $L$ under colimits by \Cref{thm:Compactly_Assembled_Omega_1}(1).
	\end{proof}

	% Source: Mathew, Lecture 7 (MPIM Bonn, 12.09.2024), recorded in Chowdhury, Remark 6.7; Ramzi, Corollary 4.40. Proof added in preparing these notes.
	\begin{corollary}[{\cite[Corollary~4.40]{Ramzi_Dualizable_Presentable}}]
		\label[corollary]{cor:Finite_Dualizable_Limits_Fully_Faithful}
		Let $S$ be a finite simplicial set and $D\colon S\to\Pr^{\dual}$ a diagram. The canonical functor from its limit in $\Pr^{\dual}$ to its limit in $\PrL_{\st}$ is fully faithful.
	\end{corollary}
	\prooffor{cor:Finite_Dualizable_Limits_Fully_Faithful}
	\begin{proof}
		For the empty diagram both limits are zero, and for a simplex both are the value at its initial vertex. We proceed by induction on the dimension of $S$ and, within each dimension, on the number of top-dimensional simplices. Write $S=S'\sqcup_{\partial\Delta^n}\Delta^n$ by removing one nondegenerate simplex of maximal dimension. In either category a cone over $S$ is a pair of cones over $S'$ and $\Delta^n$ agreeing over $\partial\Delta^n$. Thus each limit is the pullback of the limits of the three restrictions. By \Cref{prop:Dualizable_Pullbacks}(1), the pullback in $\Pr^{\dual}$ embeds fully faithfully in the ordinary pullback of those same three dualizable limits. The comparison functors for $S'$, $\Delta^n$ and $\partial\Delta^n$ are fully faithful by induction and the simplex case. Their ordinary pullback is fully faithful, since limits of fully faithful functors are fully faithful. The composite is the desired functor.
	\end{proof}

	% Source: Mathew, Lecture 7 (MPIM Bonn, 12.09.2024); Ramzi, Theorem 3.7
	\begin{remark}
		\label[remark]{rem:Dualizable_Limits_Further}
		\leavevmode
		\begin{enumerate}
			\item If the functor $\lim_i C_i^{\kappa} \to \lim_i \Calk^{\cont}_{\kappa}(C_i)$ is a homological epimorphism, then $\lim^{\dual}_i C_i$ is the ordinary kernel of the induced functor on ind-completions. Indeed, the kernel of a strongly continuous localization of dualizable categories is dualizable with strongly continuous inclusion (\Cref{lem:Strong_Continuity_Localization_Sequence,prop:Localization_Sequence_Dualizable}), so it agrees with the dualizable kernel. Products are an instance of this (\Cref{cor:Products_Dualizable_Calkin}).
			\item The full faithfulness of \Cref{cor:Finite_Dualizable_Limits_Fully_Faithful} can fail for infinite products. For compactly generated categories, the functor $\prod^{\dual}_i\Ind(A_i) \simeq \Ind(\prod_iA_i) \to \prod_i\Ind(A_i)$ is in general not fully faithful. Indeed, mapping spectra out of compact objects in the source compute $\colim_j \prod_i$, while those in the target compute $\prod_i \colim_j$, and these differ by \Cref{rem:AB6_Versus_Naive_Commutation}.
			\item By a theorem of Ramzi, the adjunction of \Cref{prop:Dualizable_Ind_Adjunction} is comonadic: $\Pr^{\dual}$ is equivalent to the category of coalgebras for the comonad $D \mapsto \Ind(D^{\kappa})$ on $\PrL_{\st,\kappa}$ \cite[Theorems~3.4 and~3.7]{Ramzi_Dualizable_Presentable}.
		\end{enumerate}
	\end{remark}

	\section{Infinite products and exactness}
	\label[section]{sec:Infinite_Products_Dualizable}

	% Source: Talk 7 (Christoph Winges, 24.11.2026); Efimov, Propositions 2.99--2.106 and Lemma 2.104; proofs completed with generative AI (Claude)
	We now specialize \Cref{thm:General_Limits_Dualizable} to products. In this case the kernel in the formula turns out to be an ordinary kernel, since products of homological epimorphisms are again homological epimorphisms. As a consequence, products in $\Pr^{\dual}$ preserve short exact sequences, and $\Pr^{\dual}$ satisfies (AB6). We write $\prod^{\dual}_i C_i$ for products in $\Pr^{\dual}$.

	We first compare ind-objects in a product with products of ind-objects. Let $(A_i)_{i \in I}$ be a family of small stable categories and put $A := \prod_i A_i$. The functors $\Ind(\mathrm{pr}_i)\colon \Ind(A) \to \Ind(A_i)$ induce a colimit-preserving functor $\Ind(A) \to \prod_i \Ind(A_i)$. We denote by
	\[
		P_A\colon \prod_i \Ind(A_i) \to \Ind(A)
	\]
	its right adjoint.

	\begin{lemma}[{\cite[Proposition~2.99 and Corollary~2.100]{Efimov_Large_Categories}}]
		\label[lemma]{lem:Products_Ind_Comparison}
		\leavevmode
		\begin{enumerate}
			\item Under the identification $\Ind(A) \simeq \Fun^{\ex}(A\catop,\Sp)$ of \Cref{cor:Ind_As_Exact_Functors}, we have $P_A((X_i)_i)\big((a_i)_i\big) \simeq \prod_i X_i(a_i)$.
			\item If $X_i \simeq \quotcolim_{j \in J_i} a_{i,j}$ for filtered categories $J_i$, then $P_A((X_i)_i) \simeq \quotcolim_{(j_i)_i \in \prod_i J_i} (a_{i,j_i})_i$. In particular, $y((a_i)_i) \simeq P_A((y(a_i))_i)$.
			\item For exact functors $F_i\colon A_i \to B_i$ with product $F := \prod_i F_i$, we have $\Ind(F) \circ P_A \simeq P_B \circ \prod_i \Ind(F_i)$ and $\Ind(F)^R \circ P_B \simeq P_A \circ \prod_i \Ind(F_i)^R$.
		\end{enumerate}
	\end{lemma}
	\prooffor{lem:Products_Ind_Comparison}
	\begin{proof}
		(1) The right adjoint of $\Ind(\mathrm{pr}_i)$ is restriction along $\mathrm{pr}_i$, and the right adjoint of a functor into a product is the product of the right adjoints of its components.

		(2) Evaluating at $(b_i)_i \in A$ and using that filtered colimits in $\Fun^{\ex}(A\catop,\Sp)$ are computed pointwise, the claim becomes $\colim_{(j_i)_i} \prod_i \hom_{A_i}(b_i,a_{i,j_i}) \simeq \prod_i \colim_{j \in J_i}\hom_{A_i}(b_i,a_{i,j})$, which is (AB6) in $\Sp$ (\Cref{prop:Dualizable_AB6}). This is also a special case of \Cref{lem:Products_In_Ind}.

		(3) The first equivalence follows from (2), since $\Ind(F)$ and $\Ind(F_i)$ are computed on filtered presentations levelwise. For the second, pass to left adjoints: the left adjoint of $\Ind(F)^R \circ P_B$ is the functor $\Ind(A) \to \prod_i \Ind(B_i)$ with components $\Ind(\mathrm{pr}_i)\Ind(F) \simeq \Ind(F_i)\Ind(\mathrm{pr}_i)$, which is also the left adjoint of $P_A \circ \prod_i \Ind(F_i)^R$.
	\end{proof}

	\begin{proposition}[{\cite[Proposition~2.102]{Efimov_Large_Categories}}]
		\label[proposition]{prop:Products_Small_Projections}
		Let $(F_i\colon A_i \to B_i)_{i \in I}$ be a family of exact functors in $\Cat_{\perf}$, and put $F := \prod_i F_i$.
		\begin{enumerate}
			\item If each $K_i \to A_i \xrightarrow{F_i} B_i$ is a short exact sequence, then so is $\prod_i K_i \to \prod_i A_i \xrightarrow{F} \prod_i B_i$.
			\item If each $F_i$ is a homological epimorphism, then so is $F$. Moreover, $\ker(\Ind(F))$ is generated under colimits by the objects $P_A((X_i)_i)$ with $X_i \in \ker(\Ind(F_i))$.
		\end{enumerate}
	\end{proposition}
	\prooffor{prop:Products_Small_Projections}
	\begin{proof}
		(1) This is \cite[Lemmas~5.2 and~5.3]{Kasprowski_Winges_Products}; we give the argument. Put $A := \prod_i A_i$ and $K := \prod_i K_i$, the kernel of $F$. For $x = (x_i)_i$ and $y = (y_i)_i$ in $A$, the category $K_{/y}$ is the product of the categories $(K_i)_{/y_i}$, which are filtered. By \Cref{rem:Mapping_Spectra_Quotient} and (AB6) in $\Sp$, we get
		\[
			\hom_{A/K}(x,y) \simeq \colim_{(k_i \to y_i)_i} \prod_i \hom_{A_i}\big(x_i,\cofib(k_i \to y_i)\big) \simeq \prod_i \hom_{A_i/K_i}(x_i,y_i) \simeq \hom_{\prod_i B_i}(Fx,Fy),
		\]
		so $A/K \to \prod_i B_i$ is fully faithful. Every object $(b_i)_i$ of $\prod_i B_i$ is a retract of an object in the image of $F$, choosing retractions componentwise. Hence $(A/K)^{\idem} \simeq \prod_i B_i$.

		(2) By \Cref{lem:Products_Ind_Comparison}(3), we have
		\[
			\Ind(F)\Ind(F)^RP_B \simeq \Ind(F)P_A\prod_i\Ind(F_i)^R \simeq P_B\prod_i\Ind(F_i)\Ind(F_i)^R,
		\]
		and under these identifications the counit of $\Ind(F) \dashv \Ind(F)^R$ at $P_B((Y_i)_i)$ is $P_B$ applied to the componentwise counits. These are equivalences since the $F_i$ are homological epimorphisms. So the counit is an equivalence on the essential image of $P_B$, which contains all compact objects of $\Ind(B)$ by \Cref{lem:Products_Ind_Comparison}(2). Since $\Ind(F)$ preserves compact objects, $\Ind(F)\Ind(F)^R$ preserves colimits, so the counit is an equivalence, i.e.\ $F$ is a homological epimorphism.

		For the last assertion, note that $P_A((X_i)_i) \in \ker(\Ind(F))$ if each $X_i \in \ker(\Ind(F_i))$, by \Cref{lem:Products_Ind_Comparison}(3). Conversely, the colimit-preserving functor $\fib(\id \to \Ind(F)^R\Ind(F))$ is the identity on $\ker(\Ind(F))$, so $\ker(\Ind(F))$ is generated under colimits by its values on the compact objects $y(a)$, $a = (a_i)_i \in A$. By \Cref{lem:Products_Ind_Comparison}, its value at $y(a) \simeq P_A((y(a_i))_i)$ is $P_A\big((\fib(y(a_i) \to \Ind(F_i)^R\Ind(F_i)y(a_i)))_i\big)$, whose components lie in $\ker(\Ind(F_i))$.
	\end{proof}

	\begin{corollary}[{\cite[Proposition~2.103]{Efimov_Large_Categories}}]
		\label[corollary]{cor:Products_Dualizable_Calkin}
		Let $(C_i)_{i \in I}$ be a family of dualizable categories. Then there is a short exact sequence
		\[
			\prod\nolimits^{\dual}_i C_i \to \Ind\Big(\prod_i C_i^{\kappa}\Big) \to \Ind\Big(\prod_i \Calk^{\cont}_{\kappa}(C_i)\Big).
		\]
		If each $C_i$ is compactly generated, then $\prod^{\dual}_i C_i \simeq \Ind(\prod_i C_i^{\omega})$. In other words, $\Ind\colon \Cat_{\perf} \to \Pr^{\dual}$ preserves products.
	\end{corollary}
	\prooffor{cor:Products_Dualizable_Calkin}
	\begin{proof}
		The quotient functors $p_i\colon C_i^{\kappa} \to \Calk^{\cont}_{\kappa}(C_i)$ are homological epimorphisms, since $\Ind(p_i)$ is the localization in the Calkin sequence of $C_i$. By \Cref{prop:Products_Small_Projections}(2), the functor $\Ind(\prod_i p_i)$ is a strongly continuous localization. By \Cref{lem:Strong_Continuity_Localization_Sequence} and \Cref{prop:Localization_Sequence_Dualizable}, its kernel is dualizable with strongly continuous inclusion, so it agrees with the dualizable kernel, which is $\prod^{\dual}_i C_i$ by \Cref{thm:General_Limits_Dualizable}.

		If each $C_i$ is compactly generated, then $C_i^{\omega} \to C_i^{\kappa} \to \Calk^{\cont}_{\kappa}(C_i)$ is short exact by \Cref{rem:Calkin_Small_Vs_Continuous}. By \Cref{prop:Products_Small_Projections}(1), their product is short exact, and applying \Cref{thm:Thomason_Neeman}(1) identifies the kernel of $\Ind(\prod_i p_i)$ with $\Ind(\prod_i C_i^{\omega})$.
	\end{proof}

	To show that products preserve short exact sequences, we use duality. For a dualizable category $C$ we use the model $C^{\vee} = \Fun^{\colim}(C,\Sp)$ of \Cref{sec:Dualizable_Categories}, and for a colimit-preserving functor $F\colon C \to D$ we write $F^{\vee} := - \circ F\colon D^{\vee} \to C^{\vee}$. For a small stable category $A$ we have $\Ind(A)^{\vee} \simeq \Fun^{\ex}(A,\Sp) \simeq \Ind(A\catop)$ by \Cref{thm:Ind_Kappa_Universal_Property}(3) and \Cref{cor:Ind_As_Exact_Functors}. For an exact functor $f\colon A \to A'$, the functor $\Ind(f)^{\vee}$ is restriction along $f$, i.e.\ the right adjoint of $\Ind(f\catop)$.

	\begin{lemma}[{\cite[Section~2.5]{Efimov_Large_Categories}}]
		\label[lemma]{lem:Duality_Localizations}
		Let $F\colon C \to D$ be a strongly continuous functor between dualizable categories. Then $(F^R)^{\vee}$ is left adjoint to $F^{\vee}$, and $F$ is a localization if and only if $F^{\vee}$ is fully faithful. If $F$ is fully faithful, then so is $(F^R)^{\vee}$.
	\end{lemma}
	\prooffor{lem:Duality_Localizations}
	\begin{proof}
		For an adjunction $L \dashv R$, precomposition with $R$ is left adjoint to precomposition with $L$, with unit and counit induced by those of $L \dashv R$. Applied to $F \dashv F^R$, this shows $(F^R)^{\vee} \dashv F^{\vee}$, with counit $\varphi \circ FF^R \to \varphi$ induced by the counit of $F \dashv F^R$. If $F$ is a localization, this counit is an equivalence, so $F^{\vee}$ is fully faithful. Conversely, if $F^{\vee}$ is fully faithful, then $\varphi(FF^Rd) \to \varphi(d)$ is an equivalence for all $\varphi \in D^{\vee}$ and $d \in D$. Since $D$ is dualizable, the functor $D \to \Fun^{\colim}(D^{\vee},\Sp)$, $d \mapsto (\varphi \mapsto \varphi(d))$, is an equivalence, so $FF^Rd \to d$ is an equivalence. Similarly, if $F$ is fully faithful, then the unit $\varphi \to \varphi F^RF$ of $(F^R)^{\vee} \dashv F^{\vee}$ is an equivalence.
	\end{proof}

	\begin{lemma}[{\cite[Lemma~2.104]{Efimov_Large_Categories}}]
		\label[lemma]{lem:Dual_Product_Embedding}
		Let $(C_i)_{i \in I}$ be a family of dualizable categories, and let $j\colon \prod^{\dual}_i C_i \hookrightarrow \Ind(\prod_i C_i^{\kappa})$ be the inclusion of \Cref{cor:Products_Dualizable_Calkin}. Then the functor
		\[
			\Phi := (j^R)^{\vee}\colon \Big(\prod\nolimits^{\dual}_i C_i\Big)^{\vee} \to \Ind\Big(\prod_i C_i^{\kappa}\Big)^{\vee} \simeq \Ind\Big(\prod_i (C_i^{\kappa})\catop\Big)
		\]
		is fully faithful. Its essential image is generated under colimits by the exact functors $E_x\colon \prod_i C_i^{\kappa}\to\Sp$ given by
		\[
			E_x((a_i)_i) = \prod_i x_i(a_i), \qquad x=(x_i)_i,\quad x_i\in C_i^{\vee}.
		\]
	\end{lemma}
	\prooffor{lem:Dual_Product_Embedding}
	\begin{proof}
		The functor $\Phi$ is fully faithful by \Cref{lem:Duality_Localizations}. Let $p := \Ind(\prod_i p_i)$ be the localization of \Cref{cor:Products_Dualizable_Calkin}. For $X \in \Ind(\prod_i C_i^{\kappa})$ we have a cofiber sequence $jj^RX \to X \to p^Rp X$. Since $\Phi$ is fully faithful with right adjoint $j^{\vee}$, its essential image consists of those $\varphi$ for which the counit $\varphi jj^R \to \varphi$ is an equivalence, i.e.\ for which $\varphi p^R \simeq 0$. By the discussion preceding \Cref{lem:Duality_Localizations}, $(p^R)^{\vee}$, as the left adjoint of $p^{\vee}$, is identified with $\Ind(\prod_i p_i\catop)$. So the essential image of $\Phi$ is $\ker(\Ind(\prod_i p_i\catop))$. The same argument for a single $C_i$ identifies $C_i^{\vee}$, via $(\hat{y}^R)^{\vee} = - \circ \colim$, with $\ker(\Ind(p_i\catop))$; this sends $x$ to $x \circ \colim$, whose restriction to $C_i^{\kappa}$ is $x|_{C_i^{\kappa}}$.

		A functor $f$ in $\Cat_{\perf}$ is a homological epimorphism if and only if it is an epimorphism in $\Cat_{\perf}$: indeed, $\Ind\colon \Cat_{\perf} \to \PrL_{\omega,\st}$ is an equivalence, pushouts in $\PrL_{\omega,\st}$ are computed in $\PrL_{\st}$ (as in \Cref{prop:Prdual_Colimits}), and we can apply \Cref{prop:Epimorphisms_Localizations}. Since $(-)\catop$ is an automorphism of $\Cat_{\perf}$, the functors $p_i\catop$ are homological epimorphisms as well \cite[Remark~1.10]{Efimov_Large_Categories}. By \Cref{prop:Products_Small_Projections}(2), the kernel of $\Ind(\prod_i p_i\catop)$ is generated by applying the product comparison functor of \Cref{lem:Products_Ind_Comparison} to families of these restrictions $x_i|_{C_i^{\kappa}}$. Its pointwise formula gives exactly the functors $E_x$.
	\end{proof}

	\begin{proposition}[{\cite[Propositions~2.105 and~2.106]{Efimov_Large_Categories}}]
		\label[proposition]{prop:Exact_Products_AB6_Dualizable}
		Products in $\Pr^{\dual}$ of short exact sequences of dualizable categories are short exact, and $\Pr^{\dual}$ satisfies (AB6): for filtered categories $J_i$, $i \in I$, and functors $J_i \to \Pr^{\dual}$, $j \mapsto C_{i,j}$, the comparison functor
		\[
			\colim_{(j_i)_i \in \prod_i J_i} \prod\nolimits^{\dual}_i C_{i,j_i} \to \prod\nolimits^{\dual}_i \colim_{j \in J_i} C_{i,j}
		\]
		is an equivalence, where the colimits are computed in $\Pr^{\dual}$.
	\end{proposition}
	\prooffor{prop:Exact_Products_AB6_Dualizable}
	\begin{proof}
		Let $K_i \to D_i \xrightarrow{F_i} E_i$ be short exact sequences of dualizable categories. Since limits commute with limits, $\prod^{\dual}_i K_i$ is the fiber of $F := \prod^{\dual}_i F_i$ in $\Pr^{\dual}$. So, by \Cref{rem:Short_Exact_Fiber_Cofiber}, it suffices to show that $F$ is a localization; then $\prod^{\dual}_i K_i \simeq \ker^{\dual}(F) = \ker(F)$. By \Cref{lem:Duality_Localizations}, it suffices to show that $F^{\vee}$ is fully faithful.

		Let $j_D$ and $j_E$ be the inclusions of \Cref{cor:Products_Dualizable_Calkin}, with associated fully faithful functors $\Phi_D$ and $\Phi_E$ from \Cref{lem:Dual_Product_Embedding}, and put $G := \Ind(\prod_i F_i^{\kappa})$. By construction of the induced map of limits, $F$ is the restriction of $G$, i.e.\ $j_EF \simeq Gj_D$. The mate of this equivalence is a natural transformation $Fj_D^R \to j_E^RG$, which induces a natural transformation
		\[
			\alpha\colon \Phi_DF^{\vee} = - \circ Fj_D^R \longrightarrow - \circ j_E^RG = G^{\vee}\Phi_E
		\]
		of colimit-preserving functors. We show that $\alpha$ is an equivalence. By \Cref{lem:Dual_Product_Embedding}, it suffices to check this on objects $\psi$ with $\Phi_E\psi \simeq E_x$ for some $x_i \in E_i^{\vee}$. Then $G^{\vee}\Phi_E\psi \simeq E_x \circ G$ is the functor $(a_i)_i \mapsto \prod_i x_i(F_ia_i)$ on $\prod_i D_i^{\kappa}$, i.e.\ $E_x \circ G \simeq E_{F^{\vee}x}$ with $F^{\vee}x := (x_i \circ F_i)_i$, which lies in the essential image of $\Phi_D$. Unwinding the definition of the mate, $\alpha_{\psi}$ is identified with the map $E_{F^{\vee}x} \circ j_Dj_D^R \to E_{F^{\vee}x}$ induced by the counit, which is an equivalence on the essential image of $\Phi_D$ (see the proof of \Cref{lem:Dual_Product_Embedding}).

		So $\Phi_DF^{\vee} \simeq G^{\vee}\Phi_E$. The functor $G^{\vee}$ is the right adjoint of $\Ind(\prod_i (F_i^{\kappa})\catop)$. As in the proof of \Cref{prop:Calkin_Short_Exact}, the sequences $K_i^{\kappa} \to D_i^{\kappa} \to E_i^{\kappa}$ are short exact \cite[Proposition~A.26]{Ramzi_Dualizable_Presentable}, hence so are their opposites, and their product is short exact by \Cref{prop:Products_Small_Projections}(1). By \Cref{thm:Thomason_Neeman}(1), $\Ind(\prod_i (F_i^{\kappa})\catop)$ is a localization, so $G^{\vee}$ is fully faithful. Since $\Phi_D$ and $\Phi_E$ are fully faithful, so is $F^{\vee}$.

		For (AB6), we first treat the case where all $C_{i,j}$ are compactly generated. The category $\Cat_{\perf}$ satisfies (AB6): it is compactly generated \cite[Proposition~1.4(3)]{Efimov_Large_Categories}, and every compactly generated category $\mathcal{E}$ satisfies (AB6), since mapping out of a compact object $c$ commutes with filtered colimits and products, so the comparison map induces on $\Map(c,-)$ the comparison map for (AB6) in $\An$ (\Cref{ex:AB6_Sets_And_Animae}), and the compact objects jointly detect equivalences. The functor $\Ind\colon \Cat_{\perf} \to \Pr^{\dual}$ preserves products by \Cref{cor:Products_Dualizable_Calkin} and filtered colimits by \Cref{prop:Prdual_Colimits} and the proof of \Cref{prop:Compact_Objects_Filtered_Colimits}. This proves the compactly generated case.

		In general, apply both sides of the comparison functor to the Calkin sequences $C_{i,j} \to \Ind(C_{i,j}^{\kappa}) \to \Ind(\Calk^{\cont}_{\kappa}(C_{i,j}))$, which are functorial by \Cref{constr:Calkin_Functoriality}. By the first part and \Cref{prop:Filtered_Colimits_Fully_Faithful}(2), both resulting sequences are short exact. The comparison functors on the middle and right terms are equivalences by the compactly generated case. Passing to kernels, so is the comparison functor on the left terms.
	\end{proof}


		\chapter{Localizing invariants of dualizable categories}
	\label[chapter]{chap:Localizing_Invariants_Dualizable}

	% TODO

	\noindent\emph{This chapter is planned. No proofs are included yet.}
		\chapter{Sheaves on continuous posets and locally compact Hausdorff spaces}
	\label[chapter]{chap:Sheaves}

	% TODO

	\noindent\emph{This chapter is planned. No proofs are included yet.}
		\chapter{Dualizability and flatness}
	\label[chapter]{chap:Flatness}

	% TODO

	\noindent\emph{This chapter is planned. No proofs are included yet.}
		\chapter{\texorpdfstring{$\aleph_1$}{aleph\_1}-presentability of \texorpdfstring{$\Pr^{\dual}$}{Pr\textasciicircum dual}}
	\label[chapter]{chap:Aleph_1_Presentability}

	% TODO

	\noindent\emph{This chapter is planned. No proofs are included yet.}

	\appendix

	\titleformat{\chapter}{\normalfont\huge\bfseries}{\Alph{chapter}}{20pt}{\Huge}
	\titleformat{\section}{\Large\bfseries}{\Alph{chapter}.\arabic{section}}{10pt}{\Large}
	\titleformat{\subsection}{\large\bfseries}{\Alph{chapter}.\arabic{section}.\arabic{subsection}}{6pt}{\large}
	\titleformat{\subsubsection}{\large\bfseries}{}{6pt}{\large}
	\addtocontents{toc}{\vskip 15pt}
	\renewcommand{\thechapter}{\Alph{chapter}}
	\renewcommand{\thesection}{\Alph{chapter}.\arabic{section}}
	\renewcommand{\thesubsection}{\Alph{chapter}.\arabic{section}.\arabic{subsection}}
	\renewcommand{\thetheorem}{\Alph{chapter}.\arabic{theorem}}

		\chapter{Categorical background}
	\label[appendix]{app:Categorical_Background}

	In this appendix we collect background material on accessible and presentable categories that is used throughout the notes. Most of this material is standard, see \cite{lurie2009HTT,lurie2016HA}, and much of it is also treated in \cite{Cnossen_StableHomotopy,CCNW_Formalization}. We give complete proofs of the statements that are less easy to find in the literature or that are used in a specific form in the main text, and we refer to the literature for the foundational results on accessible categories.

	Throughout, \emph{category} means $\infty$-category, and $\An$ denotes the category of animae. We fix a regular cardinal $\kappa$ whenever it appears. For a small category $C$ we write $\PSh(C) = \Fun(C\catop,\An)$ and $y\colon C \hookrightarrow \PSh(C)$ for the Yoneda embedding.

	\section{Filtered and \texorpdfstring{$\kappa$}{kappa}-filtered categories}
	\label[section]{sec:Filtered_Categories}

	Filtered colimits are the basic tool throughout these notes: compactness, Ind-categories and accessibility are all defined in terms of them. We recall the relevant definitions, and prove some technical results on $\kappa$-filtered partially ordered sets and sequential diagrams, which are used in the study of $\omega_1$-presentability in \Cref{sec:Omega_1_Presentability}.

	\begin{definition}
		\label[definition]{def:Kappa_Filtered}
		A category $K$ is \emph{$\kappa$-small} if it is equivalent to (the nerve of) a simplicial set with fewer than $\kappa$ non-degenerate simplices. A category $I$ is \emph{$\kappa$-filtered} if every functor $K \to I$ with $K$ a $\kappa$-small category admits a cocone. We say \emph{filtered} for $\omega$-filtered. A colimit indexed by a (small) $\kappa$-filtered category is called a \emph{$\kappa$-filtered colimit}.
	\end{definition}

	\begin{definition}
		\label[definition]{def:Kappa_Compact}
		Let $C$ be a category which admits $\kappa$-filtered colimits. An object $x \in C$ is called \emph{$\kappa$-compact} if the functor $\Map_C(x,-)\colon C \to \An$ preserves $\kappa$-filtered colimits. We write $C^{\kappa} \subseteq C$ for the full subcategory of $\kappa$-compact objects, and say \emph{compact} for $\omega$-compact.
	\end{definition}

	We will freely use the following fundamental fact.

	\begin{theorem}[{\cite[Proposition~5.3.3.3]{lurie2009HTT}}]
		\label[theorem]{thm:Filtered_Colimits_Commute_With_Small_Limits}
		In $\An$, $\kappa$-filtered colimits commute with $\kappa$-small limits.
	\end{theorem}

	\begin{corollary}
		\label[corollary]{cor:Kappa_Small_Colimits_Of_Kappa_Compacts}
		Let $C$ be a category with $\kappa$-filtered colimits.
		\begin{enumerate}
			\item If $\kappa \leq \lambda$, then every $\kappa$-compact object of $C$ is $\lambda$-compact.
			\item If $C$ admits $\kappa$-small colimits, then $C^{\kappa} \subseteq C$ is closed under $\kappa$-small colimits.
			\item Let $x \in C$ and let $u\colon \Map_C(x,-) \to \Map_C(x',-)$ exhibit $\Map_C(x,-)$ as a retract of $\Map_C(x',-)$ in $\Fun(C,\An)$. If $x'$ is $\kappa$-compact, then so is $x$. In particular, $C^{\kappa}$ is closed under retracts.
		\end{enumerate}
	\end{corollary}
	\prooffor{cor:Kappa_Small_Colimits_Of_Kappa_Compacts}
	\begin{proof}
		(1) A $\lambda$-filtered category is $\kappa$-filtered.

		(2) If $x = \colim_{k \in K} x_k$ with $K$ $\kappa$-small and each $x_k$ $\kappa$-compact, then $\Map_C(x,-) = \lim_{k \in K\catop} \Map_C(x_k,-)$ is a $\kappa$-small limit of functors preserving $\kappa$-filtered colimits, hence preserves $\kappa$-filtered colimits by \Cref{thm:Filtered_Colimits_Commute_With_Small_Limits}.

		(3) A retract of a functor preserving $\kappa$-filtered colimits again preserves them, since a retract of an equivalence in $\An$ (namely the comparison map from the colimit) is an equivalence.
	\end{proof}

	For partially ordered sets, $\kappa$-filteredness admits the expected elementary description.

	\begin{lemma}
		\label[lemma]{lem:Kappa_Filtered_Posets}
		A partially ordered set $P$ is $\kappa$-filtered if and only if every subset $S \subseteq P$ of cardinality $< \kappa$ admits an upper bound in $P$. In particular, $P$ is filtered if and only if it is non-empty and every two elements admit a common upper bound.
	\end{lemma}
	\prooffor{lem:Kappa_Filtered_Posets}
	\begin{proof}
		Suppose $P$ is $\kappa$-filtered and let $S \subseteq P$ have cardinality $< \kappa$. Regarding $S$ as a discrete category, the inclusion $S \to P$ is a $\kappa$-small diagram, and a cocone on it is precisely an upper bound of $S$.

		Conversely, let $F\colon K \to P$ be a functor with $K$ $\kappa$-small. Then $K$ has fewer than $\kappa$ objects up to equivalence, so the image $S$ of $F$ on objects has cardinality $< \kappa$. Choose an upper bound $u$ of $S$. Since $\Fun(K,P)$ is again a partially ordered set, in which $F \leq G$ if and only if $F(k) \leq G(k)$ for all $k$, there is a (unique) natural transformation from $F$ to the constant functor with value $u$, i.e.\ a cocone on $F$.

		For the last statement, note that every finite subset of $P$ has an upper bound if and only if $P$ is non-empty and every two elements have an upper bound, by induction on the size of the subset.
	\end{proof}

	\begin{lemma}
		\label[lemma]{lem:Omega_1_Filtered_Via_Chains}
		A filtered partially ordered set $P$ is $\omega_1$-filtered if and only if every increasing sequence $p_0 \leq p_1 \leq p_2 \leq \dots$ in $P$ admits an upper bound.
	\end{lemma}
	\prooffor{lem:Omega_1_Filtered_Via_Chains}
	\begin{proof}
		The condition is clearly necessary. Conversely, let $S = \{s_0,s_1,s_2,\dots\} \subseteq P$ be a countable subset (a finite subset has an upper bound since $P$ is filtered). Using that $P$ is filtered, choose inductively $p_0 \geq s_0$ and $p_n \geq p_{n-1}, s_n$ for $n \geq 1$. Then an upper bound for the increasing sequence $(p_n)_n$ is an upper bound for $S$. We conclude by \Cref{lem:Kappa_Filtered_Posets}.
	\end{proof}

	We will often use the following criterion for cofinality, see \cite[Theorem~4.1.3.1]{lurie2009HTT}: a functor $f\colon A \to B$ is cofinal if and only if for every $b \in B$ the category $A \times_B B_{b/}$ is weakly contractible. If $A$ and $B$ are partially ordered sets, then $A \times_B B_{b/}$ is the partially ordered set $\{a \in A \mid b \leq f(a)\}$. Since filtered categories are weakly contractible \cite[Lemma~5.3.1.20]{lurie2009HTT}, we obtain:

	\begin{lemma}
		\label[lemma]{lem:Cofinality_Posets}
		Let $f\colon A \to B$ be a map of partially ordered sets such that for every $b \in B$ the subset $\{a \in A \mid b \leq f(a)\} \subseteq A$ is non-empty and filtered. Then $f$ is cofinal. In particular:
		\begin{enumerate}
			\item A map $f\colon \N \to B$ into a partially ordered set is cofinal if for every $b \in B$ there exists $n$ with $b \leq f(n)$.
			\item Every countable filtered partially ordered set $Q$ admits a cofinal map $\N \to Q$.
			\item The diagonal $\N \to \N \times \N$ is cofinal, and so is every map $f\colon \N \to \N$ with $f(n) \geq n$ for all $n$.
		\end{enumerate}
	\end{lemma}
	\prooffor{lem:Cofinality_Posets}
	\begin{proof}
		The first statement was explained above. For (1), the subset $\{n \mid b \leq f(n)\}$ is non-empty and upwards closed in $\N$ (as $f$ is order preserving), hence of the form $\{n \geq n_0\}$, which is filtered. For (2), if $Q$ is finite it has a largest element (\Cref{lem:Kappa_Filtered_Posets}) and we may take the constant map. Otherwise enumerate $Q = \{q_0,q_1,\dots\}$ and choose inductively $f(0) = q_0$ and $f(n) \geq f(n-1), q_n$; then $f$ is order preserving and cofinal by (1). Statement (3) follows directly from (1).
	\end{proof}

	\begin{theorem}[{\cite[Proposition~5.3.1.18]{lurie2009HTT}}]
		\label[theorem]{thm:Filtered_Posets_Cofinal}
		For every $\kappa$-filtered category $I$ there exist a $\kappa$-filtered partially ordered set $P$ and a cofinal functor $P \to I$.
	\end{theorem}

	\begin{proposition}
		\label[proposition]{prop:Colimits_Over_Countable_Subposets}
		Let $P$ be a filtered partially ordered set, and let $\Sigma$ be the set of countable filtered subsets $Q \subseteq P$, ordered by inclusion. Then $\Sigma$ is $\omega_1$-filtered and $P = \bigcup_{Q \in \Sigma} Q$. Moreover, for every functor $F\colon P \to D$ into a category with small colimits, there is a natural equivalence
		\[
			\colim_{p \in P} F(p) \simeq \colim_{Q \in \Sigma} \colim_{q \in Q} F(q).
		\]
	\end{proposition}
	\prooffor{prop:Colimits_Over_Countable_Subposets}
	\begin{proof}
		Every $p \in P$ lies in the countable filtered subset $\{p\}$. Let $Q_0,Q_1,\dots \in \Sigma$ be countably many elements. Put $U_0 = \bigcup_n Q_n$, and inductively let $U_{k+1} \supseteq U_k$ be obtained by adding a chosen common upper bound in $P$ for every pair of elements of $U_k$. Then $U = \bigcup_k U_k$ is countable, and any two elements of $U$ have an upper bound in $U$; hence $U \in \Sigma$ by \Cref{lem:Kappa_Filtered_Posets}, and $U$ is an upper bound of the $Q_n$. So $\Sigma$ is $\omega_1$-filtered by \Cref{lem:Kappa_Filtered_Posets}.

		Consider the partially ordered set $E = \{(Q,q) \mid Q \in \Sigma, q \in Q\}$, with $(Q,q) \leq (Q',q')$ if $Q \subseteq Q'$ and $q \leq q'$, together with the projections $\pi\colon E \to \Sigma$ and $\rho\colon E \to P$. The map $\rho$ is cofinal by \Cref{lem:Cofinality_Posets}: for $p \in P$, the subset $\{(Q,q) \mid p \leq q\}$ is non-empty, and it is filtered since two elements $(Q,q),(Q',q')$ are bounded by $(Q'',q'')$ where $q''$ is an upper bound of $q,q'$ and $Q'' \in \Sigma$ contains $Q \cup Q' \cup \{q''\}$. Hence $\colim_P F \simeq \colim_E F\rho \simeq \colim_{\Sigma} \pi_!(F\rho)$, where $\pi_!$ denotes left Kan extension along $\pi$. For $Q \in \Sigma$ we have $\pi_!(F\rho)(Q) = \colim_{(Q',q') \in E,\, Q' \subseteq Q} F(q')$, and the map $Q \to \{(Q',q') \in E \mid Q' \subseteq Q\}$, $q \mapsto (Q,q)$, is cofinal by \Cref{lem:Cofinality_Posets}: for $(Q',q')$ with $Q' \subseteq Q$, the subset $\{q \in Q \mid q' \leq q\}$ is non-empty and filtered. Hence $\pi_!(F\rho)(Q) \simeq \colim_{q \in Q} F(q)$, which proves the claim.
	\end{proof}

	\begin{remark}
		\label[remark]{rem:Sequential_Diagrams}
		Since the nerve of $\N$ contains the simplicial subset $\bigcup_n \Delta^{\{n,n+1\}}$ via an inner anodyne map, a functor $\N \to D$ amounts to a sequence of objects and morphisms $x_0 \to x_1 \to x_2 \to \dots$ in $D$, without further coherence data. Similarly, since $\Delta^1 \times \bigcup_n \Delta^{\{n,n+1\}} \subseteq \Delta^1 \times \N$ is inner anodyne \cite[Corollary~2.3.2.4]{lurie2009HTT}, a natural transformation $\eta\colon x_{\bullet} \to x'_{\bullet}$ of sequences amounts to morphisms $\eta_n\colon x_n \to x'_n$ together with homotopies witnessing the commutativity of each square
		\[
			\begin{tikzcd}
				x_n \rar \dar["\eta_n"'] & x_{n+1} \dar["\eta_{n+1}"] \\
				x'_n \rar & x'_{n+1}.
			\end{tikzcd}
		\]
		We will use this freely.
	\end{remark}

	\begin{lemma}[Telescopes]
		\label[lemma]{lem:Retract_As_Telescope}
		Let $D$ be a category with sequential colimits, and let $X \xrightarrow{s} Y \xrightarrow{r} X$ be maps with $rs \simeq \id_X$. Then $X \simeq \colim(Y \xrightarrow{e} Y \xrightarrow{e} Y \to \dots)$, where $e = sr$.
	\end{lemma}
	\prooffor{lem:Retract_As_Telescope}
	\begin{proof}
		Consider the sequence $Y \xrightarrow{r} X \xrightarrow{s} Y \xrightarrow{r} X \xrightarrow{s} \dots$. The maps $\N \to \N$ given by $n \mapsto 2n$ and $n \mapsto 2n+1$ are cofinal by \Cref{lem:Cofinality_Posets}. Restricting along the first gives the sequence $Y \xrightarrow{e} Y \xrightarrow{e} \dots$, and restricting along the second gives a sequence $X \to X \to \dots$ all of whose maps are equivalences $rs \simeq \id_X$; the colimit of the latter is $X$. Hence both colimits agree with the colimit of the interleaved sequence.
	\end{proof}

	We now describe the $\omega_1$-compact objects of $\Ind(C)$. Recall that for a small category $C$, the category $\Ind(C) \subseteq \PSh(C)$ is the full subcategory spanned by the filtered colimits of representable presheaves. It admits filtered colimits, computed in $\PSh(C)$, and the representable presheaves are compact in $\Ind(C)$ \cite[Section~5.3.5]{lurie2009HTT}. We usually suppress the Yoneda embedding $y$ when no confusion can arise, and write $\quotcolim_i c_i$ for the colimit in $\Ind(C)$ of a diagram of representables when we want to stress that it is a formal colimit.

	\begin{lemma}[Strictification of maps between sequential ind-objects]
		\label[lemma]{lem:Strictification_Sequential_Maps}
		Let $C$ be a small category, let $a_{\bullet}, b_{\bullet}\colon \N \to C$ be sequences, and put $Y = \colim_n y(a_n)$ and $Z = \colim_n y(b_n)$ in $\Ind(C)$. For every map $g\colon Y \to Z$ there exist a strictly increasing map $f\colon \N \to \N$ with $f(n) \geq n$ and a natural transformation $\eta\colon a_{\bullet} \to b_{f(\bullet)}$ of functors $\N \to C$ such that the induced map
		\[
			Y = \colim_n y(a_n) \xrightarrow{\colim \eta} \colim_n y(b_{f(n)}) \simeq Z
		\]
		is homotopic to $g$. Here the last equivalence comes from cofinality of $f$ (\Cref{lem:Cofinality_Posets}).
	\end{lemma}
	\prooffor{lem:Strictification_Sequential_Maps}
	\begin{proof}
		Denote by $\alpha_n\colon a_n \to a_{n+1}$ and $\beta_{m,m'}\colon b_m \to b_{m'}$ the transition maps, and by $\iota_m\colon b_m \to Z$ the canonical maps. Since $\Map(Y,Z) \simeq \lim_n \Map(a_n,Z)$, the map $g$ amounts to maps $g_n\colon a_n \to Z$ together with paths $\gamma_n\colon g_{n+1}\alpha_n \simeq g_n$ in $\Map(a_n,Z)$. Since $a_n$ is compact in $\Ind(C)$, we have $\Map(a_n,Z) \simeq \colim_m \Map(a_n,b_m)$. By \Cref{thm:Filtered_Colimits_Commute_With_Small_Limits}, path spaces in this filtered colimit are computed as filtered colimits of path spaces; in particular every point and every path in $\Map(a_n,Z)$ between images of points of $\Map(a_n,b_m)$ lifts, up to homotopy, to some stage $m' \geq m$.

		By \Cref{rem:Sequential_Diagrams}, the data of $f$ and $\eta$ amount to maps $\eta_n\colon a_n \to b_{f(n)}$ and homotopies $\theta_n\colon \eta_{n+1}\alpha_n \simeq \beta_{f(n),f(n+1)}\eta_n$. The induced map $\colim \eta$ corresponds to the point $(\iota_{f(n)}\eta_n)_n$ of $\lim_n \Map(a_n,Z)$, with the paths $\iota_{f(n+1)}\eta_{n+1}\alpha_n \simeq \iota_{f(n+1)}\beta_{f(n),f(n+1)}\eta_n \simeq \iota_{f(n)}\eta_n$ induced by $\theta_n$. To exhibit a homotopy $\colim\eta \simeq g$, it therefore suffices to choose paths $\sigma_n\colon \iota_{f(n)}\eta_n \simeq g_n$ in $\Map(a_n,Z)$ such that for each $n$ the two resulting paths from $\iota_{f(n+1)}\eta_{n+1}\alpha_n$ to $g_n$, namely $\sigma_n \circ \iota\theta_n$ and $\gamma_n \circ \sigma_{n+1}\alpha_n$, are homotopic.

		We construct these data inductively. First choose $f(0)$, $\eta_0\colon a_0 \to b_{f(0)}$ and $\sigma_0\colon \iota_{f(0)}\eta_0 \simeq g_0$. Suppose $f(0),\dots,f(n)$, $\eta_0,\dots,\eta_n$, $\theta_0,\dots,\theta_{n-1}$ and $\sigma_0,\dots,\sigma_n$ have been constructed. Choose $m > f(n)$, a map $\eta'\colon a_{n+1} \to b_m$ and a path $\sigma'\colon \iota_m\eta' \simeq g_{n+1}$. The maps $\eta'\alpha_n$ and $\beta_{f(n),m}\eta_n$ in $\Map(a_n,b_m)$ have images in $\Map(a_n,Z)$ connected by the path
		\[
			\iota_m\eta'\alpha_n \overset{\sigma'\alpha_n}{\simeq} g_{n+1}\alpha_n \overset{\gamma_n}{\simeq} g_n \overset{\sigma_n^{-1}}{\simeq} \iota_{f(n)}\eta_n.
		\]
		Hence there exist $m' \geq m$ and a path $\theta$ from $\beta_{m,m'}\eta'\alpha_n$ to $\beta_{f(n),m'}\eta_n$ in $\Map(a_n,b_{m'})$ whose image in $\Map(a_n,Z)$ is homotopic to this path. Put $f(n+1) = \max(m',n+1)$, $\eta_{n+1} = \beta_{m,f(n+1)}\eta'$, let $\theta_n$ be the image of $\theta$, and let $\sigma_{n+1}$ be the path $\iota_{f(n+1)}\eta_{n+1} \simeq \iota_m\eta' \overset{\sigma'}{\simeq} g_{n+1}$. By construction, the compatibility condition above holds. This completes the induction.
	\end{proof}

	\begin{proposition}
		\label[proposition]{prop:Omega_1_Compact_In_Ind}
		Let $C$ be a small category and let $X \in \Ind(C)$. Then $X$ is $\omega_1$-compact if and only if $X \simeq \colim_n y(c_n)$ for some sequence $c_0 \to c_1 \to c_2 \to \dots$ in $C$.
	\end{proposition}
	\prooffor{prop:Omega_1_Compact_In_Ind}
	\begin{proof}
		If $X$ is a sequential colimit of representables, then $X$ is a countable colimit of compact objects, hence $\omega_1$-compact by \Cref{cor:Kappa_Small_Colimits_Of_Kappa_Compacts}.

		Conversely, let $X$ be $\omega_1$-compact. By definition of $\Ind(C)$ and \Cref{thm:Filtered_Posets_Cofinal}, we may write $X = \colim_{p \in P} y(c_p)$ for a filtered partially ordered set $P$. By \Cref{prop:Colimits_Over_Countable_Subposets}, we have $X \simeq \colim_{Q \in \Sigma} X_Q$ with $X_Q = \colim_{q \in Q} y(c_q)$ and $\Sigma$ $\omega_1$-filtered. As $X$ is $\omega_1$-compact, $\Map(X,X) \simeq \colim_{Q} \Map(X,X_Q)$, so the identity of $X$ factors as $X \xrightarrow{s} X_Q \xrightarrow{r} X$ for some $Q \in \Sigma$. By \Cref{lem:Cofinality_Posets}(2), $X_Q$ is a sequential colimit $Y = \colim_n y(a_n)$ of representables. By \Cref{lem:Retract_As_Telescope}, $X \simeq \colim(Y \xrightarrow{e} Y \xrightarrow{e} \dots)$ with $e = sr$.

		By \Cref{lem:Strictification_Sequential_Maps}, there exist $f\colon \N \to \N$ with $f(n) \geq n$ and a natural transformation $\eta\colon a_{\bullet} \to a_{f(\bullet)}$ inducing $e$ on colimits. Define a functor $\Phi\colon \N \times \N \to C$ by $\Phi(k,n) = a_{f^k(n)}$, with horizontal maps the transition maps of $a_{\bullet}$ and vertical maps $\eta_{f^k(n)}\colon a_{f^k(n)} \to a_{f^{k+1}(n)}$; by \Cref{rem:Sequential_Diagrams} this is the functor $\N \to \Fun(\N,C)$ given by the sequence $a_{\bullet} \to a_{f(\bullet)} \to a_{f^2(\bullet)} \to \dots$ of whiskerings of $\eta$. Since each $f^k$ is cofinal, $\colim_n y\Phi(k,n) \simeq Y$, and under these identifications the induced map in the $k$-direction is $\colim\eta \simeq e$. Hence
		\[
			X \simeq \colim_k \colim_n y\Phi(k,n) \simeq \colim_{(k,n) \in \N\times\N} y\Phi(k,n) \simeq \colim_n y\Phi(n,n),
		\]
		where the last equivalence uses cofinality of the diagonal (\Cref{lem:Cofinality_Posets}(3)). So $X$ is a sequential colimit of representables.
	\end{proof}

	\section{Accessible and presentable categories}
	\label[section]{sec:Accessible_Presentable}

	We recall the basic notions; see \cite[Chapter~5]{lurie2009HTT} for proofs of the results quoted in this section.

	\begin{definition}
		\label[definition]{def:Ind_Kappa}
		For a small category $C_0$, we let $\Ind_{\kappa}(C_0) \subseteq \PSh(C_0)$ denote the full subcategory spanned by the $\kappa$-filtered colimits of representable presheaves. We write $\Ind = \Ind_{\omega}$.
	\end{definition}

	\begin{theorem}[{\cite[Section~5.3.5]{lurie2009HTT}}]
		\label[theorem]{thm:Ind_Kappa_Universal_Property}
		Let $C_0$ be a small category.
		\begin{enumerate}
			\item The category $\Ind_{\kappa}(C_0)$ admits $\kappa$-filtered colimits, computed in $\PSh(C_0)$, and the Yoneda embedding $y\colon C_0 \hookrightarrow \Ind_{\kappa}(C_0)$ lands in $\kappa$-compact objects.
			\item For every category $D$ admitting $\kappa$-filtered colimits, restriction along $y$ induces an equivalence between the category of functors $\Ind_{\kappa}(C_0) \to D$ preserving $\kappa$-filtered colimits and $\Fun(C_0,D)$. The inverse is given by left Kan extension along $y$.
			\item If $C_0$ admits $\kappa$-small colimits, then $\Ind_{\kappa}(C_0)$ admits all small colimits, $y$ preserves $\kappa$-small colimits, and for every category $D$ admitting small colimits, restriction along $y$ induces an equivalence between colimit-preserving functors $\Ind_{\kappa}(C_0) \to D$ and functors $C_0 \to D$ preserving $\kappa$-small colimits.
		\end{enumerate}
	\end{theorem}

	\begin{definition}
		\label[definition]{def:Accessible}
		A category $C$ is \emph{$\kappa$-accessible} if it is locally small, admits $\kappa$-filtered colimits, the full subcategory $C^{\kappa}$ is essentially small, and every object of $C$ is a $\kappa$-filtered colimit of $\kappa$-compact objects. It is \emph{accessible} if it is $\kappa$-accessible for some $\kappa$. A functor between accessible categories is \emph{accessible} if it preserves $\kappa$-filtered colimits for some $\kappa$. A category is \emph{presentable} if it is accessible and admits all small colimits.
	\end{definition}

	\begin{theorem}[{\cite[Section~5.4.2]{lurie2009HTT}}]
		\label[theorem]{thm:Accessible_Recollections}
		\leavevmode
		\begin{enumerate}
			\item A category $C$ is $\kappa$-accessible if and only if $C \simeq \Ind_{\kappa}(C_0)$ for some small category $C_0$. In this case, the inclusion $C^{\kappa} \subseteq C$ induces an equivalence $\Ind_{\kappa}(C^{\kappa}) \simeq C$, and $C^{\kappa}$ is the idempotent completion of $C_0$.
			\item If $C$ is $\kappa$-accessible, then $C^{\lambda}$ is essentially small for every regular cardinal $\lambda$, and $C$ is $\lambda$-accessible for arbitrarily large $\lambda$.
		\end{enumerate}
	\end{theorem}

	\begin{theorem}[{\cite[Proposition~5.4.7.3]{lurie2009HTT}}]
		\label[theorem]{thm:Limits_Of_Accessible_Categories}
		The subcategory $\mathrm{Acc} \subseteq \widehat{\Cat}$ of accessible categories and accessible functors admits small limits, and the inclusion $\mathrm{Acc} \subseteq \widehat{\Cat}$ preserves them.
	\end{theorem}

	\begin{theorem}[Adjoint functor theorem, {\cite[Corollary~5.5.2.9]{lurie2009HTT}}]
		\label[theorem]{thm:Adjoint_Functor_Theorem}
		Let $F\colon C \to D$ be a functor between presentable categories. Then $F$ admits a right adjoint if and only if it preserves small colimits, and $F$ admits a left adjoint if and only if it is accessible and preserves small limits.
	\end{theorem}

	\begin{lemma}
		\label[lemma]{lem:Accessible_Objects_Lambda_Compact}
		Let $C$ be an accessible category. Then every object of $C$ is $\lambda$-compact for some regular cardinal $\lambda$. Consequently, $C = \bigcup_{\lambda} C^{\lambda}$, where each $C^{\lambda}$ is essentially small.
	\end{lemma}
	\prooffor{lem:Accessible_Objects_Lambda_Compact}
	\begin{proof}
		Let $C$ be $\kappa$-accessible and $x \in C$. Write $x = \colim_{i \in I} x_i$ with $I$ small and $\kappa$-filtered and each $x_i$ $\kappa$-compact. Choose a regular cardinal $\lambda \geq \kappa$ such that $I$ is $\lambda$-small. Each $x_i$ is $\lambda$-compact by \Cref{cor:Kappa_Small_Colimits_Of_Kappa_Compacts}(1). Moreover $\Map(x,-) \simeq \lim_{i \in I\catop}\Map(x_i,-)$ is a $\lambda$-small limit of functors preserving $\lambda$-filtered colimits, hence preserves them by \Cref{thm:Filtered_Colimits_Commute_With_Small_Limits}. The last statement follows from \Cref{thm:Accessible_Recollections}(2).
	\end{proof}

	\begin{construction}
		\label[construction]{constr:Ind_Large}
		Let $C$ be a locally small, possibly large, category. We let $\Ind(C) \subseteq \Fun(C\catop,\widehat{\An})$ denote the full subcategory spanned by the small filtered colimits of representable functors, where $\widehat{\An}$ denotes the category of large animae. If $C' \subseteq C$ is a small full subcategory, then left Kan extension along $C' \subseteq C$ restricts to a fully faithful functor $\Ind(C') \hookrightarrow \Ind(C)$ which preserves filtered colimits: indeed, for filtered diagrams $(c_i)$ and $(d_j)$ in $C'$ both categories compute
		\[
			\Map(\quotcolim_i c_i, \quotcolim_j d_j) \simeq \lim_i \colim_j \Map_C(c_i,d_j).
		\]
		If $C$ is accessible, then \Cref{lem:Accessible_Objects_Lambda_Compact} shows that
		\[
			\Ind(C) = \bigcup_{\lambda} \Ind(C^{\lambda}),
		\]
		since a small filtered diagram in $C$ takes values in $C^{\lambda}$ for $\lambda$ large enough.

		If $C$ admits filtered colimits, then the functor $\colim\colon \Ind(C) \to C$ sending $\quotcolim_i c_i$ to $\colim_i c_i$ (the unique filtered-colimit-preserving extension of $\id_C$, by the analogue of \Cref{thm:Ind_Kappa_Universal_Property}(2)) is left adjoint to the Yoneda embedding $y\colon C \hookrightarrow \Ind(C)$, since
		\[
			\Map_C(\colim_i c_i, c) \simeq \lim_i \Map_C(c_i,c) \simeq \Map_{\Ind(C)}(\quotcolim_i c_i, y(c)).
		\]
	\end{construction}

	\section{Idempotent completeness}
	\label[section]{sec:Idempotent_Completeness}

	Recall from \cite[Section~4.4.5]{lurie2009HTT} that a category $C$ is \emph{idempotent complete} if every idempotent in $C$ splits. We will use the following characterization, see \cite[Section~5.1.4]{lurie2009HTT}: a locally small category $C$ is idempotent complete if and only if the essential image of the Yoneda embedding $C \hookrightarrow \Fun(C\catop,\widehat{\An})$ is closed under retracts, if and only if the essential image of $C\catop \hookrightarrow \Fun(C,\widehat{\An})$ is closed under retracts. Every category admitting sequential colimits is idempotent complete \cite[Section~4.4.5]{lurie2009HTT} (compare \Cref{lem:Retract_As_Telescope}); this applies in particular to presentable categories and to all categories admitting filtered colimits. Moreover, every accessible category is idempotent complete \cite[Chapter~5]{lurie2009HTT}, and every small category $C_0$ admits an idempotent completion $C_0 \hookrightarrow C_0^{\idem}$, which may be computed as the full subcategory of $\PSh(C_0)$ spanned by retracts of representables.


	% Source: Talk 2 (Marc Hoyois, 20.10.2026)
	\begin{lemma}
		\label[lemma]{lem:Retract_Left_Adjoint}
		Let $g\colon C \to D$ be a retract of $g'\colon C' \to D'$ in $\Fun(\Delta^1,\widehat{\Cat})$, where $C,C',D,D'$ are locally small. Suppose that $g'$ has a left adjoint $f'$ and that $C$ is idempotent complete. Then $g$ has a left adjoint.
	\end{lemma}
	\prooffor{lem:Retract_Left_Adjoint}
	\begin{proof}
		Consider a retract diagram
		\[
			\begin{tikzcd}
				C \rar["u"] \dar["g"'] & C' \rar["v"] \dar["g'"] & C \dar["g"] \\
				D \rar["\hat{u}"] & D' \rar["\hat{v}"] & D,
			\end{tikzcd}
		\]
		with $vu \simeq \id_C$ and $\hat{v}\hat{u} \simeq \id_D$, compatibly with the squares. By \cite[Proposition~5.2.4.2]{lurie2009HTT}, it suffices to show that for every $d \in D$ the functor $F_d = \Map_D(d,g(-))\colon C \to \An$ is corepresentable. Since $C$ is idempotent complete, it suffices to exhibit $F_d$ as a retract of the corepresentable functor $\Map_C(x,-)$, where $x = vf'\hat{u}(d)$. Denote by $\eta$ the unit of the adjunction $f' \dashv g'$ and by $\phi \mapsto \phi^{\flat}$ the adjunction bijection $\Map_{D'}(e,g'(c')) \simeq \Map_{C'}(f'(e),c')$. We define natural transformations
		\[
			F_d \xrightarrow{\;i\;} \Map_C(x,-) \xrightarrow{\;p\;} F_d
		\]
		as follows. For $c \in C$ and $\phi\colon d \to g(c)$, let $i(\phi)$ be the composite
		\[
			x = vf'\hat{u}(d) \xrightarrow{v((\hat{u}\phi)^{\flat})} vu(c) \simeq c,
		\]
		where $\hat{u}\phi\colon \hat{u}(d) \to \hat{u}g(c) \simeq g'u(c)$. For $\chi\colon x \to c$, let $p(\chi)$ be the composite
		\[
			d \simeq \hat{v}\hat{u}(d) \xrightarrow{\hat{v}(\eta_{\hat{u}d})} \hat{v}g'f'\hat{u}(d) \simeq gvf'\hat{u}(d) = g(x) \xrightarrow{g(\chi)} g(c).
		\]
		Both constructions are natural in $c$. Using $gv \simeq \hat{v}g'$ and the triangle identity $g'((\hat u \phi)^{\flat}) \circ \eta_{\hat{u}d} = \hat{u}\phi$, we compute
		\[
			p(i(\phi)) = gv((\hat u\phi)^{\flat}) \circ \hat{v}(\eta_{\hat u d}) \simeq \hat{v}\big(g'((\hat u\phi)^{\flat})\circ \eta_{\hat u d}\big) = \hat{v}\hat{u}(\phi) \simeq \phi,
		\]
		so $pi \simeq \id$. Hence $F_d$ is a retract of $\Map_C(x,-)$, as desired.
	\end{proof}

	\section{Mates}
	\label[section]{sec:Mates}

	We freely use the calculus of mates in the homotopy $2$-category of categories, see \cite{KellyStreet1974Mates} and \cite{Riehl_Verity}. Since invertibility of a natural transformation can be checked objectwise in homotopy categories, all statements in this section only concern natural transformations up to homotopy.

	\begin{construction}
		\label[construction]{constr:Mates}
		Consider a square of categories
		\[
			\begin{tikzcd}
				A \rar["G"] \dar["U"'] & B \dar["V"] \\
				A' \rar["G'"'] & B'
			\end{tikzcd}
		\]
		commuting via a natural isomorphism $\alpha\colon VG \simeq G'U$, and suppose that $G$ and $G'$ admit left adjoints $F$ and $F'$, with units $\eta,\eta'$ and counits $\varepsilon,\varepsilon'$. The \emph{mate} of $\alpha$ is the natural transformation $\beta\colon F'V \to UF$ given by the composite
		\[
			F'V \xrightarrow{F'V\eta} F'VGF \overset{\alpha}{\simeq} F'G'UF \xrightarrow{\varepsilon' UF} UF.
		\]
		If instead $U$ and $V$ admit right adjoints $U^R$ and $V^R$, the \emph{vertical mate} of $\alpha$ is the natural transformation $\gamma\colon GU^R \to V^RG'$ given by
		\[
			GU^R \to V^RVGU^R \overset{\alpha}{\simeq} V^RG'UU^R \to V^RG'.
		\]
	\end{construction}

	\begin{lemma}[Conjugate transformations]
		\label[lemma]{lem:Conjugate_Transformations}
		Let $F_1 \dashv G_1$ and $F_2 \dashv G_2$ be adjunctions between categories $A$ and $B$, with $F_i\colon A \to B$. For every natural transformation $\phi\colon F_1 \to F_2$ there is a unique (up to homotopy) natural transformation $\phi^{\vee}\colon G_2 \to G_1$, its \emph{conjugate}, such that the square
		\[
			\begin{tikzcd}
				\Map_B(F_2a,b) \rar["\phi^*"] \dar["\simeq"'] & \Map_B(F_1a,b) \dar["\simeq"] \\
				\Map_A(a,G_2b) \rar["(\phi^{\vee})_*"] & \Map_A(a,G_1b)
			\end{tikzcd}
		\]
		commutes naturally in $a$ and $b$. Moreover, $\phi$ is an equivalence if and only if $\phi^{\vee}$ is.
	\end{lemma}
	\prooffor{lem:Conjugate_Transformations}
	\begin{proof}
		The composite of the top map with the vertical equivalences is a natural transformation $\Map_A(-,G_2(-)) \to \Map_A(-,G_1(-))$, which by the Yoneda lemma is induced by a unique natural transformation $\phi^{\vee}\colon G_2 \to G_1$. Again by the Yoneda lemma, $\phi$ is an equivalence if and only if $\phi^*$ is an equivalence for all $a,b$, if and only if $(\phi^{\vee})_*$ is an equivalence for all $a,b$, if and only if $\phi^{\vee}$ is an equivalence.
	\end{proof}

	\begin{lemma}[{\cite[Section~2.2, Propositions~2.1 and~2.2]{KellyStreet1974Mates}}]
		\label[lemma]{lem:Mates_Conjugates}
		In the situation of \Cref{constr:Mates}, assume that $G,G'$ admit left adjoints $F,F'$ and that $U,V$ admit right adjoints $U^R,V^R$. Then the mate $\beta\colon F'V \to UF$ (a transformation between left adjoints of $V^RG'$ and $GU^R$) is conjugate to the vertical mate $\gamma\colon GU^R \to V^RG'$. In particular, $\beta$ is an equivalence if and only if $\gamma$ is.
	\end{lemma}

	\begin{lemma}[Adjoint triples]
		\label[lemma]{lem:Adjoint_Triples}
		Let $F \dashv G \dashv H$ be an adjoint triple, with $F,H\colon A \to B$ and $G\colon B \to A$. Then $F$ is fully faithful if and only if $H$ is fully faithful.
	\end{lemma}
	\prooffor{lem:Adjoint_Triples}
	\begin{proof}
		We have adjunctions $GF \dashv GH$ and $\id_A \dashv \id_A$. The unit $\eta\colon \id_A \to GF$ of $F \dashv G$ is a transformation between the left adjoints, and its conjugate is the counit $\varepsilon\colon GH \to \id_A$ of $G \dashv H$: indeed, under the identifications $\Map_A(GFa,a') \simeq \Map_B(Fa,Ha') \simeq \Map_A(a,GHa')$, precomposition with $\eta_a$ corresponds to postcomposition with $\varepsilon_{a'}$, by the triangle identities. Now $F$ is fully faithful if and only if $\eta$ is an equivalence, and $H$ is fully faithful if and only if $\varepsilon$ is an equivalence, so the claim follows from \Cref{lem:Conjugate_Transformations}.
	\end{proof}

	The following was used in the definition of strongly continuous functors between compactly assembled categories, see \Cref{def:Strongly_Continuous_Compactly_Assembled} and \Cref{rem:Strongly_Continuous_Right_Adjoint}.

	\begin{proposition}
		\label[proposition]{prop:Strongly_Continuous_Via_Right_Adjoint}
		Let $C$ and $D$ be compactly assembled categories, with left adjoints $\hat{y}_C$ and $\hat{y}_D$ of the colimit functors, and let $f\colon C \to D$ be a functor with a right adjoint $g$. Then the mate
		\[
			\beta\colon \hat{y}_D f \to \Ind(f)\hat{y}_C
		\]
		of the commutative square $\colim_D \circ \Ind(f) \simeq f \circ \colim_C$ is an equivalence if and only if $g$ preserves filtered colimits.
	\end{proposition}
	\prooffor{prop:Strongly_Continuous_Via_Right_Adjoint}
	\begin{proof}
		Applying $\Ind$ to the unit and counit of $f \dashv g$ shows that $\Ind(f) \dashv \Ind(g)$. By \Cref{lem:Mates_Conjugates}, $\beta$ is an equivalence if and only if the vertical mate
		\[
			\gamma\colon \colim_C\Ind(g) \to gf\colim_C\Ind(g) \simeq g\colim_D\Ind(f)\Ind(g) \to g\colim_D
		\]
		is an equivalence. Evaluating at an object $\quotcolim_i d_i \in \Ind(D)$, the map $\gamma$ becomes
		\[
			\colim_i g(d_i) \to gf\big(\colim_i g(d_i)\big) \simeq g\big(\colim_i fg(d_i)\big) \to g\big(\colim_i d_i\big),
		\]
		whose restriction to each $g(d_j)$ is $g(d_j) \to gfg(d_j) \to g(d_j)$ followed by the canonical map $g(d_j) \to g(\colim_i d_i)$. By the triangle identity, $\gamma$ is therefore the canonical comparison map $\colim_i g(d_i) \to g(\colim_i d_i)$. Hence $\gamma$ is an equivalence if and only if $g$ preserves filtered colimits.
	\end{proof}

	\section{Tensor products of presentable categories}
	\label[section]{sec:PrL_Tensor_Products}

	We denote by $\PrL \subseteq \widehat{\Cat}$ the subcategory of presentable categories and colimit-preserving functors, and by $\PrR \subseteq \widehat{\Cat}$ the subcategory of presentable categories and accessible limit-preserving functors. By \Cref{thm:Adjoint_Functor_Theorem}, passing to adjoints gives an equivalence $(\PrL)\catop \simeq \PrR$ \cite[Corollary~5.5.3.4]{lurie2009HTT}. Both $\PrL$ and $\PrR$ admit small limits, computed in $\widehat{\Cat}$ \cite[Proposition~5.5.3.13 and Theorem~5.5.3.18]{lurie2009HTT}; in particular, colimits in $\PrL$ may be computed as limits in $\widehat{\Cat}$ along the right adjoints.

	% Source: Talk 2 (Marc Hoyois, 20.10.2026)
	\begin{theorem}[{\cite[Section~4.8.1]{lurie2016HA}}]
		\label[theorem]{thm:Lurie_Tensor_Product}
		The category $\PrL$ admits a symmetric monoidal structure $- \otimes -$ with unit $\An$, characterized by the property that colimit-preserving functors $C \otimes D \to E$ correspond to functors $C \times D \to E$ preserving colimits in each variable separately:
		\[
			\Fun^{\colim}(C \otimes D, E) \simeq \Fun^{\colim}(C,\Fun^{\colim}(D,E)).
		\]
		In particular, $\PrL$ is closed, with internal hom given by $\Fun^{\colim}(-,-)$. Moreover, there is a natural equivalence
		\[
			C \otimes D \simeq \Fun^{\lim}(C\catop,D),
		\]
		where the right-hand side denotes the full subcategory of $\Fun(C\catop,D)$ spanned by the functors preserving small limits.
	\end{theorem}

	\begin{remark}
		\label[remark]{rem:Tensor_Product_Description}
		The second description arises as follows: one first shows that $\Fun^{\lim}(C\catop,D)$ is presentable, and then one computes that
		\[
			\Fun^{\colim}(\Fun^{\lim}(C\catop,D),E) \simeq \Fun^{\colim}(C,\Fun^{\colim}(D,E))
		\]
		for every presentable $E$.
	\end{remark}

	The full subcategory $\PrL_{\st} \subseteq \PrL$ of stable presentable categories is closed under the tensor product, and becomes a symmetric monoidal category with unit $\Sp$; moreover, a presentable category $C$ is stable if and only if the map $C \to \Sp \otimes C$ is an equivalence \cite[Section~4.8.2]{lurie2016HA}.

	\begin{proposition}
		\label[proposition]{prop:Tensor_With_Ind}
		Let $C_0$ be a small category with finite colimits and let $D$ be a presentable category. There is an equivalence, natural in $D$,
		\[
			\Ind(C_0) \otimes D \simeq \Fun^{\lex}(C_0\catop,D),
		\]
		where the right-hand side denotes the full subcategory of $\Fun(C_0\catop,D)$ spanned by the functors preserving finite limits. If $C_0$ and $D$ are stable, this gives natural equivalences
		\[
			\Ind(C_0) \otimes D \simeq \Fun^{\ex}(C_0\catop,D) \simeq \Fun^{\colim}(\Ind(C_0\catop),D).
		\]
	\end{proposition}
	\prooffor{prop:Tensor_With_Ind}
	\begin{proof}
		By \Cref{thm:Lurie_Tensor_Product}, we have
		\[
			\Ind(C_0) \otimes D \simeq \Fun^{\lim}(\Ind(C_0)\catop,D) \simeq \Fun^{\colim}(\Ind(C_0),D\catop)\catop.
		\] Since $D$ admits small limits, $D\catop$ admits small colimits, so by \Cref{thm:Ind_Kappa_Universal_Property}(3) restriction along $y$ gives an equivalence
		\[
			\Fun^{\colim}(\Ind(C_0),D\catop) \simeq \Fun^{\rex}(C_0,D\catop) = \Fun^{\lex}(C_0\catop,D)\catop.
		\] This proves the first claim. In the stable case, a functor between stable categories preserves finite limits if and only if it is exact, and applying \Cref{thm:Ind_Kappa_Universal_Property}(3) to $C_0\catop$ (which again has finite colimits) and the cocomplete category $D$ gives $\Fun^{\ex}(C_0\catop,D) \simeq \Fun^{\colim}(\Ind(C_0\catop),D)$.
	\end{proof}

	\begin{corollary}
		\label[corollary]{cor:Ind_As_Exact_Functors}
		For a small stable category $A$, the restricted Yoneda functor induces an equivalence $\Ind(A) \simeq \Fun^{\ex}(A\catop,\Sp)$.
	\end{corollary}
	\prooffor{cor:Ind_As_Exact_Functors}
	\begin{proof}
		Apply \Cref{prop:Tensor_With_Ind} with $D = \Sp$, using that $\Ind(A) \otimes \Sp \simeq \Ind(A)$ since $\Ind(A)$ is stable.
	\end{proof}

	% Source: Talk 2 (Marc Hoyois, 20.10.2026)
	\begin{corollary}
		\label[corollary]{cor:Compactly_Generated_Dualizable}
		Let $C_0$ be a small stable category. Then $C = \Ind(C_0)$ is dualizable in $\PrL_{\st}$, with dual
		\[
			C^{\vee} = \Ind(C_0\catop).
		\]
		The evaluation map $\Ind(C_0\catop) \otimes \Ind(C_0) \to \Sp$ is the colimit-preserving functor which restricts to the mapping spectrum functor
		\[
			\hom_{C_0}\colon C_0\catop \times C_0 \to \Sp,
		\]
		and the coevaluation $\Sp \to \Ind(C_0) \otimes \Ind(C_0\catop) \simeq \Fun^{\ex}(C_0,\Ind(C_0))$ sends the sphere spectrum to the Yoneda embedding $C_0 \hookrightarrow \Ind(C_0)$.
	\end{corollary}
	\prooffor{cor:Compactly_Generated_Dualizable}
	\begin{proof}
		We use the criterion \cite[Section~``Symmetric monoidal duality'']{Cnossen_StableHomotopy}: an object $X$ of a closed symmetric monoidal category is dualizable if and only if for every $Z$ the canonical map $Z \otimes \iHom(X,\unit) \to \iHom(X,Z)$, adjoint to $Z \otimes \iHom(X,\unit) \otimes X \xrightarrow{\id \otimes \ev} Z$, is an equivalence. For $X = \Ind(C_0)$ we have $\iHom(X,Z) = \Fun^{\colim}(\Ind(C_0),Z) \simeq \Fun^{\ex}(C_0,Z)$ by \Cref{thm:Ind_Kappa_Universal_Property}(3), and in particular $\iHom(X,\Sp) \simeq \Fun^{\ex}(C_0,\Sp) \simeq \Ind(C_0\catop)$ by \Cref{cor:Ind_As_Exact_Functors}. Under these identifications, $\ev\colon \Ind(C_0\catop) \otimes \Ind(C_0) \to \Sp$ is the colimit-preserving functor restricting on $C_0\catop \times C_0$ to $(a,c) \mapsto \hom_{C_0}(a,c)$, and the canonical map $Z \otimes \Ind(C_0\catop) \to \Fun^{\ex}(C_0,Z)$ agrees with the equivalence of \Cref{prop:Tensor_With_Ind}. Indeed, both are colimit-preserving functors out of $Z \otimes \Ind(C_0\catop)$, hence determined by their restrictions to $Z \times C_0\catop$ by \Cref{thm:Lurie_Tensor_Product} and \Cref{thm:Ind_Kappa_Universal_Property}(3); the canonical map sends $(z,a)$ to $z \otimes \ev(a,-) = z \otimes \hom_{C_0}(a,-)$, and so does the equivalence of \Cref{prop:Tensor_With_Ind}, since the equivalence $C \otimes D \simeq \Fun^{\lim}(C\catop,D)$ of \Cref{thm:Lurie_Tensor_Product} sends an elementary tensor $c \otimes d$ to the functor $c' \mapsto \Map_C(c',c) \otimes d$ \cite[Proof of Proposition~4.8.1.17]{lurie2016HA}.
		Hence $X$ is dualizable with dual $\Ind(C_0\catop)$ and evaluation as stated. The coevaluation corresponds, under $\iHom(X,X) \simeq X \otimes X^{\vee}$, to the identity of $X$, which under $\Fun^{\colim}(\Ind(C_0),\Ind(C_0)) \simeq \Fun^{\ex}(C_0,\Ind(C_0))$ corresponds to the Yoneda embedding.
	\end{proof}

	\begin{remark}
		\label[remark]{rem:Tensor_With_Compactly_Generated}
		In other words, tensoring with $\Ind(C_0)$ in $\PrL_{\st}$ is equivalent to forming the internal hom out of $\Ind(C_0\catop)$.
	\end{remark}


	\addtocontents{toc}{\vskip 10pt}
%	\addcontentsline{toc}{chapter}{Bibliography}
	\printbibliography[heading=bibintoc]

\end{document}
